% =====================================================================
%  Harald Høffding, OPLEVELSE OG TYDNING.  Religionsfilosofiske Studier.
%  København: Nyt Nordisk Forlag (Otto Halling & Chr. Nielsen), 1918.
%  Krohns Bogtrykkeri 1918.  238 pp.
%
%  GENERATED; DO NOT EDIT BY HAND.  Built by ebook-method/oplevelse-og-tydning/
%  (rebuild.sh) from the scan, the collation decisions and the patches there.
%
%  SOURCE: the National Library of Norway's scan (nb.no item
%  363c131e7e2e2c6e3b0602f49dae5268, URN:NBN:no-nb_digibok_2012041108082; the
%  Stavanger kommunebibliotek copy), ~/bibliotek/Høffding, Harald/
%  1918-oplevelse-og-tydning.pdf (SCANS.tsv).  PDF page = printed page + 4.
%
%  EDITION CHECK (2026-10-14): the year is 1918 (title verso "KROHNS
%  BOGTRYKKERI 1918"; the Fortale is dated 19 August 1918; NB's record 1918);
%  the title page itself carries no year.  The SAGA / Lindhardt og Ringhof
%  e-book (ISBN 9788726363029, "Copyright 1918, 2020") is the same text
%  (same contents, same Fortale, orthography kept); it was used only as a
%  voting witness in the collation.
%
%  HOW ACCURACY WAS ESTABLISHED (2026-10-14; ebook-method/oplevelse-og-tydning/, README entry there):
%  (a) Draft: NB's own OCR (ALTO) laid over the page images (nb_scan.py), turned into LaTeX by rule (draft.py);
%      no agent re-typed the text.  Headings from the Indhold and the e-book; the notes' boundaries checked by
%      script against the e-book's 141 numbered notes (notecmp.py: 141 = 141, after 5 merged notes were split).
%  (b) Word collation (collate/check.py, punctuation included) against an INDEPENDENT tesseract reading of the
%      images: 1637 disagreements.  The SAGA / Lindhardt og Ringhof e-book (2020) of this 1918 text voted where the
%      two OCRs differed (rules.py; the vote was first scored by script on Etik's 204 checks decided at the image:
%      35 right, 2 wrong, both marks only).  413 checks went to the image: blind crop sheets read by five readers
%      (b1, 88/88 planted controls right), 36 of the sheets' crops an audit of the e-book vote (1 wrong of 26
%      confirmed: „Bakehantismen“, mended; every e-book decision whose word is unattested in the corpus was then
%      re-read).  OPEN 0, italics OPEN 0 (NB's italic flags + stroke slant; 16 read blind, i1, 4/4 controls).
%  (c) The footnotes (small type) read WHOLE and blind at 300 dpi (n1, three readers, 100 note strips; 13 of 16
%      control lines right, the 3 others controls whose "known" text was NB's own error): 62 word differences, chiefly
%      French accents both OCRs had missed (Térèse, Genève, très), mended; the readers' letterspacing (authors'
%      names: \emph) and italics (titles: \textit) applied from their reading (notefix.py -> postfix.tsv).
%  (d) The whole text compared word by word with the e-book (ebookdiff.py): 191 differences; the e-book's own
%      OCR faults set aside, 62 read at the image (b2, 14/14 controls) and the rest settled as plain OCR faults.
%  (e) A blind sampled page check (pagecheck.py): 10 pages drawn at random from five strata (pp. 43, 55, 69, 101,
%      126, 147, 195, 200, 220, 230; 2948 words, footnotes included) transcribed whole by three readers who saw no
%      text: 0 word, 0 punctuation and 0 markup errors (95 % upper bound 1.0 per 1000 words); the 3 differences
%      were the readers' (a heading numeral's stop, a word divided at the page foot).
%  Title page, verso, Indhold and the subtitle on p. 9 read at the image by the main session.
%  (f) The translators' odd spots (128 that touch the Danish text): 109 read blind at the image (b3, two readers,
%      24/24 controls; the readers' 2 departures, „Cécilie“ p. 111 and „Eckchardt“ p. 179, re-read at 600 dpi: printed
%      so), 19 settled against the e-book and the image.  7 mended: „Modstykke“ (p. 33), a speck read as ~ (p. 36)
%      and a worn hyphen read as ~ (hver-|ken, p. 170), a speck read as a stop (p. 171), „Enten---Eller“ (p. 72), the
%      extent of an italic title (p. 130), a word space (p. 183).  The rest are the print's own, kept.
%  NOT done: a whole-book reading.
%
%  EDITORIAL: diplomatic.  Printed as is: „RELIGIONPSYKOLOGISK“ (p. 9), „Ekssempel“ (p. 132), „Nigthingale“,
%  „Dcouments“, „frane.“, „philosophiea«“, „positivs“, „Edwad Cooke“ (p. 106), „Eckchardt“, „Cécilie“ (p. 111),
%  „Grnndtvig“, „Welhausen“, „undergaaaede“, „Udvik-|viklingsmaade“, „en en“ (p. 207), „IX“ for Plotinos' Enneads,
%  the Indhold's „Bakchantisme“ and „Djævlebesættelse“ against the headings' „Bakchantismen.“ and
%  „Djævlebesættelser“ (both kept), the notes' labels *) **) ***) and, on pp. 25 and 182, †) ††) (the print's
%  sequence).  „Austrálien“ (p. 138 note): an accent-like mark over the a, read at the image.  The ø rule applied
%  where a plain o is no word (oefix.py: no site).  Running heads are the folio only; sheet signatures dropped.  Page markers
%  \opage{N} at the exact page turns (NB's ALTO lines); a word divided across a page turn is joined with %.
% =====================================================================
\documentclass[12pt,a4paper,openany]{book}
\usepackage[utf8]{inputenc}
\usepackage[T1]{fontenc}
\usepackage{libertinus}
\usepackage{textalpha}
\usepackage[danish]{babel}
\usepackage{setspace}
\usepackage[margin=1.2in]{geometry}
\usepackage{fancyhdr}
\setlength{\headheight}{28pt}
\addtolength{\topmargin}{-13.5pt}
\usepackage{hyperref}
\hypersetup{pdftitle={Oplevelse og Tydning (1918)}, pdfauthor={Harald Høffding}, hidelinks}
\pagestyle{fancy}
\fancyhf{}
\fancyhead[LE,RO]{\thepage}
\fancyhead[RE]{\textit{Harald Høffding: Oplevelse og Tydning}}
\fancyhead[LO]{\textit{\leftmark}}
\fancypagestyle{plain}{\fancyhf{}\fancyhead[LE,RO]{\thepage}}
\setstretch{1.3}
\usepackage{marginnote}
\usepackage{graphicx}
\DeclareUnicodeCharacter{0254}{\reflectbox{c}}  % ɔ: the old Danish id-est mark (a reversed c)
\newcommand{\starmark}[1]{\ifnum#1>10 \ifcase\numexpr#1-10\relax\or\dag\or\dag\dag\or\dag\dag\dag\fi\else\ifcase#1\or*\or**\or***\or****\or*****\fi\fi)}   % \footnote[11], [12]: the print's †) ††)
\renewcommand{\thefootnote}{\starmark{\value{footnote}}}   % the print marks its notes *) **) ***) per page
\newcommand{\opage}[1]{\marginnote{\footnotesize\textit{[#1]}}}
\newcommand{\apage}[1]{\marginnote{\footnotesize\textit{[#1]}}}
\newcommand{\whiteline}{\par\vspace{\baselineskip}\par}
% the printed Indhold: an essay or chapter line (title, page), a section line
\newcommand{\ientry}[2]{\par\smallskip\noindent #1\dotfill\,#2\par}
\newcommand{\isub}[2]{\par\noindent\hspace*{2em}#1\dotfill\,#2\par}
% the headings (no \addcontentsline: the print has its own Indhold, and a footnote in a heading breaks a toc entry)
\pdfstringdefDisableCommands{\renewcommand{\footnote}[2][]{}}
\newcommand{\prefacehead}[1]{\chapter*{#1}\markboth{Fortale}{Fortale}}
\newcommand{\booktitle}[2]{\clearpage\begin{center}{\Large #1}\\[6pt]{\large #2}\end{center}\markboth{Oplevelse og Tydning}{Oplevelse og Tydning}\medskip}
\newcommand{\introhead}[1]{\begin{center}{\large #1}\end{center}\medskip}
\newcommand{\chapterhead}[2]{\bigskip\begin{center}#1\\[3pt]{\large #2}\end{center}\medskip}
\newcommand{\conclhead}[1]{\bigskip\begin{center}{\large #1}\end{center}\markboth{#1}{#1}\medskip}
\newcommand{\sectionhead}[1]{\par\medskip\begin{center}#1\end{center}\smallskip}
\newcommand{\essayhead}[1]{\clearpage\begin{center}{\Large #1}\end{center}\markboth{#1}{#1}\medskip}
\newcommand{\essaysource}[1]{\begin{center}{\small #1}\end{center}\medskip}
\newcommand{\subhead}[1]{\par\medskip\begin{center}#1\end{center}\smallskip}
\newcommand{\romanhead}[1]{\par\medskip\begin{center}#1.\end{center}\smallskip}

\begin{document}

% ---------------------------------------------------------------------
%  FRONT MATTER (scan PDF 5-11).  Title page and verso read at the image
%  2026-10-14; the Indhold read at the image 2026-10-14.
% ---------------------------------------------------------------------
\frontmatter
\begin{titlepage}
\centering
\vspace*{2cm}
{\large HARALD HØFFDING}\par
\vspace{0.6cm}
\rule{2cm}{0.4pt}\par
\vspace{0.6cm}
{\Huge OPLEVELSE\\[4pt] OG TYDNING}\par
\vspace{1cm}
{\large RELIGIONSFILOSOFISKE\\ STUDIER}\par
\vfill
{[\textit{publisher's device}]}\par
\vspace{1cm}
{\large NYT NORDISK FORLAG}\par
{\small (OTTO HALLING \& CHR. NIELSEN)}\par
{\large KØBENHAVN}\par
\end{titlepage}
\thispagestyle{empty}
\vspace*{\fill}
\begin{center}
{\footnotesize KROHNS BOGTRYKKERI 1918}
\end{center}
\clearpage

% --- p. 5 ---
\setcounter{footnote}{0}%

\opage{5}\prefacehead{FORTALE}

Den Afhandling, som har givet denne Bog Navn,
 blev udarbejdet paa Foranledning af, at min udmærkede
Ven og Kollega Théodore Flournoy i Genève i
sine »Archives de Psychologie« i Aaret 1915 udgav Optegnelser
af en Dame, i hvilke hun beskrev sine ekstatiske
Tilstande og den Indflydelse, de havde paa hendes
religiøse Liv. Denne Skildring forekom mig ikke blot
at have stor psykologisk Interesse i og for sig, men ogsaa
at opfordre til Sammenligning med lignende Fænomener,
særlig hvad Forholdet mellem selve det Oplevede
og de uvilkaarlige eller fuldt bevidste, maaske villede
Tydninger angaar. Jeg genoptog derfor tidligere Studier
angaaende den mærkelige spanske Nonne Santa Teresa.
Som Indledning har jeg givet kortere Skildringer af
beslægtede Fænomener fra forskellige Tider.

Paa Bogens andre Dele passer Overskriften »Oplevelse
og Tydning« ogsaa nogenlunde. Med god Villie
vil man ogsaa kunne lade den gælde for Stykket om
Madvig som Filosof, da det for en Del handler om hans
paa egen Erfaring grundede Stilling til det religiøse
Problem. Men jeg indrømmer, at Afhandlingen om Aristoteles'
Opfattelse af Venskabet kun kan indordnes under
den fælles Overskrift, naar man ser den som beskrivende,
hvorledes den store Græker tydede sine Oplevelser angaaende
et betydningsfuldt menneskeligt Livsforhold.

Carlsberg, d. 19. August 1918.
HARALD HØFFDING

% the Indhold (p. 7, PDF 11; p. 8 blank): read at the image 2026-10-14 (the four chapter titles letterspaced: \emph)
\chapter*{INDHOLD}
\markboth{Indhold}{Indhold}
\opage{7}%
\noindent\hfill{\small Side}\par
\ientry{OPLEVELSE OG TYDNING}{9}
\isub{Indledning}{9}
\isub{I. \emph{Det primitive Stadium.}}{}
\isub{\hspace*{1em}1. Totemisme}{17}
\isub{\hspace*{1em}2. Bakchantisme}{22}
\isub{\hspace*{1em}3. Hekseprocesser og Djævlebesættelse}{23}
\isub{\hspace*{1em}4. Spiritisme}{25}
\isub{II. \emph{Profetisme og Mystik.}}{}
\isub{\hspace*{1em}1. Israelitiske Profeter}{34}
\isub{\hspace*{1em}2. Urkristendommen}{37}
\isub{\hspace*{1em}3. Nyplatonismen}{42}
\isub{\hspace*{1em}4. Augustinus}{46}
\isub{\hspace*{1em}5. Middelalderens Mystik}{48}
\isub{III. \emph{Den hellige Teresa.}}{}
\isub{\hspace*{1em}1. De første Stadier}{58}
\isub{\hspace*{1em}2. Det højeste Stadium}{63}
\isub{\hspace*{1em}3. Teresa og Kirkelæren}{67}
\isub{\hspace*{1em}4. Selvbiografi og Systematik}{72}
\isub{IV. \emph{En moderne Mystiker.}}{}
\isub{\hspace*{1em}1. Psykologi og Mystik}{80}
\isub{\hspace*{1em}2. En moderne Mystikers Ekstaser}{86}
\isub{\hspace*{1em}3. Ekstase og Religion}{96}
\isub{\hspace*{1em}4. Personlighedsprincip og Menneskekærlighed}{102}
\isub{\hspace*{1em}5. Cécile’s »syvende Bolig«}{107}
\isub{Slutning}{113}
\ientry{SOCIOLOGI OG RELIGIONSFILOSOFI}{117}
\ientry{OM TRØST}{142}
\ientry{TRÆK AF BØNNENS HISTORIE}{163}
\ientry{TO PRÆSTER}{189}
\ientry{MADVIG SOM FILOSOF}{210}
\ientry{ARISTOTELES OM VENSKAB}{225}

\mainmatter

% --- p. 9 ---
\setcounter{footnote}{0}%
\opage{9}\booktitle{OPLEVELSE OG TYDNING}{EN SAMMENLIGNENDE RELIGIONPSYKOLOGISK STUDIE}

\introhead{INDLEDNING}

D\textsc{er} gives i Ordets strengeste Forstand ingen ganske
simpel og umiddelbar Oplevelse. Vi er aldrig nogensinde
fuldstændig passive. Selv de simpleste Sansefornemmelser
faar deres Beskaffenhed bestemt ved deres
Forhold til andre Fornemmelser, der optræder samtidig
eller er optraadte forud, saavel som ved hele vor Sindstilstand.
Stærk Sindsbevægelse kan saaledes være en
Hindring for Sansefornemmelsers Optræden. Afset herfra
væver Genkendelsesprocesser sig straks sammen med,
hvad der sanses eller føles. Baade ved umiddelbar og
ved middelbar Genkendelse bringes det Oplevede straks
i Forbindelse med Noget, vi forud kendte eller mente
at kende. Og denne Forbindelse kan indtræde saa hurtig,
at vi ikke mærker den Proces, der her gaar for sig,
men mener at forholde os absolut modtagende. Dette
kan være Tilfældet, baade naar Genkendelsen er gyldig,
og naar den (som ved Erindringsillusion) er ugyldig.
At vi i Virkeligheden har lagt noget ind deri, opdager vi
først, naar Genkendelsen viser sig at være falsk. Saa
kan vi skelne mellem Oplevelse og Tydning. Men de
forekommer i første Øjeblik stedse i uvilkaarlig Forbindelse.
Hvad der fremtræder i vor Bevidsthed, er
stedse Oplevelse og Tydning paa én Gang. Hvis ved
Erfaring forstaas Oplevelser, der er bragte i bestemt
Sammenhæng, kan vi sige, at Oplevelse er et Grænse%
% --- p. 10 ---
\setcounter{footnote}{0}%
\opage{10}tilfælde af Erfaring. Tydningsmomentet er her ved sit
Minimum.

Ved mere specielle og sammensatte Tilfælde er det
let at skelne mellem de to Momenter. Naar en Genstand
staar for en Person som lille, for en anden som
stor, ligger det lige for, at Forskellen beror paa den
Maalestok, enhver af dem medbringer eller anvender,
uden at være sig det bevidst. Naar en Farve i et Tilfælde
staar for mig som mættet, i et andet Tilfælde ikke,
kan det opdages, at Forskellen beror paa, hvilken Farve
jeg saa sammen med den eller forud for den. Naar en
og samme historisk Person eller historisk Begivenhed
bedømmes forskellig, trods lige stor Kundskab om den,
maa Forskellen udspringe af Bedømmernes forskellige
personlige Forudsætninger. Man er da straks rede med
den Sætning, at det Sete beror paa Øjnene, som ser.

I en lille Skizze om Jordskælvet i St. Francisco (1906),
som han oplevede paa nært Hold, har William James
givet en Analyse af, hvad der ved denne pludselige Begivenhed
er foregaaet i ham. Jordskælvet varede kun
48 Sekunder, og han kunde ved at tænke tilbage paa
dette Øjeblik --- hvad han gjorde, saasnart han kom til
Besindelse --- ikke spore nogen egentlig Eftertanke (little
thought, and no reflection); Fornemmelse og Sindsbevægelse
beherskede Øjeblikket. Sindsbevægelsen var Fornøjelse
og Beundring, og de forudsætter en vis Genkendelse
(gleeful recognition). Sagen var, at den, der drager
til St. Francisco, maa være sig Muligheden af Jordskælv
bevidst. James har da straks været klar over, hvad han
oplevede, og været optagen af at komme til at opleve
det. Der var ingen Frygt tilstede. En af de første Tanker
var den, hvor godt det lille Træhus, han boede i,
havde holdt. Naar der var en Genkendelse, gik den ikke
ud paa Jordskælv som en Proces, i hvilken Jordlag
% --- p. 11 ---
\setcounter{footnote}{0}%
\opage{11}styrter sammen; det stod som et Noget med en vis Personification
og kom over ham som et samlet Hele. Og
denne Opfattelse af Begivenheden som et Individ meldte
sig aldeles uvilkaarlig og saa at sige paa uundgaaelig
og uimodstaaelig Maade. Mere bestemt kunde ingen
menneskelig Handling tyde paa et levende Væsen som
Kilde og Udspring. Ogsaa paa Andre gjorde Begivenheden
et saadant Indtryk af individuel Handling. Men
ikke Alle stod saa intellektuelt interesserede overfor
Oplevelsen og var saa glade ved at være med i den,
som Filosofen. For Nogle stod det som et dæmonisk
Væsen. For en Dame stod det som Verdens Ende og
Begyndelsen til Dommen, og denne teologiske Tydning
optog hendes Sind ganske, saa at det først, efter at hun
var kommen ned paa Gaden, gik op for hende, at det
var et Jordskælv.\footnote[1]{\emph{William James}: \textit{On some mental effects of the earth-quake}. (Memories and Studies. 1911).}

Ikke mindst i vore Drømme væver Tydning og Oplevelser
sig sammen. Ved Overgangen fra Vaagen til
Søvn kan man undertiden iagttage, hvorledes Sansefornemmelser
og elementære Hallucinationer uvilkaarlig
forbindes til Genstande og Begivenheder. Saadanne hypnagogiske
Tilstande giver Mulighed for at belure vort
Sjælelivs Arbejde paa at forbinde sine Elementer og
bringe forstaaelig Sammenhæng mellem dem. Og ved
dette Arbejde virker naturligvis Erindringer fra tidligere
Oplevelser med. Ogsaa indenfor selve Drømmen kan
dette Arbejde undertiden spores, især i de saakaldte dramatiske
Drømme, hvor vi er i Nød eller Forlegenhed
og saa finder Hjælp og Trøst. En afdød Kollega fortalte
mig, at han drømte at gaa nøgen paa Gaden og var
meget ulykkelig derover; men saa havde han mødt
% --- p. 12 ---
\setcounter{footnote}{0}%
\opage{12}mig, og jeg havde trøstet ham med, at han drømte.
Denne Trøst maa han selv have udfundet, skønt den syntes
at komme til ham gennem en Anden. Under Opvaagningen
(i den hypnopompiske Tilstand) foregaar der ofte
et Arbejde, der bringer mere Sammenhæng i de Drømme,
man har haft, end der har været. Man har herved konstateret,
at naar en Drøm pludselig afbrydes ved Opvaagnen,
og man saa straks noterer dens Indhold, frembyder
dette en Række af Scener uden Sammenhæng.
Jo længere man opsætter Noteringen, des mere Sammenhæng
frembyder Drømmen, og denne Sammenhæng
bliver i Reglen fremdeles større gennem de senere Beretninger.
Jo mere man saa at sige kan overraske Drømmen,
des mere kaotisk en Karakter har den.\footnote[1]{Smlgn. \emph{Foncault}: L'évolution du rêve pendant le réveil. (Revue philosophique. 1904).}

Sindsbevægelser og Sindsændringer kan opstaa pludselig
og tilsyneladende umotiverede og bliver da ofte
meget pinlige for den, der oplever dem, fordi de ikke
synes at hænge sammen med tidligere Oplevelser. Det
viser sig da ofte, at det er en eller anden smertelig Oplevelse,
maaske langt tilbage i Tiden, og længe uvilkaarlig
tilbagetrængt og udestængt af Erindringen, --- et
»Trauma«, --- som saa at sige paa underjordisk Maade
øver sine Eftervirkninger gennem Forstemning og Depression.
Ved indtrængende Udspørgen og Drøftelse, ---
saakaldt Psykoanalyse, --- søger moderne Sindssygelæger
at drage den skjulte Aarsag frem og at ophæve dens
Indflydelse ved at stille den i Bevidsthedens klare
Dagslys.

Hallucinationer, hvad enten de optræder som Visioner
eller som Auditioner eller rører sig paa Berøringsprocessens
Omraader, gør ofte stor Virkning ved deres plud%
% --- p. 13 ---
\setcounter{footnote}{0}%
\opage{13}selige Optræden, og dog kan det i mange Tilfælde spores,
at de har udviklet sig successivt, ikke med et Slag
staar i deres fuldendte Form, men vokser fra det Ubestemte
og Enkelte til bestemtere, klarere og rigere udstyrede
Former. Og denne Vækst viser sig da at være
betinget ved Individernes Erindringsverden og Tankeverden.
Derved forklares det, at Visionerne gennemgaaende
stemmer med hele den aandelige Horizont,
indenfor hvilken de visionære Individer lever. Foruden
de Eksempler, jeg har henvist til i min Psykologi
(6. Udg. p. 192), nemlig Benevenutos Vision i Fængslet
og Bernadettes Vision i Klipperiften ved Lourdes, skal
jeg her henvise til en Vision i Johannes' Aabenbaring
(I, 10 ff), hvor Auditionen gaar forud for Visionen, idet
Seeren hører en Stemme bag sig, og Synet først kommer,
da han vender sig om. Hos senere Mystikere, hvis
Beskrivelse er nøjagtigere, gaar ofte den ubestemte Fornemmelse
af en Nærværelse forud for Auditionen, og
denne igen forud for Visionen. Den spanske Nonne
Sta Teresa mærkede i ekstatisk Tilstand, at der var Nogen
nærværende; den uvante Tilstand, hvor Intet hørtes
eller saas, indgød hende Frygt, og hun gav sig til at
græde. Men da hørte hun beroligende og trøstende Ord
og blev saa forvisset om, at det var Jesus, der var tilstede.
Skønt hun blev plaget af sin Skriftefader med det
Spørgsmaal, hvorledes hun kunde vide, at det var Jesus,
naar hun Intet saa, holdt hun dog fast ved sin Tydning.
Fornemmelsen af Nærværelse holdt sig i flere Dage, og
derpaa viste Skikkelsens Hænder sig for hende, nogle
Dage senere hele Skikkelsen. At Jesus saaledes lidt efter
lidt aabenbarede sig for hende, og ikke ligestraks helt,
udleder hun af sin Svaghed og sin ufuldkomne Tro.\footnote[1]{\textit{Vie de Ste Terèse écrite par elle-même}. Trad. Bruix. Chap. 27-28.}
% --- p. 14 ---
\setcounter{footnote}{0}%
\opage{14}--- Hvad der fremtraadte i den spanske Nonnes ophøjede
Visioner, findes der ogsaa Eksempler paa i mere vulgære
Hallucinationer.\footnote[1]{Se min \textit{Religionsfilosofi} § 57.} Og hvad der især maa mærkes, de
Forestillinger, der indvæves i Oplevelserne, behøver ikke
at være Genstand for Genkendelse, selv om de kan paavises
at skyldes Individets tidligere Erfaringer. En Dame
blandt Morton Prince's Patienter saa sin afdøde Husbond
for sig og hørte ham tale trøstende Ord til sig.
Det viste sig siden, at disse Ord havde hun nogle Maaneder
forud hørt af en Ven, til hvem den Afdøde i sin
Tid havde udtalt dem. Da de lød for hende i Hallucinationen,
kendte hun dem ikke igen, men de stod for
hende som nye. Da hun var bleven hypnotiseret, erindrede
hun baade Samtalen med Vennen, og at det var
denne, der havde hørt de trøstende Ord af den Afdøde.\footnote[2]{\emph{Morton Prince}: \textit{The Unconscious}. The fundamentals of human personality normal and abnormal. New York 1914. p. 40.}

Fra de nu omtalte Oplevelser adskiller de rent ekstatiske
Tilstande sig ved at være i den Grad koncentrerede,
at der Intet kan erindres fra dem. Enhver Sindsbevægelse
har en Tendens til at hæmme eller udelukke
alle Fornemmelser og Forestillinger; i Ekstasen optræder
denne Tendens paa sit Højdepunkt. I den ekstatiske
Tilstand kan der opstaa Hallucinationer, der maaske
kan erindres senere, ligesom Drømme kan erindres.
Men den rene Ekstase erklæres af dem, der have oplevet
den, for ubeskrivelig, skønt den tillige erklæres at
være den højeste Salighed. Som vi i Løbet af denne
Undersøgelse vil se, formaar en nyere Tids Mystikere
bedre end tidligere Tiders at skelne mellem selve Ekstasen
og de Visioner eller Billeder, den kan føre med
sig, og det er netop et Hovedøjemed med nærværende
% --- p. 15 ---
\setcounter{footnote}{0}%
% CHECK p. 15: notes paired with the marks in order (the layer misread a number)
\opage{15}Afhandling at paavise dette. --- Spørges der om den
fysiologiske Side ved Ekstasen, kan der næppe gives
bedre Svar end det, der indeholdes i følgende Udtalelser
af Flournoy: »Hvad angaar det, der foregaar i
Nervesystemet i de faa Øjeblikke, i hvilke Krisen er
paa sit Højdepunkt, maa man tilstaa, at vi Intet véd
derom. At slutte fra de sjælelige Virkninger, der følge
efter den, --- Vished, Glæde, Fred, Mod, Følelse og Udøvelse
af større Energi og Hengivelse --- maa det være
en Reorganisation af Hjernecentrene, Oprettelse af nye
Koordinationer af sensorisk og motorisk Art, der giver
Individet en ny Holdning overfor Omgivelserne\dots{}.
Man kunde fremstille Sagen ved et af de anatomofysiologiske
Skemaer af den Slags, man har foreslaaet for
lignende Tilfælde. Men alt det er endnu saa vagt, at jeg
ikke vil opholde mig derved.«\footnote[1]{Th. \emph{Flournoy}: \textit{Une Mystique moderne}. (Archives de Psychologie. Genève 1915). p. 181.}

Ved sin Karakter af uvilkaarlig Koncentration\footnote[2]{Den ekstatiske Koncentration kan undertiden forberedes ved en med Anstrengelse forberedt og fastholdt Koncentration af Opmærksomhed og Tanke. Men dette er ikke altid muligt, og paa Højdepunktet har den en Karakter af Betagethed og Grebethed, der netop betinger det Gaadefulde ved Fænomenet.}
træder Ekstasen i skarp Modsætning til det sædvanlige
Bevidsthedsliv, der netop er betinget ved Modsætning
og Skiften. Derfor egner ekstatiske Tilstande sig særlig
godt til at belyse Forholdet mellem Oplevelse og Tydning.
Tydningen har ved saadanne Tilstande saare lidet
at bygge paa. Den Maade, paa hvilken ekstatiske Tilstande
bringes i Forhold til Menneskers andre Oplevelser
og til deres Forestillingsliv, vil variere med den
aandelige Udviklings Trin og Former.

Der kan skelnes mellem tre historiske Typer, hvad
% --- p. 16 ---
\setcounter{footnote}{0}%
\opage{16}dette Forhold angaar, --- en animistisk Type, en dogmatisk
Type og en empirisk Type.

I Animismen findes Forklaringen af ekstatiske Tilstande
ved Indflydelse af, eller snarere Iboen af et fremmed
Sjæleliv. Denne Tydning foretages aldeles uvilkaarlig,
idet den overleverede Tro paa Forfædres Aander
(eller lignende Væsener) umiddelbart behersker Sindet.
Oplevelse og Overlevering smelter fra første Færd af
sammen, og det er deres fælles Produkt, der kundgør sig
for Bevidstheden, uden at denne formaar at analysere det.

Paa det dogmatiske Trin er der en udformet Kultus
og en med Eftertanke gennemført Troslære, der af Individet
anerkendes som Autoriteter, der bør lægges til
Grund ved Tydningen. Forskellen mellem Oplevelse og
Tydning begynder at gøre sig gældende, men ophæves
igen gennem Ærefrygt for den overleverede Norm. I den
israelitiske Profetisme og den kristelige Mystik haves
betydningsfulde Eksempler herpaa, og vi vil især dvæle
ved den hellige Teresa, hvis Mystik har stor psykologisk
Interesse.

Naar Eftertanken har udviklet sig saaledes, at Forholdet
mellem Oplevelse og Tydning kan blive Genstand
for kritisk Prøvelse, kommer Oplevelsen til at staa med
større Selvstændighed end i de to foregaaende Former,
og den kan endog hævde sin Overmagt overfor ethvert
Tydningsforsøg. Der finder hverken Sammensmeltning
eller dogmatisk Tillempelse Sted, og det umiddelbart
givne kommer mere til sin Ret. I den nyeste Tid har
Flournoy i sin Afhandling »Une Mystique Moderne« (der
har Undertitlen »Documents pour la psychologie religieuse«)
givet el overordentlig interessant Eksempel paa
denne Type.
% CHECK p. 16: note mark without a note

% --- p. 17 ---
\setcounter{footnote}{0}%
\opage{17}\chapterhead{I.}{DET PRIMITIVE STADIUM}

\sectionhead{1. \textit{Totemisme.}\footnote[1]{Smlgn. I. G. \emph{Frazer}: \textit{The Golden Bough}\textsuperscript{3}. (1913). Især Part VII, 2. p. 225---278. --- Émile \emph{Durkheim}: \textit{Les formes élémentaires de la vie religieuse}. Le système totémique en Australie. (1912). p. 141---199; 223—233.}}

Rundt om paa Jorden, hos de mest forskellige Folkeslag,
finder der ved Overgangen til den voksne Alder
en Indvielse af de unge Mænd Sted, der skal gøre dem
delagtige i Stammens Overlevering og Tro. Hos de mest
primitive Folk, man kender nærmere til, Centralaustralierne,
staar denne Indvielse i nøje Sammenhæng med
Totemismen. Et Totem er en Genstand, for det meste
en Plante eller et Dyr, der staar som Udtryk for Stammens
Væsen og som Betingelse for dens Velfærd. Mytologisk
forklares det ved, at Stammens Urfædre i Tidernes
Morgen har forvandlet sig til en saadan Genstand.
Det gælder nu at føre de Unge til at føle sig Et, eller rettere
blive Et med Stammens i Totem udtrykte Væsen. Dette
opnaas gennem en Indvielse, efter hvilken der overrækkes
den Unge et helligt Redskab, en churinga, en hellig
Stok eller Sten, der er forsynet med Tegn og gælder for
Menneskets og dets Totemurfaders fælles Legeme, saa
at den Unge derigennem forbindes paa én Gang med
Stamfaderen og med Stammens Totem, til hvilket Stam%
% --- p. 18 ---
\setcounter{footnote}{0}%
\opage{18}faderen i sin Tid havde forvandlet sig.\footnote[1]{Se \emph{Strehlow}: \textit{Die Arunda- und Loritza-Stämme in Zentralaustralien}. Frankfurt 1911. II p. 75 ff. Churinga skal betyde »det egne Hemmelige«. --- Den fortræffelige Missionær, hvem nogle af de bedste Efterretni ger om Centalaustralierne skyldes, bevidner den Alvor, med hvilken Indvielserne foretoges. Ceremonierne betragtes som en Art Gudstjeneste, ligesom de Kristne betragter deres egne Religionsøvelser. »Dette,« føjer han til (III p. G.), »forsikredes mig baade af kristne og hedenske Indfødte.«} Enhedsfølelsen
med Stammens Totem bliver efter Indvielsen saa stærk
og inderlig, at ethvert enkelt Individ erklærer sig at
være det Dyr eller den Plante, som er Totem. En Centralaustralier
kan derfor sige: »Jeg er en Kænguruh!«
»Jeg er en Emu!«

Forskellen mellem Væsen og Symbol eksisterer paa
dette Stade ikke. Naar denne Forskel fremhæves, er
Tvivlen allerede begyndt at røre sig. Den store psykologiske
Kamp, der strækker sig gennem hele Religionens
Historie, angaar netop denne Forskel mellem Væsen og
Symbol (logisk set: mellem Identitet og Analogi). Naar
Luther under Religionssamtalen i Marburg hævdede, at
i Sætningen: »Dette er mit Legeme og Blod« maa Ordet
»er« tages bogstavelig (som Kopula), medens Zwingli
fortolker »er« som »betyder«\footnote[2]{Se om Diskussionen mellem Luther og Zwingli: \emph{Eduard Zeller}: \textit{Das theologische System Zwingli's}. Tübingen 1853. p. 129.} --- saa taler den oprindelige,
primitive Religions Aand utvivlsomt gennem
Luther. Han taler som hin Totemist, der »var« en Kænguruh.
---

Et Totem er Udtryk for Stammens Histore. Urfædrene,
i de gamle Dage, følte deres Skæbne knyttet til
en bestemt Genstand, som derved blev Stammens Totem;
disse Fædres Sjæle fødes efter Folkets Tro paany
i de følgende Slægter. Men for ret at blive delagtig i
Forfædrenes Væsen og blive Et med dem, er den natur%
% --- p. 19 ---
\setcounter{footnote}{0}%
\opage{19}lige Fødsel ikke nok. Der maa en ny Fødsel til, og den
opnaar de Unge gennem den Indvielse, de gennemgaar
ved Overgangen til den voksne Alder. Ved strænge Øvelser
og Spægelser, dels Lemlæstelse (som Omskærelse og
Afhuggen af Tænder), dels Savn og Selvpinsler, beredes
deres Sind til at modtage den hellige Overlevering. De
bringes i en Tilstand af Passivitet, som er beslægtet med
Drøm og Ekstase, og derpaa bringes de, under meget
dramatiske Former, i Forbindelse med det Hellige,
hvorefter Overrækkelsen af de hellige Redskaber finder
Sted. Undertiden efterlignes ved Indvielsen, saa godt
som muligt, en virkelig Fødsel. Fra de Paroxysmer,
Indvielsen og Forberedelserne til den har fremkaldt,
vaagner Individet op til nyt Liv, og der kan gaa nogen
Tid, før Erindringen om det, der ligger forud for den
hele Prøve, vaagner igen.

Ikke altid er den Enkeltes Totem det samme som
Stammens. I de Visioner, der kan opstaa hos en ung
Mand under den ved strenge Øvelser og Selvpinsler
vakte Exaltation, ser han maaske et Dyr eller en Plante,
der kommer til at staa for ham som knyttet til hans
Væsen og Skæbne. Men saadan Totemisme er afledet af
den sociale Totemisme, hvilket viser sig ved, at en ældre
Slægtning vejleder den unge Mand under den Forberedelse,
der fører ham til at vælge sit Totem.

Paa lignende Maade, kun endnu mere voldsomt,
gaar det til, naar særegne Medicinmænd, Shamaner
eller Præster indvies. Ogsaa her gælder det at frembringe
en Tilstand, i hvilken Sindet aabnes for den hellige
Overlevering. Paroxysmerne er her stærkere, antager,
hvad man nu kalder en hysterisk Karakter, og det
er vel ogsaa Individer med Anlæg til Hysteri, der snarest
bliver Medicinmænd, og gennem hvilke de Afdødes
Aander lettest kan komme til at tale.

% --- p. 20 ---
\setcounter{footnote}{0}%
\opage{20}En Analogi til Centralaustraliernes Indvielsesmagi, altsaa
et lignende Forhold mellem Ekstase og Overlevering,
forekommer hos Eskimoerne. Angakokkerne trænes
og indvies, under Vejledning af betalte Lærere, til at se
Aanderne, faa Svar fra dem og faa dem til at hjælpe.
Uddannelsen varer omtrent i 10 Aar. Efter en Slags
Selvhypnotisering (ved i timevis at gnide en lille Sten
mod en større) henleves en Tid i Faste og Spænding.
Saa fortæres den unge Mand af et Søuhyre, men lever
op igen, --- genfødes altsaa. Ekstasen, som indtræder
paa Højdepunktet af hele denne Proces, opfattes som
Sjælens Bortgang fra Legemet.\footnote[1]{W. \emph{Thalbitzer}: \textit{The Heathen Priests of East Greenland} (10. Amerikanistkongres. Wien 1909).}

Paa et Stade som det i disse Eksempler fremtrædende
gør Forskellen mellem Oplevelse og Tydning sig
ikke gældende. Efter sin Indvielse føler den unge Centralaustralier
sig umiddelbart Et med sit Totem. Han skelner
ikke mellem sin Tilstand og det, der i denne Tilstand
fremtræder for ham. En saadan Skelnen vilde
forudsætte, at han kunde opleve lignende Tilstande uden
at tyde dem, og at han paa den anden Side kunde staa
overfor Stammens Overlevering som en særlig Kendsgerning.
Men ingen af Delene er Tilfældet. Jo mere tilbøjelig
til Ekstase han er, des hurtigere melder de overleverede
Forestillinger sig, og des mere er han overtydet
om, at hans egen Sjæl har været borte, og at hans
Stamfaders Sjæl har indtaget dens Plads. Ekstatiske Tilstande
betragtes ikke blot som særlig modtagelige for
Aanders Indgriben, men saadanne Tilstande staar som
uforklarlige, naar man ikke antager en saadan Indgriben.
Selve Ekstasen er Tegn paa en Aands Virken.\footnote[2]{I. G. \emph{Frazer}: \textit{The Belief in Immortality and the Worship of the Dead}. 1913. p. 15. 444.} I
% --- p. 21 ---
\setcounter{footnote}{0}%
\opage{21}Grønland udvælger ældre Angakokker til deres Elever
særlig Drenge, der har vist sig at være synske.\footnote[1]{W. \emph{Thalbitzer}: \textit{The Heathen Priests of East Greenland}. p. 454.}

Der er jo noget hemmelighedsfuldt baade med de ekstatiske
Tilstande og ved de overleverede Forestillinger,
der anvendes til at tyde dem. Hverken Ekstasen eller
Overleveringen synes at kunne føres tilbage til bestemte
Aarsager. Man véd ikke, og man spørger ikke, hvorfra
de kommer. De virker begge ved deres Hemmelighedsfuldhed,
peger begge ud over Individets begrænsede Erfaring
og øver derved en gensidig Tiltrækning paa hinanden,
umærkelig for de enkelte Individer. Gennem den
højtidelige Indvielse de Unge gennemgaar, forenes de
to Elementer en Gang for alle for dem.

Det er fuldstændig konsekvent, naar det primitive
Menneske betragter sine ekstatiske Oplevelser dels som
Virkning af Aander, dels som Udtryk for et Fællesskab
med Aander. Aandeverdenen er for ham en Realitet, en
Del af hans Verden. Tvivl rører sig maaske hos Enkelte;
selv hos Papuaerne skal der findes Skeptikere,
som ikke deler Stammens Tro. Men gennemgaaende staar
Overleveringen som hellig, d. v. s. urørlig, ubestridelig.
Det er da fuldt konsekvent, naar det, der berettes om
Aandeverdenen, uvilkaarlig anvendes til Forklaring af
ellers uforklarlige Fænomener. ---

Hvad der saaledes er karakteristisk for det Stade, vi
kalder det primitive (uden derfor at forudsætte, at der
ikke skulde være gaaet andre Stader forud)\footnote[2]{I min Kritik af den ovenfor omtale Bog af Durkheim (se nedenfor Afsnittet »Sociologi og Religionsfilosofi«) har jeg lagt Vægt paa, at den Religion, der anses for den mest primitive, vi kender, den centralaustralske, kender vi kun som fuldt færdig, saa at sige systematiseret. Den maa være naaet til denne faste Form gennem en lang Række af individuelle Initiativer og psykiske Processer (Kombinationer, Udskydninger og Forskydninger).} gør sig
% --- p. 22 ---
\setcounter{footnote}{0}%
\opage{22}gældende senere hen under andre Former, helt op i vor
egen Tid, I det nærmest følgende anfører jeg nogle Eksemplarer
herpaa.

\sectionhead{2. \textit{Bakchantismen.}}

Fra Orienten er der til forskellige Tider gaaet stærke
Strømninger af ekstatisk Begejstring ind over Europa.
Saaledes kom Dionysosdyrkelsen fra Lilleasien gennem
Trakien ind i den græske Verden paa et meget tidligt
Trin af dennes Kultur. Vel undergik den orientalske
Dionysos en gennemgaaende Hellenisering, og Ekstaserne
bragtes, særlig ved det delfiske Orakels kloge Optræden,
ind i harmoniske Former. Men det oprindelige
ekstatiske Element viser sig stadig, ogsaa efterat Apollon
havde indgaaet Pagt med Dionysos. I Delfi besjæles
den ekstatiske Præstindes Legeme af Guden, eller hendes
Sjæl udløses af hendes Legeme, for at modtage guddommelige
Aabenbaringer; man synes ikke skarpt at
have skelnet mellem disse to Forklaringer, idet Mennesket
ad begge Veje kan siges at blive Et med Guden,
hævet over sin sædvænlige Bevidsthed. Det oprindelig
græske Ord Entusiasme, der er gaaet over i nyere Sprog
til at betegne al Art Begejstring, betyder egentlig den
Tilstand at være i Gud. I Tiden fra det ottende til det
sjette Aarhundrede f. Kr. opstod i Grækenland en hel
Række Seere og Profeter, i Sagnene under Navn af Sibyller
og Bakider. I dem fandt en Trang og Higen Luft,
der ikke tilfredsstilledes ved den klart udformede græske
Kultus. Det var enkelte Personligheder, som ved deres
Begejstring og ekstatiske Væsen vandt Autoritet i større
eller mindre Kredse og i følgende Generationer maatte
give Navn til forskellige nedskrevne Spaadomsbøger.\footnote[1]{Erwin \emph{Rhode}: \textit{Psyche}\textsuperscript{2}. II. p 59—65. }

% --- p. 23 ---
\setcounter{footnote}{0}%
% CHECK p. 23: notes paired with the marks in order (the layer misread a number)
\opage{23}Den orientalske Aandsretning, der ytrede sig i Dionysosdyrkelsen,
breder sig i Oldtidens Slutning i en
Række af Former.\footnote[1]{\emph{Sam Wide}: \textit{De græsk-orientalske Mysteriereligioner} (1915).}

\sectionhead{3. \textit{Hekseprocesser og Djævlebesættelser}\footnote[2]{\emph{Bror Gadelius}: \textit{Tro och öfvertro i gångna tider}. Stockholm 1912—13.}}

Middelalderens Tro paa Djævlebesættelser og de dermed
sammenhængende Hekseprocesser hviler, som den
svenske Sindssygelæge Bror Gadelius har vist i et interessant
Skrift, paa samme psykologiske Forudsætninger
som det primitive Menneskes Trosforestillinger. I
tidligere Tid var, hævdede han, Tendensen til Ekstase
og Hallucination langt stærkere end nutildags, og de
dunkle Sider af Sjælelivet spillede en langt større Rolle
end nu i Forhold til Eftertanke og Selvbesindelse. Hvad
der nu vilde staa som abnormt, kunde dengang meget
vel stemme med normalt Sjæleliv. Vi tør ikke overføre
den moderne Skelnen mellem sundt og sygt Sjæleliv
paa de ældre Tider. Endnu den Dag i Dag er denne
Forskel ikke let at gøre. Tilbøjeligheden til Personlighedsspaltning
var forhen langt almindeligere end nu, og
pludselige Tilstandsændringer, Syner og Drifter, der
traadte i Modsætning til Individets sædvanlige Færd
betragtedes som Udslag af Djævelens Virksomhed; fik
de Overhaand betragtede baade Individet selv og Andre
Tilstanden som en Djævlebesættelse.

Af Gadelius' mange Eksempler hensætter jeg følgende,
som er hentet fra den svenske Adelsmand Niels Gyllenstierna's
utrykte Selvbiografi og handler om hans
aandelige Hjemsøgelse under et Stockholmerbesøg, da
han var 22 Aar gl. (1732). Han hensattes i Ektasen i
% --- p. 24 ---
\setcounter{footnote}{0}%
\opage{24}»den mørke Evighed«; »det evige Mørke, der var i hans
Sjæl og udenfor den i hele den sorte Afgrund«, skjulte
næsten den sædvanlige Verden for ham. Idet han saa
beskriver, hvad han i denne Tilstand undertiden »saa,
hørte, forstod, følte og erfarede«, omtaler han en Mængde
mørke Aander, der var Ophavsmænd til alle Laster
og Misgærninger paa Jorden. Han saa den røde Drage
fra Johannes Aabenbaring o. s. v. --- Gadelius bemærker,
at hvad der her viste sig for Gyllenstierna, stod i umiddelbar
Sammenhæng med de religiøse Bevægelser, der
gennembævede Pietismens Tidsalder. »Han er en visionær
Hysteriker af den folkepsykologiske Type, der er
saa hyppig i Religionshistorien, og man maa bedømme
ethvert saadant Tilfælde efter en særegen Maalestok, som
er afpasset efter Tidsomstændighederne og andre Forudsætninger.
Især er det værdt at lægge Mærke til den
Lethed, med hvilken den indre Overbevisning legemliggøres
af Sanseillusioner.«

Atter her foreligger altsaa et Møde af indre Tilstande
og overleverede Forestillinger, uden at Individet formaar
at tænke sig Muligheden af andre Forklaringer af sin
umiddelbare Oplevelse, der her bestaar i Sindets Formørkelse.
---

Et mærkeligt Modstykke til Gyllenstiernas »Oplevelser«
danner Swedenborgs Visioner, der snarest havde
en lys Karakter. »Ogsaa Swedenborg er«, siger Gadelius
»et Eksempel paa den hysteriske Udvikling af underbevidste
(visionære) Evner, og det er utivlsomt ikke rigtigt
at betragte ham som sindssyg i moderne klinisk
Betydning«.\footnote[1]{Ifølge (\emph{Kleen}: \textit{Swedenborg}. En lefnadsskildring. Stockholm 1917. I. p. XIX.) skal Gadelius senere have erkendt, at \textit{Swedenborg} ikke blot led af Hysteri, men ogsaa af Paranoia (Forrykthed). Det karakteristiske Forhold mellem Stemningslivet og Visionernes Art bevarer sin Interesse, og Gadelius' betydningsfulde Advarsel mod uden videre at anvende den moderne Sindssygelæres Forskel mellem sundt og sygt Sjæleliv paa tidligere Tiders »folkepsykologiske Typer« staar stadig ved Magt.} Det var lysere Tanker, tildels fremkomne
% --- p. 25 ---
\setcounter{footnote}{0}%
\opage{25}under den gryende Rationalisme\footnote[1]{Som det tydelig ses af den Swedenborgske Vision, jeg har anført i min \textit{Religionsfilosofi}. § 57.}, der her bestemte Swedenborgs
Visioner. Selv hvor han skildrer Helvede, er
denne Skildring abstrakt og farveløs i Sammenligning
med hvad Folketroen i det attende Aarhundrede krævede.

En saadan Modsætning viser, hvorledes Tydningen
af det Oplevede, selv naar den kommer uvilkaarlig og i
Visionens Form, er betinget ved den Oplevendes Karakter,
hans Fortid og hans Fantasi, og tillige hvorledes
Tydning og Oplevelse umiddelbart kan gaa i Et.\footnote[2]{Først efter Affattelsen af dette Kapitel kom jeg til at læse \emph{Henrik Schück}'s fine Skildring af Petrus de Dacia, Forfatter af den hellige Kristinas Legende \textit{(Vår første Forfattere}. En själshistoria från medeltiden. Stockholm 1917). Den giver klassiske Eksempler paa, hvorledes ekstatiske og hallucinatoriske Oplevelser tydes med stor Inderlighed, men i Overensstemmelse med Tidens Forestillinger. Det er hos Kristina dels djævelske, dels guddommelige Kræfter, der staar bag hendes overordentlige Oplevelser. »Alle hendes Drifter legemliggjordes og antog Skikkelser af Engle og Djævle, der førte en morderisk Kamp om hendes Sjæl. De Drømme, der for et moderne Menneske fremtræder i Digtets Form, de fremtræder for hende i Virkelighedens Form.« (p. 101). --- Af særlig Interesse vilde det være at sammenligne hendes ekstatiske Tilstande med den moderne Mystikers, vi siden kommer til at tale om.}

\sectionhead{4. \textit{Spiritisme.}}

Om Swedenborg siger den Mand, som i det nittende
Aarhundrede stræbte at naa en videnskabelig Basis for
Aandetroen,\footnote[11]{\emph{Frederic W. H. Myers}: \textit{Human Personality and its survival of bodily death}. (1904). I. p. 6.} at han førte Filosofien tilbage til Him%
% --- p. 26 ---
\setcounter{footnote}{0}%
\opage{26}melen, i Modsætning til Sokrates, der i sin Tid havde
ført den fra Himmelen til Jorden. Det var Frederic
Myers' Livsopgave at finde en videnskabelig Basis for
Troen paa menneskelige Sjæles Evighed, hvilket hverken
de hidtil bestaaende Religioner eller de hidtil forsøgte
Spekulationer havde formaaet. Og Opgaven skulde
løses gennem et Studium af de ekstatiske Tilstande, et
Studium, der skulde gennemføres efter Erfaringsvidenskabens
strenge Metode. Idet han mener at være naat
til det Resultat, at de Tilkendegivelser, som forekommer
i saadanne Tilstande, ikke kan forklares ad psykologisk
og fysiologisk Vej, men nødvendigvis forudsætter
afdøde Menneskers stadige Tilværelse og Muligheden
af deres Indgriben, finder han, hvad han søger, i Spiritismen,
med fuld Bevidsthed om, at han derved vender
tilbage til en gammel Folketro.

»Jeg indrømmer«, siger han (II. p. 218), »saare gerne,
at mine sande Forgængere ved mange af de Tanker,
som hævdes i denne Bog, er Stenalderens arme Tænkere,
Hexevæsenets og Shamanismens troende Tilhængere.«
Og han finder her en stor Kontinuitet gennem
Tiderne. »Af alle subjektive religiøse Erfaringer er Ekstasen
den, som mest paatrængende og overbevisende
gør sig gældende, og den er ikke begrænset til nogen
enkelt Religion\dots{} Ligefra de laveste Vildes Medicinmænd
op til Apostlene Johannes, Petrus og Paulus,
Buddha og Muhammed ikke at forglemme, finder vi
Beretninger, der i deres psykologiske Væsen er ens,
hvor forskellige de end kan være i moralsk og intellektuel
Henseende. Vi finder tillige paa alle Trin, at Sjælen
menes at forlade Legemet, eller, om ikke forlade det,
dog udvide sit Erfaringsomraade (range of perception)
i en med Ekstase beslægtet Tilstand\dots{}. Medicinmanden
har (naar han ikke er en ren Bedrager) ligesaavel Ad%
% --- p. 27 ---
\setcounter{footnote}{0}%
\opage{27}gang til en aandelig Verden som Petrus eller Paulus;
--- kun er den Del af Aandeverdenen, han faar Adgang
til, forskellig fra deres: i Stedet for at føle sig hævet op
i en lysere Verden, slaas han af Rædsel ved sine taagede
og mørke Syner. Men for os hjælper den blotte Kendsgerning,
at vi tror paa hans Vision, til at styrke og støtte
vor Tro paa Apostlens Vision af »den syvende Himmel«.«
(II. p. 260).

De mediumistiske Fænomener tyder ifølge Myers
paa, at Mennesker kan miste deres sædvanlige Bevidsthed
(den supraliminale Sjæl), og at Underbevidstheden
(den subliminale Sjæl) kan træde i dens Sted. I Underbevidstheden
kan Afdødes Sjæle opslaa deres Bo og
sige eller skrive, hvad Mennesket i supraliminal Tilstand
ikke kunde vide eller tænke. Eller ogsaa kan den
subliminale Sjæl selv gaa paa Vandring, færdes i Aandeverdenens
Regioner og bringe Kundskab med tilbage.
I første Tilfælde finder der en Besættelse (possession)
Sted; i andet Tilfælde foreligger der Ekstase i egentlig
Betydning (II. p. 259).\footnote[1]{Ekstase betyder ordret det »at være ude af sig selv«.}

Naar Myers, den »videnskabelige« Spiritismes Teoretiker,
saaledes resolut stiller sig paa Linie med »Stenalderens
arme Tænkere«, saa er der dog den betydelige
Forskel, at medens de primitive Menneskers Aandetro
beror paa umiddelbar og uvilkaarlig Sammensmeltning
af Oplevelse og Tydning, af individuel Trang og social
Overlevering, bygger Myers sin Tro paa, hvad han anser
for en erfaringsvidenskabelig Undersøgelse. Det var hans
personlige Trang til Lys og Trøst, der førte ham til at
prøve, om der ikke ad Iagttagelsens Vej skulde kunne
opnaas en Viden om Aandeverdenen, som hverken positiv
Religion eller spekulativ Tænkning havde formaaet
% --- p. 28 ---
\setcounter{footnote}{0}%
% CHECK p. 28: notes paired with the marks in order (the layer misread a number)
\opage{28}at bringe. Han hylder da ogsaa paa sin Vis Erfaringsvidenskabens
vigtigste Princip, nemlig at de Aarsager,
ved hvilke foreliggende Fænomener forklares, selv skal
kunne paavises som Fænomener indenfor vor Erfaring.
Herved træder han i Modsætning til tidligere Tiders
Overtro. Kun Aander, hvis Eksistens er bleven erfaret
uafhængig af de Oplevelser, der skal forklares, kan betragtes
som Aarsager til disse. Derfor er der i Spiritismen
ikke mere Tale om Engle eller Djævle\footnote[1]{ A devil is not a creature whose existence is independently known to science. (II. p. 199).}, kun om
afdøde Menneskers Sjæle. I hvert Tilfælde kan der ikke
være Tale om Djævles Indflydelse; thi i ingen af de saare
mange Iagttagelser, Myers mener at raade over, har han
truffet paa nogen Forbindelse af Aand og Ondskab.
Den gamle Djævletro smelter bort, jo mere man studerer
det Bevismateriale, som foreligger (the evidence before
us) (II. p. 78-200).

Der viser sig her en rationalistisk Tendens hos Myers,
som øver en kritisk udrensende Virkning. Den skiller
ham fra saadanne Spiritister, for hvem Spiritismen fra
første Færd er en Religion, og som derfor udtrykkelig
advarer mod metodisk Forsken i disse Sager.\footnote[2]{Se om \emph{Stainton Moses}: \textit{Human Personality} II, 223---237.} Myers
vil forene Tro og Viden. Han og hans Tilhængere har
ogsaa (ligesom saa mange Mystikere) havt at kæmpe
med den Tro, at alle Aandeaabenbarelser, især saadanne,
der optræde i ekstatiske Tilstande, er Djævelens Værk.\footnote[3]{Præsten \emph{Sanders} var saa uheldig at have uautoriserede Aabenbaringer. (II, p. 17).}
Enhver positiv Religion maa begrænse Omfang og Beskaffenheden
af de Aabenbaringer, der skal kunne anerkendes
som saadanne. Og den nye Adgang til Aandeverdenen,
der ifølge Myers og hans Tilhængere nu er
% --- p. 29 ---
\setcounter{footnote}{0}%
\opage{29}aabnet, efterat den siden det primitive Menneskes Tid
saa ofte var bleven miskendt, gør jo ogsaa Fordring
paa at træde i Dogmets og Kirketroens Sted, hvis Afmagt
menes at staa fast. »Jeg paastaar«, siger Myers
(II, 279), »at der nu er gjort Begyndelse til Anvendelse
af en Metode til at vinde guddommelig Kundskab med
samme Vished og rolige Sikkerhed som den, med hvilken
vi stadig gør Fremskridt i Kundskaben om jordiske
Ting. Iagttagelsens og Forsøgets Autoritet vil træde
istedetfor Dogmers og Kirkers Autoritet« (II. p. 279).
Kristus staar her som en stor Forgænger. Men det er
gaaet med hans Opdagelse af Aandeverdenen som med
Nordboernes Opdagelse af Amerika. Den tilhører Legenden
og kan ikke nu vise os Vej. Landet maa opdages
paany, ikke af nogen Enkelt, men gennem et stort
menneskeligt Fællesarbejde paa at tyde Iagttagelser, man
hidtil har miskendt.\footnote[1]{\textit{Fragments of Inner Life}. p. 86. (I »Fragments of Prose and Poetry by F. W. H. Myers. Edited by his wife Eveleen Myers«. 1904).}

Overfor den Indvending, der ofte rettes mod Spiritismen,
at de Ting, »Aanderne« meddele, i Regelen er
trivielle og ubetydelige, om ikke taabelige, erklærer
Myers (II. p. 202), at man ikke maa møde med for store
Forventninger overfor, hvad de Afdødes Aander kan
meddele. De kan naturligvis kun sige os, hvad de véd,
og vi kan ikke vente, at de hurtig har kunnet indsamle
Viden om de nye Forhold, i hvilke de lever. Og tillige
er deres Meddelelser stedse betingede ved Mediets Opfattelsesevne
og kan ikke hæve sig højt over dettes Forestillingskreds.

Medens Totemismen og Hexetroen umiddelbart beherskes
af Folketroen, kommer den moderne Spiritisme
til at træde i Modsætning til to selvstændige, fuldt ud
% --- p. 30 ---
\setcounter{footnote}{0}%
\opage{30}organiserede aandelige Myndigheder. Den ene er Kirkens
Dogmatik, den anden er den moderne Videnskab.
Mod den første vender Spiritismen sig principielt, dog
saaledes, at den søger at vise, at de store Personligheder,
hvis Værk Kirkens Opstaaen og Bestaaen skyldes, hentede
deres Tanker fra ekstatisk Forbindelse med Aandeverdenen.
Overfor Videnskaben mener Spiritismen selv
derimod ikke at staa i nogen Modsætning; den mener
netop at gaa frem efter videnskabelig Metode. Som vi
har set, mener Myers at bygge paa det Princip, at Aarsagen
til et Fænomen skal søges i Noget, der selv foreligger
indenfor Erfaring. Appellen til Afdødes Aander
fyldestgør efter hans Mening dette Princip, thi det er jo
en Kendsgerning, at de Afdøde virkelig have levet! Men
--- Spørgsmaalet er netop, om de lever endnu, og derom
haves ingen Erfaring. Og Spørgsmaalet er dernæst, om
de er, hvad de i Synerne udgive sig for at være, --- om
deres Identitet kan paavises. Og til begge Dele mangler
Spiritismen ethvert Middel.

Hvad det sidste Spørgsmaal, Genkendelsen af Aanderne,
angaar, er det allerede kommet frem i den gamle
Kirke. Medens man i den ortodoxe Kirke ved Aar 200
e. Kr. var begyndt at slaa sig til Ro med, at Aabenbaringen
nu engang var sket og nedlagt i de hellige Bøger,
fastholdt Montanisterne, at der stadig kom nye Aabenbaringer
gennem ekstatiske Oplevelser. Der blev da spurgt
dem, hvorledes de kunde genkende Aanderne, være forvissede
om, at det var den sande Gud eller den virkelige
Jesus, der viste sig for dem. Hertil svarede Tertullian,
den mest fremragende Repræsentant for Montanismen,
med det Modspørgsmaal, som unægtelig rammede
hans bibeltroende Modstandere: »Hvorledes har
Petrus paa Forklarelsens Bjerg kunnet genkende Moses
og Elias? Det maa have været i Aanden (in spiritu).
% --- p. 31 ---
\setcounter{footnote}{0}%
\opage{31}Fra Billeder eller Statuer har han ikke kunnet genkende
dem, thi saadanne var forbudte i Loven. Han maa
have været i oversanselig Tilstand (in spiritu, non in
sensu constitutus), siden han har genkendt dem!« Altsaa
Genkendelsen af, hvad der viser sig i de ekstatiske
Tilstande, er for Tertullian saa langt fra at være et Problem,
at han deri, at Petrus genkendte Moses og Elias,
netop ser et Bevis paa, at han har været i Ekstase
(amentia), overskygget af guddommelig Kraft (obumbratus
virtute divina).\footnote[1]{Se F. C. A. \emph{Schwegler}: \textit{Der Montanismus}. Tübingen 1841, p. 25.}

Men Spiritisterne kan ikke afvise Problemet paa
denne Maade, naar de vil stille sig paa Erfaringsvidenskabens
Grund. Det gælder i Erfaringsvidenskaben ikke
mindst om at kunne identificere Fænomener. Som vi
vil se, dukker Genkendelsesproblemet igen op i den katolske
og i den moderne Mystik.

Hovedsagen er imidlertid den, at førend man griber
til overnaturlige (eller dog overnormale) Forklaringer,
maa man paavise Umuligheden af, at Visionernes Indhold
kunde skyldes tidligere Oplevelser, som Individerne
i deres sædvanlige Tilstande ikke kan fremdrage i Erindringen.
Studiet af hypnotiske Tilstande og af Dobbeltbevidsthed
har i det ene Tilfælde efter det andet godtgjort,
at Forestillinger, et Individ i sin sædvanlige Tilstand
ikke har til Raadighed, kan komme frem i Hypnose
eller i Trance. Man vil idethele ved nøje Kendskab
til et Mediums tidligere Historie meget ofte finde
Nøglen til Forestillinger og Tendenser, der i Øjeblikket
staar som hemmelighedsfulde, fordi de ikke synes forenelige
med de forhaandenværende Tilstande og Situationer.
Selv Fantasien frembringer jo ofte, især i sine
højeste Former og Grader, Billeder og Tanker, der staa
% --- p. 32 ---
\setcounter{footnote}{0}%
\opage{32}som uforklarlige for Individet selv og for Andre. Der
er foregaaet et ubevidst Arbejde, som dog meget vel har
kunnet benytte, hvad Individet har lært at kende i sit
tidligere Liv. Det Ubevidste i Sjælelivet (hvad det saa
ellers er) kan altsaa ikke blot bevare Oplevelser og Forestillinger,
Individet i sædvanlig Tilstand ikke vedkender
sig, men ogsaa frembringe nye Tanker og Billeder, uden
at der er Grund til at antage andre Kræfter virkende
end saadanne, som Historie, Psykologi og Fysiologi kan
paavise. Saalænge en Visions Indhold kan forklares ad
»endogen« Vej, det vil sige, ud fra Individets tidligere
Historie, er der ingen Grund til at appellere til de Afdøde.
Myers indrømmer jo selv, at Visionens Indhold
betinges ved Individets Opfattelsesevne og Forudsætninger.
Men kan der sættes nogen Grænse for denne Betingethed?
Hvor den endogene Forklaring ikke kan
fuldendes, staa vi overfor et Problem, og Beviset for,
at dette Problem kun kan løses ved Optræden af Afdøde,
paahviler Spiritismen, ikke dens Modstandere.
Deraf at Individet under Visionen mener at være aldeles
passiv, kan Intet sluttes. Thi der kan gives en Aktivitet
i os, som ikke bliver Genstand for vor Bevidsthed, men
først i sine Resultater kommer op over Bevidsthedens
Tærskel.

Det er Genferpsykologen Théodore Flournoy, som
med størst Energi og med mest aaben Sans har studeret
de saakaldte spiritistiske Fænomener. Han har særlig
lagt Vægt paa Kundskab til Mediernes hele Fortid
og har anvendt dette Princip med stort Held i sin mesterlige
Studie over Sonnambulen Mlle Smith.\footnote[1]{ \textit{Des Indes à la Planète Mars}. Essai sur un cas de sonnambulisme avec Glossolalie. Genève 1900.} Senere har
han foruden forskellige nye »Documents d'Enquête«
% --- p. 33 ---
\setcounter{footnote}{0}%
\opage{33}givet en fortrinlig Udvikling af de psykologiske og erkendelsesteoretiske
Principer, der er af Betydning overfor
de Fænomener, Spiritismen beraaber sig paa, og
overfor den Tydning, den giver af dem.\footnote[1]{\textit{Esprits et Médiums}. Genève 1911.} I den nyeste
Tid har han fremdraget »en moderne Mystiker«, der i
et senere Afsnit vil blive undersøgt som et karakteristisk
Modstykke til tidligere Tiders Mystik.

Overfor den spiritistiske Tydning staar Flournoy afgjort
tvivlende. Der er ikke en eneste Kendsgerning,
der har overbevist ham, skønt han indrømmer, at flere
henstaa uforklarede. Han er her kommen til samme
negative Resultat, som Henry Sidgwick og William James,
der efter nogle og tyve Aars Beskæftigelse med de
nævnte Fænomener erklærede, at de ikke var komne
videre, end de fra først af var.\footnote[2]{\emph{William James}: \textit{Memories and Studies}, p. 174.}

% --- p. 34 ---
\setcounter{footnote}{0}%
\opage{34}\chapterhead{II.}{PROFETISME OG MYSTIK}

\sectionhead{1. \textit{Israelitiske Profeter.}}

Jødefolkets Profeter er mærkelige i Religionshistorien,
ikke blot ved deres Syners Vælde, deres etiske Kraft
og deres ophøjede Poesi, men --- hvad der for vort Emne
har særlig Interesse --- ved, at Forholdet mellem Oplevelse
og Tydning fremtræder hos dem med større Klarhed
end der, hvor Aandetroen straks væves ind i Oplevelserne.
Indenfor det gamle Testamente viser sig en
skarp Modsætning mellem den dæmoniske og den egentlig
profetiske Inspiration.\footnote[1]{Gustav \emph{Hølscher}: \textit{Die Profeten}. Untersuchungen zur Religionsgeschichte Israëls. Leipzig 1914, p. 94, 100, 197.} I den dæmoniske Inspiration
er det en dunkel Magt, der betager Individet. Paa uklar
Maade smelter det guddommelige Væsen og den Inspireredes
Personlighed sammen til Et. Under Ekstasen
fremstammedes uforstaaelige Lyde, og ved exalterede
Handlinger, Dans og Musik, virkedes der betagende paa
Mængden. Den egentlige Profeti, saaledes som den i skriftlig
Overlevering først fremtræder hos Amos, tolker Ekstasen
som Aabenbaring af en Gud, hvis Villie Profeten
kundgør. Selve de indre Oplevelser var for Profeten
Vidnesbyrd om en fra ham selv forskellig Magts Herredømme
i hans Sind; men uklar Sammensmeltning fin%
% --- p. 35 ---
\setcounter{footnote}{0}%
\opage{35}der ikke Sted. Naar Profeterne bragte Jahve's Ord til
Folket, og naar de optraadte med Myndighed, var det,
fordi der var sket Noget i deres Sind, som ikke var »af
dem selv«, og som de maaske endog havde gjort Modstand
imod. Den, der taler ud af sig selv, erklærer de
for en Løgnprofet. Ofte har Profeten gjort Modstand
imod, hvad han dreves til at forkynde; han ser deri
(se Jeremia I, 6-11) et Vidnesbyrd om, at han virker i
en Højeres Tjeneste. Vi genfinder ofte dette Træk af
Modstræben, før den fulde Hengivelse til den ekstatiske
Tilstand kan finde Sted, i Mystikens Historie, og ofte
netop anvendt som Bevis for en overnaturlig Indgriben.
Fra en senere Tids Mystik adskiller Profetismen sig
derved, at det for Profeterne ikke først og fremmest er
Fred og Hvile i Sindet, der skal opnaas gennem de ekstatiske
Tilstande, men disse Tilstande ansporer til kraftig,
ofte voldsom Optræden udadtil. Profeten staar paa
en Gang midt i sit Folk og hævet over det. Dets Færd
og dets Skæbne betinger hans Optræden. Han retter sin
Tale baade mod Præst og mod Konge. Profetismen er
en Modbevægelse mod den stivnede Kultus, især mod
de Elementer af hedensk Kultus, der efter Kongedømmets
Oprettelse var blandede med Mosaismen.\footnote[1]{Hermann \emph{Gunkel}: \textit{Israëls Profeter}. (Norsk Overs.) Kristiania 1916 p. 33---37.} Amos
taler saaledes mod Ofre, Valfarter, Sange og Fester og
kræver i deres Sted Ret og Retfærdighed (V, 21—24).
Paa lignende Maade Jeremia (Kap. 7) og Jesaja (Kap. 8).
Profeterne hævder, at Kultus ikke hørte med til Folkets
oprindelige Religion og ikke fandtes i den Tid, da
Jahve ledede det gennem Ørkenen. Overfor det Bestaaende
hævder Profeterne det Oprindelige, ligesom de overfor
det Ydre hævder det Indre.

% --- p. 36 ---
\setcounter{footnote}{0}%
\opage{36}Den samme Forblindelse, der lod Folket henfalde
til magisk Kultus, hildede det i Tryghed og ubegrundet
Selvtillid overfor de Farer, der truede fra de store Verdensriger,
der Skridt for Skridt udviklede sig paa de
smaa Nabolandes Bekostning. De store Profeter stod
her som »Vanheldsprofeter«, atter og atter overvældede
af de Farer, der truede Land og Folk, og forgæves advarende
det kortsynede Folk og dets Fyrster. Netop
under saadanne Forhold krævede de af hver Enkelt
Renhed og Redelighed i Sindet. Ligesom den enkelte
Profet traadte op overfor de bestaaende nationale
Myndigheder i Kirke og Stat og derigennem hævdede den
enkelte Personligheds Myndighed til at følge, hvad hans
egne Oplevelser havde ført ham til at anerkende som
Ideal, saaledes var det ogsaa et personligt Ideal, de
krævede fyldestgjort af de Enkelte. Hvis der --- som
Jesaja varsler --- gennem den Lidelse, Folket vil komme
til at gennemgaa, kan øjnes Mulighed for Frelse, saa
vil Betingelsen for hver Enkelt være den, at han i Tro
og Tillid bøjer sit Sind under Renhedens og Retfærdighedens
Krav. Baade i Profetens eget Sind og i det
Krav, han stiller til sit Folk, kundgør sig en Frigørelse
for Personligheden fra Overleveringens og de ydre Autoriteters
Herredømme. Og det er det frie Forhold mellem
Oplevelse og Tydning, der gør denne Frigørelsesproces
mulig, hvorved Israels Profetisme betegner et
stort Vendepunkt i det menneskelige Aandslivs Historie.
Efterhaanden som Profetismen ebbede ud, kom Kultus
og Tempeltjeneste igen i Forgrunden, og Folkets aandelige
Ledere blev Præsterne, ikke Profeterne, selv om de
profetiske Former endnu holdt sig længe.

Baade hos Amos, Jesaja og Jeremia fremtræder i
Beskrivelsen af Synerne en fremskridende Klarhed. Ofte
kommer Anelse og Syn overvældende, uden at der kan
% --- p. 37 ---
\setcounter{footnote}{0}%
\opage{37}dannes nogen Forestilling om deres Betydning. Det er
da en guddommelig Stemme (Audition), der hjælper til
at forstaa Visionen. Og ved nærmere Eftertanke kan
saa Synets Betydning og Anvendelse gennemføres i det
Enkelte. Det følger af sig selv, at naar Oplevelserne
skulde nedskrives, maatte Eftertanken træde i den umiddelbare
Betagetheds Sted. Hvad der saaledes gælder for den
enkelte Profets Vedkommende i Forhold til hans
Oplevelser, det fremtræder i Profetiens Historie, der
viser en fremskridende Udvikling fra Ekstase til Eftertanke
og til Poesi. »Ekstatikeren bliver Forfatter, Profeten
Poet«(Hølscher p. 317). Ekstaser og Hallucinationer
var jo heller ikke i sig selv Vidnesbyrd om Sandhed;
de forekom hos »Løgneprofeter« saavel som hos de
»rette« Profeter. Vægten kom da, især fra Jeremia's Tid
(Hølscher p. 291 f), mere og mere til at ligge paa Erfaringens
og Eftertankens Bekræftelse af, hvad der var
forkyndt efter de ekstatiske Tilstande.\footnote[1]{Naar der her, foruden selve Kilderne, især er gjort Brug af \emph{Hølscher}'s og \emph{Gunkel}'s religionshistoriske Arbejder i religionspsykologisk Øjemed, saa har dette for en Del kun været muligt, fordi Forskere som de nævnte selv er vel orienterede i vor Tids psykologiske Forskning.}

\sectionhead{2. \textit{Urkristendommen.}}

Naar Forholdet mellem Oplevelse og Tydning skal
drøftes for Kristendommens Vedkommende, særlig med
Hensyn til ekstatiske Oplevelser, vil det, da Striden netop
paa dette Omraade endnu er staaende, være af Vigtighed
at være klar over de Forudsætninger, som lægges til
Grund. De er allerede udtalte i det foregaaende, især i
Anledning af den moderne Spiritisme. Den væsentlige
Forudsætning for al psykologisk og historisk Forsken
% --- p. 38 ---
\setcounter{footnote}{0}%
\opage{38}er, at den Tydning, som Menneskene hylder af deres
Oplevelser, ikke mindst af de stærkest indgribende, stedse
vil være betinget ved de Erfaringer, de har gjort før
Oplevelserne, og ved de Forestillinger, der ad forskellige
Veje er komne til at udgøre deres aandelige Indhold.
Det er, hvad man har kaldt den endogene Forklaring.
Oplevelsen og dens Tydning staa da begge som Led i
Menneskets hele Liv og forstaas indenfor dette Livs hele
Sammenhæng. For Kristendommens Vedkommende maa
der her lægges Vægt paa Forestillingerne om en kommende
Forløser, der gennem Lidelse og Fornedrelse
skulde bringe Frelse og Ophøjelse, Forestillinger, der
fra Denterojesaja's Tid havde udviklet sig i den jødiske
Profetisme og senere indgik Forbindelse med orientalske
Forestillinger om Opstandelse og Verdensafslutning,
der især stammer fra den persiske Indflydelse. I Lyset
af saadanne Forestillinger maatte den Enkelte, der ved
usædvanlige indre Oplevelser følte sig kaldet til en stor
Opgave, naturlig betragte sig selv, og i Lyset af dem
maatte de, der sluttede sig til ham, naturlig komme til
at betragte ham. Om Anvendelsen i det Enkelte kan
der være Strid, men en endogen Forklaring maa holde
sig til dette Synspunkt. Det udelukker ikke stor og
maaske enestaaende Oprindelighed; ti denne ytrer sig
i Maaden, hvorpaa de forskellige historisk givne Elementer
forbindes --- dels indbyrdes, dels med de egne
Oplevelser. Og den beror særlig paa de etisk-humane
Ideer, der drages frem i Lyset af den vundne Tydning
og hævdes som afgørende for Livet.

Overfor det nye Testamente og Urkristendommen er
der for Forskningen to Hovedproblemer, der begge beror
paa Forholdet mellem Oplevelse og Tydning.

For det første: hvorledes har Jesus af Nazareth selv
tydet, hvad der rørte sig hos ham og særlig kom til
% --- p. 39 ---
\setcounter{footnote}{0}%
\opage{39}Gennembrud hos ham i ekstatiske Tilstande? Har han
kaldt sig den bebudede Messias, eller Menneskens Søn,
eller den menneskeblevne Gud? Han har i hvert Tilfælde
følt sig som den, hvis Komme betød den nære
Afslutning af Menneskelivets jordiske Bestaaen, saa at
Tro paa ham og hans Betydning maatte blive et Spørgsmaal
om Liv eller Død. Hans Liv, som det foreligger i
Legenderne, er et Liv i Forventning, i Henblik til en
nær forestaaende Forherligelse gennem Lidelse, en Forherligelse,
der var Et med Indtræden af Verdenslivet i
nye, forklarede Former. Derfor er ogsaa den nytestamentlige
Etik en Forventningens Etik, en Spændingens
og Haabets Etik, der lader alle rent menneskelige Forhold
og Opgaver ligge som ligegyldige, om ikke hindrende;
disse Forhold og Opgaver hører jo den gamle
Tingenes Orden, ikke det kommende Rige til. I denne
Tanke, som var historisk forberedt baade i Jødedommens
Historie og i andre orientalske Religioner, laa et
Grundlag for Tydning af de ekstatiske Tilstande, som
atter og atter har fundet Sted hos Jesus selv og hans
Nærmeste.\footnote[1]{Om de forskellige Forsøg, der paa psykologisk-historisk Grundlag er gjorte paa Udredning af Jesu Liv og Personlighed, se \emph{Albert Schweitzer}: \textit{Von Reimarus zu Wrede}. Eine Geschichte der Leben-Jesu-Forschung. Tübingen 1906.}

Det andet Spørgsmaal er, hvorledes den første Menighed
har opfattet Jesus fra Nazareth, --- især hvorledes
den har tydet de Sindsbevægelser, de ekstatiske
Oplevelser, som hans Liv og Død har fremkaldt hos
dem selv. Saa meget staar fast, at hos de nytestamentlige
Forfattere henføres alle indgribende sjælelige Erfaringer,
alle overordentlige Oplevelser til Jesus af Nazareth som
Livets Kilde, og at alt Haab og al Frygt tilsidst samles
i Tanken om ham og hans Genkomst. --- Da nu de
% --- p. 40 ---
\setcounter{footnote}{0}%
\opage{40}nytestamentlige Forfattere selv hørte til den første Menighed
og delte dennes Forudsætninger, er det første
Spørgsmaals Besvarelse tilsidst afhængig af Besvarelsen
af det andet, og derpaa beror det store og vanskelige
Problem, som her er stillet for historisk Forsken. Men
tilsidst er og bliver Problemet et psykologisk og
angaar Forholdet mellem Oplevelse og Tydning. Og dette Problem
bliver saa meget vanskeligere, som der endnu
ikke bliver skelnet bestemt mellem Oplevelse og Tydning,
men de smelter meget hurtig sammen. »Jesu begejstrede
Diciple udsagde om ham det Højeste, som Jødedommen
forstod at sige. Og det ikke for at udgrunde
hans Persons Hemmelighed, som om Jesus var det
først Givne og Kristologien det Andet, men de Sjæle,
som længtes efter Guddommens Nærhed og følte Trang
til en fra Himmelen kommen Gudssøn, overførte deres
Hjerters Idealer paa ham. Saa er da den nytestamentlige
Kristologi en mægtig Hymne, som Historien synger
om Jesus«.\footnote[1]{\emph{Hermann Gunkel}: \textit{Zum religionsgeschichtlichen Verständnis des neuen Testaments} (1903), p. 94.}

Den kristne Kirkes Troslære er, fra det psykologiske
Synspunkt, en mere eller mindre systematisk gennemført
Tydning af Beretningerne om Jesu og den første
Menigheds Oplevelser, men saaledes, at senere Tiders
Oplevelser og Bestræbelser søger at finde deres Tydning
i og med Tydningen af Menighedens ældste Overleveringer.
Gennem en Række af Kombinationer, ofte blotte
Kompromiser, skrider gennem Dogmehistorien disse
Tydningsforsøg frem. Et Hovedproblem, der stiller al
Teologi overfor store Vanskeligheder, er det, hvorledes
den helt nye Kultur, intellektuel, emotionel og materiel,
der har udviklet sig siden den første Menigheds Tid,
% --- p. 41 ---
\setcounter{footnote}{0}%
\opage{41}kan bringes i Sammenhæng med den nytestamentlige
Tankekreds, for hvem hele dette Kulturarbejde dels ikke
eksisterede, dels ikke kunde have nogen Værdi.

Men hvad Forholdet mellem Oplevelse og Tydning
angaar, bliver der en afgørende Forskel mellem Kirkens
Grundlæggelse og det kirkelige Livs Fortsættelse,
efterat den kirkelige Overlevering og den dogmatiske
Tydning havde fæstnet sig. Hos de store Profeter, hos
Jesus og endnu tildels i den første Menighed kunde
Livet endnu paa afgørende Maade gribe ind og sætte
nye Begyndelser, saa at Valget af Tydning, af Forbilleder
og Symboler, betingedes ved Oplevelserne. Men
for dem, der anerkender Kirken og de af Kirken hævdede
Normers Myndighed, bliver selve Oplevelsen bestemt
ved de normerede Tydninger. Der tør ikke lægges
Vægt paa Oplevelser, som ikke falder indenfor de en
Gang for alle givne Tydningers Ramme. Man anvises til
at gøre visse bestemte Erfaringer, faa visse Oplevelser,
helst i en bestemt Orden (ordo salutis). Derved er
religiøs Erfaring i kirkelig Betydning forskellig fra den
Art Erfaring, der ellers i Liv og i Videnskab faar Indflydelse
paa Menneskers Tanker\footnote[1]{Se nærmere herom \textit{Religionsfilosofi}, § 29 og 59---60.}.

Der stod i den gamle Kirke en lang Kamp mellem
nye Oplevelser og overleverede Tydninger. De ekstatiske
Tilstande, Visioner og Aabenbaringer hørte ikke
op, fordi visse Bøger anerkendtes som Normer. Montanismen,
som ovenfor er omtalt, var et af de stærkeste
Indlæg i denne Brydning, som vel for en Tid lagde
sig ved den katolske Kirkes Konsolidering, men som ---
hvad der i det følgende vil blive omtalt --- fortsattes i
Middelalderen i Modsætningen mellem Mystik og Skolastik.
Det var nu ikke blot indenfor Kirken, at eksta%
% --- p. 42 ---
\setcounter{footnote}{0}%
\opage{42}tiske Tilstande optraadte, ogsaa efter Apostlenes Tid.
Ogsaa Kristendommens Modstandere og Kætterne kunde
beraabe sig paa Tilstande, der tolkedes som Virkninger
af Aanders Indgriben. Et gribende og rørende Eksempel
er den saakaldte fjerde Esra Bog, som er skreven af
en from Jøde, hvis Syner bestemmes ved Rædslen over,
at Jerusalem ødelagdes, før Messias kom. Naar nu saaledes
Aander talte paa begge Sider, maatte der være
Trang til et Kriterium til at afgøre, paa hvis Side Retten
og Sandheden var. Ekstasen er i og for sig væsentlig
af samme Natur, hvorsomhelst den optræder, og
indeholder derfor i sig selv ingen Vejledning til at afgøre,
hvilken Tydning der er den rette. Det gjaldt om at »prøve
Aanderne«, og Kriteriet, der anvendtes ved denne Prøvelse
maatte tilsidst blive de etiske Virkninger, som de
ekstatiske Tilstande voldte. Man sluttede da, som Weinel,
der meget klart har belyst dette Punkt, bemærker,
ikke fra det Hemmelighedsfulde eller Underfulde til, at
det er en hellig Aand, der virker; men man sluttede fra
det etiske Indhold, fra den Handlen, der fremkaldtes,
til, at Tilstandene var af overnaturlig Karakter\footnote[1]{\emph{Heinrich Weinel}: \textit{Die Wirkungen des Geistes und der Geister im nachapostolischen Zeitalter} (1899), p. 149.}. Det
er da ikke Intensiteten, men Kvaliteten, som bliver afgørende.
Og Intensiteten tabte sig da ogsaa efterhaanden,
og med den den levende Forventning, som bar den
urkristelige Etik.

\sectionhead{3. \textit{Nyplatonismen.}}

I Mystikens Historie blev det af stor Betydning, at
den sidste store Tankeretning, der udviklede sig paa
græsk Grund, Nyplatonismen, lagde stor Vægt paa ekstatiske
Oplevelser.

% --- p. 43 ---
\setcounter{footnote}{0}%
\opage{43}Allerede hos Platon er der et mystisk Element, i
ejendommelig Forening med den intellektuelle og etiske
Bevægelse, der stammede fra Sokrates. Det synes at
have været orfisk Indflydelse, som saaledes er gaaet ved
Siden af den sokratiske. Platon henviser ofte til Mysteriekultusen,
dens Udtryk og dens Ceremonier. Selv om
han afgjort tydede disse paa symbolsk Maade, --- saaledes
at f. Eks. »Renselse« (Katharsis) hos ham ikke betyder
Frigørelse for liturgisk Urenhed, men for Tankens og
og Sindets Hildethed i lav Sanselighed, --- saa er det
dog karakteristisk for ham, at han føler Trang til Mysteriernes
Billedkreds, netop naar hans Tanke naar sine
Højdepunkter. Naar Tanken under sin Opstigen fra de
 endelige og sanselige Fænomener kommer til de afsluttende
Begreber, de sidste Forudsætninger og søger at
samle dem i én Idé, saa afløses Dialektiken af Entusiasme,
og i et pludseligt Syn afsløres det Guddommelige
for den (Fajdros 249, Symposion 210 E). Og da er det
Højeste naat; da først er Livet værdt at leve. Men Platon
bliver dog ikke selv paa disse Højder. Straks tager
han paany Tankearbejdet op, uden hvilket hint Syn
ingen Betydning vilde have, skønt det selv staar over al
Tænken og Videnskab (Symp. 211 A). Det Højeste er ikke
blot det, der gør Erkendelse mulig, men er tillige det, der
gør, at de erkendende Væsener eksisterer, altsaa den
sidste Grund for al Tanke og for al Væren (Staten VII
p. 508 E; 509 B).

I Forhold til, hvad der saaledes kan fremtræde paa
Tankelivets Højdepunkter, skønt det selv ligger ud over
al Tænken, er Troen paa Guder for Platon noget underordnet.
Gudernes Guddommelighed beror kun paa deres
Delagtighed i det, der ogsaa for Mennesker er det Højeste.
Den sande Tænker lever i »det, som det er Guds
Guddommelighed at leve i« (Fajdros 249 C); Forskellen 

% --- p. 44 ---
\setcounter{footnote}{0}%
% CHECK p. 44: centred line: heading or verse
\opage{44}er kun, at hvad der for Guderne er en stadig Besiddelse,
er for Mennesket Genstand for stadig Stræben (Fajdros
278 D).

Naar vi gaar fra Platon til Plotinos, fremtræder det
mystiske Element afgjort i større Fylde, optager og fyldestgør
Livstrangen paa en langt stærkere Maade end
hos Platon. Metoden er den samme: Opstigen ad Tankens
Vej til det Højeste, der er mere end Tanke. Plotins
Biograf Porfyrios siger udtrykkelig om ham, at han
fulgte den Vej, Platon havde anvist i Symposion. Men
skønt denne Vej kunde være lang, var dog det Maal, der
naaedes, ikke et fjernt og hinsidigt; Maalet, »Enhed med
og Nærhed ved den i alt virkende Gud«, var stedse nært
(Vita Plotini. Cap. 23). Denne umiddelbare Nærværelse,
som er langt værdifuldere end Videnskab (παρουσία κρείττων. Ennead. VI 19, 4), kan der kun
naas til gennem Tankearbejde. Den er altsaa paa én
Gang fjern og nær. Men Nærheden, den umiddelbare
Enhed med det Guddommelige, forudsætter en sjælelig
Tilstand, der er hævet over det sædvanlige Bevidsthedsliv,
--- en Ekstase.

Den ekstatiske Tilstand beskrives af Plotin, ikke blot,
som af Platon, som Begejstring, Entusiasme, men som
Stilhed og Hvile. Paa Højdepunktet er der ikke mere,
som hos Platon, Tale om Skuen. Sjælen er »ude af sig
selv« (ἔκστασις), og den kommer i en Tilstand, der paa
én Gang er Enfold (ἅπλωσις) og Udfoldelse (ἐπίδοσις).
Modsætningen mellem Seende og Syn falder bort. Det
er ligesom to Centre, der falder sammen, og det er, som
om Mennesket er blevet til en anden. Derfor er det saa
vanskeligt at give en Beskrivelse af, hvad der opleves i
saadanne Tilstande (Ennead. VI, 9, 10---11---11).

Af stor Interesse er den Konsekvens, Plotinos drager
af den Omstændighed, at det Højeste ikke naas gennem
% --- p. 45 ---
\setcounter{footnote}{0}%
\opage{45}Videnskab, men kundgør sig gennem umiddelbar Oplevelse
af ubeskrivelig Karakter. Naar vi vil meddele
Andre saadan Oplevelse gennem Tale eller Skrift, saa
kan vi kun middelbart lede dem til den. Saalænge vi
er paa Vejen, kan der være Tale om Belæring; men
selve Maalet forudsætter personlig Trang, uden hvilken
Sjælen ikke gennemstraales af det højeste Lys (Ennead.
IX, 6, 41). Det viser sig her allerede, hvorledes Mystiken
fører til at lægge Vægt paa individuelle Erfaringer, i Modsætning
baade til videnskabelig Forsken og overleveret
Dogmatik. Da en af Plotin's Venner, som var en
ivrig Deltager i Tempelfesterne, vilde have Plotin til at
følge med til en saadan Fest, svarede denne: »Det er
Guderne, der skal komme til mig, ikke mig, der skal
gaa til Guderne«. (Vita Plotini. Cap. 10). Plotin er enig
med Platon i, at Mennesker kan have Adgang ad deres
egen Vej til det, hvorpaa Gudernes Guddommelighed
beror, eller, som han udtrykker det (Ennead. IX, 1, 2),
»Aarsagen til, at Guder er Guder, maa være en Gud,
der staar over dem i Værdighed«. Begrebet Guddommelighed
fremtræder i saadanne Ytringer tydelig som en
Værdibestemmelse.

Grunden til, at den højeste Oplevelse ifølge Plotin
faar Ekstasens Karakter, og ikke kan være Genstand for
Beskrivelse, er for det første, at det, som skulde erkendes,
ligger ud over alle Modsætninger, medens Erkendelsen
kun bliver til ved, at Tanken gaar fra det ene
til det andet. Den absolute Enhed kan kun erfares
gennem Ro og fuldstændig Samling. Men for det andet
er det, fordi selve denne Enhed i sig indeholder
en uendelig Fylde. De to Bestemmelser hører nøje sammen,
ti hvis der var Noget udenfor Guddommen, vilde
den være begrænset og kunde ikke være en absolut Enhed.
(VI, 9, 4---6). Af begge Grunde er det Højeste uud%
% --- p. 46 ---
\setcounter{footnote}{0}%
\opage{46}sigeligt. Ved ethvert Prædikat, man vilde anvende paa
det, vilde man gøre det til noget enkelt bestemt og
begrænset. (V, 3, 13). Plotinos opstiller her en Hovedsætning
i al Mystik\footnote[1]{Se herom min \textit{Religionsfilosofi}, § 21 og Note 57.}.

Nyplatonismen bliver grundlæggende for Middelalderens
Mystik, som den faar Indflydelse paa dels gennem
Augustinus, dels gennem Dionysius Areopagita. Medens
Middelalderens Tænkning væsentlig støttede sig til Aristoteles'
encyklopædiske Produktion, hvad Verdenserkendelsen
angaar, blev det nyplatoniske Ideer, der betingede
det afsluttende Aandsindhold. Den middelalderlige
Mystiks Forhold til Skolastiken betinges ved, at den
energiske Fordybelse i og Betagethed af dette Aandsindhold
kom til at danne en Modsætning til den skolastiske
Dialektik. For at bruge Plotin's Udtryk: det
Guddommeliges umiddelbare Nærværelse (Parusi) stilledes
over Videnskaben (Episteme) om Guddommen.

\sectionhead{4. \textit{Augustinus.}}

Ogsaa den store Kirkefader har haft ekstatiske Oplevelser,
og selv om de ikke fremkaldes ved nyplatonisk
Indflydelse, saa har han dog, hvad der i denne
Sammenhæng er af særlig Interesse, tydet dem i nyplatoniske
Former. Fra denne Tid af brydes idethele to
Tydningsmaader overfor overordentlige, især ekstatiske
Oplevelser, --- den animistiske (i videste Forstand), der
finder Forklaringen i Sjæles eller personlige Væseners
Indgriben, og den platoniske, der paa Højdepunktet antager
en Aabenbaring af noget, som ikke kan finde
Udtryk i noget Billede eller noget Begreb. Vendepunktet
betegnes med Plotinos og Augustinus.

Augustinus skylder, det har han gentagne Gange selv
% --- p. 47 ---
\setcounter{footnote}{0}%
\opage{47}erklæret, Plotinos som Repræsentant for Platonismen,
overordentlig meget. Platonismen førte ham bort fra
materialistiske Forestillinger. Efter en haard Kamp opgav
han at forestille sig Gud i legemlig Form. Dernæst
førtes han ved sine platoniske (nyplatoniske) Studier
til Selverkendelse, til Studiet af sit eget Indre. Og endelig
førtes han i Plotin's Fodspor Trin for Trin bort fra
de tilvante Forestillinger og fra Sansningens Verden op
til det Uforanderlige, det sande Væsen, der kun kan
gribes i mystisk Skuen, --- eller, som han selv udtrykker
det, »i et skælvende Syns Glimt« (in ictu trepidantis
aspectus). Fastholde dette Syn kunde han ikke, og
da han igen var i sin daglige Tilstand (redditus solitis),
var det kun en »elskende Erindring« (amans memoria),
en Stemning af Glæde og Tak, han beholdt tilbage ---
en Trang til Duften af den Føde, han ikke formaaede
at tilegne sig\footnote[1]{\textit{Confessiones} VII, 23 (17). --- Ved Udtrykket amans memoria kan ikke forstaas Erindring i egentlig Forstand, men den Stemning, der bliver tilbage efter et Syn, hvis Indhold ikke kan fastholdes.}. J
efter hans Omvendelse, oplevede han igen i
Samtale med sin Moder paa hendes Dødsdag en Ekstase,
idet Moder og Søn efterhaanden førtes i Tale og
Tanke ud over alt Sanseligt --- og de opnaaede en sagte
Berøring med det Evige under Hjertets fulde Begejstring
(attingimus [sapientiam] modice toto ictu cordis); men
den Aandens Førstegrøde, de var blevne delagtige i,
maatte de forlade og atter vende tilbage til, hvad det
sædvanlige Sprog kunde udtrykke. Og han forsøger en
Oversættelse deraf, der, naar man ser nærmere til, viser
sig at være en Oversættelse af Steder hos Plotin: Naar
Alt tier, inden i os og udenfor os, og selve Sjælen tier
og skrider ud over sig selv ved Ikketænken, saa aabenbarer
Gud sig ved sig selv, uden Røst og uden Syn, og
% --- p. 48 ---
\setcounter{footnote}{0}%
\opage{48}der kommer et Forstaaelsens Øjeblik (momentum intelligentiæ),
i hvilket en hurtig Tankeakt berører den evige
Visdom. --- Augustinus føjer til, at omtrent saaledes
havde han talt, skønt ikke med de samme Ord. Og han
føjer fremdeles til: skulde det, der opleves i et saadant
Øjeblik, ikke være det, der udtrykkes i de Ord: »Gaa
ind til din Herres Glæde?«\footnote[1]{\textit{Confessiones} IX, 23—26 (10). Smlg. \emph{Plotinos}: \textit{Ennead}. V, 1, 2 og III, 8, 9. --- Se \textit{Religionsfilosofi}, § 61.} --- Dette Sidste er karakteristisk
for den Sammensmeltning af Platonisme og Animisme,
hvorpaa al Teologi beror, saavel som for Mystiken,
saaledes som den under Indflydelse af Augustinus
udviklede sig i den middelalderlige Kirke. Til Belysning
af Forholdet mellem Oplevelse og Tydning er de her
benyttede Afsnit af Augustin's Konfessioner af stor Interesse.


Naturligvis identificerer Augustin ikke Platonisme og
Kristendom. Platonismen kunde kun se Maalet, men
ikke gaa Vejen til Ende\footnote[2]{\textit{Confessiones} VII, 27 (21).}. Paa sine gamle Dage bebrejdede
han sig i sine »Retractationes«, at han havde
stillet Platonisme og Kristendom for nær hinanden.

Men Plotin har, foruden meget andet, lært ham at
skelne bestemt imellem den umiddelbare Oplevelse
og dens Tydning. Helt glemt blev denne Lærdom ikke
i Mystikens Historie, skønt den strengere Udformning
af Dogmerne under kirkelig Autoritet baade af
indre og ydre Grunde kunde gøre det vanskeligt at
fastholde den.

\sectionhead{5. \textit{Middelalderens Mystik.}}

Skønt Middelalderens Teologi søgte at optage Oldtidens
Viden, særlig som den var bleven sammenfattet
% --- p. 49 ---
\setcounter{footnote}{0}%
\opage{49}i Aristoteles', Galenos' og Ptolemaios' Værker, laa det
Maal, den tilstræbte, dog ud over al Viden og Tænken
i den umiddelbare Forening med Guddommen. Man
har derfor med en vis Ret kunnet sige, at ikke Aristoteles,
men Plotinos er de middelalderlige Filosofers sande
Lærer\footnote[1]{Fr. \emph{Picavet}: \textit{Esquisse d'une Histoire comparée des philosophies médiévales}. Chap. 5.}.
Den middelalderlige Organisation af Studierne,
der begynder med Encyklopædi og Filosofi og ender i
Mystik, foreligger egentlig allerede i den nyplatoniske
Skole i Alexandria, der benyttede Aristoteles' Skrifter
som Indledning og Forberedelse til Studiet af Platon,
hvis Tankegang fortolkedes i mystisk Retning\footnote[2]{Smlgn. \emph{William Norvin}: \textit{Olympiodoros fra Alexandria}. (1915). Kap. 1.}. Mystiken
stod ikke nødvendigvis i Modsætning til Skolastiken,
men kunde staa som dennes Afslutning. Det er jo
ogsaa blevet paavist, at den dristigste af de middelalderlige
Mystikere, Eckehart, selv har forfattet skolastiske
Værker\footnote[3]{\emph{Denifle}: \textit{Meister Eckehart's lateinische Schriften und die Grundanschauungen seiner Lehre} (Archiv für Literatur- und Kirchengeschichte des Mittelalters) (1886). --- Allerede tidligere (1838) havde \emph{Engelhardt} i sin Bog om Richard a St. Victore og Johannes Ruysbrook (p. 97) vist, at Mystikere og Skolastikere ikke hørte til to helt forskellige Verdener.}. Men der var hos mystiske Aander en
Tendens til at bygge paa det indre Liv og dets Oplevelser
og sætte dem langt over den analyserende og distingverende
Tænken, en Tendens, der let kunde føre til Strid
med den kirkelig sanktionerede Skolastik. Og i det
indre Livs Skole lærte Mystikerne, hvad der blev endnu
farligere, at skelne mellem Oplevelse og Tydning. Ofte
var det kun paa krampagtig Maade, at de formaaede at
forene deres personlige Erfaringer med den kirkelige
Tydning. Modsætningen mellem Mystikere og
Skolasti%
% --- p. 50 ---
\setcounter{footnote}{0}%
\opage{50}kere beror, som jeg har paavist i min Religionsfilosofi\footnote[1]{\textit{Religionsfilosofi} § 21 (med Note 37). --- Ogsaa paa dette Punkt staar \emph{Plotinos}, og efter ham \emph{Augustinus}, som Forgænger, idet de hævder, at de videnskabelige Grundbegreber (Kategorierne) kun har Gyldighed for den sanselige Verden, ikke for Guddommen, hos hvem der ikke kan skelnes mellem Væsen og Egenskab eller mellem Egenskaber indbyrdes. Se herom \textit{Den menneskelige Tanke}. p. 143.},
i erkendelsesteoretisk Henseende paa, at disse hævder
Gyldigheden af de Analogier, ved hvilke Guddommen
ene kan tænkes, medens hine nægter denne Gyldighed.
Modsætningen mellem Mystik og Skolastik hænger altsaa
ikke blot sammen med Modsætningen mellem Oplevelse
og Tydning, men ogsaa med Modsætningen mellem
Kriticisme og Dogmatisme.

Nogle enkelte Eksempler fra den middelalderlige Mystiks
Historie skal her fremdrages til Belysning af vort
Emne.

Af Hugo a St. Victore (i Begyndelsen af det 12. Aarhundrede)
beskrives Ekstasen paa følgende Maade. Efter
at have gennemgaaet forskellige Trin gaar Sjælen ligesom
ud af sig selv og rives bort, uden at vide hvorhen.
»Saa bliver jeg«, siger den, »pludselig fornyet og forvandlet
og føler en Glæde, jeg ikke kan beskrive. Samvittigheden
opklares; al Sorg over udstaaet Lidelse forglemmes;
Hjertet jubler; Forstanden bliver klar, Sindet
oplyst. Længselen tilfredsstillet;\dots{} og der er noget,
jeg holder inderlig fast paa med Kærlighedens Arm,
uden at vide, hvad det er«. Dertil svares: »Det er din
Elskede (dilectus tuus), der besøger dig; men han kommer
usynlig, skjult og ubegribelig. Han kommer for at
berøre dig, ikke for at ses af dig«. Tydelig skelnes her,
ligesom hos Plotin, mellem ekstatisk Oplevelse og visionær
Skuen. Og den kirkelige Lære om en personlig
Gud bruges, med Antydninger af »Højsangens« erotiske
% --- p. 51 ---
\setcounter{footnote}{0}%
\opage{51}Sprog, til at tyde det Ubeskrivelige, der opleves. --- Som
Trin, der fører op til det Højdepunkt, hvor den ekstatiske
Oplevelse bliver mulig, nævner Hugo fire: Læsning
(til Andagt) (lectio); Eftertanke (meditatio); Bøn
(oratio); gode Handlinger (operatio). Højdepunktet kaldes
Kontemplation (contemplatio), der er et med Ekstasen;
den er forskellig fra Eftertanken og betyder ikke
noget intellektuelt\footnote[1]{A. \emph{Liebner}: \textit{Hugo von St. Victor und die theologischen Richtungen seiner Zeit} (1832) p. 270—276. --- De forskellige Stadier, som Hugo angiver, findes allerede i lignende Karakteristik hos tidligere Tiders Mystikere, ikke blot hos Nyplatonikerne, men ogsaa i den indiske Mystik (Yoga og Buddhismen). \emph{Friederich Heiler} \textit{(Das Gebet}. 1918 p. 253—258) giver en (vel skematisk og paa løsrevne Citater bygget) Sammenstilling (Tafel der Gebets- und Versenkungsstufen), der viser, hvorledes alle Tiders Mystikere tilnærmelsesvis kender de samme Trin af Fordybelse. Meget beror vel her paa Tradition; men at Traditionen er bleven mere end død Overlevering, kan kun forklares ved ny, selvstændig Oplevelse. Og dette bliver klart, naar man fra de enslydende Udtryk og Ytringer gaar tilbage til den Maade, hvorpaa de motiveres hos de enkelte Mystikere.}. --- Denne Trinrække viser gennem
mange Mellemled tilbage til Platon's Symposion og
kommer med forskellige Nuancer igen hos andre Mystikere.
Saaledes særlig hos Richard a St. Victore (i Midten
af 12. Aarhundrede), der ikke blot beskriver de forud
for »Kontemplationen« gaaende Trin, men ogsaa selve
det højeste Trin mere indgaaende og med flere Distinktioner.
Af Interesse for vort Emne er den Forskel, der
gøres mellem »Spekulation«, i hvilken vi se gennem et
Spejl, og »Kontemplation«, i hvilken vi skue Sandheden
uden Dække. Men det er netop denne Umiddelbarhed,
der gør Kontemplationstilstanden til Ekstase\footnote[2]{Træffende taler H. O. \emph{Taylor} \textit{(The Mediaeval Mind. A History of the development of Thought and Emotion in the Middle Ages} (1911 II p. 363) om den »irrational immediacy of conviction«, ved hvilken Mystikeren føler og kender sin Genstand.}. Denne
% --- p. 52 ---
\setcounter{footnote}{0}%
\opage{52}Tilstand optræder, efter at der er følt Længsel, Angst
og Anger, og den har Karakteren af en fredelig Stilhed,
en blid Luftning, i hvilken Mennesket føres ud over sig
selv og ser den højeste Visdoms Lys uden Dække. Men
hvad der saaledes er skuet i Henrykkelsens Tilstand,
det skal siden overvejes og gennemtænkes, for at det
kan forstaas, og for at det ved Hjælp af Grunde og Sammenligninger
kan bringes til almindelig Forstaaelse. Og
Troens Hemmeligheder, som kendes af Aabenbaring
eller læres af Teologer, dem skal man føre ind i Kontemplationen,
saa at de kan blive Genstand for Beundring
og Nydelse\footnote[1]{\emph{Engelhardt}: \textit{Richard von St. Victor und Johannes Ruysbrook}. p. 81.}.

I disse sidste Sætninger antydes det for al Mystik
skæbnesvangre Forhold mellem Oplevelse og Tydning.
Mystikeren er overbevist om, at det er Gud, han har
været i Berøring med under Ekstasen; men naar ingen
Analogier slaar til, hvorledes saa forklare sig for andre,
og hvorledes erkende Forholdet mellem det Oplevede
og det, som Kirken lærer? Richard raader ganske vist
Bod paa Vanskeligheden ved at opfordre til at lægge
Bibel- og Kirkelære ind i selve den ekstatiske Tilstand
for at kunne beundre Læren saa meget des mere (idet
den jo saa nyder godt af Entusiasmetilstanden). Men
denne Indlægning kan ikke betyde andet end en hurtig
Tydning. I selve Ekstasen er man ikke istand til at
lægge Noget ind, skønt de Forestillinger, man ellers
lever i, naturlig vil komme til at bestemme de Visioner,
der kan knytte sig til Ekstasen.

Hos Mester Eckehart (omkring Aar 1300) træder Forskellen
mellem Oplevelse og Tydning tilbage, ikke fordi
han følger Richard's Raad at lægge Teologi ind i den
umiddelbare Oplevelse, men netop omvendt, fordi han
% --- p. 53 ---
\setcounter{footnote}{0}%
\opage{53}straks og bestemt tolker de dogmatiske Forestillinger
ud fra den mystiske Oplevelse. Saalænge Sjælen er optagen
af det i Tid og Rum fjerne, og saa længe overhovedet
Forskelle og Modsætninger gør sig gældende
for den, saalænge er den ikke, som den skal være, ikke
hævet over al Sorg. Det gælder om i et evigt Nu at opleve
Tilværelsens hele Fylde. At Gud sender sin Søn til
Verden, er ikke en fjern Begivenhed, men det er noget,
der sker i Sjælens Inderste, hvor Guds Væsen og Menneskets
eget Væsen er Et. Salighed beror hverken paa
Erkendelse eller paa Kærlighed, men paa det, hvoraf
baade Erkendelse og Kærlighed udspringer, og hvor der
hverken gives et Forbigangent eller et Tilkommende. For
den, der blot en eneste Gang har anet dette Dyb, taber
de største materielle Goder deres Værdi, og der er idethele
intet ydre Formaal for Handling, end ikke Gud
eller den evige Salighed; der handles ud af Sjælens
inderste Dyb, hvor Gud bor --- nærmere ved Sjælen end
denne ved sig selv. Man skal ikke her spørge, om det
er gennem Naade eller gennem Natur, at Gud virker i
Sjælen, begge er Et. Beviset for, at Gud virker i Sjælen,
bestaar i den Lykke og Hvile, der opstaar. Der gives
en umiddelbar Vej til Enhed med Gud, og Mestrene i
Skolen gør sig med Urette store Betænkeligheder over,
hvorledes det kan blive muligt for Mennesket at blive
Et med Gud. Det kræver Afkald og Afholden; men
Gud er i os; det er kun os, der ikke er hjemme hos
os selv\footnote[1]{I denne Karakteristik har jeg benyttet den Oversættelse af udvalgte Skrifter af Eckehart, som \emph{Aage Marcus} har givet. (\textit{Mester Eckehart; Prædikener og Traktater}. Oversat og indledet. 1917).}.)

Det var intet Under, at man fra Kirkens Side blev betænkelig
ved saadanne Udtalelser. Eckehart blev anklaget
for Kætteri, og 28 Sætninger, der var uddragne af hans
% --- p. 54 ---
\setcounter{footnote}{0}%
\opage{54}Skrifter, blev fordømte. Eckehart selv oplevede dog ikke
Domfældelsen. I den katolske Verden sætter man ikke
Eckehart højt. Denifle søger at reducere ham til en
Prædikant for fromme Damer og udleder hans Mystik
af hans Ukyndighed i Dogmatiken.

Eckehart er paavirket af Augustinus og (gennem Dionysius
og Erigena) af Nyplatonismen. Men der ligger
tydelig en oprindelig sjælelig Oplevelse til Grund. Han
har haft Øjeblikke, der var opfyldt af den Fred og Hvile,
som kun er mulig gennem Fordybelse i sit Indre og i
det Liv, som dér uvilkaarlig kan fremvælde. Og saa vidt
man kan se, har disse Øjeblikke for en væsentlig Del
været betingede ved en intellektuel Interesse for at arbejde
sig ud over de Forskelle og Modsætninger, som
det mod Erfaringen rettede Tankeliv har at brydes med.
Men denne intellektuelle Interesse er for Eckehart gaaet
i Et med Stræben efter Tilknytning til, eller rettere Forening
med noget, der er hævet over Strid og Delthed.
Og naar han uden videre bruger Ordet »Gud« om det,
der inderst inde rører sig hos ham, fremtræder tydelig
Gudsbegrebet som et Prædikatsbegreb. Oplevelsen er af
den højeste Værdi; derfor er den at kalde guddommelig;
og derfor maa den betyde Guds Iboen. Vi møder
her den nyplatoniske Tanke, der ovenfor (p. 46) blev
kaldt Grundlaget for al religiøs Mystik.

Iøvrig vil det stedse vise sig, at naar Spørgsmaalet
om, hvad Gud er, for Alvor opkastes, er det stedse den
indre eller ydre Erfaring, Mennesket har, der dikterer
Svaret\footnote[1]{Sml. \textit{Upanishad}. Übers. v, Deussen. p. 204 f, 395, 799 f. --- \emph{Augustinus}: \textit{Confessiones}. 10. Bog. --- \emph{Luther}: \textit{Catechismus Major}. (Forklaringen af det første mosaiske Bud). (Libri Symbol. ed. Hase p. 403). --- \emph{Spinoza}: \textit{Tract. theol. polit}. Cap. 14. (Deus: veræ vitæ exemplar).}. Men fra det uvilkaarlig givne Svar er der igen

% --- p. 55 ---
\setcounter{footnote}{0}%
% CHECK p. 55: centred line: heading or verse
\opage{55}en Afstand til det Svar, den religiøse Tradition, indenfor
hvilken den Enkelte bevæger sig, vil give.

Hvad Kirken fordømte hos Eckehart, var den strenge
Gennemførelse af Enhedstanken paa Grundlag af, hvad
der i ekstatiske Tilstande var erfaret om Frihed for alle
Forskelle og alle Skranker. Hans Discipel Suso paadrog
sig ogsaa en kirkelig Domfældelse; men da han var
mindre dristig i sin Tankegang, og mere omhyggelig for
at blive i Overensstemmelse med Kirkelæren, undgik
han større Ulemper. Og han faar netop ved denne Stræben
efter kirkelig Udtydning Interesse for vort Emne,
idet han derved føres til at skelne bestemt mellem Oplevelse
og Tydning.

I 54. Kapitel af sit Levned beskriver han ekstatiske
Tilstande, som en Tidlang indtraf hver Morgen og hver
Aften. »Han hensank da saa ganske i Gud, den evige
Visdom, at han ikke kunde tale derom. Undertiden
havde han en kærlighedsfuld Samtale med Gud; til andre
Tiden henfaldt han i Jamren og Sukken eller i længselsfuld
Graad eller i stille Latter, ofte var det ham, som
om han svævede i Luften eller svømmede mellem Tid
og Evighed i Guds uudgrundelige Underes dybe Hav.
Undertiden var hans Hjerte saa fuldt, at han lagde Haanden
paa det og sagde: Ak, kære Hjerte, hvorledes skal
det gaa dig!« --- Denne Ekstase har nærmest Karakteren
af en Vision. I en tidligere Ekstase, den første, han
havde, træder det Visionære mere tilbage fra det Emotionelle.
I 3. Kapitel af sit Levned fortæller han, at som
han lidende og trøstesløs en Dag stod alene i Kirken,
blev hans Sjæl henrykt --- hvad enten den var i Legemet
eller udenfor Legemet --- og han saa og hørte, hvad
ingen Tunge kan udsige: »Det var uden Form og Skikkelse
og voldte dog en saa stor Lyst, som alle Former
og Skikkelser tilsammen volde. Hjertet higede og var
% --- p. 56 ---
\setcounter{footnote}{0}%
\opage{56}dog mættet; Sindet var glad og tilfreds; ethvert Ønske
og alt Begær var tilfredsstillet. Han gik helt op i Stirren
mod et straalende Lys, der lod ham glemme sig
selv og alle Ting. Var det dog ikke Nat? det vidste han
ikke. Det var et lifligt Frembrud af det evige Liv i nærværende,
stille Modtagen. Han sagde da: Er dette ikke
Himmerige, saa véd jeg ikke, hvad Himmerige er!\dots{} Da
han igen kom til sig selv, var han tilmode som et Menneske,
der kommer fra en anden Verden«. Eftervirkninger
af denne Oplevelse holdt sig længe, som en »himmelsk
Smag«, ligesom en Æske, i hvilken man har haft
et kosteligt Stof, beholder Smagen deraf. Det er Augustinus's
»elskende Erindring« (amans memoria) (se ovenfor
p. 47), vi her træffe igen.

Ligesom Augustinus erfarede Suso en stor Modsætning
mellem de ekstatiske Øjeblikke og de sædvanlige Forhold.
Og der var en stor Forskel mellem hvad han skrev
i sin sædvanlige Tilstand, og hvad der blev ham givet
i de ekstatiske Tilstande. Men selv dette Sidste underkastede
han en indgaaende Kritik, for at der ikke skulde
findes noget deri, der afveg fra de hellige Fædres Udsagn,
men at det, den guddommelige Visdom indgav ham
gennem aandelig Oplysning, kunde støttes ved den hellige
Skrifts Autoritet, helst saaledes, at det baade i Ord
og i Mening var Et med dens Udsagn\footnote[1]{Et quum ea, quæ data erant, ad tabulam concepisset, ad se reversus studiose rimatus est, ut nihil in iis esset, qvod a dictis sanctorum Patrum discreparet, sed hoc, qood divina tradente Sapientia sibi per spiritualem illustrationem infundebatur, auctoritate sacræ scripturæ\dots{} fulciretur, quum idem in sensu, aut etiam idem in verbis et sensu\dots{} reperiretur esse. (\textit{Heinrich Suso's Leben und Schriften}. Herausg. von Diepenbrock 1829. p. VIII f.)}.

Denne Udtalelse (som jeg allerede har benyttet i min
Religionsfilosofi § 62) er af klassisk Betydning, hvad For%
% --- p. 57 ---
\setcounter{footnote}{0}%
\opage{57}holdet mellem personlig Oplevelse og kirkelig Autoritet
angaar. Den belyser klart Forskellen mellem den psykologiske
og sociologiske Opfattelse af Religionens Væsen.

Socialt og kirkelig lægger man Vægt paa Regelmæssighed
og bestemte, genkommende Former, og her trues
de ekstatiske Oplevelser med Fare, naar de ikke omhyggelig
tydes efter kirkelig Dogmatik eller idetmindste
udtrykkes overensstemmende med den kirkelige Sprogbrug.
Der rører sig i saadanne Tilstande et uroligt, og
i hvert Tilfælde et afgjort individuelt Element. Kirken
maa fastholde, at alle Livserfaringer kan indgaa under
dens en Gang for alle udformede Billed- og Forestillingskreds,
og den maa se en Fare i, at de værdifuldeste Oplevelser
ikke skulde kunne tydes udtømmende ved Hjælp
af den. ---

Jeg gaar nu over til denne Afhandlings Hovedemne,
en Sammenligning mellem Mystiken, som den i det
sekstende Aarhundrede optræder hos en spansk Nonne,
og Mystiken, som den i det tyvende Aarhundrede optræder
hos en protestantisk Lærerinde. De foregaaende
Karakteristiker skulde kun tjene som Indledning hertil.

% --- p. 58 ---
\setcounter{footnote}{0}%
\opage{58}\chapterhead{III.}{DEN HELLIGE TERESA}

\sectionhead{1. \textit{De første Stadier.}}

Teresa de Cepeda blev født i 1515 i Byen Avila i Spanien,
af en adelig Slægt. I sin Selvbiografi, som hun skrev
efter sin Skriftefaders Befaling, fortæller hun om sin
religiøse Udvikling. Meget tidlig rørte religiøs Romantik
sig hos hende; syv Aar gammel vilde hun sammen med
sin Broder drage til Maurerne for at blive Martyr og
derved hurtigere komme i Himmelen. Hun hørte til en
strengt katolsk Familie; men hun bekender, at hun
havde megen Lyst til at pynte sig og stor Lyst til at
læse Ridderromaner. Fra sit fjortende Aar af blev hun
opdragen i et Kloster, men følte i Begyndelsen stor Uvillie
mod Klosterlivet, og hun valgte det kun, fordi de Afsavn
og Lidelser, det medførte, var mindre end de,
Skærsilden vilde bringe. Men efterhaanden kom hun i
et inderligere Forhold til det Liv, der levedes i Klosteret,
og da hun blev indviet til Nonne, nitten Aar gammel,
følte hun sig som omstraalet af et guddommeligt
Lys. Meget tidlig traadte Modsætningen mellem hendes
indre Liv og hendes Skriftefaders kirkelige Fremstillinger
frem. Hun havde her en lang Kamp at føre, ti hendes
ekstatiske og visionære Oplevelser blev ofte af Skriftefædrene
opfattede som Indskydelser af Djævelen, der
burde bringes for Inkvisitionens Domstol.  Men afset fra
% --- p. 59 ---
\setcounter{footnote}{0}%
\opage{59}Anfægtelsens Tider følte hun sig saa sikker paa, at hendes
indre Oplevelser var af guddommelig Oprindelse,
at hun ikke lod sig rokke ved de Advarsler, hun fik i
Skriftestolen. Hun blev ved at skifte Skriftefader, indtil
hun fandt en, som forstod hende. Hun taler skarpe Ord
om halvlærde, om det indre Liv uvidende Skriftefædre.
En Skriftefader ansaa hun for nødvendig, ti hun var
Katolik i sin inderste Sjæl og vilde leve og dø i Overensstemmelse
med den ene saliggørende Kirkes Lære.
Og hun følte sin Mangel paa Lærdom, naar det gjaldt
om at finde den rette Tydning af, hvad hun oplevede.
Der var her, mente hun, en Fare, der især kunde true
Kvinder. De maa have en Autoritet at støtte sig til for
ikke at komme paa gale Veje, men det maa være Mænd,
som af egen Erfaring kende »de aandelige Veje«. Og
man skal derfor vælge sin aandelige Fører med Omhu
og efter Prøvelse. Det var hos Jesuiterne, hun fandt,
hvad hun søgte. Deres Mester havde selv haft Ekstaser
og Visioner, og dog havde Ingen som han virket til den
katolske Kirkes Befæstelse.Teresas jesuitiske Skriftefader
var forvisset om, at hendes mærkelige Oplevelser var af
guddommelig Oprindelse. Kun mente han, at hun manglede
Selverkendelse og Ydmyghed, og han lod hende gennemgaa
St. Ignatius' »aandelige Øvelser«. Denne Skole
øvede Fantasien i at forestille sig Himmelens Salighed
og Helvedes Rædsel og at leve sig ind i Jesu Liv i Fornedrelse
og Højhed, men gik især ud paa at styrke
Lydigheden mod den hellige Kirke og Villien til at følge
det Livsmaal, den satte\footnote[1]{En indgaaende Karakteristik af Ignatius' \textit{Exercitia spiritualia} giver \emph{Boehmer}: \textit{Les Jésuites}. (Trad. frane. par Monod. p. 25---38). Kfr. ogsaa \emph{Henri Joly}: \textit{St. Ignace de Loyola}. (1899).}.

Det var gennem Bønnen, at Teresa fik Adgang til
den aandelige Verden, der stedse udfoldede sig rigere
% --- p. 60 ---
\setcounter{footnote}{0}%
\opage{60}og inderligere for hende, indtil hun i Løbet af 2 eller
3 Aar (fra hendes 44. Aar af) gennemlevede en Række
af Ekstaser og Visioner, der utvivlsomt betegner hendes
aandelige Livs Højdepunkt.

Der er tre Hovedkilder til Teresas indre Oplevelser,
--- hendes Selvbiografi, som hun nedskrev og omskrev
i Aarene fra 1562 til 1566\footnote[1]{\textit{Vie de Sainte Térèse} écrite par elle même. Trad. Brouix 14. éd. Paris 1896.}%
% CHECK p. 60: star lost in the layer: paired with the mark by order
, hendes Breve\footnote[2]{\textit{Lettres de Sainte Térèse}. Trad. Brouix. 2 éd.} og hendes
Skrifter\footnote[3]{\textit{Les Oeuvres de Sainte Térèse}. Trad. Arnauld d'Andilly. Paris 1670.}, især Bogen om »Sjæleslottet« (Château de
l'âme 1577). ---

Teresa bekender, at hun i den første Tid i Klosteret
havde stor Lyst til Adspredelse, især til Samtale med
Besøgende. Hun passede sine Øvelser og bad de reglementerede
Bønner, men det var kun Munden, ikke
Hjertet, som bad, og naar hun laa i Bøn i Kirken, tilsyneladende
fuld af Andagt, længtes hun dog i sit Indre
efter at høre Klokkens Slag, der bebudede Bedetimens
Ophør. Hun har da i denne Periode hørt til dem, om
hvem hun senere i sin Bog om »Sjæleslottet« siger, at
de kun kendte det Slot, deres egen Sjæl er, paa samme
Maade som Vagten, der gør sin Ronde udenom Slottet,
men aldrig kommer ind i det. Døren til Sjælens Slot er
den Bøn, der ikke blot bestaar i Ord, men hvor Tanken
er med, og hvor hele Sindet er med: oraison mentale,
i Modsætning til oraison verbale.

Den første Bolig i Sjæleslottet er Selverkendelsens
Hjem. Men sig selv lærer man bedst at kende ved at
kende Gud; hans Storhed faar Mennesket til at erkende
sin egen Ringhed. I de to næste Boliger taler Gud til
Mennesket, ikke blot gennem Bønnen, men ogsaa gennem
gode Menneskers Ord, gennem Bøger og gennem
% --- p. 61 ---
\setcounter{footnote}{0}%
\opage{61}Sygdom og Sorg. Trangen til Bod opstaar, og Kærligheden
til Næsten vaagner. For at undgaa falsk Sikkerhed
bør selv de, der ikke gaar i Kloster, have et Menneske,
hvis Villie de underkaster sig, saa at de ikke handler
efter deres egen Villie, men lærer fuld Ydmyghed.
Og da først kommer --- i Slottets fjerde Bolig --- det
Stadium, da Bønnen ret faar en inderlig Karakter og
gennem Eftertankens Bøn (oraison de méditation) gaar
over til at blive Hvilebøn (oraison de quiétude). Saalænge
Eftertanken endnu raader, er fuld Samling og
Fred ikke naaet. Men midt under Betragtning eller Læsning
kan det ske, at Sjælen føler sig greben af Guds
Nærværelse, hensynker i ham som i en guddommelig
Afgrund. Under saadan Fordybelse arbejder hverken
Tanken eller Fantasien; Erindringen er næsten forsvunden;
Villien gaar op i Kærlighed, og Forstanden er betagen
af den Storhed, der aabenbarer sig.

Til en Tilstand som denne kan man ikke selv hæve
sig op. Man maa vente, til man drages opad. Deri bebestaar
netop Forskellen mellem Hvilebøn og Meditationsbøn.
Meditationen er at ligne med den Vanding af
en Have, der sker ved, at Vandet ledes gennem Rør ind
over Havens Vækster, eller ved, at man møjsommelig
trækker Vandet op af en Brønd. Kilden, hvorfra Vandet
hentes, er Gud selv. Men i Hvilebønnen sker den aandelige
Vanding umiddelbart som naar en liflig Regn strømmer
ned over vor Have. Det sker ganske stille; men
naar vor egen Eftertanke og Stræben virker med, gaar
det ikke af uden Larm. ---

I Bogen om »Sjæleslottet« siger Teresa, at hun nu
ser klarere paa disse Ting end 15 Aar forud (det vil
sige, end da hun begyndte at skrive sin Selvbiografi).
Hun tænker vist især paa den psykologiske Beskrivelse
af sine mærkelige Tilstande. Hun havde tidligere sørget
% --- p. 62 ---
\setcounter{footnote}{0}%
\opage{62}over, at endog i Tilstande, i hvilke alle hendes Evner
burde være koncentrerede om Gudstanken, kunde hendes
Forestillinger flagre om i højst forskellige Regioner,
ofte blandt de urene og giftige Dyr, der tumlede sig udenfor
Slottet, i hvis indre Boliger hun mente, at hendes
Sjæl nu havde taget Ophold. Hun forstod nu, at der var
to Slags Tanke. Den ene gaar fra Emne til Emne ligesom
Fantasien gaar fra Billede til Billede; den anden
bestaar i Koncentration om et eneste Punkt og er i Grunden
Et med Villien til at fastholde en bestemt Retning
i Higen og Stræben. Al Tanke er derfor ikke, som hun
forhen havde sagt, forsvunden i »Hvilebønnen«. Det er
Forskellen mellem diskursiv Tænken og instinktiv Koncentration,
der nu staar tydeligere for hende. Efterat
hun havde faaet det bekræftet af »en meget lærd Mand«,
lod hun sig ikke mere anfægte af, at der kunde være
Larm og Uro i Sjælens indre og ydre Regioner, medens
hun i sin inderste Higen samlede sit Sind om den ene,
store, guddommelige Tanke.

Andre psykologiske Erfaringer, der ogsaa synes at
stamme fra en senere Tid end Selvbiografien, er ligeledes
af Interesse for de mystiske Tilstandes Psykologi.

Det hjælper ikke at øve Selvtvang for at fremkalde
den Tilstand, som kaldes »Hvilebøn«. Denne Bemærkning
er rettet mod nogle »Kontemplative«, som hævder,
at man med Opmærksomhed skal vende sig mod, hvad
Gud vil virke i os, medens man ikke maa lade Eftertanken
arbejde. Saadan Forberedelse hjælper ifølge Teresa
slet ikke. Førend vi har vendt Sindet i Retning af
Gud, kan Gud allerede have ført os ind i Sjæleslottet,
uden at vi vide, hvorledes vi er komne derind. Og naar
Forberedelse ikke hjælper, hvorfor saa ikke bruge vor
Eftertanke til noget nyttigt, istedetfor at anstrenge sig
for at lade være at tænke! Brug dine Evner paa bedste
% --- p. 63 ---
\setcounter{footnote}{0}%
\opage{63}Maade, indtil Gud selv fører saa høje Emner frem for
dig, at Tanken synker hen i hellig Betagethed! --- Ved
disse Bemærkninger sættes der aabenbart en større Forskel
mellem Eftertankens Bøn (oraison de méditation)
og Hvilebønnen (oraison de quiétude), end hun forhen
havde sat. Dette Punkt er af saa meget des større Betydning,
som alle Tiders Mystikere og Ekstatikere netop
i den pludselige, uforberedte og uvilkaarlige Maade, paa
hvilken deres højeste Oplevelser opstod, saa et Bevis
for, at de havde et overnaturligt Udspring, --- hvad enten
saa Gud eller Djævelen var denne overnaturlige
Aarsag.

»Hvilebønnen« er, trods dens Liflighed, ikke den
højeste aandelige Tilstand. Og den kan indeholde en
Fare, idet den kan efterlade en Periode af Passivitet, af
Hensynken, der især for melankolske Naturer har en
særlig Tiltrækning, saa at de svælge i saadanne Tilstande
og søge at fremkalde og vedligeholde dem ved
lang Tids Bøn, Vaagen og Spægelse. De mener derved
at opnaa en Fortsættelse af den »Hvilebønnen«s Tilstand,
der er Guds Værk, og er maaske endog tilbøjelige
til at forveksle, hvad de paa denne Maade oplever,
med de højeste (endnu ikke omtalte) ekstatiske Tilstande.
Nydelsen ved en naturlig Tilstand forveksles med Glæden
ved en overnaturlig Titstand. Og Djævelen har sikkert
sin Finger med i Spillet, idet han frister svage Naturer
til gennem Spægelse at naa det, der maa komme
som en Gave. Spægelser med dette Formaal bør derfor
forbydes af Skriftefaderen eller Priorinden.

\sectionhead{2. \textit{Det højeste Stadium.}}

Den højeste Form for Bøn er Enhedsbønnen (oraison
d'union). I Hvilebønnen var Sjælens Higen vel koncen%
% --- p. 64 ---
\setcounter{footnote}{0}%
\opage{64}treret om et eneste Emne, men alligevel var der endnu
Mulighed for, at Erindring og Eftertanke kunde tumle
sig. Denne Mulighed er forsvunden i Enhedsbønnen.
Enhedsfølelsen bliver her til Ekstase, der snart beskrives
som Opløftelse (élévation), snart som Henrykkelse (ravissement).
Ekstasen forholder sig til den blotte Enhedsfølelse,
som et Baal forholder sig til en sagte Ild. I den
sagte Ild ligger Jernet endnu som en sort Masse, men
i Baalet bliver det glødende. --- I Bogen om »Sjæleslottet«
beskrives Ekstasen paa lignende Maade som i
Selvbiografien; den hører der hjemme i den femte
Bolig.

Beskrivelsen af denne fører Teresa ind paa nogle
interessante Bemærkninger. Man mærker, at hendes Opmærksomhed
for Nuancer og Overgange er bleven skærpet
i Oplevelsernes Skole.

a. Ved sin pludselige Indtræden, saavel som paa sit
Højdepunkt, er den ekstatiske Tilstand ganske simpel
og enkelt. Sjælen har ingen Følelse og ingen Villie; og
den véd ikke, hvad den er opfyldt af, --- véd ikke, hvad
den elsker. Der kundgør sig en umiddelbar Nærværelse,
uden at der ses eller høres noget. Tilstanden er derfor
ubeskrivelig. Og dog er den en Tilstand af umiddelbar
Vished. Hvis man vil spørge, hvorledes der kan være
Vished, naar der hverken ses eller høres noget, er Svaret,
at det er en af Almagtens Hemmeligheder. Det er som
en Føde, vi mærker i os, uden at vi har spist den.
Bagefter, naar den stærkeste Ekstase er forbi, kan der
opstaa Visioner og Auditioner, som kan være Genstand
for Beskrivelse. Og senere igen kan man forsøge gennem
udtrykkelig Sammenligning at gøre sig Oplevelsen klar.

Den simple Vished om en Nærværelse kan kaldes en
Vision, men man maa saa skelne mellem aandelig Vision
(vision intellectuelle) (og Audition), som er uden Billeder
% --- p. 65 ---
\setcounter{footnote}{0}%
\opage{65}og uden Ord\footnote[1]{I et Brev fra 1576 beskriver Teresa vision intellectuelle som Opfattelse af en Nærværelse, uden at noget indre eller ydre Syn er forbundet dermed. Man mærker, paa hvilken Side Genstanden er, tydeligere, end om man saa den. Ogsaa den umiddelbare Genkendelse omtales i dette Brev. \textit{Lettres de Sainte Térèse}\textsuperscript{2}. Trad. par Brouix. 1882. I p. 366.}, og de egentlige Visioner (visions imaginatives)
(og Auditioner). Ofte opstaa Hørehallucinationer
før Synshallucinationer. Da Teresa spurgte, hvem det
var, der kundgjorde sig for hende i en intellektuel Vision
(Nærværelseshallucination), gav Jesus sig tilkende ved at
tale til hende. Hun havde dog i Forvejen (ved den blotte
Nærværelse) været vis paa, at det var ham.

b. De egentlige Visioner (de imaginative), i hvilke der
ses eller høres noget, udvikler sig ofte successivt. Saaledes
lod Jesus, efter forud at have virket ved sin blotte, ubestemte
Tilværelse, hende først se hans Hænder, og derefter
hans øvrige Legeme.

Hvad selve disse billedlige Visioner angaar, skelner
Teresa mellem saadanne, i hvilke hun saa Jesus selv,
og saadanne, i hvilke hun kun saa et Billede af ham.
--- Denne Forskel synes at svare til den, man har gjort
mellem Hallucination og Pseudohallucination. I Pseudohallucinationen
staar Genstanden ligesom med symbolsk
Karakter, uden direkte Realitet. ---

Den successive Opstaaen af Visioner er af stor Interesse,
fordi der --- som allerede ovenfor bemærket ---
derved aabnes større Adgang for Forestillinger, Individet
tidligere har haft, til at faa Indflydelse paa Indholdet
af Visionerne. Derved forstaas, hvorfor Visionerne gennemgaaende
stemmer med de Visionæres hele Erfarings- og
Forestillingskreds. Allerede hvor der ingen Billedvision
er, men hvor, som ved Teresa's »intellektuelle
Vision«, den umiddelbare Fornemmelse eller Følelse af
% --- p. 66 ---
\setcounter{footnote}{0}%
\opage{66}en Nærværelse udfylder hele Tilstanden, begynder en
Tydning i Form af Genkendelse: Teresa var jo vis paa,
at det var Jesus, der var nærværende, skønt hun endnu
intet hørte eller saa! --- Vi vil siden se et Eksempel paa
umiddelbar Fornemmelse eller Følelse af en Nærværelse,
der beholder sin ubestemte Form, uden nogen Genkendelse.
Det kommer an paa, hvor nær og hvor fast de
tydende Forestillinger ligger.

c. Den begejstrede og begejstrende Henrykkelse kan
være forbunden med Smerte, fordi den kommer saa
overvældende, og fordi den guddommelige Indvirken er
som en fortærende Ild. Under selve Husvalelsen kan
der lides Smerte --- ogsaa fordi Trangen er saa dyb, at
den stedse paany gør sig gældende. Engang, da Teresa
var sunken hen, betagen af Savn og Længsel, saa hun
ved sin Side en Kerub, der gennemborede hende med
en glødende Jernpil, der stødtes gennem hendes Hjerte
ned i hendes Indvolde, som syntes at følge med, da
Pilen blev dragen ud. »Smerten ved dette Saar var«,
siger hun, »saa levende, at den fremkaldte Stønnen;
men dette uudsigelige Martyrium lod mig til samme Tid
føle den lifligste Fryd, og jeg kunde ikke ønske, at den
skulde holde op. Det er ikke en legemlig Lidelse, men
helt aandelig, skønt Legemet er med deri.«

Det vilde vistnok være urigtigt at opfatte denne Beskrivelse
som Vidnesbyrd om religiøse Ekstasers erotiskseksuelle
Udspring. Den Freud'ske Skole kan ikke uden
videre tage sig det til Indtægt. Forholdet er snarere det,
at den voldsomme Betagethed er ligesom strømmet ud
over de centrale Organer og har taget hele Legemet i
Beslag, saa at ogsaa »Indvoldene« er komne med. Jeg
har i min Psykologi kaldt dette Fænomen Følelsernes
Diffusion. I en saadan Diffusion, der kan iagttages ved
alle stærke Sindsbevægelser af inderlig og gennemtræn%
% --- p. 67 ---
\setcounter{footnote}{0}%
\opage{67}gende Art, ligger en af Grundene til, at det erotisk-seksuelle
Sprog ligger saa nær, naar Mystikerne søger Udtryk
for deres Oplevelser\footnote[1]{Smlgn. \textit{Religionsfilosofi} § 62.}. Spørgsmaalet er dog indviklet, og
et andet mystisk Fænomen vil i et senere Afsnit give
Anledning til at komme tilbage til det.

\sectionhead{3. \textit{Teresa og Kirkelæren.}}

Kun enkelte Gange antyder Teresa en Forskel mellem
Oplevelse og Tydning. Som Regel var der med Henrykkelsen
forbundet en umiddelbar Vished, der stod for
hende selv bagefter som et Under. Og ligeledes gjorde
der sig, selv om intet saaes eller hørtes, en umiddelbar
Genkendelse gældende, som nærmere kunde udfoldes og
motiveres gennem Auditioner og Visioner. Naar hun
kommer til at opkaste Spørgsmaalet om, hvad der rørte
sig hos hende i de ekstatiske og visionære Tilstande,
var det snarere et psykologisk end et teologisk Problem,
hun stillede. Det angik de indre Oplevelser i Sjælen i
saadanne Tilstande (les sentiments intérieurs dans cet
état), --- det, som Sjælen tog sig til, medens de varede
(ce que l'âme pouvait faire pendant ce temps). Spørgsmaalet
stillede sig for hende, naar hun, efter at Ekstasen
var ophørt, og hun havde kommuniceret, beredte sig til
at skrive derom. Men hun kunde ikke blive klar derover
og erklærede det for en guddommelig Hemmelighed.
Især var Oplevelsen af en Nærværelse uden Syn
og Hørelse gaadefuld. Men --- føjer hun til --- noget
maa hun dog have kunnet gribe fra disse Tilstande,
ellers havde hun end ikke kunnet udtrykke sig om dem
ved Hjælp af Sammenligninger.

For at udtrykke, hvorledes Sjælen undertiden kan
hæve sig over sig selv uden dog at ophøre at være sig
% --- p. 68 ---
\setcounter{footnote}{0}%
\opage{68}selv, siger hun, at det er ligesom en Arneild, der flammende
hæver sig i Vejret: jo voldsommere Ilden er, des
højere hæver Flammen sig over Arnen, uden dog at
ophøre at være Arnens Ild. Dette Billede, der skal anskueliggøre
Forholdet mellem »Sjælen« og »Aanden«,
bærer, saaledes som det fremkommer i ottende Kapitel
af Selvbiografien, Præget af et uvilkaarligt Symbol. Derimod
fremtræder en mere vilkaarlig Sammenligning i
det Billede, der udføres baade i Bogen om »Sjæleslottet«
og i Selvbiografien, med Haven og de forskellige Maader,
paa hvilke den kan vandes. Det er en vidt udspunden
Allegori. En anden Allegori er den, der er udtrykt i
Titelen paa Bogen om »Sjæleslottet«. Sjælens forskellige
Udviklingstrin i Retning af Fordybelse i og Enhed med
Guddommen sammenlignes med de forskellige Boliger
i Slottet. Der skelnes mellem syv saadanne Boliger, og
der kan raade Fred og dyb Stemning i et Værelse, skønt
der høres Støj og Uro fra et andet.

Som Allegorien forholder sig til Symbolet, den vilkaarlige
Sammenligning til den uvilkaarlige, saaledes
forholder --- naar hun saa at sige gaar over fra Psykologien
til Teologien --- den bevidste teologiske Tydning af
Oplevelserne sig til den umiddelbare Forvisning om guddommelige
Væseners Nærværelse eller Aabenbarelse i
Ekstasen. Skønt hun ydmygt indrømmer, at hun ikke
er Teolog, er hun dog overbevist om, at hun i nogle af
sine Visioner har skuet Treenighedens Hemmelighed ---
og det netop saaledes, som den atanasianske Bekendelse
fremsætter den! Alt stod klart og anskueligt for hende.
Jesus viste hende, hvorledes én Gud kan fremtræde i
tre Personer. Men hun var (desværre) ikke istand til
bagefter at give en nærmere Forklaring af, hvad hun
havde skuet. Den højere Oplysning, der i saadanne Tilfælde
blev hende tildel, kommer undertiden gennem
% --- p. 69 ---
\setcounter{footnote}{0}%
\opage{69}Ord, hun hører, men undertiden ogsaa »paa en ubegribelig
Maade, der uden Ord gav hende en endnu
klarere Forstaaelse end den, mange Lærde i Forening
vilde have formaaet at bibringe hende«. Saaledes saa
hun i en Vision, hvorledes alle Ting er i Gud; men hun
kan ikke beskrive, hvorledes dette blev vist hende. »Jeg
saa dem ikke i deres særegne Skikkelser, og dog var
det Syn, jeg havde, af en overordentlig Klarhed.«

Naar hun teologiserede med fuld Bevidsthed, drog
hun undertiden Slutninger af sine Oplevelser, der forargede
hendes Skriftefædre. At Gud er nærværende i alle
Væsener, blev hende klart derved, at hun i »Enhedsbønnen«
opdagede Guds Nærværelse i hendes Sjæl.
»Nogle Folk, som ikke var lærde«, sagde til hende, at
Gud kun gennem Naaden bliver nærværende i os; men
en lærd Dominikaner forklarede hende, at Gud i Virkeligheden
var nærværende i alle Ting, hvad der var hende
en stor Trøst. --- Den »lærde Dominikaner« Mester Eckehart
havde ogsaa denne Tro, der hos ham hang nøje
sammen med meget kætterske Anskuelser (se ovenfor p.
53), som sikkert vilde have vakt Teresa's Forargelse.

Teresa følte sig, som man ser, sikker i sin Sag, men
det var hende dog en Trang at faa en kirkelig Bekræftelse
paa sine Tydninger. Ved dette dobbelte Forhold
til religiøs Autoritet bestemtes særlig Forholdet til hendes
Skriftefædre, der allerede er berørt ovenfor, men
fortjener en nærmere Udredning.

I 20 Aar havde hun, siger hun, forgæves søgt en
Skriftefader, der forstod hende. En Skriftefader skal have
en god Dømmekraft, personlig Erfaring og grundig Lærdom.
Der er stor Fare ved at blive ledet af en Mand,
der er ukyndig i Kirkelæren. Hellere opgive den aandelige
Bøn end fra Begyndelsen af komme paa Afveje!
Men de Skriftefædre, der vel har teologisk Lærdom, men
% --- p. 70 ---
\setcounter{footnote}{0}%
\opage{70}ikke kender de indre Oplevelser, er tilbøjelige til at mene,
at Ekstaser og Visioner stammer fra Djævelen. Teresa
selv blev en Tidlang af alle, med Undtagelse af hendes
jesuitiske Skriftefader, betragtet som et Offer for Djævelens
Indskydelser. Man raadede Skriftefaderen til at
have Mistro til hende, saa hun frygtede, at hun tilsidst
ingen Skriftefader vilde kunne faa. Skriftefædre maa
være ydmyge, ikke bilde sig ind at forstaa, hvad der nu
engang er en lukket Verden for dem. Teresa har den
største Medlidenhed med ekstatiske Sjæle, der lider under
ukyndige Skriftefædre, som ikke forstaar dem. Og hun
klager over, at naar man beretter om Visioner, Djævelen
har indskudt, har Skriftefædrene ingen Tvivl, hvorimod
de straks er betænkelige, naar man beretter om
Visioner, der har Gud til Ophav. (Fondation de Medine
des Champs. chap. 8). Hun blev ofte truet med Inkvisitionen
paa Grund af sine Visioner.

Saavidt man kan se af Teresas Fremstillinger, stammede
Skriftefædrenes Betænkeligheder ikke fra Indholdet
af hendes Visioner. Hun var ogsaa for katolsk præget
og sindet fra Barndommen af til at indlade sig paa
Kætterier. Hun siger om sig selv: »Hun er saa urokkelig
knyttet til Trosbekendelsen, at alle tænkelige Aabenbaringer,
selv om de viste hende Himlene aabne, ikke
vilde rokke hendes Tro, end ikke paa den mindste af
de Artikler, Kirken lærer. Dersom det sker, at Sjælen
vakler i sin Tro angaaende en kirkelig Lære, eller dersom
den blot et lille Øjeblik stanser ved den Tanke,
at naar Gud har sagt noget til den, saa kan dette være
lige saa sandt, som det, han har sagt Helgenerne, --- saa
vilde denne Vaklen eller denne Tanke stamme fra Djævelen\dots{}
Naar et Menneske ikke finder denne mandige
Styrke i sin Tro, og naar de fromme Stemningers
Inderlighed og de Visioner, der bliver det tildel, ikke
% --- p. 71 ---
\setcounter{footnote}{0}%
\opage{71}bidrager til at styrke Troen, saa maa det ikke betragte
dem som tilforladelige\dots{} Jeg ser klart, og jeg véd af
Erfaring, at man ikke bør være sikker paa, at noget
stammer fra Guds Aand, uden forsaavidt det stemmer
med den hellige Skrift. Var der den mindste Afvigelse,
vilde jeg tro, at Visionerne stammede fra Løgnens Fader.«
(Selvbiografi. Kap. 25)\footnote[1]{Teresa havde undertiden Visioner af Djævelen, i nogle Tilfælde »intellektuel Vision«, i andre Tilfælde »imaginativ Vision«; desuden havde hun Anfald af Tomhed, Frivolitet og Ligegyldighed, som hun var vis paa, direkte vaktes af Djævelen. Ogsaa Helvede viste sig undertiden for hende, ganske stemmende med de Beretninger, hun havde læst og hørt om det. Ja, hun saa i en saadan Vision endog den Plads i Helvede, der stod beredt til hende, hvis hun ikke skulde blive frelst! (Vie de Sainte Térèse. Chap. 30—31).}.

Betænkelighederne maa have skrevet sig fra, at hun
overhovedet havde anden Adgang til at komme i Forhold
til Gud end gennem Kirke, Præster og Sakramenterne.
Hun er ogsaa selv undertiden betænkelig ved det
umiddelbare Forhold, der er blevet hende tildel. Fra et
kirkeligt Synspunkt ligger der en Fare i ethvert saadant
umiddelbart Forhold, i hvilket Individet bliver uafhængigt
af Overlevering og Kultus. Denne Fare søgte Kirken
at afværge ved sin Fordømmelse af to saa fremragende
Mystikere som Eckehart og Suso; men hos dem var det
jo heller ikke blot Kilden, men selve Indholdet, der østes
af Kilden, man vilde tillivs. Hos en Mystiker i det tyvende
Aarhundrede vil vi se Modsætningen mellem Oplevelse
og Overlevering træde meget skarpt frem. ---

En Tvivl om den objektive, reale Karakter af, hvad
der rørte sig i Ekstasen og kundgjordes i Visionerne,
spores ikke hos Teresa, ligesaa lidt som hos hendes
betænkelige Skriftefædre. Den Mulighed, at de store Oplevelser
blot skulde skyldes hendes egen Indbildning,
afvises straks, dels paa Grund af den Modstand, Sindet
% --- p. 72 ---
\setcounter{footnote}{0}%
\opage{72}prøver at gøre mod at hengive sig til dem, dels endelig
paa Grund af deres Virkninger til at skaffe Sjælen
Fred og vække Trang til Kærlighedsgerninger. Der gør
sig intet Enten---Eller mellem Indbildning og Virkelighed
gældende, men derimod et Enten---Eller mellem to
Virkeligheder, en, der udspringer fra Gud, og en anden,
der udspringer fra Djævelen. Kriteriet til Afgørelse mellem
disse to Ophav var for Teresa de Virkninger, de
store Oplevelser førte til: Glæde, Mod, Afsky fra Verden,
Følelse af egen Intethed. Teresa staar her paa samme
Stade som den gamle Kirke (se ovenfor p. 42).

\sectionhead{4. \textit{Selvbiografi og Systematik.}}

Hidtil har vi i denne Fremstilling kunnet benytte
Selvbiografien og »Sjæleslottet« Side om Side. Men paa
det Punkt, hvortil vi nu er naaede, gør der sig en Forskel
gældende mellem den biografiske og den systematiske
Fremstilling.

I Selvbiografien fortæller Teresa, hvorledes der efter de
høje Oplevelser, hun har beskrevet, er opstaaet en usigelig
smertelig Lidelse, som hun ikke kan frigøre sig for, og som
hun har stor Vanskelighed ved at beskrive. Hun tilføjer:
»Det er denne Tilstand af usigelig Lidelse, hvori
jeg nu befinder mig. Den ligger altsaa efter den Periode,
da jeg, trofast i Bønnen, modtog saa store Glæder og
Naadesbevisninger af Herren. Af og til er han saa naadig
at ødsle saadanne paa mig; men min Sjæls sædvanligste
Tilstand er at føle sig knuget af den Smerte, jeg
nu skal tale om.« Den Beskrivelse som nu følger, støtter
sig først til en Sammenligning med den Smerte, der
--- paa Grund af Sjælens uendelige Trang og ringe Evne
--- føltes midt under selve Henrykkelsen. Forskellen
er, at nu føler Sjælen sig forladt --- »som den ensomme
Spurv paa Taget«. Gud synes at være langt borte. Og
% --- p. 73 ---
\setcounter{footnote}{0}%
\opage{73}bliver der Sjælen et Glimt tildel af Herligheden, saa
virker det kun til, at Forladtheden føles saa meget des
stærkere. Sjælen spørger sig selv: Hvor er nu din Gud?
(Selvbiografien Kap. 20—21).

Denne Skildring svarer til, hvad der i »Sjæleslottet«
anføres som karakteristisk for den sjette Bolig. Opholdet
her betyder en Tilstand af store Lidelser. Udefra møder
Miskendelsen, og Gud er fjern. Og Smerten ved at være
skilt fra Gud er des større, jo herligere Visionerne var
i fordums Tid. Hvad er der da at gøre? Fremtvinge de
høje Oplevelser kan man ikke, og derfor bør man i de
mørke Tider fordybe sig i Betragtninger af Jesu Menneskelighed,
af Madonna og Helgenerne, --- Betragtninger,
man ikke skal kaste Vrag paa, fordi de ikke naa op til
Kontemplationens Højder. Vi er nu engang ikke Engle,
i hvem Kærligheden aldrig slukkes! --- Der hører stort
Mod til at udholde en saadan Tilstand.

I »Sjæleslottet« er dette Stadium, der i Selvbiografien
er det afsluttende, Indledning til et sidste Stadium, den
syvende Bolig, om hvilken Intet meldes i Selvbiografien,
der jo ogsaa er 15 Aar ældre end det mere systematiske
Værk.

To Ting siges at være karakteristiske for Opholdet
i den syvende Bolig, nemlig for det første en inderligere
Fordybelse i de religiøse Forestillinger, ikke længere
blot Forstaaelse, men aandelig Skuen af Mysteriet,
--- og for det andet en større Kontinuitet mellem de
forskellige Tilstande, der i den sjette Bolig stod i saa
pinlig Modsætning til hverandre; Ekstase og Aandsflugt
er sjeldne, men den aandelige Udtørrethed er forbi;
naar Gud undertiden overlader Sjælen til den selv, er
det kun, for at den saa meget mere skal skønne paa,
hvad det betyder at være i Samfund med Gud. ---

Af disse to Ejendommeligheder er det nu kun den
% --- p. 74 ---
\setcounter{footnote}{0}%
\opage{74}sidste, der er afgørende. Ti, som ovenfor (p. 68) omtalt,
kunde Teresa allerede paa tidligere Stadier --- 15 Aar
før Affattelsen af »Sjæleslottet« --- i ét Syn opfatte Treenighedens
Hemmelighed. I et Brev fra Aarene 1560—1562
(de samme Aar, i hvilke det første Udkast til Selvbiografien
udarbejdedes) meddeler Teresa sin Skriftefader
en imaginativ Vision af Treenigheden, som stemmer
med, hvad hun havde hørt lærde Mænd sige. Og denne
Vision er aabenbart den samme som den, der berettes
i sidste Kapitel af Selvbiografien, og som den, hvis Virkninger
paa Sindet beskrives ved »Sjæleslottet«s syvende
Bolig\footnote[1]{\textit{Lettres de Sainte Térèse}. I p. 37 ff. (med Udgiverens Noter).}. Det er derfor ikke rigtigt, naar man har ment,
at Teresa først i den sidste Periode af sit Liv har formaaet
at forene Bibelens og den mystiske Oplevelses
Gud. Denne Evne har hun --- hvad hele den foregaaende
Fremstilling viser --- allerede tidligere besiddet,
og paa intet af sine Stadier har hun haft den mindste
Tvivl om Kirkelærens Sandhed, saa imprægneret som
hun fra først til sidst har været af Katolicismens Aand.

Beskrivelsen af den syvende Bolig i »Sjæleslottet«
angaar altsaa ikke et Stadium, paa hvilket der opnaaedes
Intuitioner, der ikke tidligere skulde have været mulige.
Derimod gengiver den en Periode, hvor de store Modsætninger
paa det indre Livs Omraade var faldne bort
og afløste af en Kontinuitet, indenfor hvilken der vel
kunde fremkomme Stigen og Falden, men hvor dog
Billedet med den ensomme Spurv paa Taget ikke paa
noget Punkt kom til at gælde. Der forekommer ikke
mere de overvældende Ekstaser, heller ikke den smertelige
Længsel, den brændende Tørst efter at føle sig Et
med Gud. Der raadede gennemgaaende en Ligevægt, ---
en Ligevægt, som Afbrydelserne lærte at skønne saa
% --- p. 75 ---
\setcounter{footnote}{0}%
\opage{75}meget des mere paa. Dette Karaktermærke skiller den
sidste af de Tilstande, der beskrives i »Sjæleslottet«, fra
den Tilstand, med hvilken Selvbiografien slutter.

I en interessant Diskussion om Teresa's Mystik i det
franske filosofiske Selskab hævdede Delacroix, at Ekstasen
ikke betegner Højdepunktet i Teresa's mystiske Liv,
men at det maa søges i »Sjæleslottet«s syvende Bolig,
hvor Modsætningerne overvindes, og hvor en rolig og
kontinuerlig Harmoni raader, --- en Syntese eller Systematisering
af tidligere Perioders Modsætninger. »Vi ser,«
siger han\footnote[1]{\textit{Bulletin de la Société Française de Philosophie}. 26. Oktober 1905. (Le développement des états mystiques chez Sainte Térèse).}, »den hellige Teresa naa til den syvende
Bolig, denne sidste og definitive Tilstand, en Syntese af
de foregaaende. Den er Ækvivalent til Ekstasen, idet
den over hele Sjælen og over hele Livet udbreder det,
som Ekstasen havde samlet i et evigt Øjeblik, der saa
at sige levedes højt over Livet. Den syvende Bolig er
Livets Apoteose.« --- En lignende Opfattelse er fra katolsk
Side gjort gældende af Abbé Pacheu i en Bog, i hvilken
han søger at forene psykologisk Metode øg skolastisk
Teologi i sin Teori om Mystiken. Forskellen mellem
Pacheu og Delacroix er den, at medens Delacroix vil
gennemføre en rent psykologisk Metode paa alle Punkter,
hævder Pacheu, at hvor de naturlige Aarsager ikke slaar
til, maa der appelleres til Guds Indgriben, saa at den
psykologiske Mekanisme bliver et Instrument, paa hvilket
Gud spiller, --- og han gør meget hurtig Overgangen
fra Instrumentet til den guddommelige Organist. Men
deri er Pacheu enig med Delacroix, at Ekstasen ikke
er Højdepunktet i Udviklingen af Teresa's Mystik; Højdepunktet
er en Skuen, som medfører Fred i Sindet, idet
de mystiske Oplevelser og den teologiske Lære smelter
% --- p. 76 ---
\setcounter{footnote}{0}%
% CHECK p. 76: notes paired with the marks in order (the layer misread a number)
\opage{76}sammen i en Slags Krystallisation\footnote[1]{\textit{L'expérience mystique et l'activité subconsciente}. (1911). p. 76.}. Men for at begrunde
denne Teori maa den katolske Forfatter gøre Vold paa
Psykologien. Han paastaar nemlig, at der er en Erindringssammenhæng
mellem Ekstasen og den senere Besindelse,
der fører til Tydning; derved skal Ekstase være forskellig
fra Hypnose. Naar man er kommen til sig selv
efter Ekstasen, skal Mystikeren, mener Pacheu, have en
fuldkommen Viden af, hvad der er gaaet for sig\footnote[2]{Après être revenu à soi ou sait parfaitement ce qui c'est passé. \textit{L'expérience mystique et l'activité subconsciente}. p. 142. cfr. p. 176.}. Men
at dette altid skulde være muligt, strider mod de udtrykkelige
Vidnesbyrd af Paulus, Plotinos, Suso og Teresa.
Ekstasen udelukker paa sit Højdepunkt enhver Tydning,
ogsaa enhver egentlig Skuen. Som vi har set, er Teresa
henvist til uvilkaarlige og vilkaarlige Sammenligninger,
til Symboler og Allegorier for at udtrykke, hvad hun
har oplevet. Og selv de imaginative Visioner udfolder
sig i Regelen successivt, saa at der bliver Mulighed for
tidligere vundne Forestillingers uvilkaarlige Indflydelse.
Endelig er Teresa, ligesom Suso, meget omhyggelig for
at udelukke alle ikke-ortodoxe Elementer fra den Tydning,
hun fastholder. Der foregaar her et kritisk Arbejde,
der ikke bygger paa Analyse af selve Oplevelserne, men
skal sikre Overensstemmelsen med Overleveringen. ---

For at vende tilbage til Teresa selv, saa staar utvivlsomt
Ekstaserne for hende som de højeste Tilstande, og
hun jubler over den Herlighed, der i dem er bleven
hende tildel, saa lidet hun end kan tydeliggøre den for
andre. Med Vemod erkender hun i Selvbiografien, at
deres Tid er forbi, og at hun nu lever i Længselens og
de aandelige Svingningers Tid --- svarende til, hvad hun
15 Aar efter kalder at leve i den sjette Bolig. Alligevel
er Skildringen af Livet i den syvende Bolig i »Sjæleslottet«
% --- p. 77 ---
\setcounter{footnote}{0}%
\opage{77}næppe en ren Konstruktion. Den hviler paa den Kendsgerning,
at Teresa har forsonet sig med, at de store Oplevelsers,
de mægtige Syners Tid er forbi, og at hun nu
paaskønner den Harmoni og Kontinuitet, som hendes
Liv har vundet paa hendes gamle Dage. Om der, som
Delacroix paastaar, er Ækvivalens mellem Ekstasernes
mægtige momentane Bevægelse og den senere Tilstands
dæmpede Harmoni, er et meget vanskeligt Spørgsmaal
at afgøre. Det er i Grunden det samme Problem som
det, hvorvidt der i Sindelagets Form kan komme fuld
Erstatning for, hvad der oplevedes i Sindsbevægelsens
Form. I min Psykologi har jeg drøftet dette Problem.
Men hvad Teresa angaar, saa hævder hun ikke nogen
Ækvivalens her. Hun anstiller ingen Sammenligning
mellem det Stade, paa hvilket hun skrev Selvbiografien,
og det, paa hvilket hun forfattede »Sjæleslottet«. Ellers
kunde det ikke være undgaaet hende, at hun allerede
i, hvad der svarede til den femte Bolig havde kunnet
skue de teologiske Mysterier i fuld Klarhed, saa det ikke
var nogen Ejendommelighed for »den syvende Bolig«.
Hun siger selv flere Gange, at naar praktiske Hverv
havde afbrudt hendes Fremstilling, læser hun ikke det
foregaaende igennem, før hun fortsætter.

I et Brev fra Maj 1884 (Aaret før hendes Død) gør
hun Rede for sin aandelige Tilstand. Denne Beskrivelse
svarer til Skildringen af »den syvende Bolig«, som hun
forfattede fire Aar forud. Der er dog et vist Vemod tilstede,
og Modsætningen til de foregaaende Perioder betones
meget stærkere end i »Sjæleslottet«. Hun føler sig
nu i Sikkerhed og har ikke mere Trang til stærke legemlige
Spægelser. Imaginative Visioner er hørte op, medens
intellektuelle Visioner stadig kan komme; Treenigheden
bor til Stadighed i hendes Sjæl. Saadanne flammende
Bevægelser, som hun tidligere kunde stedes i
% --- p. 78 ---
\setcounter{footnote}{0}%
\opage{78}opstod nu ikke, men heller ikke saa stærke Smerter. Det
er en rolig og forsonet Alderdomsstemning, der her raader,
--- men ingen »højere Enhed«, ingen Ækvivalens i
Forhold til, hvad de store Oplevelsers Tid havde bragt.
Og at det er »den syvende Bolig«, hun nu mener at
leve i, fremgaar af et Brev fra November s. A. Hun beder
her en Veninde læse Kapitlet om »den syvende Bolig«
for Pater Rodrigo og sige ham, at »den, han véd
nok«, er naat til den beskrevne Tilstand og nyder den
fuldkomne Ro og Fred, hvorom der tales i det nævnte
Afsnit\footnote[1]{\textit{Lettres de Sainte-Térèse}. III p. 367.}. ---

Ak, Alderdommen kan bringe sin Fred og sin Resignation;
men den bringer ikke altid, vel endog meget
sjælden, fuld Erstatning for, hvad der i sin Tid rørte
sig mægtig i os. Selv om der med dette kraftigere Liv
ogsaa fulgte Uro, Tvivl og Spænding, saa var dog ogsaa
dette Vidnesbyrd om Livets Kraft. Hvor man tror at
have faaet fuld Ækvivalens, end sige at have naat en
højere Enhed, beror det ofte paa en Illusion. --- Dog
havde den hellige Teresa en lykkeligere Alderdom end
de to danske Præster, der paa deres gamle Dage bekendte
deres Fattigdom for hinanden\footnote[2]{Se nedenfor min lille Studie \textit{To Præster}.}.

Den Tid, hvor Livet rører sig kraftig, og hvor det
egentlige Arbejde gøres, skal have sin Ret, og beholde
den. ---

Sammenligner man Teresa med Middelalderens Mystikere,
faar man Indtrykket af større psykologisk Evne
ved Beskrivelsen af de mystiske Oplevelser i deres successive
Udvikling, mere rent historisk i Selvbiografien,
mere systematisk og paavirket af Alderdommens Ro i
andre Skrifter (især »Sjæleslottet«). Derved er det for en
% --- p. 79 ---
\setcounter{footnote}{0}%
\opage{79}undersøgende Betragtning lettere at opdage Forholdet
mellem Oplevelse og Tydning. Hos hende selv er naturligvis
dette Forhold ikke kommet til klar Bevidsthed,
dertil er hun fra først af for indlevet i den kirkelige
Overlevering, skønt hendes Selvstændighed viser sig i
den kritiske Holdning overfor Skriftefædrene. Selv ved
de mest ubestemte, uartikulerede og usædvanlige Tilstande
optræder umiddelbar Genkendelse og tillidsfuld
Forklaring. Hos moderne Mystikere var det at formode,
at det paa dette Punkt vilde stille sig anderledes, og der
foreligger et højst interessant Eksempel paa, at det ogsaa
er Tilfældet.

% --- p. 80 ---
\setcounter{footnote}{0}%

\opage{80}\chapterhead{IV.}{EN MODERNE MYSTIKER}

\sectionhead{1. \textit{Psykologi og Mystik.}}

Den udmærkede Psykolog Théodore Flournoy i Genève,
som vi allerede ovenfor (p. 32f) har haft Lejlighed
til at omtale, fik i December 1910 Besøg af Forstanderinden
for et i ortodox protestantisk Aand ledet Opdragelsesinstitut
for unge Piger, Mlle Cécile Vé, som han
kalder hende i sin Beretning\footnote[1]{\textit{Une Mystique Moderne}. Dcouments pour la Psychologie religieuse. (Archives de Psychologie. Ed. Flournoy et Claparède. Tome XV. Mai 1915.) --- Det er Cécile's egne Oplevelser og Meddelelser, der her foreligger, for største Delen fra Tiden, før hun vidste, at Flournoy tænkte paa at udgive dem. Jeg har i det følgende paa Grundlag af dem forsøgt en Karakteristik af hendes mærkelige Oplevelser og sammenligner den senere med den Opfattelse, Flournoy har gjort gældende.}. Hun kom for at spørge
ham tilraads om sin hele psykiske Tilstand. Denne var
en egen Art Dobbeltbevidsthed. Hun følte sig delt mellem
to Jeger. Det ene var sanselig-lidenskabeligt og ytrede
sig i stærke erotiske Impulser, der antog en dæmonisk
Karakter og gjorde hendes Nætter urolige og pinlige.
Det andet fremtraadte i hendes Stræben efter at arbejde
i sin Lærervirksomhed, besjælet af positiv-religiøse Forestillinger,
saaledes som de raadede i den Kreds, indenfor
hvilken hun virkede. Hendes Tilstand faldt dog hverken
ind under egentlig Dobbeltbevisthed, hverken den
% --- p. 81 ---
\setcounter{footnote}{0}%
\opage{81}simultane eller den successive, ti hun oplevede begge
Tilstande som sine egne, og hendes to »Jeger« havde
gensidig Bevidsthed om hinanden. Spaltningen var ikke
fuldstændig, og den blev netop derved saa meget des
pinligere. »Je suis à la fois A et B, et l'une aussi sincèrement
que l'autre.«

De to »Jeger« viser tilbage til de Modsætninger, mellem
hvilke hun havde været omtumlet i sit tidligere
Liv. Hun var 47 Aar, da hun henvente sig til Flournoy.
18 Aar gammel var hun pludselig bleven reven ud af
sin barnlige Enfoldighed om seksuelle Forhold gennem
brutal Voldtægt. Denne Episode kunde hun aldrig glemme.
Erindringsbilledet dukkede atter og atter op med
en Blanding af Rædsel, Afsky og Vellyst. Senere havde
hun følt Tilbøjelighed for en gift Mand, men havde
brudt med ham. Forholdet synes at have haft en legende
og letfærdig Karakter. Men dette, at hun ikke
havde faaet Tilfredsstillelse for sin erotiske Trang, og
ikke mindst for sin Trang til at opleve Moderforholdet,
fremkaldte voldsomme Kriser hos hende, i hvilke hun
syntes, at dæmoniske Kræfter overvældede hende. Hvad
hun oplevede, betegnede hun selv som »natlige genitale
Hallucinationer«, »Mareridt«, »Halvsonnambulismer«,
»Besættelse af en lavere Personlighed«. Flournoy siger
med Rette, at hvis hun havde levet i tidligere Tider,
vilde hun være bleven betragtet som besat. Den »Halvsonnambulisme«,
hvoraf hun led, kom til at staa i et
saa meget grellere og mere afskrækkende Lys for hende,
som hendes Religiøsitet, efter hendes Hjemkomst fra et
Ophold i Udlandet, havde lagt sig til Hvile i den ortodoxe
Protestantismes Overleveringer. I Tilslutning til
dem virkede hun med Alvor og Energi som Leder og
Lærerinde ved et Institut for unge Piger. Modsætningen
mellem, hvad der rørte sig hos hende i »Besathedens«
% --- p. 82 ---
\setcounter{footnote}{0}%
\opage{82}Tider, og hvad der baade var hendes personlige Tro og
besjælede hendes Undervisningsarbejde, blev hende tilsidst
utaalelig.

Det er betegnende, at hun gik til en Psykolog og ikke
til en Præst. Hun havde ikke blot psykologisk Sans,
men var heller ikke ubekendt med moderne Psykologi.
Hun saa idethele paa Forhaand anderledes paa sine
Oplevelser, end hun vilde have gjort, hvis hun havde levet
i tidligere Tider eller var mødt med andre Forudsætninger.
Hun blev ikke skuffet ved at henvende sig til
Flournoy. Han forestillede hende ikke blot, at selv meget
troende Mennesker ikke slap for Anfægtelser, og at
selv Paulus af egen Erfaring kendte en »Satans Engel«,
men han indskærpede ogsaa, at seksuelle Synder ikke
er af saa væsentlig Betydning som Synd mod Retfærhed
og Kærlighed. Og især udviklede han for hende,
at Kønsdriften har saa dyb en Grund i alt Levendes
Natur, at den stedse paany rejser sig mod de Skranker,
det sociale Liv sætter for den, medens den paa den
anden Side hænger saa nøje sammen med Energi i Livsførelsen,
at dens Udryddelse ikke vilde medføre Renhed,
men Sløvhed. Endelig raadede han hende til nøjagtig
Iagttagelse og Undersøgelse af sine Tilstande, idet Studiet
af vore urolige og ubændige Tilstande ofte er det
bedste Middel til at forminske deres Indflydelse. Disse
Belæringer supplerede han med klog Benyttelse af Hypnose
til at befæste Indsigten og styrke Villien til at anvende
den. --- Flournoy's Maade til at tage hele denne
Sag paa minder i høj Grad om de Midler, Spinoza,
(Ethica V. 1—20) anviser overfor »Affekterne« (remedia
affectuum).
Cécile har selv udtalt, hvor velgørende Anlæggelsen
af det psykologiske Synspunkt var for hende. Det befriede
hende fra de mange Selvanklager og fra »prote%
% --- p. 83 ---
\setcounter{footnote}{0}%
\opage{83}stantiske Skrupler«. »Min virkelige guddommelige Inspiration«,
siger hun, »kom til mig den Dag, da jeg bestemte
mig til at betragte hele den endeløse Kamp fra
et medicinsk Synspunkt, og da psykologiske Forklaringer
bragte mig flot (m'ont mise au large), idet de i høj
Grad lettede min moralske Ansvarlighed\dots{} De sygelige
Tilstande kan og bør jeg kæmpe imod, men de angriber
ikke min moralske Renhed«. »Det, at jeg er bragt
flot, straaler ud i hele mit Liv. Hvis jeg i en vis Forstand
er bleven mere overbærende overfor mig selv, saa
kan jeg sige, at jeg ogsaa er bleven det overfor andre,
og at jeg formaar at elske dem paa en bedre Maade«. ---

Kampen var dog ikke forbi. Men godt halvandet Aar
efter, at hun havde henvendt sig til Flournoy, fik hun en
indre Hjælp, der gjorde, at hun kunde undvære Hypnosen.
Naar de halvsonnambule Oplevelser, som hun
kaldte dæmoniske, havde en Tendens til at personificere
sig som uhyggelige Væsener, saa optraadte der nu som
Modvægt ogsaa en begyndende Personification af de modstaaende
Tendenser i hendes Natur. Hun havde ofte en
ubestemt Fornemmelse af en usynlig Vens Nærværelse.
Alt hvad der foregik hos hende og pinte hende, kunde
hun betro ham. Synsfornemmelser var der ikke forbundet
med den »Nærværelse«, hun mærkede. Og naar
»Vennen« talte, kunde hun ikke afgøre, om Ordene klang
i hendes Indre eller kom til hende udefra. En af disse
Oplevelser beskriver hun saaledes: »Jeg mærkede, at han
kom fra den anden Kant af Værelset og nærmede sig
min Seng, saa at jeg havde Fornemmelsen af, at jeg
kunde naa ham, naar jeg udstrakte min Arm. Men den
materielle Fornemmelse spillede den mindste Rolle ved
denne Aabenbaring (apparition); det herskende Indtryk
var det af venlig Nærværelse, og Nærværelse af noget,
som var mægtigere end jeg selv. Hele mit Liv igennem
% --- p. 84 ---
\setcounter{footnote}{0}%
\opage{84}har jeg følt en stærk Trang til at støttes, til at være beskyttet,
ja, til at lyde en Villie, der er stærkere end min
egen. Og det er denne saa ofte følte Trang, der nu har
personificeret sig«. Undertiden kunde den usynlige Gæst
minde hende om hendes Fader, som hun i høj Grad
havde set op til og i høj Grad savnede. Maaske ogsaa
hendes Hypnotisør og psykologiske Skriftefader har afgivet
Elementer til Skikkelsen.

Naar hun talte til »Vennen«, fik hun ikke altid Svar;
men hun kunde mærke, om han billigede eller misbilligede,
hvad hun sagde. Naar han svarede, kom Svaret
langsomt, automatisk, og Indholdet var kun, hvad hun
ofte selv havde tænkt paa, men ikke tænkte paa i Øjeblikket.
Et Digt, hun havde lært i sin Ungdom og holdt
meget af, men af hvilket hun kun kunde huske det første
Vers, reciterede han for hende næsten fuldstændig.

Fænomenet havde udviklet sig gradvis, ligesom det
efter Flournoy's Erfaring ofte sker hos Sonnambuler.
Cécile kunde ikke angive nogen bestemt Datum for dets
første Indtræden, men hun kunde genkende det, naar
det kom, allerede før det havde antaget en bestemtere
Skikkelse. Undertiden kunde hun mærke Gæsten, før
han virkelig kom; det syntes hende da, som om der
hinsides Tærskelen til hendes Bevidsthed var noget, der
ikke kunde trænge ind, fordi hvert af hendes Øjeblikke
var optaget af hendes Arbejde eller dets Forberedelse.
Hun kunde da prøve paa at skaffe Plads til den skjulte
Erfaring (l'expérience latente), omtrent som man prøver
paa at fremkalde noget i sin Erindring, der er flygtet
bort. ---

Den Bevidsthed om en Nærværelse, Cécile her beskriver,
kender vi allerede fra Teresa. Den er, hvad den
spanske Nonne har kaldt intellektuel Vision (i Modsætning
til imaginativ Vision). Det er et Fænomen, som er
% --- p. 85 ---
\setcounter{footnote}{0}%
\opage{85}vel kendt, ogsaa udenfor det religiøse Livs Psykologi.
Og ligeledes foreligger der Eksempler udenfor denne
Gren af Psykologi paa, at Nærværelsesfornemmelse gaar
over til visuel Hallucination, ligesom det ovenfor viste
sig hos Teresa\footnote[1]{Smlg. \emph{William James}: \textit{Varieties of Religious Experience}. Chap. 3.}. Selve Nærværelsesfornemmelsen beskriver
Flournoy som en Hallucination paa Livsfornemmelsernes
Omraade (hallucination cénéstétique), i Regelen
forbunden med organisk Velbefindende (Eufori),
Ro og Hvile, og undertiden forbunden med Muskelsøvn
(engourdissement musculaire) og med kolde Ben og
Arme, --- Symptoner, han ofte har fundet ved Undersøgelse
af Sonnambuler (Medier).

Hos Cécile udviklede Nærværelseshallucinationen sig
ikke videre, som hos Teresa. Hun kom aldrig til at se
sin Ven. I det højeste havde hun, hvad hun selv kalder en
Berøringsillusion, foruden Nærværelsesfornemmelserne,
men saaledes, at hun var fuldstændig klar over, at intet
berørte hende. Hun tilføjer, at alligevel var Overbevisningen
om en Nærværelse utvivlsom; hun vidste, naar
»Vennen« kom, og naar han gik, og Besøget fremkaldte
stedse en Følelse af moralsk Hjælp, af Ro og af Renhed.
Hun vil nødig indrømme, at »Vennen« kun skulde
være et Produkt af hendes Indbildningskraft, og støtter
sig til, at hun ikke er Herre over hans Kommen og
Gaaen. Men hun er klar over, at det er en dyb Trang
i hendes Væsen, der her finder Udtryk og Tilfredsstillelse.
Hun vilde ønske, at hun ikke var saa sikker paa,
at der kun forelaa en Fordobling af noget i hendes egen
Personlighed --- ligesom andre Elementer i hendes Natur
tenderede til Personifikation gennem seksuelle Hallucinationer.
Hun vilde, mener hun, tage Fænomenet
med større Alvor, naar hun ikke var saa vis paa dets
% --- p. 86 ---
\setcounter{footnote}{0}%
\opage{86}mystiske Udspring; nu ser hun altfor tydelig de Traade,
ved hvilke Dukken ledes (je vois trop les ficelles).

Den betydningsfuldeste Virkning af »Vennen«s Besøg
er ifølge Cécile den, at der kommer Sammenhæng i
hendes splittede Væsen. Skønt hans Besøg er en Illusion,
hjælper de hende dog fremad; hun mener at kunne konstatere
større Enhed i sin Natur (plus d'unification)\footnote[1]{Smlgn. hermed \emph{Hannah Joël}'s Beskrivelse af, hvorledes Overgangen for hendes Vedkommende fra sjælesyg til sjælesund Tilstand erfaredes som en Kamp for Enhed og Sammenhæng. (Se mine \textit{Mindre Arbejder}. III. p. 81.)}.
Denne Enhed kunde hos hende arbejde sig frem med
større Lethed, end hvis hun havde lidt af egentlig Dobbeltbevidsthed.
Det var ikke nødvendigt (som hos Morton
Prince's Patient B A C) at fremelske en ny Personlighed,
der kunde fortrænge begge de to stridende Jeger. Det
Jeg, der hos hende rørte sig i begge de modsatte Tilstande,
blev ved »Vennen«s Hjælp klaret og styrket.

Det er sandsynligvis en Reaktion mod den dæmoniske
Side i hendes Natur, der arbejder sig frem i
»Vennen«. Det er hendes Idealer fra forskellige Tider
og paa forskellige Omraader, der her danner et ejendommeligt
Kompleks. Cécile ser, trods sin stærke Betagethed
og trods den Fred, hun skylder Fænomenet,
rent psykologisk paa det. »Vennen« identificeres hverken
med Jesus eller med nogen anden bibelsk Personlighed,
heller ikke med hendes afdøde Fader eller med hendes
psykologiske Skriftefader. Hun mener ikke, som en Spiritist
vilde, at have Besøg af en Afdøds Aand. Hun frygter
ikke, at Fænomenet ved en rent psykologisk Tydning
skulde miste sin Værdi for hende.

\sectionhead{2. \textit{En moderne Mystikers Ekstaser.}}

Efter at Cécile et halvt Aar igennem havde haft sit
aandelige Tilhold hos »Vennen«, blev der hende en ny
% --- p. 87 ---
\setcounter{footnote}{0}%
\opage{87}Art Oplevelse tildel gennem ekstatiske Tilstande. Medens
hun ved »Vennen«s Aabenbarelser havde haft et sanseligt
Grundlag i Nærværelseshallucinationen, faldt i Ekstasen
ethvert saadant Grundlag bort. Der var intet at
beskrive. I hendes Beretninger betegnes derfor de Øjeblikke,
i hvilke Ekstasen indtræder, kun ved Tankestreger.
Medens Teresa's højeste Tilstande (se ovenfor
p. 64) var Nærværelseshallucinationer, har Cécile kendt
Tilstande, som var endnu mere ubestemte. Teresa kunde
straks genkende, hvad der var nærværende for hende,
og til den blotte Nærværelsesvision (den »intellektuelle
Vision«) sluttede sig snart andre Visioner (»imaginative«),
der gjorde udførligere Beskrivelse mulig. Hvad der stod
saa gaadefuldt for Teresa, var blot Sjælens egen Tilstand
under Visionerne; det er for denne Tilstand, at hun
leder efter Symboler og Allegorier.

Ligesom for Teresa Nærværelsesvisionen kunde gaa
kontinuerlig over til imaginativ Vision, saaledes hævder
Cécile, at der --- saa forskellige end Oplevelserne med
»Vennen« syntes at være fra de ekstatiske Oplevelser ---
var en kontinuerlig Overgang mellem dem. Overgangen
foregik gennem Tilbagetrængen af den Tendens til Personifikation,
som ytrede sig i »Vennen«; den satte en
Grænse, som der var Trang til at overskride. Med »Vennen«
kunde hun endnu tale; men hun mærkede en
dybere Trang, en Trang til noget, der kunde faa alt til
at tie i hende. Deraf udsprang, hvad hun kaldte den
guddommelige Oplevelse (l'expérienec divine).

Gjorde Forskellen mellem Oplevelse og Tydning sig
allerede gældende hos Cécile ved »Vennen«s Aabenbarelser,
maatte den --- for hendes intellektuelle Redelighed
og psykologiske Blik --- gøre sig endnu mere gældende
overfor de ekstatiske Oplevelser, der slet ingen egentlige
Erindringer efterlod.

% --- p. 88 ---
\setcounter{footnote}{0}%
\opage{88}I de Beretninger, Cécile sendte sin psykologiske Skriftefader,
optræder Ekstaserne som Lakuner. Kun de foregaaende
og de efterfølgende Tilstande kunde beskrives.
Forud kan der gaa Længsel og Forventning med uvilkaarlig
Koncentration, Hjertebanken, Muskelsøvn og
Kulde i Hænder og Fødder. Naar den ekstatiske Tilstand
er ved at bryde frem, kan der rejse sig en vis Modstand
mod den; Sjælen klamrer sig ligesom om sit Jeg og vil
ikke hengive sig, ikke blive helt passiv, --- alt imedens
den dog længes efter at blive greben af en mægtig,
fremvæltende Bølge. Hvad der i selve Ekstasen oplevedes,
stod siden for hende som en Oplevelse af en Inderlighed,
en Dybde, en Uendelighed og en Simpelhed, der stillede
den i Modsætning til alle andre Oplevelser, ogsaa til
den Halvsonnambulisme, i hvilken »Vennen« optraadte,
og som nu forekom hende kun at have bevæget sig i
mere overfladiske Lag af hendes Personlighed. Ekstasen
kom hvergang som noget Oprindeligt, Elementært; først
senere kunde hun genkende den som Ekstase, som noget,
hun ikke oplevede for første Gang.

Først efter Ekstasens Ophør var Overvejelse, Analyse
og Dom mulige. Og spørges der, hvorledes det var muligt
at tænke over den, naar den var forsvunden og ikke
kunde genkaldes, svares der, at der efter Opvaagningen
gik en Flodbølge af Undren og glad Begejstring gennem
Sjælen. »Jeg følte«, siger hun engang, »at jeg havde været
gennemtrængt af det Guddommelige, --- at jeg havde
været ude over Sansernes og Eftertankens Verden, ude
over al Tid, og at jeg nu igen var mig selv, uvidende
om, hvad der var foregaaet, men bevarende en Art Udstraaling
(rayonnement), der gjorde mig Livet lettere.«
Det er, hvad jeg har kaldet Følelsens Ekspansion, der
her gør sig gældende. Den gør sig uafbrudt gældende i
højere eller lavere Grader, idet vor Sindstilstand ikke
% --- p. 89 ---
\setcounter{footnote}{0}%
\opage{89}blot bestemmes ved, hvad vi oplever i Øjeblikket, men
ogsaa ved, hvad vi mere eller mindre umiddelbart forud
har oplevet, selv om vi ikke mere tænker paa det eller
kan tænke paa det. Efter Ekstaser og Visioner, der ikke
kan genkaldes, er den den eneste paaviselige Eftervirkning.
Det er, hvad Augustinus (se ovenfor p. 47) kaldte amans
memoria, det eneste, der blev tilbage efter den store
Vision. I sidste Sang af »Paradiset« siger Dante, efterat
det sidste store Syn er endt, at medens Synets Minder
er svundne, vedvarer »Virkningen den søde«. Hvad der
gælder Visioner, gælder endnu mere Ekstasen i snevrere
Betydning.

Men Cécile kalder Oplevelsen »guddommelig«, og
expérience divine bliver i den følgende Tid et teknisk
Udtryk for hende. Det er karakteristisk for hendes Skarpsindighed,
saavel som for hendes Redelighed, at hun
udtrykkelig søger at blive klar over, hvorfor hun kalder
Oplevelsen guddommelig. Det er en Værdibestemmelse,
men det er dog et Prædikat. Og med hvilken
Ret bruges dette Prædikat her? Hun gør straks den Bemærkning,
at hun havde tænkt sig det Guddommelige
anderledes. Hun var vant til at forestille sig det i personlig
Form, altsaa som en Genstand. Naar hun nu
brugte Ordet guddommelig, var det, fordi det, der havde
gjort sig gældende i hende under den ekstatiske Oplevelse,
stod for hende som en Norm, som det, der skal
være, og som man ubetinget maa anerkende. Eller, som
hun nærmere forklarer det: »Denne Oplevelse virker paa
mig som en Kraft, der driver mig frem, og som tillige
holder mig oppe paa et Niveau af Tanker og Forudsætninger,
der er højere end det, jeg forud indtog. Jeg mærker
dens Indflydelse selv paa Omraader, der ikke synes
mig at have noget at gøre med den Region, hvori Oplevelsen
udfolder sig. Derfor er det, tror jeg, at den for
% --- p. 90 ---
\setcounter{footnote}{0}%
\opage{90}for mig er guddommelig, d. v. s. gennemtrængende hele
mit Væsen og overtydende mig om en Realitet, der staar
over den, i hvilken mit daglige Liv foregaar.« Undertiden
kalder hun ligefrem det højere Liv, der omslutter hende
i de ekstatiske Øjeblikke, for Gud: cette vie-là, c'est
dieu! --- Faust's Trosbekendelse, der jo ogsaa fremtræder
som Udslag af en entusiastisk Tilstand, passer godt til
hendes Skildring, skønt hun ikke henviser til den. ---
Ogsaa andre Udtryk fra hendes vante religiøse Forestillingers
Kreds finder hun Anvendelse for, f. Eks. Naade.
(Ce n'est que maintenant que je comprends le sens du
mot »Grâce«). ---

Det er karakteristisk, at det første og vigtigste Prædikat,
der melder sig for den moderne Mystiker, er en
Værdibestemmelse, Udtryk for en Følelsesvurdering. Mystikeren
beskriver her noget, som vi alle kan kende i
andre Former og Grader, og som Livet ikke kan undvære.
Vore Oplevelser staar i personlig Værdi ikke paa
lige Trin. Vi skelner mellem vore bedste Øjeblikke, Højdepunkterne
i vort Liv, og det jævne Forløb, der optager
største Delen af Livet, afset fra de Øjeblikke, vi kalder
tørre, lave, maaske bestialske. Og vi betragter hine første
som »Norm«, d. v. s. som den Realitet, vi vil give saa
stor en Plads som muligt i vort Liv. Vi finder i de
jævne og lave Øjeblikke --- saa snart vi kan hæve os
over dem --- en Trøst i, at vi have haft Oplevelser af den
høje Art. Saaledes staar de ekstatiske Øjeblikke for Mystikerne.
Jeg ser her bort fra de psykologiske og fysiologiske
Betingelser for de høje Øjeblikke, som ogsaa vi med
Rette kalder det en Naade at have, fordi de kommer
uvilkaarlig, selv om vi kan stræbe at gøre os værdige
til at have dem. Disse Betingelser synes mere gaadefulde
ved Ekstase end ved almindelige Menneskers høje
Øjeblikke; og dog er det et Spørgsmaal, om vi vide mere
% --- p. 91 ---
\setcounter{footnote}{0}%
\opage{91}om disse end om hine. --- Et Spørgsmaal for sig er det,
om den Vurdering, der ligger i Udtryk som »guddommelig«
eller »høj« kan begrundes rationelt. Dette Spørgsmaal
er der her ingen Grund til at gaa ind paa\footnote[1]{Se herom \textit{Den menneskelige Tanke} § 145—146. --- \textit{Totalitet som Kategori}. (Vid. Selsk. Skrifter 1917). Kap. 5.}. Selve
Vurderingen, der faar Udtryk i disse Prædikater, er et
Faktum, og det er, som sagt, betegnende, at den moderne
Mystiker klart ser, at et Værdiprædikat er det første
Prædikat, der kan bruges om de mærkelige Oplevelser. ---

I Teresa's Forsøg paa at udtrykke sine ekstatiske Oplevelser
findes der egentlig intet, der svarer til Cécile's
Forklaring af, hvorfor hun kalder sin Oplevelse guddommelig.
Hun kender ikke til rene Værdibestemmelser.
Derimod er der noget beslægtet mellem de to kvindelige
Mystikere i deres Benyttelse af Analogier til at udtrykke,
hvad de har oplevet. Ligesom hos Teresa opdukker
uvilkaarlig Sammenligning ogsaa hos Cécile.
Ogsaa hun trænger til Symboler. Der er et Billede, som
atter og atter trænger sig frem hos Cécile, og som skyldes
en Erindring fra et Ophold ved den skotske Kyst i
hendes unge Dage. Hun havde der en Aften set en sort
Planke ligge i Strandkanten. Bølgerne rokkede stadig
ved den, indtil den tilsidst blev reven med og ført ud
paa Havet, hvor den blev gennemstraalet af Aftensolens
Glans. Dette Erindringsbillede kom uvilkaarlig frem hos
hende straks efter den første Ekstase: »Det forekommer
mig, at ogsaa jeg er bleven greben af Strømmen og ført
ud i Solen, af hviken jeg nu er helt gennemstraalet«.
Senere kommer dette Billede ofte frem efter hendes Ekstaser,
og det staar for hende, som det eneste, der i nogen
Grad svarer til, hvad hun mente at have oplevet.
En enkelt Gang ændres Billedet saaledes, at en lille
% --- p. 92 ---
\setcounter{footnote}{0}%
\opage{92}Baad optræder istedetfor Planken, og det optræder ogsaa
før en Ekstases Indtræden, hvor da Baadens eller Plankens
Modstand mod Bølgernes Forsøg paa at rive den
med symboliserer den Modstand, der uvilkaarlig gjordes
mod at overvældes af den ekstatiske Betagelse.

Hos Teresa svarer til dette Symbol Billedet om Jernet,
der i sagte Ild ligger som en sort Masse, men i et Baal
bliver glødende (ovenfor p. 64), ligesom Planken gennemstraales
af Aftensolen. I den uvilkaarlige Brug af et
beslægtet Billede lægger Ligheden mellem de to Mystikeres
Oplevelser sig for Dagen trods al den Forskel, der
ellers fremtræder mellem dem\footnote[1]{Et Eksempel paa, hvorledes et og samme Billede paa højst forskellige Kulturtrin kan bruges som Udtryk for samme religiøse Type, har jeg anført i min Religionsfilosofi § 36 (Fisken paa Landjorden).}. ---

En anden Kreds af Symboler trænger sig frem hos
Cécile, Billeder, der dog ikke blot har Betydning som
Analogier, men efter hendes Forvisning peger hen paa
en virkelig Sammenhæng. Ekstasen belyses ved Sammenligning
med den seksuelle Hengivelse. Hvad der opleves
under denne kan ligesaa lidt som den ekstatiske
Oplevelse beskrives. Cécile finder Slægtskabet mellem
Erotik og Religion at bestaa i, at paa begge Omraader
tilfredsstilles Trangen til inderlig og umiddelbar Forbindelse
og fuldstændig Gennemtrængen. Begge er Udtryk
for Livets stærke Trang, og at Livets store Kræfter
udtrykkes paa samme Sprog, hvad enten der er Tale
om guddommeligt eller menneskeligt Liv, har for hende
intet forunderligt. Mystikerne bruger atter og atter den
menneskelige Elskovs Sprog for at udtrykke deres høje
Oplevelser (Smlg. ovenfor p. 66f).

Det var dog for Cécile mere end billedlig Tale og
psykologisk Analogi, her forelaa. Hun maatte jo selv
den ene Gang efter den anden erfare, hvor dybt Køns%
% --- p. 93 ---
\setcounter{footnote}{0}%
\opage{93}instinktet er rodfæstet i vor Natur, og hendes psykologiske
Skriftefader havde jo ogsaa søgt at styrke og trøste
hende over hendes Kampe med Mareridt ved Henvisning
til de uundgaaelige Naturdrifter, vi tager med os,
selv om vi stiger nok saa højt i aandelig Kultur. Deraf
følger, som Cécile klart har set, at jo dybere en Oplevelse,
endog af helt anden Art, griber ind i vort Sjæleliv,
des mere vil ogsaa den seksuelle Side af vor Natur
blive sat i Bevægelse. Ofte maatte hun erfare, at ekstatiske
Oplevelser æggede de seksuelle Drifter frem. Hvad
der for Eftertanken kom til at staa i den grelleste Kontrast,
havde saaledes en real, underjordisk Sammenhæng.
Selve Modsætningen forudsætter Kræfter, der virker i
samme Plan. Efter paany at have erfaret den grelle
Modsætning mellem det Høje og det Lave siger Cécile:
»J'ai peine à réaliser que je suis la même qu'il y a 24
heures, alors que mon être tout entier vibrait encore de
t'attouchement du divin\dots{} Pourtant c'est bon d'arriver
à voir quelques-unes des ficelles de notre pantin,
comme cette connexion, qui semble si étrange, entre
l'émotion religieuse très pure, très haute, et l'excitation
sexuelle«.

Forholdet mellen Seksualitet og Religion er nu ikke
blot det, at den af religiøse Motiver vakte Sindsbevægelse
formedelst organisk Diffusion forplanter sig til det seksuelle
Omraade, men ogsaa det omvendte, at seksuel Trang
kan undergaa en Metamorfose, idet den Energi, den har
udløst, overføres til psykiske Processer af helt anden
Art, Muligheden af en saadan Metamorfose ligger i det
Hengivelsesmoment, som er tilstede, saasnart Kønstrangen
og Kønsforholdet har mere end en øjeblikkelig og
isoleret Karakter. Vendepunktet er der, naar den Genstand,
der skal tilfredsstille Trangen, staar som mere
end blot Middel, saa at der kan danne sig et Billede af
% --- p. 94 ---
\setcounter{footnote}{0}%
\opage{94}den, der peger ud over det enkelte Øjeblik og den isolerede
Nydelse. Drift og Interesse begynder da at være
knyttet til noget, der har Værdi ud over den momentane
Tilfredsstillelse, og som maaske endog kan kræve
Opofrelse af denne. Det er en saadan Metamorfose, eller
Række af Metamorfoser, Platon lader Sokrates skildre
i »Symposion«. Ad denne Vej kan Erotik gaa over til
Religion, især naar den hele Idealiseringsproces støttes
ved den ophøjede Karakter af religiøse Forestillinger og
Billeder, der træder Individet imøde gennem Overlevering
og Kultus. Erotik kan ikke fra først af have frembragt
Religion, men den Trang, der rører sig i hin, kan ---
især gennem Skuffelse --- omformes til religiøs Trang,
idet Hengivelsen nu rettes mod noget mere Omfattende,
maaske et Uendeligt. Energien kan bevares gennem Forskydningerne.

For Cécile's Vedkommende var det »Trangen til at
være elsket og blive Moder«, der krævede Tilfredsstillelse,
og hun mente selv, at hvis denne Trang var bleven tilfredsstillet,
havde hun fundet Ro indenfor den Kreds,
som da havde dannet sig omkring hende. Hun kunde
oprøres over, at saadan Lykke var nægtet hende. Men
saa havde netop den Kamp med utilfredsstillet Trang,
hun førtes ind i, givet hendes Liv større Værdi. »Vilde
Menneskesjælen,« spørger hun, »uden Lidelse mærke
den Tørst, som kan pine, men som betinger dens kongelige
Værdighed? Hvad véd de i Grunden, de, som ikke
har lidt?« Det er gennem Lidelsen, at Trangen idealiseres,
selv om der bliver Slakker tilbage. Og her har
Ekstasen, »den guddommelige Oplevelse«, bragt hende
positiv Hjælp: »Hvergang den guddommelige Oplevelse
har udfoldet sig hos mig, har det glødende Jern berørt
min Sjæl og har renset den.« Allerede det, at hun kunde
opleve begærfrie Tilstande, som dog betog og begejstrede
% --- p. 95 ---
\setcounter{footnote}{0}%
\opage{95}hende, maatte være af stor Virkning. Og hvad hun oplevede
i Ekstaserne, var for hende en Fortsættelse af,
hvad hun havde oplevet i de halvsonnambule Tilstande,
i hvilke hun mærkede »Vennen«s Nærværelse. I den
første af hendes Ekstaser gik de to Oplevelser ligefrem
over i hinanden. Ogsaa i de halvsonnambule Tilstande
havde hun erfaret, at der var andre Elementer i hendes
Natur end den ustyrlige Trang.

Under alt dette har ogsaa Erindringsbilledet af hendes
Fader øvet stor Indflydelse. Han har været en energisk,
dybsindig og selvstændig Mand, der stod for hende som
Norm og Forbillede, saa at det var hende en Trang at
føle sig som »sin Faders Datter« ved Selvbeherskelse
indadtil og Selvstændighed udadtil. Flournoy har dog
maaske lagt for stor Vægt paa Virkningen af dette »imago
paternelle«, tildels fordi han er for stærkt paavirket af
den Freud'ske Skoles Teorier. Jeg formoder desuden,
at Flournoy ikke har Ret, naar han mener, at det Ideal
af en Mand, Cécile har set i sin Fader, var Aarsag til,
at Cécile ikke giftede sig. Der var Friere nok; men Ingen
skulde have kunnet staa Maal med det faderlige Ideal.
Det er dog lidet sandsynligt, at Cécile med hendes Trang
til »at være elsket af en Mand og have et Barn i sine
Arme« ikke skulde have taget denne Side af Livet op,
hvis den Rette, d. v. s. den, hun kunde ønske at være
elsket af, virkelig var kommen, og der er ikke Spor
deraf i hendes Optegnelser. Men den Eros, der kunde
have fyldt hende, kom ikke. ---

Den hele Side af Livet, der --- ikke blot som naturligt
Symbol --- men dels som Modstand, der maatte
overvindes, dels som Kontrastvirkning, der kunde ydmyge,
dels som Motiv, der kunde metamorfoseres, spiller saa
stor en Rolle hos Cécile, synes saa godt som ganske
at mangle hos Teresa. Det kan maaske forklares ved,
% --- p. 96 ---
\setcounter{footnote}{0}%
\opage{96}at hun saa tidlig kom ind under Klosterlivets Tugt, saa
at de naturlige Tendenser ikke fik Tid til at udfolde
sig. De har ført en mere underjordisk Tilværelse, og
naar de rørte sig stærkere, har Bekendelse i Skriftestolen
dannet tilstrækkelig Modvægt. Hverken i Selvbiografien
eller i Skrifterne mærker man dem direkte.
Men i Symboliken kan de spores, ti erotiske Billeder
forudsætter dog forudgaaende Erfaring eller Bevidsthed
om, hvad Billedet betyder. Kun i en enkelt levende
Beskrivelse fremtræder den ovenomtalte organiske Diffusion
tydelig. Den Mulighed staar dog stedse tilbage, at
Teresa, trods hendes store Evne til Iagttagelse og Analyse,
ikke har kunnet maale sig med Cécile, hvad denne
Evne angaar. Cécile noterer saaledes ofte Bebudelser og
Varsler om Ekstasens Indtræden, medens denne hos
Teresa mest kommer uforberedt og uformodet.

Den vigtigste Forskel mellem den spanske Nonne fra
det sekstende Aarhundrede og den moderne Mystiker
fra det tyvende Aarhundrede er dog, at hin meget hurtig,
og uden at mærke Overgangen, gaar over fra Oplevelse
til teologisk Tydning, medens Cécile foreløbig bliver
staaende ved uvilkaarlig Symbolik.

\sectionhead{3. \textit{Ekstase og Religion.}}

Man kunde undre sig over, at Cécile, der hyldede en
udpræget konfessionel Religiøsitet, og som havde paataget
sig den religiøse Ledelse af en Opdragelsesanstalt,
ikke ligestraks i de stærkt betagende Oplevelser, der
faldt i hendes Lod, saa Kundgørelser fra Bibelens og
Kirkens Gud eller Aabenbarelser af Kristus. Som vi har
set, rørte der sig først og fremmest hos hende en Trang
til simpel Værdsættelse af det Oplevede i dets Betydning
for hendes indre Liv, og dernæst søgte hun uvilkaarlig
efter Billeder til at anskueliggøre denne Betydning. Et
% --- p. 97 ---
\setcounter{footnote}{0}%
\opage{97}Billede som det med Planken i det solbelyste bølgende
Hav er et godt Billede paa den Livspoesi, der er et
umiddelbart Udtryk for, hvad der er erfaret om Forholdet
mellem det Mørke og det Lyse, det Tyngende og
det Løftende indenfor Sjælelivet. Men enhver dogmatisk
Anvendelse ligger hende fjernt. Medens Teresa straks
genkender Jesus i den Skikkelse, der fremtræder for
hende, først som blot nærværende, senere som Genstand
for Hørelse og Syn, erklærer Cécile udtrykkelig, at hun
aldrig har opfattet »Vennen« som Et med Jesus, og at
hun heller ikke, naar hun bad, havde Indtrykket af, at
det var Vennen, hun kom i Forhold til. De bibelske
Forestillinger bragtes overhovedet slet ikke i Forbindelse
med det, der skaffede hende den bedste Trøst og den
dybeste Fred. Hun tager Oplevelsen som den kommer.
Paa samme Maade som overfor de bibelske Forestillinger
stillede hun sig overfor de spiritistiske Tydninger.
Det vilde jo, siger hun, have været en ganske simpel
Forklaring at antage Aander som Ophav, og det vilde
have været hende en stor Glæde, om hun havde kunnet
tro, at det var hendes Faders Aand, der besøgte
hende. Men hun havde ingen Anledning dertil.

Forskellen mellem Teresa og Cécile i denne Henseende
hænger sammen med, at hos Cécile spiller det
Uvilkaarlige en langt større Rolle end hos Teresa, og
Ekstasen træder fra først af i bestemt Modsætning til
hendes bevidste Religiøsitet; den Betydning, Ekstasen faar
for hende, er maaske, som vi vil komme til at se, betinget
ved, at hendes bevidste Religiøsitet i Grunden ikke
tilfredstillede hende. Hos Teresa var Ekstasen derimod
forberedt gennem Bønnen. »Døren til Sjæleslottet« er
Bønnen, baade den i Ord og den i det Indre, ledsaget
af Sindets Henvendelse (l'oraison tant vocale tant mentale,
accompagnée d'attention). Bønnens forskellige For%
% --- p. 98 ---
\setcounter{footnote}{0}%
\opage{98}mer er ligesaa mange Trin op til Ekstasen. Efter den
vilkaarlige Henvendelse gennem Bøn kommer saa den
højeste Tilstand uvilkaarlig; »førend vi tænker paa at
hæve vor Tanke til Gud, er vor Sjæls Kræfter allerede
inde i Slottet, uden at vide, hvorledes de er komne derind.« Det
Vilkaarlige hører altsaa kun hjemme paa de
lavere Trin; senere gælder det netop om ikke at øve Vold
mod sig selv; men vilkaarlig Forbereden spiller dog en
Rolle, der ikke har noget Tilsvarende hos Cécile.

Ligesom Teresa --- til Trods for nogle af sine Skriftefædre
--- i sine mystiske Oplevelser fandt Bekræftelse
paa Tanken om Guds Iboen i alle Ting, saa at det ikke
blot er gennem Naaden, men paa endnu inderligere
Maade, han virker i Menneskesjælen, saaledes er ogsaa
Cécile overbevidst om, at »den guddommelige Oplevelse«
ikke forudsætter et Brud paa Naturlovene, men er Udslag
af en naturlig Udfoldelse af det indre Liv. Men
for Cécile fik --- i hvert Tilfælde i en Periode --- denne
Overbevisning en principiel Betydning, som Teresa ikke
kunde tillægge den. Sagen var, at Cécile i den Periode,
i hvilken »den guddommelige Oplevelse« var af afgørende
Betydning for hende, ikke kunde anvende de bibelske
og kirkelige Forestillinger, saa gerne hun end
vilde.

»Jeg ønskede,« siger hun, »at jeg kunde fremkalde
min gamle kristelige Forvisning igen i den oprindelige
Form, tage min Tilflugt til Tanken om mine Synders
Forsoning og om en udenfor mig selv tilvejebragt Forløsning;
men jeg fandt ikke længere Styrke eller Trøst
derved; det var, som om det ikke længere naaede ind
til mig\dots{} Jeg har forhen kunnet »føle Gud« i mit Liv
eller i min Bøn, men det var en Gud, jeg var fortrolig
med, hvis Historie jeg saa at sige kendte, og hvis Indvirken
mere eller mindre var, som jeg ventede den. Det
% --- p. 99 ---
\setcounter{footnote}{0}%
\opage{99}var en Gud, der var traadt frem i Historien, aabenbaret
gennem Jesu Lære og Liv. Men det, som jeg nu er
kommen i Berøring med, overgaar ethvert Udtryk og
al Historie, --- og dog er det ikke blot en Abstraktion.
\dots{} Nu forekommer mig enhver formuleret Bøn, eller
i hvert Tilfælde enhver Bøn, som forud formes for at
forrettes i et bestemt Øjeblik, at være Helligbrøde. Jeg
gaar sandsynligvis for vidt her, men det er nu min
Overbevisning, mod hviken jeg ikke kan handle uden
at have Følelsen af at krænke min Samvittighed\dots{} Til
syvende og sidst maa det dog være af liden Betydning
for mig, hvor den guddommelige Oplevelse har sin Plads
blandt de Realiteter, der er indregistrerede under en
bestemt Etikette. Den er til, og dens Værdi for mit aandelige
Liv er ligesaa utvivlsom som dette aandelige Liv
selv\dots{} For et Menneske, der som jeg er bleven opfostret
med Fromhed og Teologi, er det vanskeligt at
antage en helt ny Form for Aabenbaring; det sker ikke
uden Vaklen. Men man maa lære ikke at være bange
for Sandheden, hvilken denne saa maatte være!« ---

Der laa i »den guddommelige Oplevelse« for Cécile
Spiren til et nyt Gudsbegreb. »Først nu,« siger hun,
»gaar det ret op for mig, hvor snevert, dogmatisk og
menneskeligt mit Begreb om det Guddommelige var.
Ledet af en bydende Trang til overalt at genfinde Karakteren
af noget personligt, var jeg aabenbart kommen
til at smedde mig en Gud, der helt indeholdes i den
moralske Sanktion og helt og holdent var aabenbaret i
den Fader, Evangeliernes Kristus forkyndte\dots{} Jeg føler
mig paa Grænsen af et stort, forjættet, men endnu ikke
udforsket Land\dots{} I Højhed og i Umiddelbarhed overgaar
den guddommelige Oplevelse alt, hvad jeg hidtil
har kunnet forestille mig\dots{} Jeg føler en saadan aandelig
Fornyelse efter de korte Øjeblikke af »guddommelig
% --- p. 100 ---
\setcounter{footnote}{0}%
\opage{100}Berøring«, at det forekommer mig barnagtigt og meningsløst
at blive ved at søge Opbyggelse ved at høre
Prædikener og deltage i Møder\dots{} Hvad betyder overfor
denne direkte Aabenbaring til en Menneskejæl hele
den kristelige Prædiken og den uhyre Bygning af Dogmer?\dots{}
Det forekommer mig nu, at jeg staar ved
Bredden af et uhyre Ocean, og at Gud gaar forbi i de
mægtige Bølger og i den nedgaaende Sols Straaler.
Men det er ikke den smaatskaarne og modellerede
Gud. Det er det højeste Virkelige, det sidste Endemaal
for alt, den levende og triumferende Kraft, i hvilken
jeg tilsidst finder det Grundvæsen, der er hævet over
Tid og Rum og over Godt og Ondt!« ---

Hvis man vilde sige, at Cécile her glider over fra Beskrivelse
til Tydning og endog lægger en hel panteistisk
Metafysik ind i sine Oplevelser, maa dog dertil svares,
at hun egentlig kun har udtalt, at hvad hun har oplevet,
ikke tillader hende at lade nogen Modsætning gælde
indenfor det, der i hendes højeste, mest styrkende Øjeblikke
er fremtraadt for hende. Hun hævder en docta
ignorantia, der fører hende til at se med Humor paa
Dogmer og teologiske Lærebygninger. »Hvis dette virkelig
er den Maade, paa hvilken Gud aabenbarer sig for
en Menneskesjæl, hvorledes har Teologerne saa kunnet
skrive saa mange tykke Bøger derom?« »For Fremtiden
paalægger Tanken om Gud mig fremfor alt Tavshed.«
--- Hun har gjort den Erfaring, at »vor Trang, især vor
aandelige Trang, altid overgaar det Udtryk, vi kan give
den.«

Af hvad der virkede i »den guddommelige Oplevelse«,
kunde Disharmonier, Ulykker og Forbrydelser, saaledes
som den sædvanlige Erfaring viser os dem, ikke udledes.
Gud kan ikke være Aarsag til det Onde. Hun fandt
Trøst i denne Overbevisning overfor de Disharmonier,
% --- p. 101 ---
\setcounter{footnote}{0}%
\opage{101}hun maatte erfare i sit eget Indre. Hun havde Øjeblikke,
hvor hun oprørtes over, at hun var skabt med saa stærke
Instinkter, naar de ikke skulde finde Tilfredsstillelse.
Men især da Verdenskrigen udbrød, blev den Tanke
hende en Trøst, at »Gud ikke er i denne Verden«.\footnote[1]{Sandsynligvis har hun fundet denne Sætning i et posthumt Skrift af Mme Flournoy, hvor den hævdes med stor Inderlighed og Eftertryk. Se om dette Skrift \textit{Den store Humor} p. 91.}
»Si je me sentais encore tenue de croire qu'il joue un
rôle dans ce drame, j'arriverais au blasphème.« ---

Den Vending i Céciles religiøse Udvikling, som indtraadte
med »den guddommelige Oplevelse«, var betinget
ved Modsætningen til den forudgaaende Periode i hendes
Liv, i hvilken hun mere og mere var kommen ind paa
en Religiøsitet, der havde sit Udtryk i stive dogmatiske
Rubriker. »Oplevelsen« udsprang da, mener hun, ud af
en dyb Trang, som denne Dogmatisme ikke kunde tilfredsstille,
og af den aandelige Nødstilstand, som hendes
konfessionelle Tro ikke havde kunnet hjælpe hende ud
over. Aldrig havde Følelseslivet rørt sig saa stærkt og
inderlig hos hende. Den Inderliggørelse, hendes religiøse
Liv derved vandt, udstraalede over hele hendes Livsførelse.
Skønt hun ikke kunde meddele selve Oplevelsen
til andre, gjorde den hende det dog muligt at virke med
større Kraft og Kærlighed. Da det, hun oplevede, var
saa højt forskelligt fra, hvad andre oplevede, følte hun
sig mere selvstændig end før og fik mere Tro til sig
selv end før. Nu var hun endelig bleven en Person
(quelqu'un), og ikke blot et Kompleks af andres Tanker.
Hun forstod sig selv bedre, og forstod derfor ogsaa
andre bedre. Med større Taalmodighed og inderligere
Sympati stod hun nu overfor andres Lidelser, selv om
de var af anden Art end hendes egne. Hun havde hidtil
ikke haft tilstrækkelig Respekt for sin egen aandelige
% --- p. 102 ---
\setcounter{footnote}{0}%
\opage{102}Udvikling, naar den ikke foregik i de vante dogmatiske
Rubriker. Hun saa, at hun ikke havde Ret til at mishandle
sin Sjæl for at faa den ind i et bestemt Skema.
Selv den mest autentiske Kristendom er --- ser hun nu
--- ikke et aandeligt Lægemiddel for alle. En Prædiken,
hun hører, forøger kun hendes Forvirring. Hun fik Lyst
til at blive Katolik for i Tavshed at kunne knæle i en
Kirke uden at forstyrres af, hvad Mennesker siger\footnote[1]{Det er vist ikke faa, hvem det gaar som Cécile, at den liturgiske Del af Kultus kan betage dem, medens Prædikenen kun forstyrrer. I en Optegnelse fra 1847 siger Kierkegaard: »Mig beskæftiger ubetinget Psalmesangen mest af hele Gudstjenesten.« William James fortalte mig, at han ikke hørte til nogen Konfession, men holdt af at høre Psalmesang; Prædikener kunde han ikke holde ud at høre. J. N. Madvig siger i en efterladt Optegnelse, at Orgelspil og Sang vækker dyb og levende Stemning, Længsel, Kærlighed og Mod, --- „men saa kommer Prædiken med dens saa ofte haltende og modsigende Dogmatik, dens Uklarhed og skæve Slutninger.« En afdød Ven af mig, en berømt Naturforsker, der var en stadig Kirkegænger, var meget optagen af Psalmesangen, men forholdt sig i Regelen overvejende kritisk overfor Prædikenen.}.

\sectionhead{4. \textit{Personlighedsprincip og Menneskekærlighed.}}

De mærkelige Tilstande, vor moderne Mystiker havde
oplevet, syntes at maatte føre til en afgørende Ændring
i hendes aandelige Liv, navnlig til en uddybet Forstaaelse
af sig selv og af andre. Hun følte sig nu som
»sin Faders Datter«. Faderen var en fremragende Pædagog,
der i religiøs Henseende var stærkt paavirket af
Vinet og i Tilslutning til ham hævdede den individuelle
Personligheds Ret overfor alle Konfessioner og Dogmer.
Hans Slagord til sine Børn var: Følg Eders Samvittighed,
og fordyb Eder med Alvor i Jesu Liv og Værk.
Cécile havde en Tid fornægtet dette Raad. Efter sin
Hjemkomst fra Udlandet, hvor hun havde følt sig me%
% --- p. 103 ---
\setcounter{footnote}{0}%
\opage{103}get omtumlet, var hun bleven stærkt paavirket af den
i det romanske Schweiz herskende Ortodoksi og Pietisme
og havde med Iver kastet sig ind i Arbejde for
denne Retning. Hun havde da ikke nok i sin Faders
Stade, men mente at hæve sig over dette ved hovedkulds
(à corps perdu) at kaste sig ind »i hele denne
noget forhastede og overfladiske, men berusende kristelige
Virksomhed«. Hendes Iver gjorde, at hun blev ansat
som Lærerinde og Forstanderinde for en i ortodox
Aand ledet Opdragelsesanstalt for unge Piger. Hun fandt
det ortodoxe Grundlag bekvemt for Udøvelse af aandelig
Paavirkning overfor sine Elever. Men nu følte hun
i stigende Grad Modsætningen mellem sit indre Liv og
de ydre Former og Formler. Ekstasen havde den Betydning
at styrke hende i Troen paa det indre Livs Ret.
Hun kunde nu trygt lade hele Dogmatiken ramle sammen.
Hendes ortodoxe Omgivelser mærkede vel, at hun
ikke mere som før gik op i Troens Former og Formler;
men man lod hende fortsætte sit Arbejde, da det
indre Livs forhøjede Energi bragte ny Inderlighed og
Kraft i det.

Men alt imedens denne aandelige Frigørelsesproces
stod paa, gjorde sig ogsaa en anden Strømning gældende
hos hende. Hun frygtede for, at det intensive indre Liv
og dets Erfaringer skulde isolere hende og efterhaanden
svække Interessen for at paavirke sine Elever. »I en vis
Forstand,« siger hun, »er min egen Sjæls Liv mig værdifuldere
end de Sjæle, der omgiver mig og behøver min
Kærlighed«. Den »guddommelige Oplevelse« var nu bleven
Midtpunkt og Sol i hendes Liv. Hun sagde sig selv, at
hvis hendes aandelige Kraft og Rigdom havde været
noget større, vilde hun kunne lade baade det indre og
det ydre Liv komme til deres fulde Ret. Men her følte
hun sin Fattigdom. Hun kunde ikke føle sig tilfreds%
% --- p. 104 ---
\setcounter{footnote}{0}%
\opage{104}stillet ved noget, hun maatte beholde for sig selv, og
som ikke kunde meddeles til andre. Hun føler stadig,
hvor meget hun skylder »den guddommelige Oplevelse«,
og hun vil ikke kaste Sten i den Kilde, af hvilken hun
har læsket sig. Hun har vundet ny Kraft, og hendes
religiøse Forestillinger er blevne frigjorte for Sneverhed,
for Dogmetro og Teologi. Men hun følte sin nye Rigdom,
som om hun havde faaet en 100 frcs. Seddel, hun
ikke kunde faa byttet, skønt det netop var Smaapenge,
hun nu havde Brug for overfor andre. ---

Denne Brydning er ikke det, der mindst interesserer
hos denne moderne Mystiker. Fra et religionspsykologisk
Standpunkt kan Ekstase og Mystik betragtes som
stærkt udprægede (i manges Øjne halvt patologiske) Former
for, hvad der i ethvert rigt indre Liv kan erfares
om høje Øjeblikke, lyse og varmende Stunder, der staar
i skarp Modsætning til det daglige Livs Spredthed og
Udvorteshed. Og den Vanskelighed, Cécile følte sig stillet
overfor, er derfor ikke noget for Mystikeren eget, men
kan fremkomme for Enhver, hos hvem Forskellen mellem
et indre og et ydre Liv overhovedet gør sig gældende.
I Grunden er fuldstændig direkte Meddelelse umulig.
Cécile har med Rette udtalt, at det indre Livs Fylde
langt kan overgaa de Meddelelsesmidler, som staar til
Raadighed. For dem, der kender dansk filosofisk Literatur,
behøves der ingen nærmere Paavisning af, at indirekte
Meddelelse er den eneste Bro, der kan bygges
mellem selvstændige og ejendommelige Personligheder.
I hvert Tilfælde er det at øve Vold imod Personligheden
at kræve samme ydre Formler anerkendte og samme
sjælelige Udvikling gennemløben af alle. Cécile klager
over religiøse Kredses ringe Respekt for Personligheden.
Man vil trumfe sine Dogmer igennem, hvorledes det saa
gaar med de enkelte Sjæle. Og hun erfarede selv, da
% --- p. 105 ---
\setcounter{footnote}{0}%
\opage{105}hun med stor Forsigtighed drøftede sit Problem om
Forholdet mellem Dogmer og indre Oplevelser med en
Præst (»Pastor Z.«), med hvilken Mistænksomhed man
ogsaa i protestantiske kirkelige Kredse, betragter Oplevelser
af den Art, der var faldet i hendes Lod. Pastor
Z. slog to Ting fast, for det første, at enhver Oplevelse,
der udfolder sig i Bevidsthedens Grænseregioner, er af
lavere, flygtig og illusorisk Art, --- for det andet, at enhver
Oplevelse, der ikke fører til Antagelse af en personlig
Gud, er sygelig og ikke kan lægges til Grund for
Livsførelse. I Modsætning hertil stillede hun --- i sit
Sind, ti hun vilde ikke udtale sig aabent for en Præst
--- den Erfaring, at intet havde grebet saa styrkende og
rensende ind i hendes Liv som »den guddommelige
Oplevelse«, skønt den ikke kunde bringes ind under de
vedtagne dogmatiske Former (clichés très connus et
catalogués).

Det var hendes egen Menneskekærlighed og Samvittighedsfuldhed,
som voldte hendes Problem. Hun fordrede,
at den Planke, hun saa ofte saa vugget i det solbelyste
Hav, skulde blive til en Redningsbaad ogsaa for
andre end hende selv. Og man kan ikke lade være at
spørge, om hun ikke i Virkeligheden var paa Vej til
at bygge en saadan Redningsbaad af Planken. Hun fastholder
stadig, at trods Usigeligheden af de indre Oplevelser,
havde de sat hende istand til »at udbrede lidt
mere aandelig Kraft og guddommelig Kærlighed omkring
sig«. Kan man i Grunden naa videre? Skal man nødvendigvis
slaa af paa sine inderligste Erfaringer og derved
nedværdige sig selv for at fyldestgøre Menneskekærlighedens
Krav? Vil ikke netop en fuld og gennemført
Anerkendelse af Personlighedsprincipet hos andre som
hos sig selv, føre til en sandere Menneskekærlighed end
% --- p. 106 ---
\setcounter{footnote}{0}%
\opage{106}den, der overfor andre taler det »katalogiserede« Sprog,
men et helt andet Sprog overfor sig selv?

Undertiden havde Cécile det Indtryk, at der kunde
komme det Øjeblik, da hun ikke havde Ret til at beholde
for sig selv det, som hun maatte kunne meddele,
nemlig, at det overhovedet var muligt »at komme i
Berøring med det Guddommelige uden nogetsomhelst
Mellemled af Spekulation, Dogme eller Trosform.« Men
hun turde ikke tage alle Konsekvenserne deraf. Hun
skimter Personlighedsprincipet, men hun lader det falde
igen. Trangen til Fællesskab, og til den Tryghed, Fællesskab
giver, var for stærk. »Vil vi faa den Tid at se,«
spørger hun, »da sand religiøs Frihed vil herske? Undertiden,
f. Eks. idag, ønsker jeg det, men til andre Tider
frygter jeg det og opdager, at jeg endnu er meget
traditionsbunden. Jeg er bange for at gaa for vidt. Jeg
slipper Duen ud gennem Arkens Vindue, men skynder
mig saa at lukke Vinduet igen\footnote[1]{En af de Kvinder i det 19. Aarhundrede, der har øvet en stor Gerning i Menneskekærlighedens Tjeneste, og som tillige var en selvstændig Tænker, \textit{Florence Nigthingale}, sendte en Due ud i Verden, og hun lukkede ikke Lugen for den. Tværtimod hævdede hun med religiøs Begejstring den Overbevisning, at det var ethvert Menneskes Opgave frit at bringe det bedst mulige ud af sin Personlighed. Hun kunde ikke opfatte Gud som et individuelt Væsen, men som et universelt Væsen, en Lov. Af de kirkelige Skikke benyttede hun kun, hvad hun havde Brug for. Hun var »meget sikker paa Gud«, men meget sparsom i at definere hans Væsen. (\emph{Edwad Cooke}: \textit{Florence Nigthingale}. Dansk Overs. 1917. I. p. 349—441; II. p. 188). --- En følgende Afhandling i denne Bog vil føre os til at tale igen om denne mærkelige Kvinde. (Se »Træk af Bønnens Historie«).}. Ganske vist, der er
snever Plads i Arken; man faar stive Lemmer, og man
aander en overhedet Luft uden tilstrækkelig Ilt; --- men
saa er man alligevel i Arken, sikker paa at være i Sik%
% --- p. 107 ---
\setcounter{footnote}{0}%
\opage{107}kerhed, medens udenfor Bølgerne og de stærke Vinde
raade.«

Cécile var klar over, at hvis hun talte ud helt om
sine indre Oplevelser, vilde hun miste sin halvt gejstlige
Stilling. Det var dog sikkert ikke af Frygt for et saadant
Brud, hun ikke gjorde det afgørende Skridt. Foruden
den Trang til at være noget for andre, som hun mente,
bedst fyldestgjordes paa Grundlag af de konfessionelle
Forestillinger, rørte der sig mere og mere ogsaa en
Trang til bestemte Billeder og Personifikationer, en
Trang, der saa meget des lettere fik Overmagten, som
der kom den Tid, da Ekstaserne hørte op. Dette fører
os over til det sidste Trin i den Del af hendes Udvikling,
Dokumenterne giver os Indblik i.

\sectionhead{5. \textit{Céciles »syvende Bolig«.}}

Det har ikke blot været Menneskekærlighed, Trang
til at være noget for andre, som paany førte Cécile tilbage
til kirkelige Forestillinger. Der rørte sig mere og
mere hos hende en Utilfredsstillethed ved det Uudsigelige,
ved det, der ikke kan udtrykkes i noget Begreb,
eller rettere i noget Billede. Hendes Fantasi har, for at
bruge Schleiermacher's Udtryk i et af hans Breve, været
»bestemt ved Gysen for det ubestemt Tænktes Dunkelhed,
ikke ved Gysen for den Modsigelse, der fremkommer
ved at laane det Uendelige det Endeliges Former.« Som
jeg har gjort opmærksom paa (Religionsfilosofi § 95).
rører der sig dog her snarere en Modsætning mellem
Tendensen til at fastholde en vis bestemt Art af Billeder
og Tendensen til at lade Oplevelse og Stemning uvilkaarlig
fremkalde de Billeder, der kan symbolisere dem.
Denne sidste Tendens kan være forbunden med Bevidsthed
om ethvert Billedes Begrænsning, og Mystikeren
vil kunne opleve Tilstande, der ikke tillader noget  Bil%
% --- p. 108 ---
\setcounter{footnote}{0}%
\opage{108}lede at dukke op, før Eftervirkning og Eftertanke gør sig
gældende. Naar Cécile mere og mere trængte til en vis
bestemt Billedkreds, var det maaske, fordi hun havde
Ret, naar hun erklærede, at hun egentlig ikke var Mystiker.
Rettere er det dog at sige, at hun ikke blev ved
at være det. Mere og mere forsøger hun at oversætte
Oplevelsen paa personificerende Vis. Oplevelsen vedblev
at være den samme, og hun udsletter ikke Forskellen
mellem Oplevelse og Tydning. Den personificerende Tydning
lykkedes heller ikke lige efter en Oplevelse. »Imorges,«
beretter hun, »prøvede jeg paa, af de mægtige
religiøse Sindsbevægelser, jeg havde oplevet, at udlede
en Kultus, at oversætte den i formuleret Bøn, at centralisere
den om min Kærlighed til Kristi Person ved
at tillægge ham en Indgriben i mit indre Liv. Det lykkedes
ikke synderlig. Det forekom mig at svække og
forringe Oplevelsen, naar jeg vilde gøre den til et med
det, der udgør eller udgjorde min kristelige Tro.« Dog
lykkedes Oversættelsen bedre og bedre under Indflydelse
af en Trang, hun mere og mere blev sig bevidst.

Jeg tror ikke, man forstaar denne Udviklingsproces
hos Cécile rigtig, naar man ikke lægger Mærke til, at
samtidig blev Ekstaserne sjeldnere og sjeldnere og hørte
tilsidst helt op, ligesom ogsaa hendes stærke Drømme
(af hvilke i et enkelt Tilfælde Ekstase udviklede sig)
efterhaanden blev borte. Muligvis kommer her Forbindelsen
af Mystik og Seksualitet igen, da hun nu var i
den Alder, da Menstruationen, som ofte syntes at begunstige
Ekstasen, ophørte. Selv mener hun ikke, at Alderen
har haft nogen Indflydelse. Hun mener, at Ekstaserne
hørte op, fordi de ikke mere tilfredsstillede
hende. Forholdet er snarere det omvendte: da »den guddommelige
Oplevelse«, der efter hendes eget Sigende
affødte høje Idealer og omfattende Horizonter, af en
% --- p. 109 ---
\setcounter{footnote}{0}%
\opage{109}eller anden Aarsag, faldt bort, gjorde den Maalestok for
religiøse Forestillinger, den havde ydet hende, sig ikke
mere gældende. --- Det gik, som jeg har søgt at vise,
Teresa paa lignende Maade. ---

Dertil kom, at Ekstasen --- i den Periode, da den
var mulig --- kun optraadte i Ny og Næ, og at den ikke
kunde fremkaldes ved Beslutning, skønt en vis uvilkaarlig
Samling og Koncentration kunde lette dens Indtræden.
Den var en Komet, men der trængtes til en fast
Stjerne, som man altid kunde finde paa sin Horizont.
En stadig Tro og Tillid er maaske mindre vidunderlig,
men er i Grunden --- mener Cécile nu --- en bedre Tilstand
end Oplevelsernes glimrende men flygtige Straaler.
Og hun sætter nu sin Villie ind paa at fæstne Blikket
paa Polarstjernen under Livets Sejlads. Herved træder
hendes Tilbagevenden til konfessionel Religiøsitet i en
skarpere Modsætning til »den guddommelige Oplevelse«,
end hun selv er sig bevidst. Overgangen forudsætter Resignation.
Forholdet er, som naar en stærkt opblussende
Sindsbevægelse afløses af et stadigt Sindelag. Men den
foregaar hos Cécile gennem et Valg, en vilkaarlig Akt,
ikke gennem uvilkaarlig Udligning. Hun havde i »den
guddommelige Oplevelses Tid« trukket Vejret dybt og
fuldt, og hun kommer aldrig til at gøre det mere saa.
Men hun tager sit Parti og finder efterhaanden Ro og
Fred derved. Den forfærdelige Verdenskrig, der, foruden
at gribe ind i hendes økonomiske Forhold, fremkaldte
Bekymring og Angst for de Venner, hun havde paa
begge kæmpende Parters Side, bidrog ogsaa til, at hun
trængte til et fast Tilhold, der var fælles for hende og
hendes Omgivelser. ---

Céciles Mystik har da, som hun selv siger, kun været
en forbigaaende Fase i hendes aandelige Liv. Maaske
er det saaledes hos alle Mystikere. De ekstatiske Øje%
% --- p. 110 ---
\setcounter{footnote}{0}%
\opage{110}blikke har været faa, skønt de har været tilstrækkelige
til at præge den Indsats, Mystikerne har givet i det
aandelige Livs Udvikling. Om Aarsagerne til, at Ekstaserne
er knyttede til visse bestemte Aar, kan der paa
Grund af manglende specielle Oplysninger næppe siges
noget. Cécile er efter sin egen og Flournoys Mening
gaaet styrket og beriget ud af den hele Proces, selv om
Overgangen til roligere Farvande efter de indre Storme
ikke er sket uden Resignation. ---

Religionspsykologien staar i stor Gæld til Flournoy
for hans Offentliggørelse af denne mærkelige Episode i
en begavet Kvindes Liv, saavel som for --- hvad der
var en Forudsætning for, at dette Dokument kunde blive
til --- den Tillid, hans Personlighed og hans Forsken
har kunnet vække. Han har søgt at se denne Episode i
dens Sammenhæng med Céciles hele Liv og Udvikling,
og han har knyttet oplysende Bemærkninger til Beretningen.

Det forekommer mig, at han for meget anerkender
den Freud'ske Teori om Ekstasernes erotisk-seksuelle
Udspring. Som vi har set, og som Flournoy selv fremhæver,
træder den ekstatiske Oplevelse netop i skarp
Modsætning til den seksuelle Betagethed. Den indre
Sammenhæng mellem de forskellige Kræfter, der her
maaske drager Næring fra en fælles Rod, er mere indviklet
ind Freudianerne i deres naive Psykologi mener.
Og saa kommer dertil, at Freudianerne har givet de
seksuelle Udtryk (libido f. Eks.) en i den Grad udvidet
Betydning, at den hele Teori tilsidst svømmer ud. Naar
ogsaa Platon har været Freudianer, bliver vi det i Grunden
alle sammen; men saa falder Teoriens specielle Karakter
bort.

Et i religionpsykologisk Henseende betydningsfuldere
Spørgsmaal er det, om Flournoy har Ret i at betragte
% --- p. 111 ---
\setcounter{footnote}{0}%
\opage{111}den Udvikling, Cécilie gennemløber, som harmonisk og
fremskridende. »Hun kommer«, siger han, »tilbage til
sit Udgangspunkt; men som efter en Opstigning i Spiral
er det et højere Stade, hun nu vender tilbage til«. I hendes
Tilbagevenden til de kristelige Forestillinger ser han
et symbolsk Udtryk for Kærlighedens og Hengivelsens
Sejr over den aandelige Egoisme, som blev fremkaldt
ved den Selvfordybelse og Selvoptagethed, Ekstasen bevirkede.
(Il y a eu abandon progressif de l'expérience
divine, privilège égoiste et involontaire, pour en venir
au don personnel de soi-même à autrui).

Jeg kan ikke se andet, end at Flournoy overfor Céciles
Udviklingsgang bringer for meget System ind, idet
han fremfor alt søger at paavise Udviklingens Afslutning
i en højere Enhed. Han begaar her den samme
Fejl, som Delacroix og Pacheu overfor Sainte-Térèse.

Spiralen kan jeg ikke faa Øje paa. Cécile selv har
utvivlsomt fundet Fred i Sindet, og ganske særlig Harmoni
med sine Omgivelser, ved sin Tilbagegang til sit
Udgangspunkt. Men det er en Tillempelse, der ikke i
og for sig betegner et Fremskridt. Ja, hvis det at være
Religionslærerinde for unge Damer var et Ideal, der
for enhver Pris burde holdes fast, gjaldt det ganske sikkert
at arbejde sig ud af alt, hvad der kunde hæmme
fuld Hengivelse dertil. Forhaabentlig gives der dog andre
Veje for Menneskekærlighed, som gør det muligt at være
noget for andre uden at dele deres Tro med Fornægtelse
af sine egne dybeste, mest indgribende Oplevelser.

Cécilie har utvivlsomt ogsaa senere bevaret meget af,
hvad hendes mystiske Periode har lært hende om det
indre Livs Selvstændighed overfor alle Ord og alle ydre
Former. Men hendes personlige Løsning af Problemet
om Modsætningen mellem indre Oplevelse og ydre Virksomhed
kan ikke siges at være den ene rette. Det vilde
% --- p. 112 ---
\setcounter{footnote}{0}%
\opage{112}i hvert Tilfælde være sørgeligt, om det var saa. Hun
indrømmer jo selv, at hun lukkede Lugen paa Arken,
efterat hun et Øjeblik havde været saa dristig at slippe
Duen ud. Og hun indrømmer, at der var kvalmt inde
i Arken. Har der været bedre Luft nu, da hun havde
slaaet sig til Ro derinde? Hun har sikkert kunnet bidrage
til en Forbedring af Luften paa Grundlag af sine
Erfaringer. Men havde det ikke været bedre at lade Lugen
staa aaben og saa taalmodig vente paa, om Duen
skulde komme tilbage --- maaske med et Olieblad? Der
dæmrede her et nyt Princip, men Overleveringen var
stærkere end Mindet om den Oplevelse, der havde ladet
Principet dæmre. Kun derved blev det muligt, at et helt
andet og modsat Princip kunde følges: at Kærlighed til
Mennesker skal vise sig i at antage deres Meninger.

% --- p. 113 ---
\setcounter{footnote}{0}%

\opage{113}\conclhead{SLUTNING}

Det er Hensigten med denne Afhandling at belyse
Forholdet mellem Oplevelse og Tydning ved en Undersøgelse
af ekstatiske Tilstandes religiøse Betydning til
forskellige Tider og under forskellige individuelle og
sociale Forhold. Dette er naturligvis kun et enkelt Eksempel
paa den Betydning, Oplevelse kan have for de
religiøse Forestillinger, men det er et vigtigt Eksempel.

Forholdet mellem Oplevelse og Tydning stiller baade
et erkendelsesteoretisk og et psykologisk Problem. For
Erkendelsesteorien maa, naar den ikke falder tilbage
i Dogmatisme, al Videnskab tilsidst staa som en stor
Tydning af Oplevelser, der hver for sig er beskrevne
saa nøjagtigt som muligt, helst gennem Tal, Maal og
Vægt. Ogsaa fra det erkendelsesteoretiske Synspunkt
gælder det ikke at forveksle Oplevelse og Tydning. Paa
Tydningen beror det, hvorledes den enkelte Oplevelse
indgaar som et Led i vore Tankers Verden. Gennem
Tydningen kommer den enkelte Oplevelse først til
sin fulde Ret. Men det maa stedse paany prøves, om
Tydningen udtømmer hele Oplevelsen, baade dens Kvalitet
og dens Kvantitet, baade dens successive Udfoldelse
og dens typiske Ejendommelighed. Og derfor maa
det prøves, om man ud fra Tydningen igen kan udlede
Oplevelsen med dens specielle Beskaffenhed\footnote[1]{\textit{Totalitet som Kategori}. (Vid. Selsk. Skr. 1917). p. 43---45.}. Psykolo%
% --- p. 114 ---
\setcounter{footnote}{0}%
\opage{114}gisk anvendes der et mere omfattende Begreb om Tydning.
Psykologien undersøger ikke blot de Tydninger,
der fører til en vel sammenhængende Tankebygning,
altsaa til Videnskab, ikke blot de klart bevidste og logisk
udformede Tydninger, men ogsaa, og især, den uvilkaarlige
og halvt bevidste Maade, paa hvilken en Oplevelse
kan bringes i Forbindelse med Menneskets øvrige
Forestillingskreds. De Forskelligheder, der her gør sig
gældende for en sammenlignende Betragtning, fra Totemismen
gennem Heksetro til Spiritisme, eller fra middelalderlig
Mystik til moderne Oplevelse af »guddommelig«
Art, er af stor kulturhistorisk Interesse. Naar man
følger Eksemplernes Række, fremtræder der en stedse
større Forskel mellem Oplevelse og Tydning, og en
stedse større Selvstændighed for selve Oplevelsen, --- i
Analogi med den Forskel, den kritiske Filosofi gør
mellem Iagttagelsen og dens videnskabelige Tydning. ---

Hvad vort specielle Eksempel, de ekstatiske Oplevelser,
angaar, saa ligger de til Grund for en betydningsfulde
Brydning mellem det indre Liv og de Overleveringer,
under hvis Indflydelse dette indre Liv føres. I
Mystiken har i lange Tider Personligheden haft sit Tilhold
overfor en fast udhamret Dogmatik, der forlangte at
bestemme, hvad der skulde anerkendes som religiøs
Erfaring. Det ses hos de israëlitiske Profeter, hos Middelalderens
Mystikere og i særlig udpræget Form hos
den moderne Mystiker, vi skylder Flournoy Bekendtskabet
med. Der foregaar her en stille Udvikling lige
op til den Tid, da det blev proklameret (her i Norden
af E. G. Geiger, Sibbern og Kierkegaard), at paa det
religiøse Omraade har intet Værdi, som ikke udspringer
af personlig Oplevelse, og da endog det Individuelle
blev erklæret for nutildags at være »det eneste Levende
i Religionen.« Det er ikke mindst de mystiske Oplevel%
% --- p. 115 ---
\setcounter{footnote}{0}%
\opage{115}sers direkte Umeddelelighed, som her er af Betydning,
idet den gør Individet selv og andre opmærksom paa,
at det vigtigste i Livet former sig ejendommelig hos
hvert alvorligt Individ, og at alt Fælles er af anden
Rang. Mange Ironikere --- d. v. s. Repræsentanter for
den store Ironi --- nærmer sig i deres indre Liv til
Mystikerne, forsaavidt de ligesom disse, mærker at der
er en umeddelelig Fylde, som udfolder sig i dem. Det
gælder, som jeg har sagt at vise, om Sokrates. Ogsaa
den store Humor betinges for en væsentlig Del ved
Inkommensurabiliteten mellem det Indre og det Ydre\footnote[1]{\textit{Den store Humor}. § 16—17; 19; 25.}.
Hvorledes man saa opfatter eller forklarer ekstatiske
Tilstande, --- hvad enten man kalder dem patologiske
eller ikke, --- hvad enten man betragter dem som indeklemt
Seksualitet eller tillægger dem et ædlere Udspring,
--- saa er det et Faktum, at der har været Mennesker,
og det ikke de ringeste, som i dem havde de vigtigste,
de mest afgørende Oplevelser, og det er religionspsykologisk
Hovedsagen.

For Mystikerne, som for Totemisterne, betegnede de
ekstatiske Oplevelser Højdepunkterne i deres Liv. Fra
disse Oplevelser udstraaler en Glans, der bestemmer
Rangforholdet mellem alle andre Oplevelser. Det ligger,
som jeg har søgt at vise, i Følelsens Natur, at den, naar
den er vakt ved en enkelt Oplevelse, søger at brede sig
over hele Sjælelivet og at farve alle andre Elementer i
dette, hvad enten de staar i Forbindelse med hin Begivenhed
eller ikke. Indre Tilstande, der i og for sig
ikke har noget med hverandre at gøre, kan derved faa
et fælles Præg, og der kan komme en Helhed ud af de
indre Oplevelsers Række, hvad enten denne er afsluttet
eller ikke\footnote[2]{Se, foruden min Psykologi, \textit{Religionsfilosofi} § 27.}. Det er som bestemmende Elementer i en
% --- p. 116 ---
\setcounter{footnote}{0}%
\opage{116}Totalfølelse, at de ekstatiske Oplevelser faar deres Betydning.
Selv om saadanne Oplevelsers Tid skulde være
forbi, vil det Problem, der stillede sig for Mystikerne,
stedse stille sig for os. Ti der vil stadig gøre sig en Forskel,
maaske en Modsætning, gældende mellem Øjeblikke,
i hvilke vi oplever, hvad der for os faar den højeste
Værdi, og de jevne, dagligdags Oplevelser, eller endog
de lave og dyriske Tilstande, og det kommer da an paa,
hvilket Resultat Brydningen mellem disse forskellige
Elementer fører til, og hvilken Dom om Livet, der som
Følge deraf fældes. Det er af stor Vigtighed, at Oplevelserne
faar Lov til at øve deres uvilkaarlige Indflydelse,
og at Tydningen foretages med fuld Frihed. Mystikens
Historie har lært os, at Selvbeherskelse og Selvtænkning
ikke derved udelukkes; de vil kun uvilkaarlig bestemmes
ved det Element, der fremtræder som det værdifuldeste.
Der vil her, især for udprægede Personligheder,
være meget, som ikke gaar op i noget afsluttende Begreb
og heller ikke kan »katalogiseres« i nogen Dogmatik,
og derfor vil, som vi har et smukt Eksempel
paa i Cécile, et poetisk Udtryk være det eneste mulige.
Al Livspoesi har her sit Udspring. Céciles Yndlingsbillede
om Planken, der vugges og gennemstraales i det af
Aftensolen glødende Hav, har sin store Værdi ikke blot
ved dets Poesi, men ogsaa derved, at det er selvfrembragt,
umiddelbart udsprunget af Oplevelse.
% --- p. 117 ---
\setcounter{footnote}{0}%

\opage{117}\essayhead{SOCIOLOGI OG RELIGIONSFILOSOFI}

\essaysource{LES FORMES ELÉMENTAIRES DE LA VIE RELIGIEUSE \\ Par Émile Durkheim. (Paris. Alcan). 1912.}

I dette Værk af den franske sociologiske Skoles Fører
er der, i højere Grad og med større Klarhed end i
noget tidligere Arbeide fra denne Skole, aabnet en Vej
til Forstaaelse mellem Sociologien og Filosofiens andre
Dele. Gennem Studiet af en speciel Form for religiøst
Liv kaster Forfatteren ikke blot Lys over selve det religiøse
Problem, men viser ogsaa, hvorledes man kan
anerkende det specielle sociologiske Synspunkt for menneskeligt
Aandsliv og dog fastholde Berettigelsen af andre
Synspunkter. Sociologien har hidtil ofte, ogsaa i Durkheims
Fremstillinger, gjort Fordring paa at være den
eneste Aandsvidenskab. Dels direkte, dels ogsaa, og især,
indirekte opgives denne Fordring i det foreliggende betydelige
Værk. Vel hævdes det sociale Liv, Samfundet og
dets Institutioner og Traditioner, som Kilde til alt menneskeligt
Aandsliv; men gennem den nærmere Udvikling
af de bærende Ideer føres den kritiske Læser ved Forfatterens
egen Haand ud over det eksklusivt sociologiske
Standpunkt. Og det betydningsfulde er, at Sociologien
ikke taber, men netop vinder derved. Ligesom selve
% --- p. 118 ---
\setcounter{footnote}{0}%
\opage{118}Samfundet netop godtgør sin Kraft ved, at de enkelte
Individer kan røre sig med Frihed og udvikle deres
Ejendommelighed hvert for sig, saaledes vil ogsaa Samfundsvidenskabens
Betydning stige, jo mere de forskellige
Discipliner, der fra forskellige Synspunkter undersøger
Aandslivet, hævder deres Selvstændighed. Allerede Comte
stillede, i sin sidste Tid, Etiken som en selvstændig Disciplin
ved Siden af Sociologien. Og paa samme Maade
vil Psykologi, Erkendelsesteori og Religionsfilosofi vise
sig at betegne selvstændige Udgangspunkter i Forhold
til Sociologien, hvor meget de end kan lære af denne.
Jeg véd vel, at Durkheim ikke vil anerkende dette fuldt
ud; men jeg vil søge at vise, at foruden alt det, vi positivt
kan lære af det foreliggende Værk, der paa sit
eget specielle Omraade betegner et stort Fremskridt,
lærer vi fyldigere end før at se, hvorledes Sociologien dels
faktisk gør Brug af psykologiske og erkendelsesteoretiske
Forudsætninger, dels i sine Konsekvenser fører til at
anerkende hine andre aandsvidenskabelige Discipliners
Betydning. Det er fra denne Side, en Filosof naturlig
kommer til at betragte et Værk som det foreliggende.
For en sociologisk Specialist vil det sikkert have en Betydning,
som jeg ikke er kompetent at give Udtryk for.
Men efter min Overbevisning har Værket en Betydning,
der strækker sig langt ud over det rent sociologiske Omraade.
Jeg for mit Vedkommende har i hvert Tilfælde
ikke blot lært Sociologi af det, men har fundet Oplysning
og Bestyrkelse ogsaa paa andre Punkter. Ideer og
Synspunkter, til hvilke jeg ad psykologisk Vej var kommen,
har jeg ved Studiet af dette grundige og mesterlig
gennemførte Værk, fundet baade Berigtigelse og Bekræftelse
af. Gennem en Karakteristik og gennem en
dertil knyttet Kritik vil jeg forsøge at begrunde dette
nærmere.

% --- p. 119 ---
\setcounter{footnote}{0}%
\opage{119}\romanhead{I}
% CHECK roman numeral read

Medens Historikeren og Etnografen studerer, hvad der
er foregaaet til bestemte Tider, og hvad der har været
eller er Sæd og Skik paa bestemte Steder, søger Sociologen
ifølge Durkheim dybere tilbunds. Han vil finde
de Kræfter, der virker i Menneskelivet den Dag idag, de
Realiteter, der ligger bag ved Historiens Begivenheder og
Folkeslagenes Sædvaner. Hvad Religionen angaar, gælder
det da ikke blot om at skildre Dogmer og Ceremonier,
men at finde det Reale, der ligger bag dem og giver
dem deres egentlige Betydning, --- det Evige og Menneskelige
i Religionen. Sociologen gaar ud fra, at Religionen
hænger sammen med Livets Virkelighed, med
en menneskelig Trang, der mere eller mindre rører sig
til alle Tider.

De fleste Teorier om Religionens Natur og Grundlag
staar magtesløse overfor den Kendsgerning, at Religionen
trods al Modstand og Kritik stadigt holder sig. Den maa
da svare til en bestemt Trang, en bestemt Kraft i den
menneskelige Natur. Hvis den, som den animistiske
Teori antager, blot skyldes Drømmeforestillinger, der
tages for Virkelighed, eller hvis den, som den naturistiske
Teori antager, blot skyldes falsk Naturforklaring, vilde
dens Tid forlængst have været forbi. Men det er ikke
en Trang til at forstaa, det er en Trang til at leve, til
at opnaa Styrke og Tryghed under Livskampen, der
ligger til Grund. Durkheim erklærer sig altsaa bestemt
imod den intellektualistiske Opfattelse af Religionen.

For at finde Religionens Grundlag maa man først og
fremmest være klar over, at den religiøse Bevidsthed ikke
selv formaar at tyde den Trang og Tilskyndelse, som
virker i den og driver den frem. Den er tilbøjelig til at
søge Aarsagen til, hvad der bevæger sig i den, udenfor
% --- p. 120 ---
\setcounter{footnote}{0}%
\opage{120}sig. Vil man ret blive klar herover, er det hensigtsmæssigt
at opsøge de simpleste, mest elementære Former
for Religion. Man kan vel ikke opstille som almindelig
Regel, at den mere elementære Form for Religion
forklarer den mere udviklede, ligesaa lidt som
man kan opstille den modsatte Regel. Men de mere
elementære Former frembyder den Fordel, at den folkelige
og præstelige Fantasi endnu ikke har ladet en yppig
Vækst af Myter og Dogmer blive til. I religiøs som i
social Henseende forefindes kun det strengt Nødvendige.
Og navnlig gør der sig ingen individuelle Forskelligheder
gældende i religiøs Henseende. Alle Bevidstheder
bevæger sig i samme Kølvand. Den individuelle og den
sociale Type falder sammen. En saadan elementær Religion
finder Durkheim i Centralaustraliens Religion,
der i de senere Aar er bleven saa fortræffelig skildret
især af Spencer og Gillen og af Strehlow. Det er den
simpleste Religion som kendes. Den har intet egentligt
Gudsbegreb. Hvad der gør den til Religion, er, at den
beherskes af den store Modsætning mellem Helligt og
Profant. Og denne Modsætning viser sig her at staa i
nøjeste Sammenhæng med hele den sociale Organisation.
Ligesom Religionen er ogsaa det sociale System overordentlig
simpelt. Klanens Medlemmer er hverken forbundne
ved Blodslægtskab eller ved fælles Bopæle; det
er selve Kultus'en, der forener dem. De har samme
Navn og Emblem (Totem) og foretager sig de samme
Ceremonier, nærer samme Opfattelse af, hvad der er
helligt, og hvad der er profant.

Det er Modsætningen mellem Helligt og Profant, det
gælder at forklare. Hvorfra kommer den Følelse af det
Ophøjede, som Mennesket allerede paa dette Trin har
overfor visse Ting? Hvorledes opstaar Modsætningen mellem
et Højere og et Lavere i og udenfor Mennesket?
% --- p. 121 ---
\setcounter{footnote}{0}%
\opage{121}Hvorledes føres Mennesket til at anerkende noget, der
er større, mere omfattende end det selv?

Baade Sjæledyrkelse og Naturdyrkelse forudsætter
denne Modsætning. Og det religiøse System, vi her staar
overfor, den saakalte Totemisme, ligger længere tilbage
end Sjæledyrkelse og Naturdyrkelse.

Totem er et Navn og et Emblem, der er fælles for
Klassens Medlemmer og forpligter dem til gensidig Hjælp,
til at hævne hverandre, sørge med hverandre og overholde
visse Ceremonier og visse Ægteskabsforbud. Totem
kan være et Dyr, en Plante eller en livløs Genstand.
Mennesket føler sig Et med sit Totem, og gennem det
med sin Stamfader, fra hvem den stammer, eller som
maaske selv har forvandlet sig til det. Ligesaa lidt som
Australieren skelner mellem Billedet og Tingen, ligesaa
lidt skelner han mellem sit Totem og sig selv eller sin
Stamfader. Og han føler sig herigennem solidarisk med
alle, der har samme Totem. Gennem en særegen, tildels
streng og smertefuld Indvielsesproces bliver de Unge
delagtige i Klanens Totemdyrkelse og faar derved Delagtighed
i det Hellige. Et helligt Redskab (churinga eller
tjurunga), som han ved denne Indvielse modtager, forbinder
ham med sit Totem og sin Stamfader, og der
udgaar fra det en styrkende og opmuntrende Kraft, hvorfor
Tabet af det er en stor Ulykke.

Forskellen mellem Helligt og Profant hævdes gennem
Kultus. Den negative Kultus inskærper, gennem Faste,
Vaagen, Smerte, Ensomhed, Tavshed, at ved Indvielsen
til Delagtighed i det Hellige maa det tidligere Menneske
aflægges, for at nye Egenskaber og Kræfter kan opnaas.
Den positive Kultus bringer Mennesket under
Paavirkning af nye Kræfter, der udstraaler fra de hellige
Genstande, hvor de indeholdes som i en Art permanent
Reserve. Gennem Kultus styrkes ikke blot Mennesket
% --- p. 122 ---
\setcounter{footnote}{0}%
\opage{122}selv, men selve Totemgenstanden næres derved. Det er
en Art Sekrament eller Kommunion, hvorved Mennesket
paa en Gang bliver delagtig i det Helliges Kraft og ved
sit Offer styrker og forøger denne Kraft. At Mennesket
trænger til saadan Styrkelse, kommer af dets Afhængighed
af Naturmagterne. Naturens Liv undergaar Forandringer,
stiger og synker i Kraft og Frugtbarhed, og
Mennesket trænger til at faa sin Tillid til Naturen styrket,
medens han tillige mener gennem Kultusskikkenes nøjagtige
Overholdelse selv at kunne hjælpe paa Naturen.

Bag Modsætningen mellem det Hellige og det Profane
ligger nu ifølge Durkheim Modsætningen mellem
Samfund og Individ skjult. Totem er fra først af et
socialt Begreb, et Udtryk for Menneskets Solidaritet med
Klanen og for Solidariteten indenfor Klanen. Ikke Frygt,
men tillidsfuld, og paa Højdepunkterne begejstret Tilslutning
til, hvad der kan bringe Frelse og Tryghed, en
Frelse og Tryghed, som kun Livet i Klanen kan bringe,
er det Grundlag, hvorpaa Modsætningen mellem det
Hellige og det Profane udfolder sig. Forudsætningen for
Totemismens Kultus er, at der er et Fond af Kraft, som
de enkelte Individer kan bidrage til at bevare og at
øge. Det er da ikke blot Spil af Ideassociationer og
Hallucinationer, der virker i Totemismen og i de Magter
og Symboler, der knytter sig til den. Kilden til Religionen
ligger tydelig i andre Regioner af Sjælelivet end den
rent intellektuelle. Gennem Hengivelsen til det Hellige
og Iagttagelsen af de hellige Skikke føler Mennesket sig
i Forhold til noget, hvis Rækkevidde naar langt ud
over de enkelte Individers og de enkeltes Generationers
Eksistens. Totem er i den simpleste Religion, vi kender,
det karakteristiske Symbol herfor, og det er Klanens Bestaaen
og Velfærd, der er det reale Grundlag for dette
Symbol. Gennem Kultus styrkes og genskabes periodisk
% --- p. 123 ---
\setcounter{footnote}{0}%
\opage{123}selve Samfundet. Under de i Fællesskab og i de fra
Fædrene overleverede Former fejrede Fester falder Særinteresser
og dagligdags Ønsker og Bekymringer bort.
Den enkelte føler sig greben af noget, der er større end
han selv. Gennem Kultusskikkene, som ofte har en
betagende og ophidsende Karakter, sættes Sindet i en
Bevægelse, som det sædvanlige Liv ikke kan volde. Det
føler sig genfødt, badet i den Kilde, alt Liv skyldes.

Klanens Bestaaen og Velfærd er paa mange Maader
knyttet til Naturens Gang og betinget ved dens Rytme,
afhænger af, om Frugten modnes, om der kommer Regn,
om Jagten lykkes. Og gennem Kultus mener Mennesket
at kunne øve Indflydelse herpaa; derfor er det af saa
stor Betydning, at de hellige Skikke strengt iagttages.
Denne Tro er et Udslag af den stærke indre Betagethed,
som Deltagelse i den fælles Kultus volder. De store Festers
periodiske Genkomst svarer til Naturens rytmiske
Gang, noget, der kan spores helt op i de højeste Religioner.
I Mellemrummene mellem Festerne spredes Klanens
Medlemmer og kæmper, enkeltvis eller i smaa Flokke,
for deres Selvopholdelse. Kun gennem Totembillederne
fastholder de da Erindringerne om de store Øjeblikke af
Betagethed og Begejstring, og de føler Modsætningen
skarpt. I deres eget Indre gør Modsætningen mellem
Helligt og Profant sig gældende, svarende til Modsætningen
mellem Deltagelse i det fælles Liv gennem Kultus
og den isolerede Eksistens i Mellemtiderne. Den levende
Trang til atter og atter at opnaa hine store Øjeblikke
er Religionens Grundlag. Saaledes kan man ogsaa i de
højere Religioner se Mennesker, hos hvem der raader
en mere eller mindre bevidst Tvivl om religiøse Forestillingers
og Skikkes Betydning, atter og atter deltage i
den gamle Kultus. Sagen er, at religiøse Skikkes virkelige
Begrundelse ikke er at søge i de tilsyneladende For%
% --- p. 124 ---
\setcounter{footnote}{0}%
\opage{124}maal, de tjener, men i den usynlige Virkning, de udøver
paa Sindene, den Maade, paa hvilken de stemmer hele
den indre Tilstand.

Naar Samfundet i hine store Øjeblikke gør sig gældende
som Magt i Individernes Indre, sker der en uvilkaarlig
Idealisation af det. Totemsymbolet betegner noget,
der rækker ud over det Bestaaende. Selvfølgeligt er
Mennesket sig ikke selv bevidst, hvad der er det sidste
Grundlag for Modsætningen mellem Helligt og Profant,
--- hvad der tilsidst vækker de Følelser, som overvælder
det. Aarsagen til disse Følelser finder det udenfor sig ---
ligesom det uden videre overfører sine Sansefornemmelsers
Kvaliteter paa selve Genstandene. ---

Nu angaar Religion og religiøse Forestillinger, allerede
paa det mest primitive Trin, ikke blot Menneskelivet,
men faar ogsaa kosmologisk Betydning. Den Kraft,
Mennesket føler sig betaget af i sin Kultus, og som styrker
og nærer det, genfinder det uvilkaarlig i hele Naturen.
Modsætningen mellem Helligt og Profant udstrækkes
til alle Naturgenstande, der kommer indenfor Synsvidde.
Det er en Art Smitte, der her virker, Alt, hvad
der har nogensomhelst Berøring med det Hellige eller
med det Profane, bliver selv helligt eller profant. Det
primitive Menneske kender ikke vore Distinktioner mellem
forskellige Arter af Genstande. Ligesom det ikke
finder nogen Modsigelse i at opfatte sig selv og sin
Stamfader som et med Totemgenstanden, saaledes opfatter
det overhovedet alt som et, hvoraf det paavirkes
paa samme Maade. Denne overordentlige Smitsomhed
af de religiøse Kræfter (l'extrême contagiosité des forces
religieuses) skyldes ikke blot Ideassociationer. Hellighedskarakteren
beror jo først og fremmest paa den stærke
Følelse af Afhængighed og Bestyrkelse, der er frembragt
i Sindene, og jo stærkere denne Følelse ytrer sig, des
% --- p. 125 ---
\setcounter{footnote}{0}%
\opage{125}mere vil den ogsaa brede sig til alt, hvad der, særlig
i Betagethedens Øjeblikke, kommer i Berøring med Mennesket.
Herved forstaar man den Lethed, hvormed det
primitive Menneske identificerer de mest uensartede
Ting, saa meget de end for Sanserne fremtræder med
forskellige Egenskaber.

Derved bliver en Religion ikke blot en Livsanskuelse,
men ogsaa en Verdensanskuelse, ikke blot en Psykologi,
men ogsaa en Metafysik. ---

Den religiøse Trang er altsaa et med den sociale
Selvopholdelsestrang, der rører sig hos alle og forbinder
dem til et Hele. Den fører nu uvilkaarlig til Tro
paa Aander. Det er en Misforstaaelse at mene, at Tro
paa Aander skulde ligge til Grund for Religionen, denne
har, som vi har set, et helt andet Grundlag. Men der
gives faktisk ingen positiv Religion uden Tro paa Aander.
For Autraliernes Vedkommende fremtræder Begrebet
Sjæl straks med Hellighedens Præg. Gennem sit Totem
er det enkelte Individ knyttet til Stamfaderen, er
et Fragment af den store Kraft, der igennem Totemdyrkelsen
besjæler hele Klanen. Der gives kun et vist Antal
Sjæle i Tilværelsen, og de nyfødte Børns Sjæle er
direkte Reinkarnationer af Forfædres Sjæle, eller de
skyldes Spirer, som Forfædre har udsaaet. Sjælen er
ikke andet end det totemiske Princips Inkarnation i et
enkelt Individ. --- Drømme og lignende Fænomener
har afgivet Støtte for Sjæletroen, men indeholder ikke
den egentlige Forklaring af dens Opstaaen. Da Klanen
stadigt fortsætter sit Liv paa samme Maade gennem de
nye Mennesker, som fødes, staar det som en ganske naturlig
Forestilling, at de nye Sjæle stammer fra Sjæle,
der forhen har levet. Udødelighedstroen har her endnu
ingen moralsk Betydning. Baade Sjæletroen og Troen
paa mægtige Aander, der i Fortiden lagde Grunden til
% --- p. 126 ---
\setcounter{footnote}{0}%
\opage{126}alt Godt, har udviklet sig af Totemismen. Egentlige Guder
kendte Centralaustralierne ikke. Der foregaar dog en
Udvikling i Retning af en stedse større Afstand mellem
Fortidens Store, om hvilke Myter og Legender taler, og
de andre Forfædre. Særlig ved de for hele Stammen,
ikke blot for den enkelte Klan fælles Fester, fik de mytiske
Væsener en mere ophøjet og universel Karakter.
Fra Totemismen kunde Udviklingen Skridt for Skridt
ske i Retning af Forestillinger, som vi finder i højere
Religioner. ---

Studiet af elementær Religion, hvorpaa Centralaustralien
giver saa klassisk et Eksempel, er ikke blot oplysende
med Hensyn til selve det egentlige religiøse Liv.
Det godtgør ogsaa paa lærerig Maade, hvorledes Religionen
fra først af staar som den eneste Form for aandelig
Kultur; særlig indeholder den Oprindelsen til Videnskaben.
Og da Samfundet er Moder til Religion, er det
ogsaa Moder til Videnskaben. De videnskabelige Tankers
Autoritet skriver sig fra, at de er et Værk af kollektiv
Erfaring, af Samfundets Fællesarbejde. Det er utallige
Individers og Generationers Iagttagelse og Eftertanke,
der ligger til Grund for, hvad de enkelte kan udforme
i klare Begreber. Især har Religionen ydet Tankelivet
en Tjeneste ved at muliggøre den første Forestilling
om, hvad Slægtsskabsforhold mellem Tingene vil
sige. Den har lært Menneskene at stille det sammen,
som Sanserne skiller ad. Til Ideernes Verden har Vejen
kun kunnet banes ved en Anspændelse af intellektuelle
Kræfter, som først blev mulig i og ved Samfundet, særlig
gennem dettes Kultus. De første Forsøg paa Klassifikation
foretoges med Stammernes Inddeling i Klaner
som Forbillede; derved var Grundlaget for al Logik
givet. Grundbegrebernes, de saakaldte Kategoriers Opstaaen
finder hverken sin Forklaring ved Apriorisme
% --- p. 127 ---
\setcounter{footnote}{0}%
\opage{127}eller ved Empirisme; den skyldes hverken oprindelige
indre Anlæg eller Sansepaavirkning udefra. Det er sociale
Tilstande og Forhold, der indeholder Forklaringen,
ligesom allerede ved Klassifikationen. Tidsforestillingen
udvikles saaledes gennem de Tidsinddelinger, som svarer
til Festers og Ceremoniers Rytme. Rummets Inddeling
hænger oprindelig sammen med de forskellige
religiøse Betydninger af de forskellige Retninger og Steder.
Ogsaa Klanernes Lejring indenfor Stammens fælles
Samlingsplads har afgivet et Skema. Aarsagssætningen,
specielt det Princip, at det Lige maa frembringes
af det Lige, skyldes Kultusoverleveringens strenge Fordring
om, at Ceremonierne foretages paa en aldeles bestemt
Maade; ellers kunde det ønskede Resultat ikke
opnaas. Her var altsaa Normer, efter hvilke alle maatte
forme deres Forestillinger. Den kollektive Handlens nødvendige
Krav udtryktes i kategoriske Former. Fornuftens
Autoritet er af social Oprindelse. Men Samfundet selv
er en Del af Naturen og har kun det Fortrin fremfor
andre Naturomraader, at de fundamentale Forhold fremtræder
tydeligere og ligger os nærmere ind paa Livet.
Derfor har Forestillinger, der udarbejdes efter sociale
Tings Mønster (sur le modéle des choses sociales), kunnet
hjælpe til at forestille sig Ting af en helt anden
Natur.

Hvad Religionens Fremtid angaar, er for Durkheim
især to Betragtninger af Betydning.

Religionen, saaledes som den kan studeres historisk,
har hidtil væsentlig været et socialt Fænomen. Religion
og Samfund hører historisk nøje sammen. Gennem Samfundets
Indvirken paa de Enkelte --- i deres Indre ved
utallige ubevidste Indflydelser, udefra ved utallige Bud,
Forbud, Vaner og Overleveringer --- opstaar i og for
dem den store Modsætning mellem Helligt og Profant,
% --- p. 128 ---
\setcounter{footnote}{0}%
\opage{128}mellem et Højere og et Lavere, mellem det Ideale og
det Sanselige. Og gennem de store Fælleshandlinger ved
Festerne opstaar den Begejstring, det høje Sving, som
er en stadig Kilde til religiøs Følelse, og uden hvilket
den i Længden ikke kan bestaa. Men nutildags er det
blevet anderledes. Vi lever i en Overgangstid. Individualismen
har bredet sig sammen med Arbejdsdelingen. Vi
kan ikke længere af fuldt Hjerte fejre de Fester, der
begejstrede vore Fædre. Vi gaar hver sine Veje. Naar
det var karakteristisk for al Religion hidtil, ligefra dens
mest elementære Former af, at den er fælles for en
Gruppe, som faar sin Livskraft næret og fornyet gennem
fælles Højtider, saa staar det som en stor Gaade,
hvorledes der for Fremtiden kan bestaa Religion, og
hvorledes ny Religion skal kunne danne sig. Hvis denne
Gaade vil finde sin Løsning, saa kan vi i hvert Tilfælde
ikke danne os nogen Forestilling om denne Løsning.
Der gives ganske vist en religiøs Idealisme i Form
af indre, subjektiv Religiøsitet, som former sig forskellig
i de forskellige Individer. Men Sociologien har kun
med fastslaaede og realiserede Kendsgerninger at gøre,
og det er et Fremtidsspørgsmaal, om den religiøse Idealisme
kan komme til at faa en saadan Karakter af Kendsgerning,
at den kan blive Genstand for objektivt Studium.
I Sociologien maa vi holde os til, hvad Religion har
været i Fortiden.

Foruden Forholdet mellem det Sociale og det Individuelle
er ogsaa Forholdet mellem Religion og Videnskab
af Betydning for Fremtidsudsigterne. Her maa det
staa fast, at Religionen forlængst har mistet sin intellektuelle
Betydning. Dens Kosmologi kan ikke længere bruges.
I denne Henseende er Videnskaben et fuldstændigt
Ækvivalent, --- Videnskaben, der selv har udviklet sig
af den. Men det er ogsaa kun i rent intellektuel Hen%
% --- p. 129 ---
\setcounter{footnote}{0}%
\opage{129}seende, at Videnskaben er Ækvivalent her. Den kan ikke
erstatte Religionen, forsaavidt denne er en Kraft, der
holder Mennesket aandelig i Live, et Grundlag for Mod
og Tillid. Livet kan ikke vente, til Videnskaben bliver
færdig. Ja, den bliver aldrig færdig. Og Livet vil derfor
stedse drive Tanken ud over Videnskabens Grænser.
Det vil derfor snarere blive en Omdannelse end en Tilintetgørelse
af Religionen, som Fremtiden vil bringe.

\romanhead{II}
Kun en ufuldkommen Ide har forestaaende Fremstilling
af det foreliggende Værks Indhold kunnet give
om den rige Belæring, der kan høstes af det. Det er
bygget over et meget specielt Emne, de centralaustralske
Stammers Religion, og dette Emne er drøftet med
grundigt Sagkundskab og energisk Analyse. Jeg er fra
Durkheim's Bog gaaet tilbage til nogle af Kildeskrifterne,
navnlig til Strehlow's Beretninger og har fundet fuld
Bekræftelse paa Bogens Opfattelse. Ved sin koncise
Fremstilling udmærker den sig fra andre Skrifter over
lignende Emner, der trætter ved Bredde og Gentagelser.
Frazer's berømte og udmærkede »Golden Bough« faar
saaledes ofte Læseren til at stanse af Træthed formedelst
Gentagelserne og Mangelen af Koncentration. Desuden
er der hos Durkheim en lærerig Kritik saavel
overfor Enkeltheder som overfor hastige Generalisationer.
Men alligevel interesserer det specielle Emne ham
kun som Grundlag for Forstaaelse af Religionens almindelige
Natur. Han mener i den Religion, der er den
simpleste af alle, vi kender, at finde det Væsentlige for
al Religion. Og ved at grave saa dybt, som han gør, ---
ved at føre Analysen til Ende med den Energi, der
overhovedet udmærker ham som Forsker og Forfatter,
er han ogsaa, saavidt jeg kan se, kommen i Berøring
% --- p. 130 ---
\setcounter{footnote}{0}%
\opage{130}med det Centrale i Religionen. Det har i hvert Tilfælde
for mig været af stor Interesse at se, hvorledes der her
fra et sociologisk Standpunkt er udarbejdet en Opfattelse,
jeg i min »Religionsfilosofi« er naat til ad psykologisk
Vej. Men for at nærmere Forstaaelse her kan
være mulig, er der et Mellemværende at opgøre mellem
Sociologi og Psykologi.

Fra den store Modsætning mellem Helligt og Profant,
i hvilken Durkheim finder det Væsentlige i al Religion,
og med hvilken al Religion begynder, idet den
er fundamental i Forhold til al Kultus, alle Myter og
alle Dogmer, gaar han tilbage til Samfundet, dets Bestaaen
og dets Velfærd. Men, som han selv antyder, er
der endnu et Spørgsmaal at gøre: hvorfor er det netop
Samfundet, der, uden at Mennesket bliver sig det bevidst,
er Kilde til alt, hvad der er helligt? Durkheim
giver selv Svaret: Det er, fordi Samfundet er Kilden til
de største Værdier for Mennesket, --- fordi det er den
store Støtte i Kampen for Opholdelsen, og ikke mindst,
fordi det er Ophav til alt, hvad der gør Mennesket til
mere end Dyr. Durkheims Tanke her er analog med
den, der laa til Grund for hans Udtalelser i det franske
filosofiske Selskab om »Moral som Kendsgerning« (le fait
moral).\footnote[1]{\textit{Bulletin de la Société Française de Philosophie}. 1906. --- Jevnfør min Bog \textit{Den menneskelige Tanke} p. 353---357.} Han førte der al Moral tilbage til Samfundet
og begrunder dette ved, at Samfundet er Kilde til, Bevarer
og Overleverer af al Civilisation, særlig af det
aandelige Livs Goder. Det kommer da an paa, om Samfundet
saa ogsaa virkelig udfylder denne Plads. Og her
viser det sig da, at der er Forudsætninger, der ligger
dybere end Samfundets Bestaaen. Naar det Hellige forudsætter
det Sociale, forudsætter det Sociale igen det
Værdifulde. Samfundet skal jo være Værdiernes Kilde,
% --- p. 131 ---
\setcounter{footnote}{0}%
\opage{131}og kun derved kan det ligge til Grund for Hellighedskvaliteten.
Ad denne Vej kommer jeg fra Durkheim's
Definition af Religion tilbage til den, jeg har forsøgt at
gennemføre i min »Religionsfilosofi«, nemlig at al Religion
hviler paa Forudsætningen om Værdiens Bestaaen.
Kun saalænge og saavidt Samfundets Bestaaen staar som
Del af eller Middel til Værdiens Bes\textasciitilde{}taaen, kan den faa
en religiøs Karakter. --- Jeg er ganske enig med Durkheim
i hans skarpe Skelnen mellem den sidste Forudsætning
for eller det sidste Motiv til Religion og den
Maade, paa hvilken Menneskene forestiller sig Aarsagen
til, hvad der som Følge af denne Forudsætning eller
dette Motiv rører sig i deres Sind. Analogien med Sansekvaliteterne
er her fuldstændig træffende. ---

Men heraf følger igen, at Samfundets Modsætning til
Individet og dettes Afhængighed af hint ikke er absolut Betingelse
for Religion. Fællesskabets og Sammenslutningens
Betydning kan ikke skildres mere træffende, end Durkheim
gør det, naar han viser, hvorledes Samfundet gør
sig gældende i Individets inderste Indre, uden at det
selv er sig det bevidst. Men deraf følger ikke, at al Erfaring
om Værdier og al Tro paa Værdiens Bestaaen er
knyttet til Samfundet.

Samfundsbevidstheden (la conscience collective) er
ifølge Durkheim en i sin Art enestaaende Syntese af
enkelte Bevidstheder. Og det samme kan siges om selve
Samfundet: la société est une synthèse des consciences
humaines\footnote[1]{\textit{Les Formes Élémentaires de la Vie Religieuse}. p. 605; 615.}. Men overalt, hvor der paa Aandsvidenskabens
Omraade tales om Syntese, staar vi overfor en
ejendommelig Antinomi. Syntese er Sammenfatten og
forudsætter forsaavidt en Mangfoldighed af Elementer,
som sammenfattes, en Mangfoldighed, der er given før
% --- p. 132 ---
\setcounter{footnote}{0}%
\opage{132}Sammenfatningen. Men paa det aandelige Omraade
kende vi kun Elementerne som Led i den Sammenhæng,
der skyldes Syntesen. Psykologisk kender vi saaledes
kun Fornemmelser, Forestillinger, Følelser og Drifter
som Led i det hele Bevidsthedsliv: Syntesen selv er den
stadige Forudsætning, ikke blot et Resultat. Og paa
lignende Maade i Sociologien: vi kender kun Individer
som socialt bestemte og betingede, som Led af et Samfund.
Der er her, som iøvrigt paa alle Omraader, en
stadig Brydning mellem Kontinuitet og Diskontinuitet,
som vi aldrig kan komme ud over. Vi har ikke Ret
til at lægge udelukkende Vægt paa det ene eller det
andet af Leddene i dette Modsætningsforhold. Forholdet
mellem Samfund og Individer er kun et specielt Ekssempel
paa Forholdet mellem Kontinuitet og Diskontinuitet.
Naar Samfundet er en Syntese af Individer, maa
disse dog paa hvert Punkt have en vis Selvstændighed
i Forhold til det. Uden en saadan Selvstændighed kunde
Samfundslivets Værdi ikke erfares af de enkelte Individer
eller gøre sig gældende i deres Indre, hvilket dog
ifølge Durkheim er Religionens Forudsætning. Den historiske
Sammenhæng, Individet lever sig ind i, er selv
bleven til som Resultatet af mangfoldige individuelle
(skønt stadig socialt betingede) Kræfters Samvirken gennem
Tiderne.

Durkheim svarer hertil, at Individualismen er sekundær.
Saaledes er f. Eks. den individuelle Totemisme,
hvor en enkelt har sit særegne Totem, afledet i Forhold
til den sociale. Dette gælder dog kun, naar man
tænker paa en bevidst og principiel Modsætning mellem
den enkelte og det bestaaende Samfund. Men Talen er
her blot om de individuelle Kræfter, i hvis Vekselvirkning
selve Samfundet egentlig bestaar. Disse Kræfter
kan ikke paa noget Punkt udelukkes fra Betragtning.
% --- p. 133 ---
\setcounter{footnote}{0}%
\opage{133}Ethvert nyt Element i Samfundets Former og Overleveringer,
hellige og profane, --- enhver ny Forestilling
eller Skik, --- maa skyldes et eller flere enkelte Individer,
efter hvilket det er blevet optaget af de andre. Der
er en stadig Brydning mellem det Individuelle og det
Fælles. Durkheim finder Studiet af en saa elementær
Form for Religion som den centralaustralske særlig
gunstig stillet, fordi her individuelle Forskelligheder
ikke gør sig gældende. Men det er den nu bestaaende,
fuldt og færdig udviklede centralaustralske Religion,
der beskrives i Kilderne. Hvorledes kan den være bleven
til? Den maa jo have haft sin Udviklingshistorie
ligesom alt andet i Verden, og her maa det Sporadiske,
Begyndelsen ud fra flere forskellige Udgangspunkter
ligesaa vel have gjort sig gældende som ved den
organiske Vækst og ved Arternes Udvikling. At man nu
ikke bestemt kan paavise det Sporadiske og det Individuelle
berettiger ikke til at antage, at det har manglet
forud eller helt skulde mangle nu. Vi véd ikke, hvem
der har digtet vore gamle Folkeviser, men der maa dog
have været enkelte Personer, som først har undfanget
dem, det være i nok saa ufuldkommen Form. Senere
er de blevne gentagne, ændrede og omdigtede paa mange
Maader. Individuelle Kræfter dels receptive, dels produktive,
kan ikke have manglet paa noget Punkt, hvor
langt vi end gaar tilbage. Selve den »menneskelige Trang«
(besoin humain), Religionen skylder sin Opstaaen og sin
Bestaaen, rører sig neppe i samme Grad og paa samme
Maade hos alle. Og naar den har ført til Udformning
i Kultus, Myte og Legende, maa det være sket gennem
Erfaringer, Kombinationer, Forskydninger og Udskydninger,
altsaa gennem en Række psykologiske Processer,
der maa være foregaaede i de enkelte Individers
Indre. Ligesom Guderne nærer sig af Menneskenes Bøn%
% --- p. 134 ---
\setcounter{footnote}{0}%
\opage{134}ner og Ofre, saaledes nærer Samfundet sig af de enkelte
Indskud i Samfundsbevidstheden (conscience collective).
Saadanne Indskud gøres, dels receptivt, gennem Tilslutning
til og Deltagelse i den en Gang for alle etablerede
Kultus, dels aktivt, gennem nye Erfaringer, Drifter og
Ønsker, der --- gennem den ogsaa af Durkheim bemærkede
Trang til at meddele sig og udbrede sin Tro ---
søger at faa social Betydning. Allerede den individuelle
Totemismes Opstaaen forudsætter noget saadant. Og naar
Totemismen hos nogle Indianerstammer i Amerika fremtræder
mere udformet, tildels ogsaa nærmere Opløsningen,
end i den elementære Form, der kendes fra Centralaustralierne,
saa forudsætter dette en Række Forandringer,
der maa være begyndt som individuelle Variationer,
hvor saa igen Aarsagen til disse Variationer maatte
være at finde.

Det er end ikke sikkert, at individuelle Afvigelser fra
fælles Tro ikke skulde kunne opdages allerede paa et
Trin som Australiernes. Frazer omtaler saaledes Skeptikere
blandt Papuanerne paa Øen Kiwai, der ikke
delte Stammens Tro paa »de Dødes Land«, hvor Sjælene
kommer hen efter Døden. Ligeledes omtaler han,
at der indenfor en Stamme kunde tvivles om Rigtigheden
af en anden Stammes Tro\footnote[1]{\emph{Frazer}: \textit{The Belief in Immortality and the Worship of the Dead}. I. London 1913. p. 213; 239.}. Havde vi flere og
nøjere Kundskaber om de mere primitive Folks aandelige
Liv, vilde vistnok flere Eksempler paa individuelle
Afvigelser og Initiativer fremkomme. De kan ikke nogensinde
ganske have manglet.

Hvad der saaledes ikke tør nægtes, selv for de mest
primitive Samfunds Vedkommende, træder naturligvis
mere og mere frem under Udviklingens Forløb, mere
i de Religioner, vi --- netop af denne Grund --- kalder
% --- p. 135 ---
\setcounter{footnote}{0}%
\opage{135}de højere, end i dem, vi kalder de lavere. Men saa er
det ogsaa ikke blot berettiget, men en nødvendig Opgave
for Religionsvidenskaben at følge det religiøse Liv
ind i de enkelte Individualiteters Indre, hvor det har sin
egentlige Eksistens. Det er et vigtigt Materiale for Religionsforskningen,
der kan hentes fra den Maade, paa
hvilken især fremragende Personligheder, bevidst eller
ubevidst, har udformet det religiøse Liv, selv om de
tilsidst beherskedes af historisk overleverede Former.
Det er til saadanne Dokumenter, jeg især har støttet
mig i min Religionsfilosofi\footnote[1]{I min Bog \textit{Den menneskelige Tanke} har jeg p. 385—392 kort angivet min Opfattelse af Forholdet mellem Religionshistorie og Religionsfilosofi.}. Jeg kan ikke andet end
finde det interessant, at der ad denne Vej er naat Resultater,
til hvilke ogsaa Durkheims Undersøgelser i sine
Konsekvenser fører. Paa Grund af det stadige Vekselforhold
mellem det Individuelle og det Sociale kunde
man heller ikke vente andet, end at man ud fra de to
Udgangspunkter maatte komme til beslægtede Resultater.

Ifølge Durkheim forudsætter Psykologien Sociologien.
Saaledes skal Modsætningen mellem Sansning og Forstand
saavel som den mellem Egoisme og Pligt finde
sin Forklaring ved den socialt betingede Modsætning
mellem Helligt og Profant. Men Psykologien lader sig
ikke nøje hermed. Den spørger om, hvorledes da selve
Modsætningen mellem Profant og Helligt opstaar, og
derefter fremdeles, hvorledes Modsætningen mellem Socialt
og Individuelt opstaar. Og Psykologien har her en
Allieret i Biologien. Der maa være et uvilkaarligt, instinktivt
Grundlag for denne sidste Modsætning. Det har været
skæbnesvangert for Sociologien, at dens Grundlægger
Auguste Comte eliminerede Psykologien fra Videnska%
% --- p. 136 ---
\setcounter{footnote}{0}%
\opage{136}bernes Række og fordelte de psykologiske Fænomener
mellem Biologien og Sociologien. Naar nu Sociologien
optræder paa egen Haand, udenfor den systematiske
Sammenhæng med andre Videnskaber, som Comte lagde
Vægt paa, kommer den let til at overse de psykologiske
Betingelser, som allerede Biologien fører til at indskærpe.
Comte selv førtes i Virkeligheden, som det især ses af
Udtalelser i hans Breve til Stuart Mill\footnote[1]{\textit{Lettres d'Auguste Comte à Stuart Mill}, Paris 1877. p. 275.}, til sin Lære
om Menneskets sociale Natur gennem Studiet af Hume
og Adam Smith. Her ses Sociologiens Gæld til Psykologien.
Det er ligeledes gennem psykologisk Analyse af de
nye Dokumenter, der nu stod til Raadighed, at Durkheim
er naat til de Resultater, han fremstiller i sin Bog.
Dokumenterne give i og for sig kun ydre Fakta. Kun
ad psykologisk Vej kan man naa tilbage fra den intellektualistiske
Opfattelse af Religionen til den »menneskelige
Trang«, hvori Durkheim med Rette finder Religionens
egentlige Grundlag. Det er endelig kun for en psykologisk
Opfattelse, at den religionsfilosofisk saa vigtige Forskel
mellem Oplevelse (religiøs Erfaring) og Tydning
kan gennemføres paa saa konsekvent en Maade, som
det er sket i »Les Formes Élémentaires de la Vie Religieuse«.

Flere Steder bygger Durkheim ligefrem paa bestemte
psykologiske Love. Saaledes forklarer han Totemismens
kosmologiske Anvendelse ved, at »de religiøse Kræfter,
som skyldes Forestillinger og Følelser, der vækkes ved
Samfundslivets Indflydelse gennem Kultus, paa Grund
af deres Intensitet har en Tendens til at brede sig til
alt, hvad der kommer Mennesket i Hænde eller for
Øje, medens hine Kræfter raade.« »Smitsomhed« (contagiosité)
kalder Durkheim dette. Den er et med, hvad
% --- p. 137 ---
\setcounter{footnote}{0}%
\opage{137}jeg i min »Psykologi« har kaldt »Følelsens Expansion«,
som først er beskrevet af Hume og Beneke. Durkheim
bemærker med Rette, at Forklaringen af, hvorledes enhversomhelst
Sammenhæng kan være nok til at overføre Hellighedskarakteren
til nye Genstande, ikke kan søges i
Lovene for Forestillingernes indbyrdes Association; det
er en dybere liggende Lov, man her maa gaa tilbage
til. Det er fremdeles en psykologisk Lov, Durkheim bygger
paa, naar han paaviser en Tendens til at identificere
alle Ting, der vækker en og samme Følelse.

Ved at være opmærksom paa disse psykologiske Love
kaster Durkheim Lys over, hvad Comte kaldte Erkendelsens
teologiske Stadium, og i det hele over den primitive
Erkendelses Natur. Han berigtiger tillige Levy-Bruhl's
Antagelse, at primitive Mennesker skulde have en helt
anden Logik end de civiliserede, fordi hos dem en »Delagtighedens
Lov« (loi de participation) traadte istedetfor
»Selvmodsigelsens Lov«. Man forstaar nemlig f. Eks., at
en Australier kan sige om sig selv, at han er en Kænguru,
idet han identificerer, hvad der paavirker ham
ens: Forestillingen om hans Totem (Kænguruen) og
Forestillingen om hans egen Person. Identifikationen af
at være og at betyde strækker sig gennem hele Religionens
Historie. Den var Stridspunktet mellem Luther og
Zwingli i Nadverlæren, og den er endnu stadigt Stridspunkt
i vore Dages Brydning mellem Dogmatisme og
Symbolisme.

Naar ogsaa Forskellen mellem Sansning og Forstand
skal føres tilbage til religiøs-social Indflydelse, saa maa
dertil bemærkes, at denne Forskel ikke har den skarpe
Karakter, Durkheim tillægger den. Sansningen indeholder
allerede Forstanden i Spiren, forsaavidt den er Evne
til at skelne og at genkende, og der er mange gradevis
forskellige Mellemled mellem denne i Sansningen inde%
% --- p. 138 ---
\setcounter{footnote}{0}%
\opage{138}holdte elementære Tænken og hvad man i snevrere Betydning
kalder Tænken. Psykologien søger at finde en
Kontinuitet her, ligesom paa Sjælelivets andre Omraader.
Under Diskussionen af Durkheims Teori i det franske
filosofiske Selskab gjorde Delacroix allerede opmærksom
herpaa\footnote[1]{\textit{Bulletin de la Sotiété Française de Philosophie}. 1913. p. 79.}. --- Særlig antager Durkheim, at den Forudsætning,
at det Lige frembringes af det Lige, der betegner
en væsentlig Side af Aarsagssætningen, skulde
skyldes Indskærpelsen af, at de hellige Handlinger stedse
maa foretages paa samme Maade, for at de kan frembringe
den ønskede Virkning. Men der har dog ellers i
Livet, i Mellemrummene mellem de rituelle Sammenkomster,
der saa afgørende skal have bestemt Menneskenes
Tankegang, været Lejlighed til at faa hin Forudsætning
praktisk indskærpet. Naar saaledes en Mand,
om hvem Strehlow\footnote[2]{\textit{Die Arunda und Loritja-Stämme in Zentral-Austrálien}. III, 2. Frankfurt 1911. p. 69.} fortæller, tærskede de Planter,
han havde samlet, med en Stok, rensede Frøet ved at
sigte det, rystede det, ristede og knuste det paa Stene
og derpaa spiste det, --- saa er der her en hel Række
Handlinger, som tilsigtede et bestemt Maal, og som
maatte gøres paa en bestemt Maade, for at dette Maal
kunde naas. Der er anden Praksis end den religiøse,
og Kultus er ikke den eneste Lejlighed til at erfare, at
man maa gaa aldeles bestemte Veje for at naa til bestemte
Maal. Livet giver mangfoldig Anledning til at
staa overfor den Nødvendighed, der først træder Mennesket
imøde i Form af det faste Forhold mellem Middel
og Formaal.

Durkheim gør et ejendommeligt Forsøg paa at udlede
den menneskelige Tankes Grundformer (Kategorierne)
fra religiøs-social Indflydelse. Naturligvis foregaar
% --- p. 139 ---
\setcounter{footnote}{0}%
\opage{139}Tankelivets Udvikling under stadig Indflydelse af den
Drøftelse, der kun er mulig under Samkvem og Samliv.
Og Samfundslivets Former og Rubriker, særligt Inddelingen
i Stammer og Klaner, har ganske naturlig afgivet
Rubriker, der uvilkaarlig anvendtes, naar det gjaldt
at ordne og klassificere Iagttagelser og Oplevelser. Ad
Analogiens Vej (sur le modèle des choses humaines)
overførtes de sociale Rubriker til alle Ting, man i Erfaringen
fik med at gøre. Ligesaa lidt her som paa saa
mange andre Punkter skelnes der fra først af mellem
Analogi og Identitet. Men heri ligger ikke andet og ikke
mere end, at sociale Fænomener var de første Objekter,
paa hvilke Sammenligning og Inddeling anvendtes
med særlig Interesse. Og den Udvikling, som Kategorierne
er undergaaaede, efterat hin første Analogi ophørte
at være bestemmende, er langt betydningsfuldere
end den, der fandt Sted i Klan- og Totemforestillingernes
Periode. Religionshistorien kan ganske vist yde et
interessant Bidrag til Kategoriernes Historie, men et Bidrag,
der ikke i Værdi kan maale sig med det, der kan
ydes af Videnskabernes Historie. Det er gennem Videnskabernes
Udvikling i Vekselvirkning med Erfaringen,
at Kategorierne efterhaanden faa deres rationelle Udformning,
deres Begrundelse og deres Begrænsning.

Selv om Udledningen af Kategorierne af religiøs-sociale
Indflydelser var historisk rigtig, vilde den dog ikke
afgive nogen Begrundelse af deres Gyldighed. Gyldighed
og Autoritet falder ingenlunde sammen. De første Eksempler
paa en Tankeforms Anvendelse har intet Fortrin
for de senere. Det afgørende er den Lov, der raader
i hele Rækken af Eksempler. Den religiøse Tanke
maa naturligvis frembyde samme Grundtype som den
videnskabelige, da de begge ere menneskelige Tanker.
At hin tildels gaar forud for denne, giver den ingen
% --- p. 140 ---
\setcounter{footnote}{0}%
\opage{140}særegen Interesse ud over den historiske Interesse. Der
er en Analogi mellem sociale og fysiske Forhold; men
ved selve denne Analogi lægger den menneskelige Tankes
ejendommelige Beskaffenhed sig for Dagen. Til fuld
Bevidsthed om Kategoriernes Væsen behøves der en
lang Række af Eksempler fra de mest forskellige Omraader.
Her er stadig nyt Arbejde at gøre.

Hvad tilsidst Religionens, eller rettere det religiøse Problems
Fremtid angaar, ser jeg den paa lignende Maade
som Durkheim. Jeg er overbevist om, at Religionen,
som vi kende den gennem Historien, hænger sammen
med et væsentligt Element, en dyb Trang i den menneskelige
Natur, og at den derfor ikke kan forsvinde
uden Ækvivalent, dersom Aandslivets Energi og Koncentration
ikke skal svækkes. Den aandelige Arbejdsdeling,
Sondringen mellem Videnskab, Kunst, Moral og
Tro, har gjort Religion til et Problem. Forskellen mellem
Durkheim og mig er især den, at Durkheim venter
sig den nye aandelige Livsform af kraftige sociale Sammenslutninger
og den Entusiasme, de maatte formaa at
vække, medens jeg har det stille Haab, at der allerede
nu i de individuelle Livsførelser og Livsanskuelser er
Spirer tilstede, der engang vil kunne udvikles og tage
Skikkelse som Forbilleder, repræsenterede af mægtige
Personligheder, saadanne som Legenderne skildrer ved
de højere Folkereligioners Grundlæggelse.

Allerede for 70 Aar siden udtalte en svensk Tænker,
Erik Gustaf Geijer, at det Individuelle nu var det eneste
Levende i Religionen. Hvis det er sandt, maa ogsaa
Muligheden af et koncentreret Aandsliv i Fremtiden ---
hvad enten man vil kalde det Religion eller ikke ---
bero paa Personlighedsprincipet. Man kommer ikke
udenom dette. Men ud fra det, gennem den Skærsild,
enhver Sandhed, der har Betydning for det personlige
% --- p. 141 ---
\setcounter{footnote}{0}%
\opage{141}Liv, maa gennemgaa af Uro, Tvivl og Søgen, vil der
kunne fremgaa Nydannelser af højeste Værdi. Den, der
tror paa »Værdiens Bestaaen«, maa i hvert Tilfælde haabe
dette.

Jeg skilles fra Durkheim's Bog med Følelsen af, hvor
meget jeg har lært af den, baade der, hvor mine Tanker
nærmer sig til hans, og hvor jeg har maattet tage
Afstand. Den betegner efter min Overbevisning et stort
Skridt fremad, maaske et Vendepunkt i religionsfilosofisk
Forskning.

% --- p. 142 ---
\setcounter{footnote}{0}%

\opage{142}\essayhead{OM TRØST}

\romanhead{I}

at gennemgaa, hvad der til forskellige Tider har trøstet
Mennesker, vilde være at skrive hele Religionens
Historie, ti Trøst er altid et Element i Religionen,
skønt et Element, der kan indtage en meget forskellig
Plads i Forhold til den enkelte Religions andre Elementer.
Og netop denne Forskel mellem Trøstelementet
og Religionens andre Elementer er af stor Betydning til
Karakteristik af en bestemt, foreliggende Religion. I al
Almindelighed kan det siges, at Trøstelementet i Løbet
af Religionens Historie er blevet mere og mere fremherskende.
Mere og mere gaar Menneskelivet op i
Arbejdet paa materiel og aandelig Kultur, og det er
dette Arbejdes Fremgang eller Tilbagegang, de Grænser,
der møder det, eller de nye Muligheder, der opdages for
det, som bestemmer menneskelige Følelsers og Villiers
Karakter og Retning. Kun naar Sindet føler sin Afmagt
og maa lade Hænderne synke, træder Religionen i Kraft
og kan da yde den Tilfredsstillelse eller Husvalelse, der
forgæves søgtes i selve Arbejdet. Arbejdet selv har Religionen
intet at gøre med direkte; dets Opgaver stilles
fra en helt anden Egn end den, hvor Religionen har
hjemme.

I Religionens klassiske Tider er det anderledes. I dem
% --- p. 143 ---
\setcounter{footnote}{0}%
\opage{143}træder Trøstelementet ikke i Forgrunden, men Religionen
er en positiv Faktor i Livet, stiller selv de afgørende
Opgaver, tilsidst den eneste Opgave, der anerkendes. Til
de klassiske Tider hører først og fremmest de store Begyndelsestider
(som vi kender dem i de højere Religioner).
I den oprindelige Buddhisme og i Urkristendommen
er Religionen Indbegrebet af al aandelig Kultur,
samler alt om et eneste Maal. Til de klassiske Tider
hører ogsaa de store Organisationstider, hvor Religionen
vel ikke er eneste Form for Kultur, men hvor den drager
al aandelig Kultur ind under sig, idet den staar
som højeste Værdimaaler for al menneskelig Stræben.

I Begyndelsestiden er Trøsten given i selve Tilslutningen
til det ene Fornødne. Der er kun Tale om én
Trang, og kan den tilfredsilles, er alt, som det skal
være. I Organisationstiderne er der flere Værdier at
værne om, og de vil aldrig helt være betingede ved Religionen
alene. Der er mange Fornødenheder foruden
det, man stadig mener at have Ret til at kalde det ene
Fornødne. Den middelalderlige Kristendom (i det 13.
Aarhundrede) er det største Eksempel paa en saadan
Organisationstid. Men der traadte allerede Videnskab og
Statsliv, ja, selve Kirkelivets ydre Ordning frem som
Tendenser, der var deres egne Love underlagte og havde
deres egne Skæbner. Moderne Kristendom kan man ikke
kalde nogen Organisationsform, skønt man i det 19. Aarhundredes
Romantik forsøgte at opfatte den saaledes. Den
vedkender sig mere eller mindre Spaltningen mellem de
forskellige Livsomraader. Det er da ikke saa let at finde
Trøst; er man trøstet, hvad det Væsentlige angaar, saa
er der dog saa meget, der ogsaa var en Hjertesag, og
stadig staar uafgjort og ufuldendt. I en Brevveksling
mellem to danske Præster omkring Aar 1900 træder
dette tydelig frem; de længes efter Kristi Genkomst, men
% --- p. 144 ---
\setcounter{footnote}{0}%
\opage{144}hvis nu den store Begivenhed snart indtræder, hvad
bliver der da af den kirkelige og folkelige Udvikling,
af hvilken en herlig Frugt af Menneskelivet skulde fremgaa?
(Se i det følgende Afhandlingen »To Præster«).

Hvad der saaledes fremtræder i Religionshistorien,
gælder ogsaa for Menneskelivet idethele. Der er en Trøst,
der er saa uadskillelig knyttet til, hvad der er Livets
Hovedinteresse og Hovedvirksomhed, at man egentlig
ikke længere kan tale om Trøst, naar blot denne Interesse
kan fyldestgøres og denne Virksomhed fortsættes.
Det vil da komme an paa, hvor inderlig alle andre Interesser
og Bestræbelser, der rører sig hos de enkelte,
kan blive Led eller Elementer i Hovedbestræbelsen. Jo
inderligere Forholdet er, des mindre vil de enkelte Delinteresser
kræve en særegen Trøst. Trøsten betegner da
overalt en Hvile, en Vederkvægelse, der er Betingelse
for fortsat Arbejde. Trøstens Nødvendighed betegner
Nødvendigheden af Opsamling og Udløsning af Kraft,
før der gaas videre. Trøsten har da, kan man sige, en
aktiv, ikke en passiv Karakter. Hvad der modtages, er
Grundlag for ny Selvvirksomhed. Trøsten staar her kun
som Formaal, fordi den kan blive et nødvendigt Middel.
Det sidste Maal kan den ikke blive, kun et Led i
Livets store Rytme. Selv den mest energiske Livsvirksomhed
kan ikke forløbe ud ad en eneste ret Linie,
men vil være som en snart stigende snart synkende
Kurve.

Ved en forunderlig Modsigelse staar alligevel endog
i meget energiske Livsanskuelser og Livsførelser passiv
Trøst som det Højeste, det Definitive, Maalet, til hvilket
der ses hen. I Saligprisningerne i Evangelierne raader
denne Tendens, mere udpræget hos Lukas end hos
Matthæus. »Himmeriges Rige« bringer den evige Trøst
efter den nuværende Tids Lidelser. I den populære Re%
% --- p. 145 ---
\setcounter{footnote}{0}%
\opage{145}ligiøsitet, hvor Forventningen om det store Gennembrud
ikke bestemmer hele Livet, og hvor man har glemt
Verdensdommen, der bringer Afgørelsen, er »den evige
Hvile« bleven det gængse Udtryk for det hinsidige Liv.
Herved stanser de Flestes Tanke. Selv for Kierkegaard,
denne Stridens Mand, stod det som det Højeste, at den
hele Strid med ét var forsvunden. Og i mere gemytlige
Former for Hinsidighedstro udmaler man sig det andet
Liv som en Fortsættelse af det nuværende Livs Ferier
og Vennemøder. Nogle af den nyere Tids største Aander
har udtalt sig mod den Art Trøst. Allerede Leibniz
sagde, at hvis Saligheden ikke bestod i Fremskriden,
vilde de Salige blive stumpe og sløve. (Nisi beatitudo
in progressu consisteret, stuperent beati). Kant's og
Goethe's Udødelighedstro gik i samme Retning. Efter
en saadan Opfattelse kan Trøsten aldrig blive definitiv.
Ti hvor der stadig virkes, er stadig ny Modstand mulig
(at virke er jo netop at overvinde Modstand), og nye
Skæbner melder sig alt efter Forholdet mellem Kraft og
Modstand. Gælder det absolut Trøst og Hvile, saa er
Buddhismen med sit Nirvana og med sin Stræben efter
Befrielse fra Gerningernes Verden (Karma) mest konsekvent.

Men hvis det er saa, at sand Trøst kun faas, naar
Virksomhed kan udløses, saa at kun den, der kan være
selvvirksom, kan trøstes, saa ligger her Mulighed for et
tragisk Misforhold. Ti da er der ingen sand Trøst for
dem, for hvem Arbejde er en Umulighed. Den Affældige,
hos hvem Erindringen om Virksomhedens Tid ikke er
helt forsvunden, kan da kun være modtagelig for rent
passiv Trøst. Eller han maa i hvert Tilfælde lade sig
nøje med en mere elementær og smaatskaaren Virken,
føre sit Liv paa en begrænset Skueplads, maaske i grel
Modsætning til det tidligere Livs store Arena. Saadan
% --- p. 146 ---
\setcounter{footnote}{0}%
\opage{146}Smerte og Disharmoni opstaar ofte, naar Livet hæmmes
uden dog at opløses helt. Smerten er jo overalt
Tegn paa, at der endnu er Livskraft tilstede, skønt ikke
nok til normal Fungeren. Overfor den fuldstændige Affældighed
trænger Individet ikke til Trøst, da han ikke
længere har Kraft til Sorg og Savn; der er det de andre,
der trænger til Trøst, især fordi de ingen Hjælp kan
bringe.

Selv de, der mener, at Trøsten beholder, og bør beholde,
det sidste Ord, kan dog indskærpe, at man ikke
skal trøstes i alle Tilfælde, og at Trøsten kan indeholde
Fare. Forfatteren af De imitatione Christi (Om Kristi
Efterfølgelse) skelner mellem den ydre Trøst, som f. Eks.
faas gennem Samtale med Venner, ved Rejser eller ved
Musik, og den indre Trøst, som bestaar i at tænke over
de himmelske Ting (speculari coelestia) og i indre Jubel
(jubilare interius). Man griber efter hans Mening
altfor hurtig til den ydre Trøst. »Naar et Menneske begynder
at sløves, frygter han en lille Anstrengelse og
griber gerne til en ydre Trøst«. Den indre Trøst forudsætter
altsaa et Arbejde. Men den gamle Mystiker gaar
et Skridt videre. »Det er«, siger han, »saare stort at kunne
undvære baade menneskelig og guddommelig Trøst og
for Guds Æres Skyld gerne at finde sig i aandelig Udlændighed
(exilium cordis)«. Dette er vel at forstaa saaledes,
at »Udlændigheden« siden efter viser sig at være
»til Guds Ære«, ikke saaledes, at den, der undværer al
Trøst, samtidig skulde være klar over, at det er »til
Guds Ære«; ti i saa Fald undværer han jo ikke al
Trøst, er ikke stedt i fuld »Udlændighed«. Og selv hvor
Mennesket ikke er i Udlændighed, men føler den guddommelige
Kærlighed, maa denne Følelse ikke gøres til
blot Middel for egen Trøst. Man skal elske Jesus for
hans egen Skyld, ikke for sin egen Trøsts Skyld. My%
% --- p. 147 ---
\setcounter{footnote}{0}%
\opage{147}stikeren antyder her Muligheden af en Forskydning i
sentimental og egoistisk Retning, som netop kan begunstiges,
naar Trangen til passiv Trøst faar Magt, og som
er karakteristisk for mange Former af moderne Religiøsitet.
Ofte vedkender man sig udtrykkelig, at man kun gaar
ind paa religiøse Antagelser, fordi man ikke kan undvære
den Trøst, de indeholder. Middelalderens Tankegang
er her mere storslaaet. Det Guddommelige stiller
sine Fordringer uden Hensyn til Individets private Befindende
og Sympati. Thomas af Aqvino kunde saaledes
meget vel tænke sig, at et eneste Menneske kunde blive
salig, medens alle andre fordømtes, og at hans Salighed
ikke vilde formørkes ved deres Skæbne. Salighed bestaar
nemlig i »at have sin Fuldkommenheds Fylde
i Gud«, --- og overfor denne Fylde er enhver Tanke
om endelige Væsener umulig\footnote[1]{Se min \textit{Religionsfilosofi}. Note 105.}.

Etisk set vil det være Opgaven at fastholde den aktive
Type som Forbillede, og at sørge for, at passiv
Trøst ikke behøves. Villien skal, som jeg har udtrykt
det i min Etik (X, 3), øves i at staa paa sine egne Ben.

\romanhead{II}
Drengen i »Flugten til Amerika« havde haft en Række
af Sorger, men han »drukned sin Sorg og fandt sin
Trøst paa Bunden af Sagosuppen«, og Jeppe mente sig
overfor sin megen Elendighed berettiget til at »bruge
de Midler, Naturen giver os til at bortdrive Sorg med«.
Men naar Trøsten ikke skal have denne passive Karakter,
men skal bestaa i Udløsning af ny Virksomhed,
saa forudsætter den en aandelig Vækst. Der maa være
Mulighed for ny Kraftudløsning, snart en mere vilkaarlig,
snart en mere uvilkaarlig. Der kan maaske med fuld
% --- p. 148 ---
\setcounter{footnote}{0}%
\opage{148}Bevidsthed vælges et Arbejde, som ikke i sig selv har
noget at gøre med Sorgen og dennes Genstand. Trøsten
bliver da, ligesom hos Drengen og hos Jeppe, Trøst ved Siden
af. Sorgen afledes ved Koncentration af Energien i
en helt anden Retning, og Sorgens Ekspansion i Sindet
hæmmes eller hindres. Det kan være Rejser eller det
daglige Livs Opgaver og Pligter, der lægger Beslag paa
Interesse og Energi. Eller Trøsten findes i Kunst, Videnskab
eller Religion. Oplevelsen søges fortrængt ved nye
Oplevelser af ydre eller indre Art. Man gaar uden om
det saarede Sted. Men det er ogsaa muligt, at Trøsten
vindes ved, at Eftertanken netop rettes mod selve Tabet
og Sorgen. Da gør man sig Værdien af det Tabte
levende, langtfra at ville forringe den ved Hengivelse
til andre Værdier\footnote[1]{\emph{Seneca} opfordrer Marcia, hvem han vil trøste over hendes Søns Død, til at laane Øre til Sønnens Navnkundighed og Ihukommelse. (Apertos aures præbe ad nomen memoriamque filii tui. \textit{Consolatio ad Marciam}. V, 3. (Dialogi ed. Gertz p. 170).}. Sorgen uddybes og gennemtrænger
hele Sindet, kommer til at løbe Linen ud. Det er ikke
fordi man vil nyde Sorgen (hvad Petrarca kaldte dolendi
voluptas, og Amiel kaldte douceur mélancolique),
heller ikke finder man en Ære i at være særlig ulykkelig
fremfor andre\footnote[2]{\emph{Seneca} advarer Marcia mod den forkastelige Ære, at ville staa som den mest ulykkelige af alle. (perversissima gloria, infelicissima videri). \textit{Consolatio ad Marciam}. V,4. (ed. Gertz p. 170).}. Sammenligning med andre behøver
slet ikke at spille nogen Rolle, hverken for at give
den egne Sorg en forherligende Baggrund, eller for at
indskærpe, at andre har haft større Sorger, eller for at
se Sorg som fælles Skæbne. Hvad der raader, er Trang
til at fordybe sig i, hvad der er sket. Kun derigennem
kan maaske den Forstenelse, der er Tabets første Virkning,
opløses. Der søges fuld Udladning for Tilstanden,
% --- p. 149 ---
\setcounter{footnote}{0}%
\opage{149}aabnes Vej for den Diffusion, den Betagethed af hele
Organismen, som kan være nødvendig. I Eddadigtet
»Første Sang om Gudrun« beskrives paa en gribende
Maade, hvorledes fornemme Kvinder, »Jarlers ædle Hustruer«
sad hos Gudrun og søgte at trøste hende ved
at fortælle om deres Sorger. »Hver af dem sagde sin
egen Sorg, --- den bitreste, som hun baaret havde\dots{}
Dog kunde Gudrun ikke græde, saa mod var hun i Hu
over den døde Mand, saa fuld af Smerte ved Sigurds
Ligfærd«. »Gudrun sad over den døde Sigurd. Hun græd
ikke som andre Kvinder, men hun var nær ved at briste
af Græmmelse«. Men da bebrejder »Gjukes Datter« de
andre Kvinder, at de ej forstod »den unge Viv at trøste«;
hun afdækker Sigurds Lig og lægger Gudruns Hoved
mod hans Kind. Og da græd Gudrun.

Ikke blot overfor et Tab, ogsaa overfor det, man
aldrig har ejet, men inderlig ønskede at vinde, kan
Diffusionen være en Trang. Det er det, Goethe skildrer
som »Trost in Tränen« og »Wonne der Wehmuth«.

Er baade Værdien af det tabte og selve Tabets Virkelighed
fuldt gennemlevet, kan der hos nogle Naturer være
Trang til at tænke begge sammen og tillige at tænke
dem ind i Livets store Sammenhæng. Hvad man har
ønsket og tilstræbt, hvad man har lidt af Tab eller
Skuffelser, ses da som hørende ind som naturlige, maaske
nødvendige Led i Tilværelsens store Hele. Det Totalbillede,
den »syntetiske Intuition«, som derved bliver
mulig, kan optage Sindet og fylde det med den Tilfredsstillelse,
som aandelig Virksomhed medfører, hvor
den har Fremgang. Spinoza og Goethe gik, hver paa sin
Vis, denne Vej. Livet bliver her omsat til en Forstaaelse,
der vil lade baade Værdien og dens Tab komme
til deres Ret ved at gøre dem til Elementer i en Helhed.
Spinoza betoner særlig den Glæde, selve Forståa%
% --- p. 150 ---
\setcounter{footnote}{0}%
\opage{150}elsen giver, og denne Erkendelsesglæde staar hos ham
som Ækvivalent for den Sorg, der fortones gennem
Trøsten. En saadan intellektuel eller æstetisk Trøst, som
Spinoza har skildret med Genialitet og stille Begejstring\footnote[1]{Smlgn. \textit{Spinoza's Ethica}. Analyse og Karakteristik. (Videnskabernes Selskabs Skrifter 1918). p. 73—77; 87---89.},
og som Goethe\footnote[2]{Smlgn. mine \textit{Mindre Arbejder}. Tredie Række. p. 162---164; 166—169.} har givet Udtryk for i sin store Kunst,
vil dog, skønt den ikke gaar udenom Sorgen, være Trøst
ved Siden af. Eftertankens og Fantasiens Verden er nu
engang ikke Sindets hele Verden. Og selv om der gennem
nok saa dybtgaaende sjælelige Sammensmeltninger
og Organisationer danner sig som Livsresultat en Totalfølelse,
til hvilken baade de store Sorger og de store
Glæder har ydet deres Bidrag, saa er der dog ingen
Sikkerhed for, at denne Helhed ikke igen kan sprænges,
naar Skæbnen banker paa Porten, uden at denne
Banken (som hos Beethoven) er Indledning til en Symfoni\footnote[3]{\textit{Den store Humor}. p. 104,}.
En definitiv Trøst kan ikke nogensinde begrundes,
ligesaalidt ad religiøs som ad filosofisk Vej.

Sagen er, at det Individ, som trøstes, stedse er Del
af et større Hele, hvis Love ikke falder sammen med
de Love, under hvilke hans Eksistens føres. Om enhver
Art af Trøst, --- den passive og den aktive, den vilkaarlig
valgte og den uvilkaarlig fremvoksede, den vemodige
eller den fortvivlede, den intellektuelt overlegne
(gennem syntetisk Intuition) eller den emotionelt overlegne
(gennem Totalfølelse), --- gælder det, at den har en
individuel Karakter. At noget yder Trøst, er et rent
subjektivt Fænomen, et psykologisk Faktum, der ikke
siger noget om Tilværelsens objektive Sammenhæng.
Fordi den enkelte er kommen til Ro, eller ser en Vej
% --- p. 151 ---
\setcounter{footnote}{0}%
\opage{151}aaben for sig, er det ikke sagt, at Tilværelsen er kommen
til Ro, eller at en ny Vej virkelig er aabnet.

Men derfor er det alligevel en Kendsgerning, at paa
dette bestemte Sted i Tilværelsen, i dette enkelte Menneskes
Indre er der naaet Harmoni eller Mulighed for
ny Udfoldelse. Og hvis der herved er gjort fuld og ærlig
Brug af den Kraft, der raades over, saa er der opnaaet
paa dette Sted, hvad der kan opnaas, og hvad
der er opnaaet, har da ogsaa sin objektive Betydning.
Der er sprungen en Blomst ud, hvilken dens Skæbne
saa end siden bliver. Verden har da været saaledes beskaffen,
at denne Blomst kunde udfolde sig paa dette
Sted og paa denne Tid. ---

Hvad der, som Verden nu engang er, kan trøste
Mennesker, kan kun læres af Erfaring. Det er ikke abstrakt
Teori, men arbejdende og lidende Menneskers
Livserfaringer, det gælder om. Og her er et stort og
lidet dyrket Omraade for sammenlignende Undersøgelse.
En Digter kan snarere yde Bidrag i denne Retning end
en Filosof, og i hvert Tilfælde behøver Filosofen et stort
og omfattende Materiale fra egen og andres Erfaring
for at kunne se klart her. Den sidste Opgave for Videnskaben
er ikke den, at finde en objektiv Teori, der passer
paa alle foreliggende Emner, men den at vise, hvorledes
forskellige Mennesker, --- forskellige Vandringsmænd
ad Livets labyrintiske Veje, --- maatte komme
til forskellige Livsopfattelser\footnote[1]{\textit{Den menneskelige Tanke}. p. 304.}. Og herhen hører da
ogsaa Forskellighederne i, hvad der har kunnet trøste
Mennesker. ---

Ogsaa Illusioner kan trøste. Og naar vi ser andre
trøste sig ved, hvad der ikke vilde trøste os, er vi tilbøjelige
til at mene, at deres Trøst hviler paa Illusion.
Men her gælder det, at hvis Illusionen er Frugt af ær%
% --- p. 152 ---
\setcounter{footnote}{0}%
% CHECK p. 152: notes paired with the marks in order (the layer misread a number)
\opage{152}ligt Arbejde, der under de givne Betingelser ikke kunde
føre til andet Resultat, har andre ingen Ret til at gribe
ind, med mindre de raader over Evner og Midler til
at føre den illusorisk trøstede til et objektivt vigtigere
Resultat\footnote[1]{\textit{Etik}\textsuperscript{4}. p. 149 og 197.}.

Ved Spørgsmaalet om, hvorvidt Illusioner kan trøste,
skelner Benedetto Croce med Rette mellem to Tilfælde.
Hvis Illusionen ikke er erkendt som Illusion og dog
yder fuld Tilfredsstillelse, maa der være Grunde, den
støtter sig paa, og den udgør da den Del af den virkelige
Sandhed, der kunde erkendes under de givne Betingelser.
Illusion kan den kun kaldes ud fra en fuldstændigere
Sandhed. Anderledes naar Illusionen erkendes
som Illusion af vedkommende Menneske selv. Den maa
da vække en Angst og en Iver, som driver til at søge
efter Sandheden, og hvis den alligevel fastholdes som
trøstende, maa den fremkalde Uro og slet Samvittighed,
indtil en alvorlig Stræben efter Sandhed opstaar. Iøvrigt
nægter Croce, at Sandheden tilsidst skulde kunne være
sørgelig. Den fulde Sandhed maa jo omfatte baade det
glædelige og det sørgelige og altsaa staa over begge
Dele\footnote[2]{B. \emph{Croce}: \textit{Logica}\textsuperscript{2}. Bari 1900. p. 346---347. --- Hele Kapitlet om »Consolatio philosophiea« er meget interessant.}. --- Til denne sidste Bemærkning er dog at sige,
at det vil komme an paa, om der kan dannes en Helhedsopfattelse,
der anviser baade Glæden og Sorgen
deres bestemte Plads i Tilværelsen. Vi kommer da i
hvert Tilfælde tilbage til Spinoza's og Goethe's Opfattelse.
Og ligesom før maa vi føje til, at det ogsaa kommer
an paa, om der overfor Livets stridende Erfaringer kan
danne sig en Totalfølelse, der lader dem alle komme
til deres Ret. Helt afsluttet foreligger Virkeligheden
aldrig, og Croce's Opfattelse er derfor kun berettiget
overfor en dogmatisk Pessimisme.

% --- p. 153 ---
\setcounter{footnote}{0}%
\opage{153}H. C. Ørsted gjorde i sin Tid (i »Aanden i Naturen«)\footnote[1]{Jeg har tidligere drøftet Ørsted's Ytring i min \textit{Religionsfilosofi} § 115 og i \textit{Danske Filosofer} p. 56.}
Indsigelse mod den Sætning, at vore Ønsker skal kunne
bestemme, hvad der er Sandhed, og krævede, at man
skulde give Sandheden Æren: »den fulde Sandhed bringer
altid sin Trøst med sig«. I denne Ytring fremtræder paa
én Gang den store Forskers Sandhedskærlighed og hans
Optimisme. Men han har ikke set, at Ytringen kan forstaas
paa to Maader. Lægges Vægten paa, at det er den
fulde Sandhed, der trøster, saa kan Spørgsmaalet rejse
sig, om vi nogensinde naar den. I hvert Tilfælde maa
vi, indtil vi har den, trøste os med Haabet. Istedetfor
at have noget at trøste os ved, maa vi da trøste os til,
at fortsat Stræben kan bringe Trøsten. Det er den gamle
Betydning af Ordet Trøst, --- Tillid, Tro, Tilforladelighed,
--- der her træder frem igen i Modsætning til den
nyere Betydning af Ordet, som vistnok overvejende
leder til Forestillingen om noget Passivt, om en Modtagen,
en Husvalelse. --- Lægger man derimod Vægten, ikke
paa Ordet »fuld«, men paa Ordet »sin«, bliver Meningen,
at Sandheden trøster blot ved at være Sandhed, og dette
kan igen enten forstaas saaledes, at Trøsten ligger i den
Uundgaaelighed, der afsløres for os, eller at den vindes
gennem det aandelige Arbejde, der kræves for at finde
Sandheden, altsaa i Form af Erkendelsesglæde. Uundgaaeligheden
fører os tilbage til Vemodens eller Fortvivlelsens
Trøst; Erkendelsesglæden fører os tilbage til
Spinoza og Goethe. Ørsted har neppe tænkt paa nogen
bestemt af disse forskellige Tydninger. Han har snarest
tænkt paa Sandhed som Afsløring af en Virkelighed,
hvor alt viser sig i Harmoni, saa ofte Trøstesløshed end
synes at have det sidste Ord.

% --- p. 154 ---
\setcounter{footnote}{0}%
\opage{154}\romanhead{III}
Henimod Slutningen af den klassiske Oldtid (3. Aarh.
f. Kr.) udviklede der sig en særegen Trøsteliteratur, som
fortsattes gennem den romerske Kejsertid og fandt sin
Afslutning i et Skrift af en Mand, som er bleven kaldt
»den sidste Romer«, Boëthius (d. 575 e. Kr.), en høj
Embedsmand under Østgoten Theodorik den Store.
Gennem Boëthius' Skrift om den Trøst, Filosofien kan
bringe (De consolatione philosophiæ), et Skrift, der blev
meget læst og kommenteret i Middelalderen, strækker
denne Literaturs Indflydelse sig over i den nyere Tid.

Den første navnkundige af disse Trøsteskribenter er
Platonikeren Krantor (i Slutningen af det tredie Aarhundrede
f. Kr.). Hans Skrift »om Sorg« er desværre
gaaet tabt, men dets Indhold kan spores hos senere
Forfattere over samme Emne, Cicero, Seneca og Plutarkh.
Cicero kalder det »den gyldne Bog«. Helt op i
moderne Prædikanters Trøstetaler mener man at kunne
genfinde de samme Tanker, naturligvis uden at Kilden
kendes. Hvad der søges Trøst for, er Alderdom, Landflygtighed,
Tab af de Nærmeste, Fattigdom. Filosoferne,
der dog paa dette Omraade snart overbødes af Retorerne,
satte da Eftertanken i Bevægelse for at vise, at hvad
der var sket ikke var idel Ulykke, maaske slet ikke var
Ulykke. Mange af disse Grunde forekommer os snusfornuftige;
men til given Tid og Sted kan de have gjort
deres Virkning, og ofte har vel Villien til at trøste gjort
større Indtryk end de Grunde, der førtes i Marken\footnote[1]{Smlgn. Constant Martha: Les Consolations dans l'Antiquité. (Etudes Morales sur l'Antiquité 1863. p. 135—189.)}.
I sine Tuskulaner (III, 31) siger Cicero, at han selv i en
stor Sorg (Datterens Død) havde samlet alle de forskellige
Trøstegrunde, han kendte, idet han for enhver Pris
søgte Lægedom for Sindets Brusen (tumor animi). Han
% --- p. 155 ---
\setcounter{footnote}{0}%
\opage{155}opregner: Beroligelse eller Afledning af Sorgen, Paavisning
af dens indbildte Karakter, Henvisning til, at man
har gjort sin Pligt. Og disse forskellige Midler maa hvert
for sig anvendes paa rette Tid og Sted.

Allerede det, at man kom til at tænke over andre
Sider ved det, der var sket, istedetfor at kredse om
selve Ulykken, kunde tjene til Afledning og hindre
Sindet i at stirre paa det Uoprettelige. I sit Skrift om
Alderdommen siger Cicero, at selve Udarbejdelsen af
dette Trøsteskrift (ved hvilket han vilde trøste sin Ven
Atticus og sig selv) har været ham en saadan Fornøjelse,
at den for ham har fejet Alderdommens Besværligheder
til Side og gjort hans gamle Dage blide og behagelige.
Et karakteristisk Udtryk for, at det at tænke over Sorgen
kan aflede den! Det er en Trøst, der naturligvis kun er
mulig, saalænge man endnu har Kraft til Eftertanke.
Men, siger Cicero saa, Affældighed og Svaghed kan jo
ogsaa overgaa unge Mennesker, saa at denne Side af
Sagen ikke særlig hører hjemme i et Skrift om Alderdommen.
Iøvrig henviser han til, at Alderdommen er
det sidste Akt af et Skuespil, som Naturen, der sandelig
ikke er nogen Fusker, har anlagt og gennemført efter
en bestemt Lov. »Det har jo imidlertid ikke kunnet være
anderledes, end at noget maatte være det sidste, og
ligesom det gaar med Træernes Bær, naar de har naaet
deres Modenhed i behørig Tid, saaledes maatte det sidste
Livsafsnit ogsaa blive runkent, om jeg saa maa sige, og
staa for Fald; dette maa da den fornuftige Mand finde
sig i med Sagtmodighed. Ellers fører han nemlig Krig
med Guderne paa Giganternes Vis; for hvad andet betyder
vel dette end netop det: at kæmpe med Naturen
selv\footnote[1]{\emph{Cicero}: \textit{Cato den Ældre. Om Alderdommen}. Gertz' Oversættelse p. 19.}«. Denne Udtalelse kan siges at være typisk for
% --- p. 156 ---
\setcounter{footnote}{0}%
\opage{156}antik Trøst. Den bygger paa en naturalistisk Optimisme,
i hvilken dog for mere stemningsfulde Naturer snart
Bitterhed, snart mild og vemodig Resignation optræder
som indoptagne Elementer\footnote[1]{Se min Karakteristik af Marcus Aurelius i \textit{Mindre Arbejder}. III. p. 115.}.

Hos Kejsertidens Stoikere faar Trøstegrundene en
kraftigere Karakter. Der appelleres paa én Gang til
Villien som det ene væsentlige og til den store kosmiske
Sammenhæng, indenfor hvilken Mennesket med alle sine
Glæder og Sorger kun er et forsvindende Punkt. Seneca
forestiller først Marcia, hvem han vil trøste for hendes
unge Søns Død, at det stemmer bedst med hendes ædle
og fine Karakter selv at overvinde Sorgen og ikke at
vente paa den Dag, da Sorgen mod hendes Villie forlader
hende. Sorgen, tilføjer han, bliver kun saa heftig,
fordi vi i vor Sorgløshed ikke venter Ulykken, hvor ofte
den end gaar forbi os for at ramme andre. Netop det,
der har den største Værdi, vil i Regelen være mest udsat
for Undergang. --- Disse Betragtninger fører da
Seneea til at udvikle, hvor forsvindende menneskelig
Skæbne er i Forhold til det store Univers, indenfor hvilket
Ustadighed og Omskiften raader. Hvor kan dog
nogen --- siger han i en lignende Sammenhæng i et andet
Skrift (Ad Polybium) --- driste sig til at tro, at han og
hans Nærmeste skulde kunne være Undtagelser indenfor
Naturens store Orden, som alt underligger?\footnote[2]{\emph{Seneca}: \textit{Ad Marciam}. (Dialogi ed. Gertz p. 173—201); \textit{Ad Polybium} p. 324.}

Epiktet betoner stærkere end Seneca, at Menneskets
Væsen er indeholdt i det Indre, i Villien. Guden vilde
gerne have givet Mennesket mere end Villien, men har
ikke kunnet, ligesaalidt som han kan berøve det denne!
Men saa er den Herlighed, Mennesket har i sit Indre,
saa meget des større. I betragter det som en Ulykke,
% --- p. 157 ---
\setcounter{footnote}{0}%
\opage{157}at I ikke har set Fidias' Zeus i Olympia --- men i Eders
egen Villie er der noget, som er værdt at skue, og som
ikke kræver nogen lang Rejse. Og derinde bor Guden
og Eders Genius, saa at I aldrig vil blive ene, selv om
I lukker Eders Sjæls Døre! Man skal da kun bekymre
sig om Villien, om at have Herredømme over sine Forestillinger.
Alt andet raader I ikke for. I skal elske Mennesker,
men huske, at de er dødelige. Andre Menneskers
Ulykker, selv Eders Nærmestes, beror ikke paa Eders
Villie. Suk, naar saadan Ulykke kommer, men agt paa,
at I ikke sukker i Eders Inderste! I er en Del af Naturen
og maa bøje Eder for dens Gang. I al Sorg og Jammer
ytrer der sig egentlig en dogmatisk Paastand om, hvad
der er godt og ondt. Den maa bekæmpes, ti ved Verdens
store Skuespil er der ikke Brug for en misfornøjet
Tilskuer. Er du ikke tilfreds, saa gaa din Vej. Sig, at du
ikke leger med mere!\footnote[1]{\textit{Epiktetus}. ed. Schmidt passim.}

Hos Boëthius mødes stoiske og nyplatoniske Tanker i
hans mærkelige Skrift om »Filosofiens Trøst«, om hvilket
en udmærket Literaturhistoriker\footnote[2]{\textit{Fridrich Leo} i Kultur der Gegenwart. I,8. p. 395.} har sagt, at det afspejler
Glansen af Oldtidens nedgaaende Sol som Symbol
paa dens Opgang i Fremtiden. Bogen er forfattet i det
Fængsel, Kong Theodorik havde kastet ham i, og som
han kun forlod for at lide Døden. Der har været Strid,
om han var Kristen eller Hedning. Af Bogens Indhold
kan der ikke sluttes til hans Kristendom; den giver en
rent menneskelig Betragtning i Stoicismens og Nyplatonismens
Spor. Da han skal være Forfatter til teologiske
Skrifter, hvis Uægthed der ikke synes Grund til at antage,
mener man nu i Almindelighed, at han har været
Kristen, »opvoxet i oplyst og dannelsesivrig Kristendom,
som dengang hos Aristokratiet udgjorde Grundlaget for
% --- p. 158 ---
\setcounter{footnote}{0}%
\opage{158}Verdensanskuelsen«\footnote[1]{\textit{Eduard Norden}: Kultur der Gegenwart. I,8. p. 412.}. Han har i hvert Tilfælde i Nød
og Dødsfare kunnet trøste sig uden at tage sin Tilflugt
til Kristendommen. Filosofien besøger ham i hans Fængsel
og forjager Muserne, som kun formaar at nære Sorgen,
ikke at læge den. Den sande Trøst, lærer den ham, bestaar
i Erkendelsen af Universet, der --- som Platon
har vist i »Timaios« --- gennemstrømmes af Guddommens
Aand. Fordi Ulykken kommer over os, har Verdensordenen
ikke forandret sig; Opgang og Nedgang skifter
nu som før for alle endelige Væsener. Men de Dødelige
søger Lykken udenfor, ikke i deres eget Indre, hvor det
højeste Gode, der er Et med Gud selv, har sit Bo. Ingen
Skæbnestilskikkelse kan røve denne indre Ejendom, og
den, der besidder den, behøver ingen ydre Domstol. Vi
lever i en Splittelsens og Ufuldkommenhedens Verden;
vor Jord er kun som et Punkt i Forhold til det hele
Univers, vor Levetid kun som et Øjeblik i Forhold til
den uendelige Tid. Men selv i denne ufuldkomne Tilværelse
kan Mennesket bevare sin Enhed med det Guddommelige\footnote[2]{\emph{Boëthii}: \textit{De Consolatione Philosophiæ}. ed. Obbarius. Jena 1843. (Især II, 3—7; III, 8—10).}.

\romanhead{IV}
% CHECK roman numeral read

Det var Eftertanke over Livet og Verden, de antike
Tænkere tog deres Tilflugt til, naar Trøsten skulde søges.
De krævede en Underordning af Følelseslivet under
Tankelivet, hvis Harmoni og Kraft skulde kunne bevares
under Skæbnens Gang. Naar hos Boëthius Filosofien
jager Muserne bort fra den gamle Mands Leje i
Fængslet, antydes derved en vis Modsætning mellem
antik Tænkning og antik Digtning. I græske Tragedier
faar Stemning og Lidenskab ofte Lov til at vinde fuldt
Udtryk. Man tænke paa Kassandra's lidenskabelige
% --- p. 159 ---
\setcounter{footnote}{0}%
\opage{159}Klage, Antigone's Jammer over at skulle forlade Livet
og Medeia's fortvivlede Udbrud. Herved udtales kun den
Trøst, der kan ligge i den fulde Hengivelse til Sorgen.
Men der findes ogsaa Exempel paa mild Husvalelse ved
at mærke Guddommens Nærhed. Midt i sin Dødskval
hører Hippolytos Artemis', hans Gudindes Stemme og
allerede det lindrer hans Smerte:

\begin{quote}
O Guddoms Duft, som til mig naar!\\
jeg mærker dig, selv i min Lidelse,\\
og Lindring føler jeg!\footnote[1]{\emph{Euripides}' \textit{Hippolytos}. v. 1391 f. (Gertz' Oversættelse).}
\end{quote}

Hvad Digtningen har fremført i disse Exempler, har
sikkert hyppig rørt sig i Livet. Men Tidens Ordførere
nærede Mistillid til Følelserne, fordi de sætter Fantasien
i Bevægelse og river Mennesket ud af den indre Frihed,
Menneskets faste Borg, og gør det afhængig af Tingene
ved at knytte Sindet til den. Derfor advares der mod
»at sukke i sit Inderste«, idethele mod at lade Følelsen
faa sit fulde Forløb, hvad enten dette kun førte til at
faa Luft, eller det bragte positiv Lindring. Saadanne
Betænkeligheder havde Kristendommen ikke. Den har
reddet Livspoesien ved at optage Livets mest modsatte
Stemninger, den højeste Jubel og den dybeste Sorg i
sin Kultus og give dem Udtryk i store Billeder, der
fremstiller Slægtens store Drama. Den Enkelte føres her
ind under en stor og omfattende Betragtning af Livet
og føler sin Lod som bestemt ved Gangen i dette Drama.
I den store Forventning om Fuldendelse og Dom var
der Mulighed for et rigt og dybt Følelsesliv, der kunde
hæve Sindet højt over alle ydre og indre Tilskikkelser\footnote[2]{Se min Sammenligning mellem Stoicismen og Kristendommen i Afhandlingen \textit{Hedenske Sandhedssøgere} (Mindre Arbejder I) og \textit{Epiktet og Kristendommen} (Mindre Arbejder III).}.

% --- p. 160 ---
\setcounter{footnote}{0}%
\opage{160}Uddybelse af Følelseslivet gennem ophøjede Billeder
og derved Grundlæggelsen af en stor Livspoesi er fra
et almenmenneskeligt Synspunkt Kristendommens Bidrag
til Trøstens Historie. Dog var det ikke den Bebudelse,
med hvilken Kristendommen selv optraadte,
--- som overhovedet vel aldrig nogen stor aandelig Bevægelse
har bragt netop den Frugt, den selv mente at
bringe. Kristendommen vilde bringe den højeste Sandhed,
stille Sandhedskravet i ny og streng Forstand, og
Fyldestgørelse af dette Krav stilledes som ufravigelig
Betingelse for, at Trøsten kunde komme.

Den Sandhed, hvorom Urkristendommen samler alt,
er, at Gudsrigets Gennembrud er nær, og at det gælder
at være beredt til at modtage det. Trøsten ydes kun
til dem, der med fuldt Hjerte gaar op i den store Forventning,
lader deres Liv besjæle af den og agter alt andet
ringe. Derfor kunde Kristendommens Grundlægger sige,
at han ikke var kommen for at bringe Fred paa Jorden,
men for at optænde Ild og vække Splid.

Det kunde ikke være andet, end at Kristendommen
under sin Vexelvirkning med den Kulturudvikling, som
Urkristendommen ikke havde ventet eller forudset,
maatte undergaa Ændringer og Forskydninger. Og Trøstelementet
har da afgjort skudt sig frem foran Sandhedselementet.
Og ofte kom da Trøsten, den Trøst, Kristendommen
forkyndte, til at staa langt ude i Horizonten som en
sidste Tilflugt, medens den umiddelbare Glæde og Tryghed
ved Livet raadede i Forgrunden. Og man fandt maaske
endog den Sandhed, der var aabenbaret i Kristendommen,
netop at bestaa i den Trøst, den bragte. Dersom det havde
kunnet vises, at den Sandhed, Kristendommen forkyndte
som den højeste, kunde omfatte, eller dog stemme med,
hvad menneskelig Forsken kan erkende som Sandhed,
vilde saadan Forskydning ikke have været nødvendig.
% --- p. 161 ---
\setcounter{footnote}{0}%
\opage{161}Men det kan siges i Grunden at være indrømmet fra
alle Sider, at Religionens, særlig Kristendommens Betydning,
ikke mere kan bestaa i at tilfredsstille Menneskeaandens
intellektuelle Krav. Uvilkaarlig holder man
sig da, naar man vil lægge sig Kristendommen til Rette,
til andre Sider ved den og stiller dem i første Linie.
Herom har jeg udførlig talt i min Religionsfilosofi og
gaar derfor her ikke nærmere ind derpaa. Jeg har i den
nævnte Bog (§111) sagt, at det er et Tegn paa Religionens
forandrede Stilling i Aandslivet, at den mere og
mere faar Betydning af Trøst: »Religionen var engang
Ildstøtten, der gik foran Menneskeslægten paa dens store
Tog gennem Historien og viste den Vej. Nu bliver den
mere og mere Ambulancen, som følger efter Toget og
opsamler de Trætte og Saarede.« --- Det er blevet mig
fortalt, at disse Ord er blevne anførte paa en københavnsk
Prædikestol og betragtede som en Forhaanelse
af Kristendommen. Den, der vil læse Ordene i den Sammenhæng,
i hvilken de forekommer, vil se, at dette er
en Misforstaaelse. Det er en historisk Bemærkning, som
man kan forsøge at bestride, men som man nok vil blive
nødt til at lade staa. --- Skal Sandheden være det første
og afgørende i Kristendommen, kan den kun hævdes
fra et Standpunkt som det, Pascal og Kierkegaard
indtog\footnote[1]{For \textit{Kierkegaard's} Vedkommende kan særlig henvises til hans \textit{Papirer} (udg. af Kuhr og Heiberg) VIII,1. p. 209. »Christus er ikke Kærlighed, og allermindst i menneskelig Forstand; han var \textit{Sandheden}, \textit{Sandheden}, absolut, derfor kunde han ikke blot forsvare, men han maatte lade Menneskene blive skyldige i hans Død, ɔ: gøre \textit{Sandheden} aabenbar til det Yderste.« Den her udtalte Tanke har han nærmere udviklet i \textit{Tvende ethiske-religiøse Smaa-Afhandlinger}. (Samlede Værker VI)}.

Hvad der en Tid stansede Kierkegaard og holdt ham
fra at gøre sin skarpe og strenge Opfattelse gældende,
% --- p. 162 ---
\setcounter{footnote}{0}%
\opage{162}var netop Hensynet til de mange, der uanfægtede af
intellektuelle og andre Betænkeligheder havde fundet
deres Trøst i kristelige Forestillinger og Billeder, ---
»de Svage, Enfoldige, Syge og Sorrigfulde«, »de mange
lykkelige, elskelige Existenser«, han ikke vilde forfærde.
Men Stansningen kunde for den intellektuelle og etiske
Redeligheds Profet kun blive foreløbig.

Imidlertid ligger der et virkeligt Problem i det Hensyn,
der voldte den foreløbige Stansning. Den rette Trøst
kræver, som vi har set, en aandelig Aktivitet, stiller en
Fordring til den, der skal trøstes. Overfor en saadan
Fordring vil Menneskene --- ligesom overfor etiske Fordringer\footnote[1]{Se herom min \textit{Etik} IV,2; VIII,6 hvor det paavises, at man i Grunden hverken i Art eller i Grad kan stille den samme etiske Fordring til Alle.}
--- møde med højst forskellige Forudsætninger,
og hvad der skal kunne trøste dem, maa derfor blive
af højst forskellig Art og Grad. Det vil gælde om at udløse
det Minimum af Selvvirksomhed, som maa være
muligt, saa længe Livet varer. Derfor er det en Kunst
at kunne trøste; den kræver Menneskekundskab og Fordomsfrihed.
Den sande Trøster er den, der formaar at
løse den levende Kraft og udløse den paa en Maade og
i en Retning, der stemmer med den enkeltes Individualitet.
Her er ingen udelukket fra Trøst, naar blot den
rette Trøster (hvad enten det er den enkelte selv eller
en anden) kommer. Om han altid kommer, eller kan
komme, derom kan Psykologien intet sige. Naar det
paastaas, at Trøst altid er mulig, er det maaske velment
Retorik, men det er ikke Tænkning.

% --- p. 163 ---
\setcounter{footnote}{0}%

\opage{163}\essayhead{TRÆK AF BØNNENS HISTORIE}

Det er af stor psykologisk Interesse, naar der, hvad
en betydningsfuld sjælelig Livsytring angaar, kan
paavises en vis Udviklingsgang, i hvert Tilfælde indenfor
enkelte Kulturperioder, --- en Række af Metamorfoser
eller Forskydninger i en vis Retning, saa at der i Vexelvirkning
med de forskellige historiske Forudsætninger
fremstiller sig Typer af mere eller mindre udprægede
Ejendommeligheder. Ikke som om en saadan Udviklingsgangs
Slutningsform skulde betyde et en Gang for alle
opnaaet Resultat. Der kan stadig være Mulighed for
nye Begyndelser, ogsaa for Gentagelse af Typer, der
syntes at tilhøre et tidligere, formentlig overvundet
Stade. Der kan fremdeles samtidig foregaa forskellige
Udviklingsprocesser, danne sig forskellige Typerækker,
hvis Betydning og hvis indbyrdes Forhold først en langt
senere Tid vil kunne gøre klare. En psykologisk og historisk
Betragtning maa stedse have Øjet aabent for saadanne
Muligheder. For det enkelte Menneske i Livets
Praxis bliver det derimod først og fremmest Opgaven
at finde den Type, han selv efter sin Ejendommelighed
og sin Livserfaring tilhører. Det er det store ved det af
Sokrates grundede Personlighedsprincip, at det stiller
samme Fordring paa alle Stadier og under alle Former
af menneskelig Udvikling. Det gælder overalt at kende
sig selv, sin Evne og sin Trang i Forhold til sit eget
% --- p. 164 ---
\setcounter{footnote}{0}%
\opage{164}Livs og sin egen Tids Betingelser. Saadan Selverkendelse
er den vigtigste Forudsætning for ærlig og sund Udvikling,
altsaa for Opstaaen af nye værdifulde Typer, og
den betegner tillige det Punkt, hvor Mennesker af højst
forskellig Aandsretning bedst forstaar hverandre.

Som Exempel skal her anføres nogle Træk af Bønnens
Historie indenfor forskellige Kulturperioder. Jeg
ordner Exempler, jeg har til min Raadighed, under tre
Hovedformer. Ordningen sker efter en Maalestok, som
naturlig frembyder sig, jo mere Eftertanken retter sig
mod dette sjælelige Fænomens Sammenhæng med det
personlige Liv og dets Forhold idethele. Men dermed er,
som allerede antydet, ikke sagt, at en saadan Udvikviklingsmaade
gennemløbes en Gang for alle i det menneskelige
Aandslivs Historie.

\subhead{A. Bøn og Besværgelse. Social og individuel Bøn.}

a. Der gives et Stade indenfor ellers højst forskellige
Kulturer, hvor der tillægges Bønnen, ligesom Offeret, en
dobbelt Virkning, en Virkning paa Guden, der bedes til,
og en Virkning paa selve det bedende Menneske. Guden
styrkes og næres ved Bønnen, og hans Selvfølelse tiltales
ved, at man udtrykker sin Afhængighed, sin Frygt,
sit Haab eller sin Beundring overfor ham. Bønnen kan
selv betragtes som et Offer, idet Mennesket giver Afkald
paa egen Magt og Stolthed. Men Mennesket selv faar
Oprejsning for, hvad det ofrer, idet Bønnen giver det
Magt over Guden\footnote[1]{Eksempler fra indisk, persisk, ægyptisk, græsk og romersk Religion hos \emph{Westermarck}: \textit{Origin and Development of Moral Ideas}. II. p. 653---659. \emph{Westermarck} vil endog (idet han specielt støtter sig til visse marokkanske Skikke) fortolke Ordet Religion saaledes, at det udledes af religare, at binde: Guden bindes med Menneskers Bønner, Ønsker eller Forbandelser (I. p. 584).}. Troen paa en højtidelig Bøns Virk%
% --- p. 165 ---
\setcounter{footnote}{0}%
\opage{165}somhed beror paa den Forudsætning, at man ved Ordets
Kraft kan fortrylle Guden, tvinge hans Villie --- naar
man kun nævner ham med hans rette Navn\footnote[1]{\emph{Usener}: \textit{Gøtternamen}. p. 335.}.

I to berømte Digterværker danner Bønnens besværgende
Magt et afgørende Motiv. I Aischylos' »Eumeniderne«
siger Apollon til Hævngudinderne, at han vil
udfri Orestes fra deres Forfølgelse, fordi denne har anraabt
Guden i Tillid til hans Beskyttelse: »Ti rædsom
er, for Gud og Mand, den Bedendes Fortørnelse«. (Brøndsted's
Oversættelse, p. 129). I det indiske Drama »Sakuntala«
hidføres Heltindens Ulykke ved, at hun, betagen
af sin Elskovslykke, ikke hører en from Vismands
Melding, og han lyser da sin Forbandelse over hende,
en Forbandelse, han selv ikke uden Besværgelse kan
hæve. (Hammerich's Oversættelse, p. 44).

Ogsaa paa mere primitive Trin gaa Bøn og Trolddom
over i hinanden, uden at der kan gøres nogen tydelig
Forskel. Troen paa Ordets Magt er saa stor, at den
blotte Fortælling om en Begivenhed kan anses som tilstrækkelig
Aarsag til at fremkalde selve Begivenheden.
Frazer formoder, at mange gamle Fortællinger, som vi
er vante til at betragte som blotte Myter, er blevne
brugte i magisk Ritus, idet det mentes, at naar de fremsagtes
paa en bestemt, overleveret Maade, bevirkede de,
hvad de handlede om. Der kan da skelnes mellem fortællende
Trylleformler (narrative spells) og bydende
Trylleformler (imperative spells)\footnote[2]{\emph{Frazer}: \textit{The Spirits of the Corn and of the Wild}. I. p. 105—106.}.

Bønnens Magt kan være farlig. Hos Ainoerne i Japan
maa Kvinderne ikke bede --- for at de ikke skal faa Magt
over Mændene! Og af samme Grund undgaar Santalerne
at aabenbare deres Husguders Navne for Kvinderne.

Den Magt, der i Bønnen og Besværgelsen er lagt i
% --- p. 166 ---
\setcounter{footnote}{0}%
\opage{166}Menneskets Haand, kan afføde to forskellige Konsekvenser.
Den kan gennem Bønnens Almagt gøre Mennesket
selv til Gud. Saaledes betegner det indiske Ord
Brahma baade den Kraft, ved hvilken Guderne lever og
virker, og den Kraft, ved hvilken Mennesket gennem
Bønnen fravrister Guderne sit Ønskes Opfyldelse. En
Identificeren af begge laa da nær, og kan spores i den
ejendommelige Udvikling, der i det gamle Indien førte
til Monoteisme\footnote[1]{Smlgn. Religionsfilosofi\textsuperscript{2}. p. 151.}. Men der kan ogsaa drages en anden
Konsekvens, idet man hævder, at det kun er lavere
Væsener, der er underlagte Ordets magiske Magt, ikke
de egentlige Guder. Dette fører i forskellige Religioner
til at anvende den Skelnen mellem Guder og Dæmoner,
som oprindelig\footnote[2]{Se f. Eks. \emph{Bezold}: \textit{Die babylonisch-assyrische Religion}. (Kultur der Gegenwart. II,3,1). p. 44.} skyldes en ny Religions Sejr over en
ældre.

b. Bøn og Besværgelse og den magiske Magt, der
kan øves gennem dem, maa naturligvis stedse være
Handlinger, der udføres af bestemte Mennesker. Men
udføres de af disse Mennesker paa deres egne Vegne,
som Udslag af en individuel Evne, eller er det en social
Gerning, der udføres gennem dem? Det er Forholdet
mellem Magi og Religion, vi her staar overfor. Frazer
har opstillet den Teori, at Magien gaar forud for Religionen.
Førend Mennesket erkender sin Afhængighed,
indser, at han ikke formaar alt, ikke har fuldstændig
Magt over Naturen, gaar der ifølge Frazer en Periode,
i hvilken et enkelt Individ tror sig istand til ved sine
Formler og Handlinger at gribe ind i Naturens Gang.
Men efterhaanden lærer Mennesket, at Haabet om at
lede Naturen ved magiske Midler maa opgives, og
han betragter Guderne som eneste Besiddere og Haand%
% --- p. 167 ---
\setcounter{footnote}{0}%
\opage{167}hævere af den Magt, han selv troede at besidde. Nu
først træder Forskellen mellem Guder og Mennesker afgjort
frem, og Mennesket føler sig ikke længere som
Gudernes Ligemand. Bøn og Offer bliver nu Led i et
fælles Ritual, og det betragtes som Overgreb og Helligbrøde,
naar den enkelte paa egen Haand og ad sin egen
Vej sætter disse Kræfter i Bevægelse. Magi er derfor
ifølge Frazer\footnote[1]{Se især \textit{The Magic Art}. (1911). I. p. 233 f; 374. --- Tidigere i \textit{Adonis, Attis, Osiris}. (1906).} ældre end Religion, naar man ved Religion
forstaar en af Afhængighed udsprungen Dyrkelse
af Væsener, der er mægtigere end Menneskene. Efterhaanden
som Religionen udvikler sig, bliver Magien til
»den sorte Kunst«.

Denne Teori giver en simpel psykologisk Forklaring
af Religionens Opstaaen gennem Erfaringen om menneskelig
Magts Begrænsning og en derved betinget Afhængighedsfølelse.
Jeg har derfor tidligere lagt Vægt
paa den. Men den historiske Udvikling skrider ikke
altid frem efter simple psykologiske Formler. Og man
maa give den franske sociologiske Skole Ret i, at saa
langt vi end gaar tilbage, finder vi stedse Religionen
socialt organiseret. Den enkelte frigør sig først efterhaanden
fra fuld Beherskethed af social Skik og Overlevering.
Efter Durkheim og hans Disciple gaar derfor
Religion forud for Magi, ligesom det Sociale stedse gaar
forud for det Individuelle. Der lægges her især Vægt
paa, at i de primitive Samfund uddannes og øves Magikeren
gennem Belæring og Indvielse, og de Væsener,
han gennem sin ofte ekstatiske Indvielsesperiode kommer
i Berøring med og henter sin Magt fra, er dem, Stammen
dyrker. Den unge Magiker maa hos og i sig selv gennem
Visioner erfare Sandheden af de Beretninger, hans Stammes
Overlevering indeholder. Han er da selv i sin største
% --- p. 168 ---
\setcounter{footnote}{0}%
\opage{168}Magtfylde et socialt Produkt og en social Repræsentant.
Først senere kan Magien træde i et skarpt Modsætningsforhold
til Religionen\footnote[1]{Se, foruden \emph{Durkheim}: \textit{Les Formes Elémentaires de la Vie Religieuse}. (1912). Livre I. Chap. 1. \emph{Hubert et Mauss}: \textit{L'Origine des Pouvoirs Magiques} (Mélanges d'Histoire des Religions). (1909). --- Smlgn. om min Stilling overfor den franske sociologiske Skole Afhandlingen \textit{Sociologi og Religionsfilosofi} i nærværende Værk.}.

Formedelst sin oprindelige nøje Sammenhæng med
Besværgelsen faar Bønnen en magisk Karakter og tillige
en social Karakter. Men ligesom den efterhaanden
udskiller sin magiske Karakter, saaledes faar den ogsaa
mere og mere en individuel Karakter. De sociale Former
individualiseres. Det bliver ikke nok, at den enkelte i
sin personlige Erfaring søger at opleve, hvad der berettes
om de Væsener, Stammen dyrker, men han gør
sine nye Erfaringer, og hans Bøn faar derved en mere
personlig Karakter. I primitive Religioner (d. v. s. de i
vore Øjne lavest staaende) vil der neppe kunne findes
afgørende Exempler i denne Retning, skønt der ikke
mangler Exempler paa Individer, der stiller sig tvivlende
overfor Overleveringen. Derimod træder i den ægyptiske
Religion baade Bønnens Sondring fra Besværgelsen
og Bønnens Individualisering for Dagen paa en mærkelig
Maade.

Denne Religions Historie viser os alle Bønnens Former,
ligefra det magiske Stadium, hvor Bøn og Trolddom
er et, til den inderlige og af personligt Sindelag
fremgaaede Bøn. Paa tidligere Trin kan der knap skelnes
mellem Bøn og Trolddom; det kan ikke ses af Bønnens Indhold,
om den er mere end en Trylleformel; det kommer
an paa Anvendelsen, ikke paa Indholdet. I den gamle
stive ægyptiske Statsreligion kunde alle bede den samme
Bøn, og Vægten blev lagt paa de Ord, som ytredes. Men
% --- p. 169 ---
\setcounter{footnote}{0}%
\opage{169}senere i den mærkelige Bevægelse i Retning af personlig
Inderlighed, som betegner den religiøse Udvikling i
Ægypten i det trettende og tolvte Aarhundrede f. Kr.,
blev Bønnen »en Aabenbaring af indre personlig Erfaring
eller et Udtryk for et individuelt Forhold til Gud«.
Der advares udtrykkelig at bruge for mange Ord, da
det slet ikke er Ordene, det kommer an paa. »Vær ikke
af mange Ord, ti i Tavshed vil du vinde, hvad godt er«.
»Bed da med et higende Sind, hvis Ord alle er skjulte«.
Og Gud kaldes »de Tavses Gud«. »Du liflige Kilde for
den, der tørster i Ørkenen, er lukket for den, der taler,
men aaben for den Tavse«.

Den Udvikling i Retning af Personlighed og Inderlighed,
der betegner et Højdepunkt indenfor den ægyptiske
Religion, afløstes af ny Stivnen under et sneversindet
Præsteskabs Herredømme\footnote[1]{\emph{James Henry Breasted}: \textit{Development and Thought in Ancient Egypt}. (1912). p. 95; 272 ff; 349; 356.}.

Der er en indre Forbindelse mellem Bønnens Frigørelse
fra Besværgelseskarakteren og dens Overgang til
at være et inderligt Udtryk for den enkelte Personligheds
Trang. Ved Besværgelser vil det komme an paa
at bruge aldeles bestemte Ord og Handlinger, ligesom
man jo overalt maa anvende visse bestemte Midler for
at naa til et bestemt Maal. Bønnen bliver ved den Frigørelsesproces,
vi har set Exempler paa, ført ind paa
en Inderliggørelsens Vej, der kan føre stedse videre. ---

Naar Vægten lægges paa Ordene og paa deres Brug
til Opnaaelse af ydre Formaal, er det naturligt, at der
lægges Vægt paa mange og lange Bønner. Det er karakteristisk,
at vi 1200 Aar før Bjergprædikenen hører en
Advarsel mod at bruge mange Ord; det er Tegn paa
en Inderliggørelsesproces, som jo desværre blev afbrudt
% --- p. 170 ---
\setcounter{footnote}{0}%
\opage{170}af den Magt, der har faaet saa mange inderlige og personlige
sjælelige Udviklingsprocesser til at ophøre.

Lotze anfører etsteds en tysk folkelig Sætning, som
aabenbart er fremkaldt ved Erfaringen om et omvendt
Forhold mellem Ordenes Mængde og Stemningens Inderlighed:
»De, der beder for meget, beder sig tvers igennem
Himmelen og bliver satte til at vogte Gæs paa den anden
Side«. ---

\subhead{B. Bønnens Inderliggørelse.}

Naar Bønnen ikke staar som en magisk Kraft, og
naar den ikke blot er en uvilkaarlig eller vilkaarlig Gentagelse
af en i Overleveringen foreliggende Formel, er
Vejen aaben for Inderliggørelse. Da kan de ydre Formaal
for Bønnen falde bort, og Bønnen kan mere og
mere faa sin Værdi i sig selv som en sluttet, koncentreret
Tilstand, en Enhedstilstand, i hvilken tilsidst
ogsaa Dobbeltheden af den, som beder, og den, der bedes
til, kan falde bort, ikke blot Dobbeltheden af Bønnen
selv og det Gode, der bedes om.

En Overgangsform danner Bønnen som Udtryk for
Lovprisning, Beundring eller Taknemmelighed. Takkebønnen
er ædlere end Bedebønnen. De smukkeste og
mest ophøjede af de gammeltestamentlige Psalmer hører
herhen. Den ottende Psalme med dens Lovprisning og
Tak for den herlige Plads, der er anvist Mennesket i
Skabningens Verden, har man kaldt det skønneste, alle
Tiders religiøse Literatur i denne Retning har frembragt\footnote[1]{\emph{Wundt}: \textit{Völkerpsychologie}. II,3. p. 665 f.}.
Stoikeren Epiktet kender kun Takkebøn, ikke
Bedebøn, og hans Tak gælder, i Lighed med den gamle
Psalmists, Menneskets Herlighed, --- kun at han som
Stoiker finder denne Herlighed deri, at Mennesket hver%
% --- p. 171 ---
\setcounter{footnote}{0}%
\opage{171}ken finder sin sande Lykke i sig selv, ikke i noget, som
er udenfor det\footnote[1]{Bonhöffer: Epiktet und das neue Testament. (1911). p. 65.}.

Allerede naar Bedebønnen angaar aandelige Goder,
der tilsidst beror paa Menneskets eget Sind og egen
Villie, har den mistet sin Udvorteshed. En saadan Bøn
beder Sokrates i Platon's »Fajdros«, da han efter en
dybsindig Samtale med sine Venner i den hellige Lund
foreslaar dem følgende Bøn, inden de gaar bort: »O kære
Pan og I andre Guder, som er hjemme her, giv mig at
faa et skønt Indre, og at de ydre Goder maa stemme
med de indre!« (Fajdros p. 279 B). --- I Dialogen »Alkibiades
II« anfører Sokrates med Bifald et Vers af en
ældre Digter, der beder Zeus give os det, som godt er,
hvad enten vi bede derom eller ikke, men afværge det
Onde, selv om vi bede derom. (Alkibiades II p. 143 A).
At denne Dialog neppe er af Platon, er i denne Sammenhæng
ligegyldigt. Den røber Indflydelse af stoiske
Tanker.

Platonisme og Stoicisme har haft meget stor Indflydelse
paa Bønnens senere Historie. Nyplatonismen
gennemfører, hvad Platon og Stoikerne havde indledet.
Plotinos lærte, at hvad der opleves i vore højeste Tilstande,
der hos ham fik Ekstasens Karakter, kan ikke
tænkes eller siges, saalænge de bestaar; bagefter kan
man forsøge at tyde det Oplevede, men det bliver altid
Stykkeværk. Vi kan have det Højeste i os, selv om vi
ikke kan sige, hvad det er, ligesom de Begejstrede og
Betagne kun vide, at de har noget Større i sig, men ikke
véd, hvad det er. (Ennead. V, 3, 10; 5, 9). Paa egentlig
Bøn lagde han ingen Vægt, saa lidt som paa Kultus
idethele; Guderne skulde komme til ham, ikke omvendt\footnote[2]{Se Afhandlingen Oplevelse og Tydning.}.
Men der var dem blandt hans Efterfølgere, der svingede
% --- p. 172 ---
\setcounter{footnote}{0}%
\opage{172}tilbage til Kultus og folkelig Religion, tildels i meget
kras Form. Et Vidnesbyrd herom er Skriftet Ȯgypternes
Mysterier«, som er blevet tillagt Jamblichos og i hvert
Tilfælde skriver sig fra hans Tid (c. 300 e. Kr.). Modsætningen
mellem Filosofi og Mytologi falder bort. Bestræbelsen
gaar ud paa at forny den gamle Folkereligion
under Indflydelse af de nye orientalske Strømninger og
hævde den overfor, hvad der paa den Tid var den moderne
Tvivl. Bønnen faar nu kun Værdi, naar den fremføres
i hellige, det vil her sige af Overleveringen godkendte
Ord --- og helst i Ord, der stammer fra Orientens
hellige Folk (Ægypter og Assyrer), hvem Guderne har
aabenbaret det hemmelighedsfulde Sprog. Er disse Ord
end uforstaaelige for os, saa forstaar dog Guderne dem!
Bønnen bliver altsaa igen til Magi. Men tillige hævdes,
--- hvad der vil vise sig at betegne et betydningsfuldt
Vendepunkt, --- at det ikke er Bønnen, der vender Gudernes
Sind til Mennesket, men at det er Gudernes eget
naaderige Sind, der straaler ned over de Bedende, idet
det allerede før Bøn og Henvendelse har virket paa dem.
Der skelnes mellem tre Trin i Bønnens Virkning, en
Trehed, der ofte kommer igen i Mystikens senere Historie
(saaledes hos Victorinerne og Sancta Teresa). Det
første og laveste Trin bestaar i Samling, Koncentration,
Henvendelse mod et eneste Maal. Det andet bestaar i
en indre Forbindelse med Guderne. Paa det tredie Trin
fuldendes Enheden med Guddommen (den Enhed, som
Plotinos naaede gennem Tankens Selvfordybelse)\footnote[1]{\emph{Harless}: \textit{Das Buch von den ægyptischen Mysterien}. (1857). p. 73---79.}.

Ikke alle Nyplatonikere gik med til en saadan Sammenblanding
af Tænkning og Magi. Den sidste betydeligere
Nyplatoniker, Olympiodoros fra Alexandria, lægger
Vægten paa Bønnens rent inderlige Karakter. Han ud%
% --- p. 173 ---
\setcounter{footnote}{0}%
\opage{173}vikler i sin Fajdonkommentar, i Tilslutning til en Tankegang
hos Platon, at den højeste menneskelige Stræben
tilsidst har sit Formaal i sig selv, og den, der ikke vil
indrømme dette, henviser han til, at vi jo dog hengiver
os til Bønnen for selve Bønnens Skyld. Dette er altsaa
for ham en selvfølgelig Sætning. Og i en anden Sammenhæng
kommer han tilbage dertil. Han viser, at Tænkning
ikke er noget, vi kan afgøre en Gang for alle; det
gælder ikke om at have tænkt, men om at forblive i
tænkende Virksomhed, og han oplyser dette ved Bønnen:
da Bønnen i sig selv er Lykke, er den noget, vi vil
blive i\footnote[1]{\textit{Olympiodori Philosophi in Platonis Phaedonem commentaria}. ed. W. Norvin. (1913). p. 19; 95.}.

De nyplatoniske Tanker gaar igen i Middelalderens
Mystik, som overalt søger at omsætte de ydre Forhold,
religiøse Forestillinger betegner, i indre. At »Guds Rige
er nær«, betyder for Mester Eckchardt, at det er i os,
ti Gud er os nærmere end vi selv er os. Som her Tidsforholdet
falder bort, saaledes falder Forskellen mellem
Bøn og Bønhørelse ogsaa bort. Gud bevæges saa lidet
ved alle Bønner og gode Gerninger, som om der slet
ikke gaves noget saadant. »I et første evigt Blik --- hvis
vi da kan antage et første Blik --- saa Gud alle Ting,
som de skulde ske\dots{} Han saa ogsaa den ringeste Bøn
og gode Gerning, som nogen vilde forøve, og han saa,
hvilken Bøn og hvilken Andagt han vilde bønhøre«.
Paa den anden Side er den rette Bøn den, der udspringer
af et frit Sind, som ikke er bundet til noget her i Verden,
men helt er hensunket i Guds Villie. I den inderlige Bøn
bliver Mennesket tilsidst et med den, det beder til, med
Gud\footnote[2]{\emph{Mester Eckehart}: \textit{Prædikener og Traktater}. Oversat og indledet af Aage Markus. (1917). p. 69; 78; 95.}.

% --- p. 174 ---
\setcounter{footnote}{0}%
\opage{174}Helt op i Kierkegaard's religiøse Taler møder vi de
samme Tanker, genfundne og selvoplevede, men hørende
til en stor Tradition, der hænger sammen med Religiøsitetens
inderligste Former.

Efter den naive Opfattelse --- siger Kierkegaard ---
er det Hovedsagen, at Gud faar at vide, hvad man har
paa Hjertet. Men det forholder sig omvendt. Det kommer
an paa, at den Bedende vedbliver med sin Bøn,
indtil han faar at vide, hvad Gud vil. Udførlig er denne
Tanke udviklet i en af de religiøse Taler, hvis Overskrift
er: »Den rette Bedende strider i Bønnen og sejrer ---
derved, at Gud sejrer«. Bønnens Inderlighed er Hovedsagen,
og Bønnen er ikke blot et Middel til at opnaa
noget. Allerede den Samling, Bønnen forudsætter, ---
det, at Mennesket samler hele sin Sjæl om et enkelt
Ønske, --- allerede det forædler, og det er Forberedelsen
til at opgive alt; ti kun den, der blot har et eneste Ønske,
kan opgive alt. Det er Guds uforanderlige Villie,
Mennesket bøjer sig for, og derved, --- ved ikke selv at
ville være noget, --- kan han helt gennemstraales af Gud
og komme til at ligne Gud. --- I en senere Optegnelse
hævder Kierkegaard, at man skal lære altid at takke.
 »Ellers fremkommer al denne barnagtige Vrøvlen, det
ene Øjeblik skal man takke, for nu er det godt eller det
og det er saa rart, det næste Øjeblik skal man sukke
og bede om Taalmodighed, for nu var det og det saa
galt. Det er umuligt at have et virkeligt Forhold til
Gud og blive i sin blot menneskelige og jordiske Forestilling
om Godt og Ondt, Behageligt og Ubehageligt«(\footnote[1]{\emph{Kierkegaard}: \textit{Samlede Værker}. V. p. 149—167. Smlgn. \textit{Papirer} (udg. af Heiberg og Kuhr) VII,1. 24; VIII,1 p. 123.}.

I Nyplatonisme og Mystik er det Enhedstanken, der
raader. Bestræbelsen gaar ud paa at komme ud over
Modsætningen mellem Bønnen og det, der bedes om,
% --- p. 175 ---
\setcounter{footnote}{0}%
\opage{175}og ud over Modsætningen mellem den, der beder, og
den, der bedes til. I Kierkegaard's religiøse Taler er det
Tanken om den guddommelige Uforanderlighed, altsaa
ogsaa en platonisk Tanke, der fører ud over den sædvanlige
Bedeteori. --- Det er naturligvis ikke Kierkegaard's
definitive Standpunkt, vi har for os her. I sit dybest
gaaende Arbejde (»Sygdommen til Døden«) siger han:
»At bede er at aande, og Muligheden er for Selvet, hvad
Ilten er for Aandedrættet\dots{}. Gud er det, at alt er
muligt, eller at alt er muligt, er Gud«\footnote[1]{\textit{Samlede Værker}. XI. p. 153.}. Uforanderligheden
tillader derimod ikke flere eller andre Muligheder,
end der altid har været. ---

Det kunde synes, at det blot var et terminologisk
Spørgsmaal, om man vilde bruge Ordet Bøn om en Akt,
i hvilken Sindet er ganske et med sin Genstand, med
det Uforanderlige, og ved hvilken der ikke er Forskel
paa Akten selv og det, den er rettet imod. Brøchner
hævder saaledes, at naar man ser Bønnens Virkning i
den indirekte, subjektive Indflydelse paa selve den Bedende,
saa er Bønnens Væsen i Virkeligheden undergravet
af Kritiken. Ti Bønnen gaar ud paa Underet, og
kun paa Underet. Derimod antager Brøchner, at der
gives et rationelt Ækvivalent for Bønnen, nemlig Sindets
Henvendelse til det uendelige Væsen, med hvilket det i
sin inderste Grund er et, og i den Kraft, der gennem
saadan Henvendelse til vort Væsens sidste Kilde kan
udløses hos os\footnote[2]{\emph{Brøchner}: \textit{Problemet om Tro og Viden}. (1868). p. 95ff. --- \textit{Det Religiøse i dets Enhed med det Humane}. (1869). p. 17.}. Det er Mester Eckehardt's Standpunkt,
Brøchner stiller sig paa, kun at han ikke vil kalde den
inderlige Henvendelse, der paa et saadant Stade er mulig,
for Bøn. I en fransk Disputats, hvis Emne var en psykologisk
Studie af Bønnen, havde Forfatteren (Segond)
% --- p. 176 ---
\setcounter{footnote}{0}%
\opage{176}netop fundet Bønnens Væsen udtrykt, ikke hos primitive
Mennesker, hvis Indre vi ikke kan rekonstruere, men
hos Mystikerne, hos hvem Bønnen opfattes med klar
Bevidsthed om dens Natur. Af hans Opponenter (Durkheim,
Delacroix og Delbos) blev der foreholdt ham det
vilkaarlige i denne Begrænsning og hævdet, at Forholdet
mellem Tavshedens Bøn og Ordets Bøn burde
være nøjere undersøgt\footnote[1]{Se det interessante Resumé af Diskussionen i Revue de Métaphysique et de Morale Janvier 1911. Supplément p. 30—35.}. Opponenterne havde aabenbart
Ret i, at det er vilkaarligt at kalde en enkelt Type for
den væsentlige. Der springes derved en hel Række Problemer
over.

Spørgsmaalet er dog ikke blot terminologisk. Historien
synes at vise, at der finder en successiv Overgang
Sted mellem de forskellige Typer, selv om de, naar de
hver for sig er fuldt udviklede, kan træde i skarp Modsætning
til hverandre. Denne Kontinuitet har sin store
Betydning for at undgaa Kraftspild; ti gennem den
kan en ny Type drage Næring af den gamles opsparede
Kraft, og desuden gør den Forstaaelsen lettere, selv om
denne Forstaaelse væsentlig maa vindes ad Analogiens
Vej. Brøchner har dog Ret i, at det for Klarhedens Skyld
er hensigtsmæssigt at markere Forskellen mellem de
forskellige Typer ved forskellige Benævnelser, selv om
det kan blive vanskeligt at afgøre, hvor den afgørende
Grænse skal sættes. Kirken har ofte sat den der, hvor
Dogmatiken trues af den dristige Mystik. Religionspsykologisk
maa man sætte den paa den anden Side af
Mystiken, hvor det store Problem om den psykologiske
Mulighed og Beskaffenhed af det Ækvivalent, Brøchner
omtaler, stiller sig. --- Et herfra forskelligt Spørgsmaal
er dét, hvilken Type der er den højeste. Foreløbig betragter vel
Enhver det Stade, han selv staar paa, som
% --- p. 177 ---
\setcounter{footnote}{0}%
\opage{177}det højeste\footnote[1]{Smlgn. Striden mellem Baggesen og Grundtvig om, hvis Stade der var det højeste. Det kunde ifølge Grundtvig først afgøres paa Dommedag. \textit{Danske Filosofer}. p. 70f.}. Det afgørende vil være, om en enkelt
Type kan godtgøres at indeholde i sig alt, hvad der er
værdifuldt og berettiget i andre Typer. Det ligger udenfor
nærværende Undersøgelse at gaa nærmere ind herpaa.

Som allerede udtalt i Begyndelsen af denne Afhandling
foreligger der ikke en eneste Udviklingsrække, men
flere Rækker paa én Gang, snart opstigende, snart nedstigende,
hvad Forholdet mellem Dobbelthed og Enhed
angaar. Og hos store Personligheder, der har grebet dybt
ind i positiv Religions Udvikling, vil man ofte finde
flere forskellige Typer gøre sig gældende. Hos Jesus af
Nazareth, saaledes som Evangelierne skildrer ham, forekommer
de fleste af de i det foregaaende omtalte Typer.
Overfor Dæmonerne anvender han Besværgelse, men
overfor Gud anvender han Bedebønnen og Takkebønnen,
og ogsaa her staar Takkebønnen som den mere sublime,
idet den er Udtryk for de herligste Erfaringer. I Johannesevangeliet
raader Enhedsbønnen, og det er den,
der har sat Frugt i den kirkelige Dogmeudvikling. Højdepunktet
er dog Bønnen i Gethsemane, --- det mest religiøse
Ord, der er sagt paa Jorden. Her kæmper Jesus
sig frem fra Dobbelthed til Enhed, fra Særvillie til fuld
Opgaaen i den guddommelige Villie, der udtrykker hans
Kald. --- Hos den største Personlighed indenfor Kirken,
Augustinus, træffer vi (se Confessiones) ogsaa, dels efterhaanden,
dels i samme Periode, flere forskellige Typer.
Som Barn bad han om, at han ikke maatte faa Prygl
i Skolen. I hans Ungdom, da Talekunst og Filosofi optog
ham, forandredes med hans Stemning (affectus) ogsaa
hans Bønner, som nu gik ud paa at vinde Visdom. Da
Trangen til at hæve sig over sine sanselige Lidenskaber
% --- p. 178 ---
\setcounter{footnote}{0}%
\opage{178}begyndte at gøre sig gældende, bad han om Kyskhed
og Selvbeherskelse; men han tilstaar, at han frygtede
for at blive bønhørt for hurtig. Derimod da den nye
Tro var ved at bryde igennem hos ham, bad han om
dens Sejr i hans Sind, og nu siger han til sin Gud: »Hvorfor
ikke strax?« En barnlig Bedebøn optræder ogsaa
efter Omvendelsen, da han ved sine egne og sine Venners
Bønner blev befriet for Tandpine. Men han bad forgæves
om at blive befriet fra den dybe Sorg over sin
fromme Moders Død, og han véd, hvorfor han ikke blev
bønhørt; Gud vilde indskærpe ham Betydningen af et
inderligt Forholds Bestaaen, selv efter den anden Parts
Bortgang. Højdepunktet i Bønnens Udvikling var det
for Augustinus, da han kunde bede: »Giv mig, hvad du
byder mig, og byd mig, hvad du vil!« ---

Under alle sine Former hører Bøn nøje sammen med
Trang og Drift, Ønske og Længsel. Men der gives en
stor Religion, ifølge hvilken Drift og Ønske netop er
det, der ikke skal være, da de forøger eller stedse paany
fremkalder den kvalfulde Existens, som det netop gælder
at bringe til Ophør. For Buddhismen\footnote[1]{\emph{Warren}: \textit{Buddhism in translation}. (1906). p. 177. 189. 245.} er Bøn Karma
(Existens), og omvendt er Karma Ophav til den Higen
og Tragten, der faar Udtryk i Bønnen. Det er det mest
radikale Standpunkt overfor Bønnen som sjæleligt Fænomen.
Bønnen er Udtryk for en Livstrang, en Drift til
at fastholde Existensen, selv om det skal være under
helt andre Former end de hidtil erfarede, og selv om
det ikke er den egne, individuelle Existens, der søges
hævdet. Der er Mulighed for Metamorfoser til nye sjælelige
Fænomener, som man maaske ikke længere vil kalde
Bøn, til fri aandelig og materiel Udviklingstrang. Men
Buddhismen tager hele dette Grundlag og alle de Muligheder,
det indeholder, bort fra Roden af.

% --- p. 179 ---
\setcounter{footnote}{0}%
\opage{179}I den Karakteristik, jeg i min Religionsfilosofi har
søgt at give af Buddha og Jesus, har jeg netop lagt Vægt
paa dette Punkt. Buddha forkaster al Higen og Stræben,
Jesus vækker en stor Forventning og lægger dermed
Vægten paa Stræben og Arbejde.

\subhead{C. Bønnen som Gave --- og som Symbol.}

Enkelte Gange er i det foregaaende den Tanke kommen
frem, at det er en Gave at kunne bede, saaledes
hos Nyplatonikere og hos Mester Eckchardt. Gennemføres
denne Tanke, staar vi ved et Vendepunkt i Bønnens
Historie. Den bliver da ikke, eller ikke blot Næring,
men Virkning, og det er dette, der bliver det væsentlige.
Og da bliver tilsidst Bønformen, det, at der bedes om
noget og til nogen, symbolsk, et poetisk Udtryk for
Higen og Længsel efter noget, der i Grunden allerede
ligger i vort eget Væsen. Bønnen virker da ikke blot
som kraftudløsende, men den er selv Følge af en Kraftudløsning,
som det til den givne Tid ikke stod i Individets
bevidste Magt at bringe tilveje. --- Dette Synspunkt
fremtræder igen med forskellig Begrundelse. Jeg
tror, man kan skelne mellem en supranaturalistisk, en
idealistisk og en naturalistisk Begrundelse.

a. Det fortælles,\footnote[1]{Tiele: Elements of the Science of Religion. II. p. 133.} at der til den persiske Profet Jelaled-din
Rumi kom en Mand og klagede over, at Allah
aldrig svarede paa hans Bøn, --- at han aldrig mærkede
Allah's Nærværelse. Da sagde Profeten: »Gud befaler mig
at sige til dig: Du haardt prøvede Mand, var det ikke
mig, der drev dig til at dyrke mig? Selve din Paakaldelse
af Allah var Udtryk for min Nærværelse, og din Sorg,
din Længsel, din Iver var mine Sendebud!« --- Selve Bønnen,
det at der bedes, er altsaa en Aabenbaring af Guddommen.

% --- p. 180 ---
\setcounter{footnote}{0}%
\opage{180}En lignende Tanke kommer frem, naar Paulus skriver:
»Aanden kommer vor Skrøbelighed til Hjælp; ti vi véd
ikke, hvad vi skal bede om, som det sig bør, men Aanden
selv gaar i Forbøn for os med uudsigelige Sukke. Og
han, som ransager Hjerterne, véd, hvad Aandens Higen
er.« (Romerbrevet VIII, 26). Og derefter tales der om
guddommelig Forudbestemmelse og Kaldelse. --- I Bønnen
er Mennesket altsaa ikke den Givende eller Ofrende,
men væsentlig den Modtagende. I selve den uformede
Bøn, »de uudsigelige Sukke« rører en guddommelig Aand
sig i Menneskets Indre, uden at Mennesket er sig det
bevidst\footnote[1]{Smlgn. om dette mærkelige Sted \emph{Bernhard Weis}: \textit{Lehrbuch der Theologie des neuen Testaments}.\textsuperscript{6} p. 343.}.

Bønnen synes at være rettet fremad og opad; men
der er altsaa for en dybere Opfattelse bag Bønnen som
fremadhigende en Kraft, der (som en vis a tergo) skyder
paa og driver frem. Denne Opfattelse bliver den diamentrale
Modsætning til det Trin, hvor Bøn og Besværgelse
gaar i Et, og hvor Bønnen som en magisk Kraft
giver Mennesket Magt baade over Gud og over Verden.
Vi kan ikke vente at finde denne dybere Opfattelse
uden hos enkelte fremragende religiøse Personligheder
til forskellige Tider. Men det maa være et underligt
Standpunkt, fra hvilket dette skulde forringe Opfattelsens
Værdi.

Det er i Afhandlingen om Oplevelse og Tydning vist,
hvorledes denne Opfattelse kommer frem hos Sancta
Teresa i den afgørende Modsætning, hun sætter mellem
Eftertankens Bøn (oraison de méditation) og Hvilebønnen
(oraison de quiétude). Til Hvilebønnen kan man
ikke selv hæve sig op, men maa vente, til man drages
op, ligesom en Have maa vente paa sin bedste Van%
% --- p. 181 ---
\setcounter{footnote}{0}%
\opage{181}ding, den liflige Regn, der stille strømmer ned fra
oven.

Pascal, for hvem det stod fast, at han ikke havde
kunnet søge Jesus, hvis han ikke allerede havde fundet
ham, hævder ligefrem, at Bønnen er en Gave, som Gud
giver til hvem det behager ham, og han afviser bestemt
den Antagelse, at man har Bønnen fra sig selv. Det er
en jansenistisk Sætning, Pascal her hævder. Den gaar
tilbage til Augustinus i hans Polemik mod Pelagius.
Selv Bossuet slutter sig derfor til den, idet han siger,
at Bønnen, ved hvilken alt gives os, selv er given os\footnote[1]{\emph{Pascal}: \textit{Pensées}. ed. Brunschwicg. II. p. 408; 438, med Udgiverens Note. --- Smlg. \emph{Sibbern}: \textit{Gabrielis' Breve}\textsuperscript{1}. I. p. 100: »At kunne bede, er den rette fulde Naade.«}.

b. Naar jeg i Modsætning til den overnaturlige Forklaring
af Bønnens Tilblivelse, der gaar ud fra en Dobbelthed
mellem guddommelig og menneskelig Natur, stiller
den idealistiske Forklaring, er det, fordi her den
guddommelige Natur tænkes virkende i og gennem den
menneskelige Natur, ligesom den virker gennem Naturen
idethele.

I det indiske Digt Bhagavadgita forklares baade
Troen paa Guder og Bønnen til Guder som Virkning af
en Magt, der ikke selv kommer til Menneskets Bevidsthed
og ikke kan iagttages, --- af selve det absolute, uforanderlige,
allerhøjeste Væsen, der rører sig i alt. Krishna,
der er en Inkarnation af Vishnu, siger i det nævnte
Digt: »De, hvis Erkendelse er berøvet dem formedelst
et eller andet Begær, nærmer sig til andre Guder end
mig, og bundne til deres egen Natur overtager de allehaande
religiøse Pligter. Naar nu en af dem ønsker i
Tro at tjene en af disse Guder, saa gør jeg denne Tro
urokkelig hos ham. Opfyldt af denne Tro dyrker han
da med Omhu denne Gud og opnaar det Gode, han
% --- p. 182 ---
\setcounter{footnote}{0}%
\opage{182}ønsker; --- men det er i Virkeligheden mig, der yder
ham det«\footnote[1]{\textit{Bhagavadgita} VII, 20—24. (Garbes tyske Oversættelse).}.

I den europæiske Verden er en saadan Tanke gennem
Nyplatonisme og middelalderlig Mystik gaaet over til
moderne Teologer, der ikke føler sig bundne ved konfessionelle
Forestillinger som afsluttende. Det at der
bedes, og at der opnaas noget gennem Bønnen, staar
som to Oplevelser, der begge udspringer af én Kilde,
Menneskene selv ikke aner, naar de ikke gaar Eftertankens
Vej til Ende. Efter Schleiermacher vilde en Teori om
Bønnen, som forudsatte et Vexelforhold mellem Skaberen
og det Skabte, betyde en Tilbagegang til Magiens Standpunkt.
Bøn og Bønhørelse er i Grunden ét, og i den
rette Bøn er Passivitet og Aktivitet forenede. Bønnen
og dens Opfyldelse eller Ikkeopfyldelse er kun Dele af
en og samme guddommelige Ordning\footnote[2]{\emph{Schleiermacher}: \textit{Der christliche Glaube}. § 47.}. I samme Retning
gaar senere Auguste Sabatier's Opfattelse\footnote[11]{Se min Karakteristik af Sabatiers Standpunkt i Mindre Arbejder. Anden Række. p. 75—78.}.

c. Som Exempler paa, hvad jeg mener med en naturalistisk
Opfattelse i dens Anvendelse paa Bønnen, vil
jeg fremdrage Udtalelser af to indbyrdes højst forskellige
Aander, den ene er af Grundlæggeren af den nyere Tids
Naturopfattelse, den anden en af de største Skikkelser i
Filantropiens Historie.

Thomas Hobbes søger at vise, at Bøn meget vel kan
finde sin Plads, selv om man antager, at der er bestemte
naturlige Aarsager til alt hvad der sker. »Bøn,« siger
han\footnote[12]{\emph{Hobbes}: \textit{Tripos}\textsuperscript{3}. (1684). p. 293.}, »er Guds Gave ligesaa vel som Velsignelse er det
\dots{} Og Bøn til den almægtige Gud er kun Taksigelse
for Guds Velsignelser. Selv om den gaar forud for den
% --- p. 183 ---
\setcounter{footnote}{0}%
\opage{183}specielle Ting, vi beder om, er den dog ikke Aarsag eller
Middel dertil, men en Tilkendegivelse af, at vi ikke vente
andet, end hvad Gud vil, ikke hvad vi vil.«

Florence Nightingale var ikke blot en af de største
og klogeste Menneskevenner, der nogensinde har levet,
men hun var en selvstændig Tænker. Stuart Mill's Logik,
Jowett's Platonforskning og Quételet's Statistik udgjorde
det Grundlag, ud fra hvilket hun byggede sine
egne Tanker. Det vidunderlige ved Gud er efter hendes
Opfattelse netop det, at han er ubønhørlig: hvis han
vilde forandre Mening efter Folks Bønner, vilde vi jo
være overgivne paa Naade og Unaade til dem, der bad
til ham! Den sande Ide om Gud, der ene er forenelig
baade med den dybeste Moral og med den mest omfattende
Betragtning af Naturen og Historien, er den
Idé, der viser ham, ikke som en Individualitet, men
som et universelt Væsen --- som en Lov. Hvilke de enkelte
Love er, det maa vi finde ved Videnskabens Hjælp.
Og paa Grundlag af Naturens Lovmæssighed, et Grundlag,
i hvilket Guds Væsen er givet, stiller der sig saa
for hvert Menneske den Opgave at gøre det bedste ud
af sin egen Personlighed. Skønnere end nogen Sætning i
de kirkelige Ritualer var i hendes Øjne den Bøn af Sokrates
i Platon's Fajdros, som allerede ovenfor er omtalt:
giv mig at faa et skønt Indre, og at de ydre Goder
maa stemme med de indre! Karakteristisk for Florence
Nightingale er det ogsaa, at Skikkelsen som Andromache
og Antigone i hendes Øjne stod over de Kvinder, Bibelen
omtaler\footnote[1]{\emph{Edward Cooke}: \textit{Florence Nigthingale}. Dansk Oversættelse. (1917). II. p. 432—442.}.

d. Naar nu Bønnen --- med en eller anden Motivering
--- betragtes som Gave, altsaa som Virkning i Menneskets
% --- p. 184 ---
\setcounter{footnote}{0}%
\opage{184}Sind, ikke blot som udsprungen af Menneskets Sind som
Aarsag, bliver den til en symbolsk Akt, til Poesi, og
hvis den er alvorlig og ærlig, til Livspoesi. Allerede Enhedsbønnen
medfører denne Konsekvens. Ti er man vis
paa at have Gud i sig, at være ét med ham, saa hans
Villie er vor egen Villie, saa er der intet at bede om,
ligesaa lidt som Gud selv kan have noget at bede om.
Dobbeltheden af Bønnen og det, der bedes om, falder
bort med Dobbeltheden af den Bedende og den, der
bedes til. Men naar Bøn og Bønhørelse bliver Et, viser
det sig ganske klart, at Bønnens Formulering i Ord ikke
er nødvendig, eller kun har poetisk Betydning. Trangen
til Formulering, med Henvendelse til et fra den Bedende
forskelligt Væsen, vil snarere hænge sammen med en
Trang til billedlige Forestillinger end med selve den
Trang, der søger Tilfredsstillelse. Hos nogle Mennesker
raader en Angst for det Ubestemte, en Sky for det Uoverskuelige;
de trænger til bestemt udformede Forestillinger,
helst saadanne, der gennem Slægternes Række har faaet
Karakter af aandelig Hjemlighed. Hos andre Mennesker
raader der netop en Sky for det Begrænsede, en Angst
ved snevre Forestillinger; de opdager strax, hvor enhver
Lignelse halter, og hvor ethvert Billede viser sig utilstrækkeligt.
En virkelig Livspoesi peger ud over denne
Modsætning; ti i den er ethvert Billede, der bruges, Udtryk
for en personlig Erfaring om det indre eller ydre
Livs Kaar og Skæbne og for en Stræben efter at fastholde
sit Forbillede. Det er et Livsideal, Sindet her vender sig
imod, --- exemplar vitæ, som Spinoza kaldte det, en
ideal socius, som James har kaldt det. Selve Positivismens
Grundlægger, Auguste Comte, henvendte den »solennelle
effusion des sentiments généreux«, som han kaldte sin
Bøn, til Menneskeheden, der for ham repræsenteredes
ved en Kvinde, som havde været for ham, hvad Bea%
% --- p. 185 ---
\setcounter{footnote}{0}%
\opage{185}trice var for Dante\footnote[1]{Smlg. \emph{Auguste Comte}: \textit{Catéchisme positif}. p. 105 f. La vraie nature de culte intime consiste dans l'adoration journalière des meilleures personifications que nous puissions assigner à l'Humanité --- I \textit{Politique positivs}. (IV p. 114) defineres priere som commémoration suivie d'effusion.}. Det vil for mange være naturligt
at give en saadan dyb Stemning, som fremkaldes af
Stræben efter at fastholde Livsidealet, hvilke Kaar og
Skæbner vi end brydes med, Udtryk i Ord. Naar en
med Inderlighed og Kraft frembrydende Tanke faar Luft
i Ord, naar der er andre tilstede, hvorfor skulde dette ikke
ogsaa kunne ske i Enrum? En Mester har jo sagt, at man
skal lukke sin Dør, naar man beder. I Fællesbøn, der
forud er formuleret og skal holdes paa bestemt Tid og
Sted, er der altid noget kunstlet, som frastøder mange
Naturer, ikke mindst saadanne, hvis indre Liv udfolder
sig paa ejendommelig og dyb Maade. I deres bedste
Tanker, i deres inderligste Følelser rører sig for dem
det Højeste, de kender, og de kan meget godt fastholde
Troen paa, at der ogsaa i andre Menneskers Indre rører
sig noget tilsvarende, uden at de behøver fælles Formler
for at slaa det fast.

At Bønnen er Virkning udelukker naturligvis ikke, at
den er Aarsag. Den er ikke blot selv Tegn paa Udløsning
af Energi, men den fører selv til ny Udløsning.
Det er fra denne Side set, at Bønnen, der ofte kaldes
det centrale religiøse Fænomen, ikke uden Ækvivalent
kan forsvinde fra den aandelige Økonomi. Og vi har
set Muligheden af, at den naive Bedebøn gennem Enhedsbøn
og gennem Bøn betragtet som Gave kan gaa
over til at blive et frit Udtryk for Mod, Tillid og Begejstring
eller for inderlig Fastholden af et Livsideal.
Der er sikkert ikke faa, som i deres Liv har gennemløbet
denne Typerække. Om de i den afsluttende Livs%
% --- p. 186 ---
\setcounter{footnote}{0}%
\opage{186}poesi kan finde fuld Erstatning for, hvad der paa de
tidligere Stadier blev dem tildel gennem Bønnen, kan
være en vanskelig Sag for dem at afgøre, og endnu vanskeligere
for andre. Og om en saadan Typerække, som
vi paa Grundlag af sammenlignende Betragtning kan
opstille, ogsaa vil gennemløbes af Slægten idethele, eller
om der ikke til alle Tider samtidig vil forefindes de forskellige
Typer, vi har beskrevet, det er ogsaa en vanskelig
Ting at afgøre. ---

Naar Talen er om Ækvivalens for tidligere religiøse
Livsformer, er der en Ting, man ikke maa glemme. Bedebønnen,
i dens naive Former, er vel den af de ovenfor
omtalte Typer, man vanskeligst vil kunne tænke sig en
Erstatning for. Den hører da ogsaa især hjemme i Tider,
hvor Religion stod som den eneste Aandsmagt. Der var
endnu ikke sket nogen Arbejdsdeling paa det aandelige
Omraade. Det er nu anderledes. I det delte Arbejde kan
Religion (hvad enten man tager Ordet i dets videste
eller dets snevreste Betydning) ikke træde ind som et
enkelt Led, men den maa enten være alt, eller staa over
alt, eller være den dybeste Forudsætning i alt specielt
Arbejde uden selv at have et eget Arbejdsomraade.
Udenfor disse tre Muligheder vil der ingen Vej kunne
findes. --- Ganske vist prøves der undertiden, ja maaske
ret ofte, paa en ejendommelig Arbejdsdeling mellem
Religion og Kultur, saa at de gaa Side om Side uden
at have noget med hinanden at gøre. Saaledes skal den
bekendte Astronom Pater Secchi have sagt: »Der er en
Trappe, som fører op fra mit Oratorium (Bedekammeret)
til mit Observatorium. Naar jeg er kommen op over et
vist Trin, tænker jeg kun paa Observatoriet; før jeg er
kommen til dette Trin, tænker jeg kun paa Oratoriet.«
Saa udvortes kan Forholdet i Længden ikke blive ved
at være. Man kan naturligvis ikke tænke paa alle Ting
% --- p. 187 ---
\setcounter{footnote}{0}%
\opage{187}lige energisk i samme Øjeblik. Men man maa dog sætte
sig den Opgave at forene Erfaringerne fra Oratoriet i
Forbindelse med Erfaringerne fra Observatoriet. Og selv
om man ikke kan og vil stille sig denne Opgave, saa vil
Erfaring vise, at der er en underjordisk Forbindelse mellem
de to Grupper af Erfaringer, som en skønne Dag
kan gøre sig gældende i og for den bevidste Tanke. ---

Den inderligste Samling om og Hengivelse til, hvad
der for os er Idealet, maa vise sin Kraft paa selve vor
Arbejdsplads, ligesom den selv er Frugt af et indre Arbejde,
vi maaske ikke blive os bevidst.

Efterskrift. --- Efterat denne Skizze var nedskreven,
kom Friedrich Heiler's Bog: Das Gebet.
Eine religionsgeschichtliche und religionspsychologische Untersuchung
(1918) mig i Hænde. Der er i denne Bog samlet
et stort Materiale til Belysning af Bønnen og dens
Historie. Det falder mig ikke ind at sammenligne min
lille Skizze med Heiler's udførlige Fremstilling. Jeg har
i den blot villet sammenstille nogle Exempler, jeg i
Aarenes Løb har samlet, for at paavise, hvilke Ændringer
Bønnen kan undergaa, indtil der indtræder et
Stade, hvor der ikke længere kan tales om Bøn i egentlig
Forstand. Forfatteren karakteriserer udmærket de to
store Typer, han kalder den mystiske og den profetiske
Bøn, og han anser den profetiske for den højeste. Han
kræver, at Bønnen skal maales efter dens egen (religiøse)
Norm og afviser den »filosofiske Norm«, der anlægger
etiske og rationelle Synspunkter. Men han erklærer, at
det er en »Forgrovelse« af Bønnen, naar den forbindes
med Troen paa, at Guds Villie kan forandres, en Tro,
der ikke blot ligger til Grund ved den primitive Bøn,
% --- p. 188 ---
\setcounter{footnote}{0}%
\opage{188}men ogsaa ved den profetiske Bøn, saaledes endnu ved
Luthers Bøn. Skulde det ikke være en »filosofisk Norm«,
der her gør sig gældende hos ham? Ud fra selve den
profetiske Bøns Natur kan i hvert Tilfælde Korrektivet
ikke være hentet, ti for den er Bønhørelse et væsentligt
Moment, og det betegner en stor Krise, naar Bønhørelse
ikke indtræder. Vilde Bønnen uden Indflydelse af etiske
og rationale Motiver have kunnet udvikle sig fra primitiv
til mystisk og profetisk Bøn, end sige til den meget idealistiske
Form, Forfatteren (iøvrig uden nærmere Begrundelse)
stiller højst? --- Det er stadig Indflydelse af de
nævnte Motiver, der tilsidst fører til et Stade som det,
jeg har antydet i Slutningen af min Skizze.

% --- p. 189 ---
\setcounter{footnote}{0}%

\opage{189}\essayhead{TO PRÆSTER}

\essaysource{Otto Møller og Skat Rørdam. En Brevvexling. København 1915---16.}

Den her foreliggende Brevvexling mellem to betydelige
og alvorlige kirkelige Personligheder fra vore
Dage har baade religionspsykologisk og religionshistorisk
Interesse. Til enhver Tid har en Religion et større Indhold,
et større Indbegreb af Lærdomme og overleverede
Forestillinger, end de enkelte gør Brug af, og kan
gøre Brug af. Det er da af stor Interesse at se, hvilke
Elementer i en foreliggende Religion der til given Tid
og hos enkelte bestemte Personligheder viser sig at øve
størst Indflydelse, --- hvilke der i Virkeligheden leves
paa. Hvad der i denne Brevvexling først vakte min Opmærksomhed
og gav mig Lyst til en lille Studie, var
de to Mænds Forhold til et Problem, som i mine unge
Dage greb mig saa stærkt, at det blev bestemmende for
hele min Livsgang, og som jeg senere atter og atter har
søgt at orientere mig overfor. Det er Spørgsmaalet om
Forholdet mellem »det nye Testamentes Kristendom« og
»moderne Kristendom«. Dog er der ogsaa andre Spørgsmaal,
som denne Brevvexling kan bringe til at tænke
over, saaledes Forholdet mellem Selvoplevelse og Overlevering,
og i Forbindelse hermed Forholdet mellem den
enkelte og det kirkelige Samfund, han tilhører. Men
ogsaa rent menneskelig frembyder Brevvexlingen Inter%
% --- p. 190 ---
\setcounter{footnote}{0}%
\opage{190}esse ved det smukke og varme Indtryk, man faar af de
to Mænds Venskabsforhold, der bestod med Inderlighed
og Trofasthed fra de unge Dage, indtil Døden skilte dem.
Brevvexlingen gaar fra 1859 til 1909, og Venskabet var
grundlagt forinden. Der var Skær nok, det kunde være
strandet paa, og det er tiltalende at se, hvorledes en
dyb Samstemmen i det, der for dem begge var det væsentlige,
hjalp dem ud over Vanskelighederne. Det kan
ikke nægtes, at Læserne her har mest Grund til at beundre
Otto Møllers Højsindethed og Fredsvenlighed. Den
konservative og ortodokse Rørdam udtrykker ofte sin
Misfornøjelse med sin Vens Kætterier paa en meget
skrap Maade. Man kommer til at tænke paa, hvor vanskeligt
det maa være for Mennesker, der staar paa forskellige
Stader, at komme til Forstaaelse, naar man ser,
hvor nær en Spaltning kan ligge selv for Mennesker, der
i det væsentlige vil bygge paa samme Grund.

\romanhead{I}
% CHECK roman numeral read

Otto Møller karakteriserer sig selv som en Melankoliker,
som lever i noget, der er langt borte fra alt, hvad
der rører sig i Øjeblikket, og som let føler sig forvirret,
naar Tiden og dens Rørelser trænger sig ind paa ham.
Hele sit Liv har han, afset fra Skolegang og Studentertid,
levet i Gylling Præstegaard, hvor han afløste sin Fader.
Han var en Grubler og en Lærd. »Jeg har altfor megen
Tid til at gaa og tygge paa mine egne Tanker«, siger
han selv. Og hans Studier gjorde ham til en af Danmarks
lærdeste Præster, --- hvad der ganske vist, efter hans og
Rørdams fælles Mening, ikke betød ret meget. Han har
været et aandeligt Midtpunkt, til hvilket mange søgte
hen for at vinde Ro og Klarhed over kirkelige og religiøse
Spørgsmaal. En begejstret og trofast Tilhænger af
Grundtvig var han; men det var Grundtvig fra den tid%
% --- p. 191 ---
\setcounter{footnote}{0}%
\opage{191}ligere Tid, der for ham var den egentlige Grundtvig. I
Møllers Øjne var Mesteren, da det fra Fyrrerne af gik
opad for ham, tilbøjelig til at se for lyst paa alting og
til at gøre Tingene lettere, end de var og skulde være.
Han mente ikke, at Grundtvig's sidste Tid svarede til,
hvad hans stærke Aar havde lovet. I hvert Tilfælde kom
nu de fleste Grundtvigianere for let til det. Der var for
mange Møder, for megen Sang og Klang. »Der pjankes
afskyelig med Kristendom\dots{} Gud fri os for den glade
Kristendom; den kan blive værre end aandelig Død,«
skriver han 1860. Og langt senere (1892) skriver han:
»Jo ældre jeg bliver, jo mere faar jeg Fornemmelse af
det ensidige og blødagtige ved det gamle »ægte« Grundtvigianske\dots{}
Grundtvig selv var en ejendommelig Personlighed,
helstøbt i sit Slags; men hos hans Disciple
blev det til noget manieret noget.«

I Forhold til den forjættelsesrige Foraarstid han
havde oplevet i sin Ungdom, følte han sig fattig paa
sine gamle Dage, naar han saa sig omkring i Tiden. Samtidens
Præster havde han, ligesom hans konservative
Ven, meget liden Respekt for. Bisperne er »tummelumske«,
--- og Præsterne er »nogle Skarn overalt«. »Gejstligheden
er altid sig selv lig og under ikke Menigheden
nogen virkelig Frihed«. Og idet hans politiske Konservatisme
forener sig med hans radikale Kirkepolitik, siger
han (1906): »Jeg har kun liden Smag for at regeres af
Degne, men endnu langt mindre af Præster«. --- Det er
ikke blot Grundtvigianerne, det gaar ud over. Ogsaa de
Højkirkelige (den tredie Retning) og den indre Mission
udtaler han sig stærkt imod.

Mismodet og Følelsen af skarp Modsætning overfor
Tidens Retninger, selv saadanne, der som han selv vilde
staa paa Kirkens Grund, forklares vel for en Del ved
hans Temperament, ved hans ensomme Liv og ved hans
% --- p. 192 ---
\setcounter{footnote}{0}%
\opage{192}Sygelighed; men ogsaa kun for en Del. I en kort Karakteristik
af de to Venner taler Bispinde Rørdam (Tilbageblik
paa et langt Liv, p. 118) om Otto Møller som »Manden
med den udprægede kritiske Sans og det tunge Syn«.
Han er mere Tænker, mere konsekvent og skarp i sin
Tankegang end Vennen. Og hans kritiske Holdning betinges
ikke blot af hans indre og ydre Kaar, men først
og fremmest af hans altid vaagne Eftertanke. Selvfølgelig
rører Eftertanken sig kun indenfor hans egen kirkelige
Forestillingskreds, men her fører den ham til en
Stilling overfor forskellige teologiske og kirkelige Spørgsmaal,
der medfører stadig Misbilligelse fra hans mere
begrænsede og autoritetsfangne Vens Side.

Begge Venner er udpræget kirkelige Naturer. De har
begge deres sidste Tilhold i »Menigheden«, et Begreb,
det viser sig meget vanskeligt at give et bestemt Indhold
og Omfang. De er begge inderlig overbeviste om,
at der i den Overlevering, i hvilken de i den Del af deres
Ungdom, de kaldte »Foraarstiden«, har levet sig ind, er
givet det for Mennesker ene fornødne, men ogsaa i sig
selv indholdsrige og dybe Udtryk for den højeste Sandhed.
Forskellen mellem dem kommer frem dels i deres
Opfattelse af Kirkens Forfatning, dels i deres Stilling til
den Ret, personlig Erfaring har i Forhold til Overleveringen,
dels endelig i det Omfang, i hvilket de tillader
friere Eftertanke at udfolde sig paa det religiøse Omraade.

Skat Rørdam var, efter hans Hustrus Udsagn, en
mere sangvinsk Natur end hans Ven, med lysere Syn
og større Frejdighed. Det lysere Blik ytrede sig ikke
mindst i, at han havde større Tillid til de bestaaende
kirkelige Institutioner og, især efterat han var bleven
Sjællands Biskop, havde større Tro til, at de kunde virke
til Held. Han er ogsaa dels paa Grund af sin Optimisme,
% --- p. 193 ---
\setcounter{footnote}{0}%
\opage{193}dels paa Grund af mindre kritisk Sans og Trang, mere
ortodox end sin Ven, hvad der ofte siges ham i Brevene.
Men man faar Indtrykket af en Mand, der lever helt
og fuldt i den Overlevering, han nu engang er ført ind
i. De Slutninger, hans Eftertanke fører ham til, stemmer
uvilkaarlig med de Slutninger, hans Kirke forlængst
havde draget\footnote[1]{Se min Karakteristik af »kirkelige Naturer« i min Religionsfilosofi § 96.}. Otto Møller har derimod et større Spillerum
mellem det kirkelige Grundlag og de Tankegange,
der kan anslaas ud fra det, og han baade saa og fulgte
flere Muligheder, end hans ortodoxe Ven indlod sig paa.

Meget skarpt udtalte Rørdam sig ofte, ja som oftest,
om Møllers Forfattervirksomhed. Om Møllers Bog om
Evangelierne siger Rørdam, efterat have læst den i
Manuskript, at den kun kan blive til Glæde for Herrens
Fjender, og det er idethele et meget hvast Brev, han
sender sin Ven; man forstaar, at denne, da han læste
det, havde Fornemmelsen af at faa en Spand koldt Vand
over Hovedet. Da Møller, trods Rørdams indtrængende
Formaning, alligevel udgav det, ønskede Rørdam, at det
maatte blive læst saa lidt som muligt for at gøre saa
lidt Ulykke som muligt. Sagen var, at Møller havde udtalt
sig meget stærkt om Markus og Lukas, som han
satte langt under Matthæus. En saa afgørende Modsætning
burde man efter Rørdams Mening ikke antage mellem
Forfattere, der dog alle var inspirerede af den hellige
Aand. --- I en senere Bog angreb Møller den ortodoxe
Forsoningslære, ifølge hvilken Gud kun ved et
blodigt Offer, der maatte bringes af hans egen Søn,
kunde vise Naade mod Menneskene. Istedetfor Forsoningslæren
stillede han Læren om Forløsningen, der
ikke forudsatte nogen Forandring hos Gud fra Vrede
til Kærlighed, men kun en stadig guddommelig Villie til
% --- p. 194 ---
\setcounter{footnote}{0}%
\opage{194}Hjælp og Udfrielse. Da Møller sendte Rørdam Bogen,
var han vis paa, at hans ortodoxe Ven vilde »ryste sit
gamle Hoved over den«. Og det slog ogsaa til. Rørdam
var meget forarget over, at Møller ikke vil indrømme,
at »der er ydet Gud noget ved Kristi Død«, og over,
at han i Grunden betragter den ortodoxe Forsoningslære
som hedensk. Møller viser sig ogsaa her som den mest
konsekvente. Ti en Gud, hos hvem der ikke er Forandring
eller Skygge af Omskiftelse, kan naturligvis ikke
ved nogen Begivenhed eller Handling gaa over fra Vrede
til Kærlighed. Naturligvis kan heller ikke Møller her
tage alle Konsekvenserne. --- Ogsaa en følgende Bog af
Møller over et beslægtet Emne er Rørdam »saa særdeles
utilfreds med«, Møller vil nemlig ikke indrømme, at Gud
er vred paa Synderne; han henlægger Vreden til Dommens
Dag og opfatter Kristi Død som Vidnesbyrd om,
at Gud aldrig har været vred paa Menneskene. --- Den
Modsætning, der paa disse Punkter fremtræder mellem
de to Mænd, er ikke ny i Kirkens Historie, men gaar
langt tilbage. Striden mellem de Synspunkter, der betegnes
ved Ordene Forløsning og Forsoning, kom i Oldtiden
frem i Modsætningen mellem de østerlandske og
de vesterlandske Kirkefædre (Athanasius og Augustinus),
i Middelalderen i Modsætningen mellem Abailard og Anselm,
og paa Reformationstiden i Modsætningen mellem
Zwingli og Luther.

Møller tager de gentagne Tilrettevisninger med stor
Sagtmodighed. »Jeg ser jo mere og mere, at Du ikke
kan forstaa mig, og det er Du ikke ene om; men det
kan dog ikke hjælpe mig til at se anderledes, end jeg
nu engang er kommen til at se paa disse Ting.«

Ogsaa i Opfattelsen af Jesu Personlighed og Skæbne
er de to Venner uenige. Møller har skrevet en Fremstilling
af den evangeliske Historie, i hvilken han søger
% --- p. 195 ---
\setcounter{footnote}{0}%
\opage{195}at vise, at Jesus ved sin Optræden led en hel Række
Skuffelser, idet hverken de Lærde i Jerusalem eller Menigmand
i Galilæa anerkendte ham, ja, end ikke Apostlene
forstod ham. Rørdam forarges. Skulde det da være gaaet
Jesus som en menneskelig Reformator, der maa gennemgaa
saa mange bitre Kvaler formedelst Modstand? Hvor
kan han da være Herren og Frelseren? Det er moderne
Snak, Møller har ladet sig betage af, og som maa føre
til at betragte Jesus som blot Menneske! --- Møller finder
det bedrøveligt, at han maa opgive at skrive noget,
der er hans ældre Ven tilpas. Men han erklærer, at han
har taget Sagen rent historisk, og at han betragter de
store Skuffelser som hørende med til Jesu Lidelse. Jesus
har haft Ønsker og Haab, der ikke opfyldtes. Kan man
være sandt Menneske uden det? At Rørdam kalder denne
Opfattelse moderne, rører ikke Møller: ti »hvor meget ser
ikke nyt ud, som dog er sandt?«

Overalt bekræfter det sig, hvad Møller etsteds siger,
at Rørdam har en særegen Kærlighed til det Overleverede,
medens han selv har en Drift til at gaa bag ved det Overleverede.
Dette vilde let have kunnet føre dem bort fra
hinanden, og havde maaske ogsaa gjort det, hvis ikke
mundtlig Samvær og Samtale fra Tid til anden havde
afløst Drøftelsen i Breve. ---

\romanhead{II}
I nøje Sammenhæng med de allerede omtalte Differenspunkter
mellem de to Præster staar deres forskellige
Opfattelse af Forholdet mellem Selvoplevelse og
Overlevering.

Det er en farlig Sag, siger Otto Møller, at tale meget
om, hvad man endnu ikke selv har oplevet. Men de fleste
Mennesker lever paa Autoritetstro, ikke paa Selvoplevelsens
fulde Vished. De tror paa, hvad Paven eller Luther,
% --- p. 196 ---
\setcounter{footnote}{0}%
\opage{196}Kirken eller Skriften siger, altsaa paa anden Haand. Det
gaar kun i stille Vejr. Det, man kan vide paa første
Haand, er ganske vist ikke meget, --- egentlig kun Modsætningen
mellem Synd og Naade; men der kan vise sig
at være mere indeholdt deri. Af afgørende Vigtighed er
det at begynde fra den rigtige Ende, nemlig fra Syndsbevidstheden.
Men om Selvoplevelsen kan føre til Tro,
betinges ved, at Aabenbaringen er bleven forkyndt for
Mennesket. --- Her brydes Møllers kirkelige Standpunkt
med hans Tilløb til en rent psykologisk Betragtning.
Allerede Syndsbevidstheden i den Betydning, Møller
fremsætter den, er dog et historisk betinget Fænomen
og forudsætter Forkyndelse. Det er ikke en almenmenneskelig
Oplevelse. Troen er egentlig allerede Forudsætning
for Syndsbevidsthed; der maa være Tro paa en overnaturlig
Autoritet, overfor hvilken Menneskets Tragten
kan vise sig som Synd. Møller er her selv for autoritetstro
overfor den Lutherske Lære om »Frelsens Orden«,
som der iøvrig allerede paa Reformationstiden blev opponeret
imod. Kan en Selvoplevelse ikke begynde med
Beundring og Begejstring for et Ideal, og vil Følelsen
af egen Afstand fra dette Ideal ikke blive des dybere,
naar Selverkendelsen vækkes ved Tanken om Idealet?
Naar man appellerer til Selvoplevelsen maa man ogsaa
anerkende, at baade dens Begyndelse og dens Forløb
kan være højst forskellig hos de forskellige Personligheder.
Lutherdommen har slaaet en eneste Type fast,
som alle mere eller mindre voldsomt skulde presses
ind i\footnote[1]{Se om forskellige religiøse Typer i min Religionsfilosofi § 36—44 og 94—96, og Afhandlingen om nogle religionsfilosofiske Arbejder fra den nyeste Tid (Mindre Skrifter. Anden Række).}.

Med sin Sans for Problemer anerkender Otto Møller
nu ogsaa, at der her for en udpræget kirkelig Opfattelse
% --- p. 197 ---
\setcounter{footnote}{0}%
\opage{197}ligger en uløselig Gaade. Kan der ikke være den inderligste
Trang og Stræben tilstede, uden at den Forkyndelse,
der ene kan føre til Tro, bliver Mennesket tildel?
Han indrømmer, at naar Tro paa ydre Autoritet helt er
fejet bort, er der intet Raad mere. --- Ja, Møller burde
have været mere Pessimist, end han er; ti har man ikke
nogen Tro, eller rettere noget Haab til, at der kan opstaa
nye Idealer ad rent menneskelig Vej, og at de kan
vække den Selverkendelse, uden hvilken Livet ingen
Sandhed faar, saa er alle Sunde lukte.

Skat Rørdam er nu fra første Færd betænkelig ved
sin Vens stærke Vægt paa Selvoplevelsen, fordi den synes
ham at føre til at bygge paa bar Grund. Mennesket,
siger han, begynder med Autoritetstro og bliver ved dermed
hele sit Liv. Og selv i den største menneskelige Nød
maa Mennesket dog kunne skimte og ane Troens Værdi
og Autoritetens Magt. Forudsætningen er, at Mennesket
lever i et kristent Land. Hvorledes det gaar, naar denne
Forudsætning ikke er tilstede, taler Rørdam ikke om.
Hans mere skarpt tænkende Ven ser her en uløselig
Gaade, fordi han saa bestemt betoner baade det personlige
Moment, Selvoplevelsen, og det objektive Moment,
Aabenbaringen og Forkyndelsen.

Rørdam stoler mere paa Menighedens Indflydelse paa
de enkelte end Møller. I den nu bestaaende Menighed
lever der ifølge Rørdam et Vidnesbyrd, som gaar tilbage
til den apostoliske Menighed og til Apostelskrifterne, og
denne apostoliske Menighed er den af Gud satte Begyndelse.
Men Rørdam indrømmer, at der indenfor de
apostoliske Skrifter, der kommer til de enkelte gennem
Menigheden, som »Menighedens Arvegods«, ikke kan skelnes
en Gang for alle mellem, hvad der er af Gud, og hvad
der er menneskeligt; enhver Kristen maa drage Grænsen
efter sin gennem Livet i Menigheden opfostrede kriste%
% --- p. 198 ---
\setcounter{footnote}{0}%
\opage{198}lige Bevidsthed. Og dog er, som Rørdam hævder i Modsætning
til »den moderne Teologi«, Kirken »Sandhedens
Pille og Fodfæste«.

Otto Møller synes ikke at nære den Tillid til Menighedens
Indflydelse, som Rørdam. Han mener, at »naar
Menigheden ikke anvises til det, tager den ikke Forargelse
af ret meget, men lader helst enhver feje for
sin egen Dør«. Og han er imod en Kirkeforfatning, fordi
det under en saadan i Grunden vil være Præster og Teologer,
der kommer til at raade. --- Spørger man, hvor
»Menigheden« saa bliver af, antager jeg, at Møller vilde
henvise til Sakramenterne. Han er ikke bange for at
give Tanken et vist Spillerum, naar blot det væsentlige
i Kultus, Sakramenterne, værnes. Overholdes blot dette,
kan der tillades »adskillige Krumspring« i andre Henseender.
--- En Løsning ligger naturligvis ikke heri, ti
Sakramenterne selv forudsætter, baade hvad deres Overlevering
og deres Tydning angaar, et dogmatisk Grundlag
--- og hvorfra kommer det? ---

Kun den katolske Kirke kan her for Alvor slutte af.
For den er der her ingen Gaade. Der er heller ingen uløselig
Gaade, skønt nok et Problem, for dem, der staar
udenfor al konfessionel Religion. For dem bestaar Sandhedskriteriet
i den mest omfattende og mest sammenhængende
Erfaring, som det er muligt at vinde, og der
stilles her, hvad Forholdet mellem Individ og Samfund
angaar, stedse nye Spørgsmaal, der maa finde deres Besvarelse
gennem historisk og psykologisk Undersøgelse
af, hvorledes Individet udvikler sig under Samfundets
Indflydelse til større eller mindre Selvstændighed, og
hvorledes Individet paa sin Side kan virke tilbage paa
Samfundet og blive medbestemmende for dettes videre
Udvikling. Intetsteds staar vi her overfor noget absolut
første eller sidste, skønt det kan være vanskeligt at
% --- p. 199 ---
\setcounter{footnote}{0}%
\opage{199}følge de Veje, ad hvilke den sociale Indflydelse gør sig
gældende hos de enkelte, og at efterspore individuelle
Initiativers første Udspring. En »uløselig Gaade« er der
kun, naar man gaar ud fra en absolut Autoritets Nødvendighed
og dog ikke vil være Katolik.

\romanhead{III}
Jeg har hidtil overvejende dvælet ved Spørgsmaal, i
hvilke de to Venner stod imod hinanden. Der er naturligvis
mange Punkter, ganske afset fra deres fællesgrundtvigske
Stade, i hvilke de staar sammen. Naar deres Uenigheder
har religionspsykologisk Interesse ved at vise, hvor
forskelligt et Syn paa Livet der kan findes hos Mennesker
paa samme Stade, saa har deres Overensstemmelser Interesse
ved at vise, hvorledes alvorlige Præster i vore
Dage se paa de aandelige Bevægelser, der kan faa
Følger for deres Kirke, og paa selve det historiske Grundlag
for deres Kirke.

De stemmer overens i deres Betragtning af Nygrundtvigianismen,
i hvilken de dels se en Virkning af, at man
er kommen for let til sit religiøse Stade, dels et Udtryk
for en umoden Trang til at bringe kirkelig Opfattelse i
Forbindelse med Retninger i Tidens Literatur og Politik.
De føler begge Uvillie overfor den indre Mission og lægger
en utvetydig Foragt for Dagen overfor »den tredie Retning«
og den raadende Teologi. De se ogsaa ens paa de
nyere Forskninger over de bibelske Bøgers Oprindelse,
Bibelens Literaturhistorie, eller, som man, fordi den i
mangt og meget forkaster de overleverede Meninger om
disse Bøger, plejer at kalde den, Bibelkritiken. De er
begge meget betænkelige ved hele denne Bevægelse, men
det maa siges til deres Ære, at de klart se, at der her
er Spørgsmaal, som ikke kan skydes tilside. De mener
begge, at indrømmer man først de Grunde, der føre til
% --- p. 200 ---
\setcounter{footnote}{0}%
\opage{200}at henlægge »Mosebøgerne« til en langt senere Tid end
Moses', saa vil det kun være et Tidsspørgsmaal, naar
man skrider til at anvende samme Metode paa de nytestamentlige
Bøger. Er Jødedommen ikke af guddommelig
Oprindelse, vil Kristendommen, der bygger paa den,
heller ikke kunne være det. Som Møller udtrykker det:
gaar Welhausen ind i Folk, saa gaar Kristendommen ud!
Det er et Undermineringsarbejde, der foregaar. Men
efterhaanden indser de begge, at her er noget, man ikke
kan komme fra. Deres fælles Stade udtrykker Møller paa
sin kraftige Maade: »Det pinlige er nu, at jeg er overbevist
om, at der er noget om det med den gammeltestamentlige
Kritik; men den er falden i nogle forbandede
Hænder.« De føler begge, at de mangler Kundskaber
til at tage Sagen op, og haaber paa, at den rette Skriftkloge
vil komme, der kan hjælpe Menigheden over Vanskelighederne.
--- Om det vil hjælpe, at Undersøgelserne
kommer i andre Hænder, naar Metoden er den samme,
vil Tiden vise.

Fra »Bibelkritikerne« til »Fritænkerne« er der i de to
Venners Øjne ikke noget langt Skridt, og dog falde deres
Domme her langt skarpere. Som de kirkelige Naturer,
de begge, hver paa sin Vis, er, indlevede paa alle Kanter
i en Tankegang, der aldrig frigør sig helt fra Autoritet,
saa meget end den ene af dem betoner Selvoplevelse
og Eftertanke, kan de i det frie intellektuelle Arbejde
ikke se andet end Ondskab, naar dets Resultater strider
mod deres Tro. Rørdam opstiller det Princip, at »hvad
vi som Guds Børn har at spørge om, er ikke Grunden,
men Øjemedet«. Han anvender dette Princip paa den
ulykkelige Krig 1864 og tyve Aar efter paa Slotsbranden.
Begge Ulykker skulde opfattes som Guds truende Tale
til Folket. (Hvad Grunden til Slotsbranden var, véd han
iøvrig god Besked med: deri, at Slottet --- som han selv
% --- p. 201 ---
\setcounter{footnote}{0}%
\opage{201}saa med sine Øjne --- brændte saa hurtig, laa tydelig
nok, at det var Guds Villie, at det skulde brænde!). Det
er klart, at naar man har Adgang til at kende Guds
Formaal med det, der sker, eller med visse Skrifters Tilblivelse,
bliver Interessen for at finde Grunden til disse
Ting meget ringe, eller rettere, Grundene maa være
givne med Formaalene, ligge i den samme Villie, der
sætter Formaalene. Ud fra et saadant Standpunkt kan
der vanskelig naas Forstaaelse af den Sandhedskærlighed,
som kan føre til Opgivelse af overleverede Antagelser
paa det religiøse Omraade. For Rørdam er Fritænkere
»Guds- og Samvittighedsfornægtere«, der betragter
Menigheden som deres egentlige Fjende, og al Kristendom,
der ikke er den pureste Individualisme, for Hykleri.
Naar saadanne Mennesker arbejder for borgerlig Frihed,
er det kun for at benytte den til at undertrykke Kristendommen,
som de hader, fordi den dømmer dem.
Deres sidste Øjemed er at gøre Mennesket til Træl af
Naturen, at nyde istedetfor at lyde, at gøre Mennesket
ligt de Umælende. (Man ser, at Rørdam ved god Besked,
ikke blot med de guddommelige, men ogsaa med menneskelige
Formaal). Otto Møller gør vel nogle Forbehold
overfor denne kraftige Karakteristik; men ogsaa han
mener, at hvad der egentlig er ivejen med Fritænkerne,
er, at de elske Mørket mere end Lyset, og at Fejlen stikker
i Villien, ikke i Tanken. Der gives i Grunden ingen
undskyldelige Fritænkere.

Det vilde dog vist være lettere at komme ud af det
med Præsten i Gylling end med Sjællands Biskop. Han
har Sans for Selvoplevelsens og Eftertankens Betydning,
og han har Trangen til at gaa »bag det Overleverede«.
I Grunden maatte han blive nødt til at indrømme, at
det ikke er Retten til fri Forsken og selvstændig Livsførelse,
der skal begrundes, men derimod Autoriteters
% --- p. 202 ---
\setcounter{footnote}{0}%
\opage{202}Ret til at sætte Grænse for denne Ret. Dette Princip for
Bevisbyrdens Fordeling paa det aandelige Omraade grunder
sig derpaa, at kun i fri Hengivelse og paa Grundlag
af personlig Erfaring kan Sandheden blive Menneskets
virkelige Ejendom.

\romanhead{IV}
% CHECK roman numeral read

Jeg kommer nu til et Spørgsmaal, som er af ganske
særlig Interesse for Forholdet mellem moderne Kristendom
og Urkristendommen, nemlig hvorledes Kirkens Repræsentanter
stiller sig til den Forventning om Kristi
nære Genkomst og dermed om hele Verdensløbets Afslutning,
som er en af det nye Testamentes Grundtanker,
uden hvilken man navnlig slet ikke forstaar den nytestamentlige
Etik. For en gammel Kierkegaardianer er
der særlig Grund til at henlede Opmærksomheden paa
dette Punkt. Ti naar Kierkegaard endte med at erklære,
at det nye Testamentes Kristendom ikke er til, saa er det
den nytestamentlige Etik, han har for Øje. I dogmatisk
Henseende mente han ikke, at der var noget ivejen.
Dog saa Kierkegaard ikke, at det nye Testamentes Etik
bestemmes og betinges ved hin store Forventning, der
er det egentlige Indhold af Urkristendommens Tro. Han
var endog tilbøjelig til at vende Forholdet om og at betragte
Forventningen som en psykologisk set nødvendig
Illusion, uden hvilken den strenge urkristelige Etik ikke
kunde fastholdes. Dette er naturligvis urigtigt. Det er
Forventningen om det nære »Guds Rige«, der ligger til
Grund, og af hvilken Etiken, der under disse Forhold
kun kunde blive en »Interimsetik«, som en nyere Forsker
har kaldt den, er en nødvendig Følge. Man prøve engang
at læse det nye Testamente igennem med den stadige
Tanke om »Gudsrigets« umiddelbare, bogstavelige Nær%
% --- p. 203 ---
\setcounter{footnote}{0}%
\opage{203}hed, og Alt vil blive ganske anderledes forstaaeligt, end
det ellers er. Men hele Kirken, den katolske saavel som
den protestantiske, gaar her Rationalismens Veje, og det
af en naturlig Selvopholdelsestrang, idet Kirkens egen
fortsatte Existens bliver ubegribelig og uberettiget ud
fra hin Forudsætning om snar Afslutning. Og hele den
Betydning, man tillægger Kristendommen for menneskelig
Kultur og for Folkelivets Udvikling, kommer til at
mangle et Fundament. Rationalismen begynder tidlig,
allerede i andet Peters Brev, med den Sætning, at for
Gud er tusind Aar som én Dag, og omvendt. Men hvis
Gud nu netop havde forjættet noget til et aldeles bestemt
Tidspunkt? Rationalismens Gud har Evigheden
for sig, og Tidsforskelle svinder hen til intet; men Bibelens
Gud handler i Tiden og fører sit Rige til Fuldendelse
til bestemte Tider. Man kan ikke beraabe sig paa Jesus
og Apostlene paa alle andre Punkter og forkaste deres
Udsagn paa et Punkt af saa afgørende Betydning som
dette. Det er nemlig to helt forskellige Livsanskuelser,
der her kommer til at staa overfor hinanden. For den
ene maa alt menneskeligt Kulturarbejde stilles i Bero
overfor det ene Fornødne, som er at forberede sig til
den nære Afslutning af alt. Den anden lever og aander
i Arbejdet for menneskelig Livsudfoldelse i Samfundsforhold,
i Videnskab og i Kunst; her arbejdes der
paa langt Sigt, og en Afslutning eller Afbrydelse staar
kun som en fjern Mulighed. Mellem disse to Livsanskuelser
maa der, eller maatte der, vælges. Forskellen
mellem dem blev tydelig, da den levende Forventning
blegnede eller i hvert Tilfælde mistede sin umiddelbare
praktiske Betydning. Saa fik man Tid til mange Ting,
og kunde have sit Hjerte mange Steder, som der indenfor
Urkristendommen ikke kunde være Tanke eller Tale om.
% --- p. 204 ---
\setcounter{footnote}{0}%
\opage{204}Det er denne Forskel, der gør moderne Kristendom til en
anden Religion end Urkristendommen\footnote[1]{Jeg har tidligere udtalt mig om dette Spørgsmaal, først i Slutningsafsnittet \textit{Kierkegaard som Filosof} (1892), senere i min \textit{Religionsfilosofi} (1901). (§ 14. 59. 99. 123).}.

Kommer man med dette Problem for Øje til Brevvexlingen
mellem de to Præster, finder man vel nogle
Antydninger af, at de har skimtet det, men de har ikke
stillet det paa afgørende Maade og derfor heller ikke
besvaret det.

Rørdam er i Forlegenhed med Afstanden mellem
Døden og Dommen. Denne Forlegenhed eksisterer ikke
for det nye Testamente, da Tidsafstanden forudsættes at
være ganske kort, og for mange af den dalevende Slægt
slet ikke at eksistere. Der er, siger Rørdam ganske rigtig
(og Møller er enig med ham deri), ikke sagt noget i
Bibelen om, hvad der sker i hint Mellemrum, navnlig
intet om en mulig Udvikling af, hvad der ikke blev
fuldendt i og med de enkelte før deres Død. --- Men,
maa der føjes til, der var jo heller ingen Grund til at sige
noget derom, naar der intet længere Tidsmellemrum
forudsattes at skulle være.

Undertiden rører der sig hos de to Præster, under
Indtrykket af Tidens mange Ugudeligheder, en Formodning
om, at Enden er nær. De er enige om, at Tidernes
Tegn tyde stærkt hen paa Enden, og at Verden er bleven
gammel. Men som gode Grundtvigianere foruroliges de
saa ved, hvad der skal blive af Mesterens Spaadomme
om Folkelivets mægtige Udvikling. Der synes nok at
være Grøde; men der er endnu langt frem, og det ser
i Grunden sort ud allevegne nu (1873). Herefter maa der
altsaa synes at være langt til Enden. Otto Møller erklærer
da ogsaa, at naar han er enig i, at Tidernes Tegn
peger mod Enden, saa er det kun Hovedet, der slutter
% --- p. 205 ---
\setcounter{footnote}{0}%
\opage{205}saaledes; Hjertet har ingen Tro paa Herrens snare Komme,
og det er ogsaa for hans Vedkommende den Grødetid,
Grnndtvig har indledet, som han haaber maa føre
til noget trods alle Hindringer.

Til det Folkeliv, hvis Fremtid de to Præster haaber
paa i deres Hjerter, saa at det kun er med deres Hoved,
de kan fastholde Urkristendommens store, alt betingende
Haab, regner de utvivlsomt og saa Kirkelivet. De taler dog
ikke om, at ogsaa den kirkelige Udvikling vilde blive
afbrudt, hvis Genkomsten indtraf nu. Ikke blot Kultur
og Folkeliv, ogsaa Kirken har indrettet sig paa langt
Sigt. Der vilde blive Bestyrtelse blandt Biskopper, Præster
og Teologer, hvis Enden kom nu. Thor Lange har i
en lille Fortælling (»Kommet igen«. »Ude og Hjemme«
1882) vist, hvorledes Jesu Genkomst vilde vække Forstyrrelse
i den bestaaende Kirke. Kirken er nu engang
traadt i Jesu Sted og er slet ikke indrettet paa, at han
selv umiddelbart tager Styret. Lange har henlagt Scenen
til Middelalderen, men Skildringen passer ligesaa godt
paa den protestantiske Kirke, der jo endog har indrettet
sig nok saa hyggelig som den katolske.

Henved 30 Aar efter støder de to Venner igen paa
Problemet uden at mærke dets Brod. De undrer sig over,
at der ingen særlig Etik fremsættes i det nye Testamente.
Alt drejer sig om at vække og nære Troen paa Jesus;
saa forudsættes, at dermed er ogsaa Livsførelsens Art og
Retning afgjort. Særlig Rørdam beskæftiger sig meget
hermed. »Der er her adskilligt at tænke paa, som jeg
langt fra er kommen til Ende med, men som jeg tumler
med paa mine Morgenspadsereture.«

Spørgsmaalet finder sin simple Besvarelse, naar man
fastholder, at Tro i det nye Testamente er Tro paa Jesus
som den, der snart bringer Fuldendelsen. At »Himmeriges
Rige« er nær, er at tage i bogstavelig Forstand, trods
% --- p. 206 ---
\setcounter{footnote}{0}%
\opage{206}de kirkelige Omtydninger. Er det saa, da bæres Livet
af en stor Forventning, i Henhold til hvilken alt andet
blegner. Inderligt Sammenhold mellem dem, der lever i
denne Forventning, og Stræben efter at udbrede den
er naturlige Følger. Men om Samfundsordning, om Kunst,
om Videnskab kan man ikke vente at finde noget i det
nye Testamente, og man finder heller intet. Den nytestamentlige
Etik, som Kierkegaard ikke kunde finde i
den moderne Kristendom, og som Biskop Rørdam ikke
kunde finde i selve det nye Testamente, er Forventningens
Etik, der naturligvis kun kunde bestaa, saalænge
Forventningen var aktuel. De biskoppelige Morgentures
Problem finder her sin simple, historiske Løsning.

Paa sin Vis kommer Præsten i Gylling ind paa den
samme Side af Problemet. Han har grublet over Lignelsen
med de kloge og de daarlige Jomfruer, som handler
om dem, der venter paa Jesu Genkomst, og han er stanset
ved det Træk, at alle Jomfruerne, baade de, der havde
Olie paa Lampen, og de, der ikke havde det, faldt i Søvn.
Dette udlægger han paa følgende Maade: »Det er jo
ganske afgjort, at Ingen kan holde ud at leve i uafbrudt
Forventning af Herrens Komme; det fører kun til usund
og unaturlig Overspændthed«. Otto Møller gaar her ud
fra, at det skal være en lang Tid, der skal ventes i. Selv
en kort Ventetid, naar det er noget Stort og absolut
Afgørende, der ventes, kan forekomme lang, da »trange
Tider langsomt skrider«. Men Tiden til Genkomsten forudsættes
netop i sig selv at være kort. I den korte Mellemtid,
der er Tale om i det nye Testamente, kunde Forventningen
meget vel holde sig, og den har holdt sig hos
den første Slægt. Anderledes maa det naturligvis gaa,
naar der gaar Tusinder af Aar, uoverskuelige Tidsrum,
ud over hvilke maaske Hovedet, men ikke Hjertet kan
se. Men desuden --- et er det at sove, et andet er det,
% --- p. 207 ---
\setcounter{footnote}{0}%
\opage{207}som Kirken og Menneskene have gjort i den Mellemtid,
der jo stadig varer, nemlig at sætte Kraft og Interesse
ind paa vaagent og levende Kulturarbejde og Menneskeudvikling
og have baade sit Hjerte og sit Hoved optagne
deraf. Søvn er en Hvile, som de første Kristne har trængt
til ligesom alle Mennesker; de kunde meget godt fra den,
eller fra anden Hvile, vaagne op til den Forventning,
om hvilken hele deres aandelige Liv havde samlet sig.
Vanskeligere er det at vaagne op fra det Kulturarbejde,
i hvilket Kristne og Ikkekristne ofte tage Del med lige
Iver. Dobbeltheden, Modsætningen, og tilsidst Modsigelsen
mellem oprindelig Kristendom og menneskelig Kultur
kunde ikke belyses klarere end ved Otto Møllers Udlægning
af Jomfruernes Søvn. Denne Dobbelthed er moderne
Kristendoms egentlige Problem, saa liden Bevidsthed
der end synes at være derom hos Kirkens Repræsentanter.

Den katolske Kirke fastholder paa sin Vis Forventningens
Etik gennem sine Munkes og Nonners Livsførelse.
Derfor lyder det baade fra den hellige Teresa i
det 16. Aarhundrede og fra Kardinal Newman i det 19.
Aarhundrede: »Hvad vilde Verden blive, naar der ikke
var Munke!« Og ifølge Pascal varer Jesu Dødskamp lige
til Verdens Ende, og man tør ikke sove i al den Tid!\footnote[1]{Jesus sera en agonie jusqu'à la fin du monde; il ne faut pas dormir pendant ce temps-là! \textit{Pensées}, éd Brunschwicg. Paris 1904. II. p. 435.}.
Kun afløses her den levende Forventnings uvilkaarlig
omdannende Indflydelse af bevidst Askese. ---

Alt idealt Liv er et Liv i Forventning, og det har
sikkert (som jeg allerede har bemærket i min Religionsfilosofi)
været af stor Betydning, at den Religion, der fik
saa stor Indflydelse paa Europas Aands- og Folkeliv,
fra først af var en en Forventningens Religion. Men for
at blive Grundlag for fri menneskelig Livsførelse maa
% --- p. 208 ---
\setcounter{footnote}{0}%
\opage{208}Forventningen være knyttet til Opgaver, paa hvilke der
kan arbejdes indenfor den Verden, der er tilgængelig for
menneskelig Erfaring. Det er en stor, og ofte vanskelig
Motivforskydning, der maa foregaa ved Overgang fra
kristelig til human Etik.

\begin{center}* * *\end{center}

Læsningen af de to Præsters Brevvexling har ikke
ringe psykologisk og tidshistorisk Betydning ved det
Lys, den kaster over moderne Kirkemænds Stilling, dels
til deres egen Lære, dels til dermed sammenhængende
Forhold. --- Jeg tvivler ikke om, at de moderneste Kirkemænd
frembyder andre Karaktertræk. Tiden vil vise, om
de kommer til at staa anderledes overfor Problemerne
end deres Forgængere.

Endnu skal nævnes, at de to Venner paa deres gamle
Dage erfarede noget, som efter den enes Forklaring
hang sammen med deres Livsgerning. --- I det sidste
Brev fra Rørdam hedder det: »Det er ikke gaaet mig,
som jeg i mine yngre Dage havde ventet, at alt som jeg
blev ældre, skulde min Evne til at glædes i Gud og til
at se opad og fremad med et levende Haab voxe\dots{}
Men dette erfarer jeg saa lidt, for ikke at sige intet
af.« Og Otto Møller svarer: »Jeg sidder ikke rigere i det,
end Du gør\dots{} Jeg føler mig saa fattig, træt og afmægtig,
som vist ingen kan det, og synes tilmed endda,
at jeg er sløv og føler det ikke saa ydmygende, som jeg
burde.«

Det er maaske en Erfaring, som mange gamle Mennesker
maa gøre, og som finder sin naturlige Forklaring
i aftagende Livskraft. Men Otto Møller antyder en Forklaring,
der faar særlig Betydning ved at komme fra en
Mand, der lagde stor Vægt paa Selvoplevelse og alvorlig
% --- p. 209 ---
\setcounter{footnote}{0}%
\opage{209}Tilegnelse: »Mon vi ikke har talt og skrevet os trætte?«
--- Naar det hører til Ens Profession atter og atter at
tale om Ting, der kun har Værdi som Genstand for
levende Følelse, vil let den Tid komme, da man mærker,
at det indre Liv ikke kan holde Skridt med de Former,
der er lagt til Rette for det, og Skuffelser af den omtalte
Art vil ligge nær. En græsk Forfatter har med Rette
sagt, at af Menneskene lære vi at tale, af Guderne at
tie. En Smule hellig Taushed kunde maaske ogsaa vore
moderneste Kirkemænd have godt af.

% --- p. 210 ---
\setcounter{footnote}{0}%
\opage{210}\essayhead{MADVIG SOM FILOSOF}

\essaysource{J. N. Madvig: Livserindringer (1887). \\ Supplementer til Livserindringer af J. N. Madvig. Udgivne af J. N. A. Madvig (1917).}

Det var for mig af stor Interesse at læse den store
Filologs Selvbiografi, og ikke mindre de senere udgivne
»Supplementer«, der bl. a. indeholder hans filosofiske
Ideer og karakteriserer hans Stilling til det religiøse
Problem. Jeg har aldrig hørt til hans nærmere Kreds,
men jeg er i mine yngre Dage gentagne Gange kommen
i Berøring med ham og har haft Samtaler med ham om
filosofiske og religiøse Spørgsmaal. De Antydninger, han
ved saadanne Lejligheder gav, er blevne bekræftede og
tydeliggjorte ved de nu foreliggende Optegnelser.

Det første Kendskab til Madvig fik jeg, da jeg som
saa mange Skoledisciple studerede hans latinske Grammatik.
Det var naturligvis ikke altid fornøjelig Læsning, men
jeg erindrer fra de senere Skoleaar, hvorledes jeg med
stort Udbytte og med, hvad jeg nu vilde kalde logisk og
systematisk Interesse, skaffede mig et Overblik over hans
Fremstilling af det latinske Sprogs Syntax. Det var ved
Siden af Matematiken det i formel Henseende mest lærerige,
jeg kom i Berøring med i Skolen. Og naar jeg saa
undertiden, træt af at gaa fra Paragraf til Paragraf, slog
op paa Fortalen og læste hans Ytring om, at »denne
Grammatik er skreven med et Sind, der lader det nu%
% --- p. 211 ---
\setcounter{footnote}{0}%
\opage{211}værende Liv dets fulde Ret, alt imedens det vender sig
mod Fortiden som den, hvis rige Arv vi skal bevare«,
--- saa fik jeg et Udblik til en større Sammenhæng, i
hvilken mine Skolestudier kunde gaa ind, og til hvilken
de skulde aabne mig Adgang. Det er for mig nu paa mine
gamle Dage rørende at læse, at det kun var modstræbende,
han havde paataget sig at skrive en Skolegrammatik,
»at knytte sit Navn til Skolebænken«; han har vel
vidst, hvorledes de fleste Skoledrenge betragter Forfatteren
til de Grammatiker, de maa pines med. Han overvandt
denne Modstræben i Erkendelse af den Betydning,
Arbejdet kunde faa. Men den hang sammen med noget,
som endog engang førte til en Krise i hans aandelige Liv.
Det filologiske Studium, som han dyrkede saa frugtbart
og saa genialt, stod dog for ham som »den største og dybest
følte Ensidighed«. Han kunde ogsaa paa ældre Dage
ikke »taale at dø efterladende en uberigtiget Forestilling
om sig som væsentlig exklusiv Filolog«, --- skrev han i
Aaret 1869 (65 Aar gammel). Men allerede som ung Professor
følte han (mest skærende i Aarene 1830—32) paa
smertelig Maade denne Ensidighed. Hans fattige Kaar i
Barndom og Ungdom havde ikke levnet ham Tid til fri
og almindelig Orientering paa Aandslivets Omraader,
eller, som han udtrykker det, »til Bearbejdelse af Personligheden
under Omskuen paa de forskellige Livsinteresser
og Livsbestræbelser i Forhold til det særlige Kald og
den særlige Gerning«. Der opstod hos ham en stærk Følelse
af hans »almenvidenskabelige og filosofiske Dannelses
Utilstrækkelighed«, og »kun langsomt og efterhaanden
lykkedes det ham at forene det ivrige Specialstudium
med en stille Sysselsættelse med almindeligere Spørgsmaal«.

Der er næppe Tvivl om, at denne Krise hang nøje
sammen dels med hans Opfattelse af den klassiske Old%
% --- p. 212 ---
\setcounter{footnote}{0}%
\opage{212}tid, dels med den Retning, i hvilken han efter sin Begavelse
især kunde dyrke den klassiske Filologi. Han
kunde ikke betragte den klassiske Oldtid som det store
Forbillede, »som en regelgivende, uopnaaelig Mønsterverden«,
saaledes som den forudgaaende Tidsalder havde
opfattet den, og som mange af hans Kaldsfæller senere
har betragtet den. I Modsætning dertil stillede han en
mere nøgtern Opfattelse. Den græsk-romerske Oldtids
Kultur stod for ham som primitiv i Forhold til vor; men
netop derved stod den som simpel og klar og som et
betydningsfuldt Grundlag for senere europæisk Udvikling,
et Grundlag, der burde kendes paa første Haand
af dem herhjemme, der skulde have en dybere gaaende
Dannelse. Dette Grundlag trængte i hans Øjne til Supplering,
og han var en af de første herhjemme, der krævede
Plads i Skolen for en videre gaaende naturvidenskabelig
og matematisk Undervisning, end tidligere Tider
have anset for nødvendig. --- Men alligevel var det i
hans Arbejde paa den klassiske Filologis Omraade ikke
særlig den kulturhistoriske Side, han lagde Vægt paa.
Hans epokegørende Arbejde gik i Retning af kritisk
Undersøgelse af Haandskrifter og af textkritiske Studier.
Dette Arbejde paa at fremskaffe ægte og fejlfrie Texter,
et Arbejde, der er saa nødvendigt for en frugtbar Benyttelse
af de gamle Forfattere, var hans Livs Hovedværk,
og hans Berømmelse grunder sig paa det. Kun i
forholdsvis ringe Omfang har han ydet kulturhistoriske
Undersøgelser, uagtet det væsentlig var paa denne Side
af de klassiske Studier, han grundede Nødvendigheden
af en anselig Plads for dem i Skolen. Mænd som Boekh
og Wilamowitz, der ganske særlig beskæftigede sig med
Oldtidens Literatur- og Kulturhistorie, har neppe i saa
høj Grad som Madvig følt den Ensidighed, der saa længe
pinte ham. Paa Grund af deres Studiers Retning har de
% --- p. 213 ---
\setcounter{footnote}{0}%
\opage{213}bedre kunnet finde Tilknytningspunkter i nyere Tids
Kultur, og derved mere direkte se den almene Betydning
af deres specielle Virksomhed.

Men afset herfra har der rørt sig noget hos ham,
som formedelst Arbejdets Deling paa det intellektuelle
Omraade maa gøre sig gældende hos enhver Forsker, ---
en Følelse af, at det Fag, man dyrker, kun bringer i
Forhold til en enkelt Side ved Tilværelsen, --- en Trang
til et alsidigere Tankeindhold end det, Ens eget Fag kan
yde. Naturligvis rører en saadan Trang sig stærkere, jo
mere specielt Faget er. Men den kan røre sig hos alle,
hos Filosofen saavel som hos Grammatikeren eller Entomologen.
For Madvigs Vedkommende førte denne Trang
ham først og fremmest til filosofiske Studier. Den Hegelske
Filosofi's spekulative Fremgangsmaade tilfredsstillede
ham ikke, skønt han selv paa enkelte Punkter havde
Lejlighed til at interessere sig for Hegels Synspunkter,
især retsfilosofiske og sprogfilosofiske. Men han foretog
et sammenhængende Studium af Spinoza og Kant, og i
sine senere Aar, da desuden hans svækkede Syn gjorde
filologisk Arbejde vanskeligere for ham, beskæftigede han
sig indgaaende med filosofisk Læsning, eller rettere med
at høre filosofiske Skrifter blive læste for ham. Og han
betragtede det som en Lykke, at han derved kunde »afslutte
og fuldstændiggøre sit Livs aandelige Indhold«, og
at han havde faaet »Ro til at hengive sig til Betragtning
over Livets og Tilværelsens dybere Spørgsmaal«. Ogsaa
paa andre videnskabelige Omraader søgte han at orientere
sig. Han hørte Forelæsninger over Kemi hos Zeise
og lod sig af sin Ven Matematikeren Jürgensen indføre
i Grundlaget for den højere Matematik. Disse Studier
kom ham til Gode ved hans Virksomhed som Undervisningsinspektør
ved de højere Skoler og gav ham et alsidigt
Blik paa Skolespørgsmaal. Jeg hørte engang en af
% --- p. 214 ---
\setcounter{footnote}{0}%
\opage{214}hans matematiske Kolleger udtale, at »kun en Kolos
som Madvig« kunde have en Dom om alle Sider af den
højere Skoles Virksomhed.

Det var nu dog ikke blot Trang til et mere omfattende
Tankeindhold, der rørte sig hos ham. Hvad der
maaske er mest interessant hos denne universelle Forsker,
er Trangen til en Aandsdannelse, der kunde gennemtrænge
hele hans Personlighed --- Trangen til, hvad han
selv har kaldt »Bearbejdelse af Personligheden under
Omskuen paa de forskellige Livsinteresser«. Hos saare
mange Videnskabsmænd staar ikke blot det specielle Fag
isoleret overfor andre Fag, men hele Tankeindholdet staar
fjernt fra det egne personlige Liv, saaledes som dette
føres udenfor Studerekammeret eller Laboratoriet. Denne
Dobbelthed vilde Madvig ud over. Han vilde, at hans
Tankeliv ogsaa skulde udvikle hans Personlighed. Han
stræbte efter »at forene Tankens Trang til Klarhed med
Følelsens Inderlighed og Kærlighedskravet«. Og denne
Stræben satte da ogsaa sit Præg paa ham, i det milde
og sandhedskærlige Sind, der aabenbarede sig i hans
Færd. ---

Lad os forsøge at faa et Indtryk af den Retning, i
hvilken hans Eftertanke har bevæget sig. Hans Optegnelser
er korte og fragmentariske, affattede i hans tunge,
men i Regelen klare Sprog. De give hans foreløbige Udkast.
Men en bestemt Tankeretning spores tydelig\footnote[1]{I det følgende Resumé har jeg saa vidt muligt beholdt Madvigs egne Udtryk.}.

Madvig's Filosoferen førte ham tidlig til at tage Afstand
fra den spekulative Filosofis Forsøg paa at udrede
alt af den rene Tanke. Filosofiens Opgave var for ham
% --- p. 215 ---
\setcounter{footnote}{0}%
\opage{215}at forme det i Iagttagelsens Verden Fundne i sammenhængende
Begreber, saaledes at vor egen Bevidsthed
danner Udgangspunktet. Filosofi er Klaring af vor Bevidstheds
uafviselige Indhold og Forfølgelse deraf i alle
dets Konsekvenser. Og vi finder aldrig Bevidstheden
uden et bestemt Indhold. Forskellen mellem vort Jeg
og dets Indhold er ikke oprindelig. Vi finder fra første
Færd Bevidstheden som dobbeltsidig, paa én Gang subjektiv
og objektiv, og Opgaven er da den at gaa tilbage
til den almindeligste Form ad Analysens Vej.

Vi kan intet fornemme eller opfatte, uden at det
fremtræder i Forskel eller Modsætning til noget andet
indenfor det Givne, og det er netop formedelst saadanne
Forskelligheder i det givne Indhold, at vi kan blive os
bevidste som omfattende disse Forskelligheder i Enhed.
Det herskende Subjekt bliver sig bevidst som en Enhed
overfor en bevægelig og foranderlig Begrænsning. Men
indenfor det foranderlige Indhold kan Intet opfattes som
virkelig sket uden i Forbindelse med noget Andet, som
forud er sket, og i kontinuerlig Forbindelse med dette.
En saadan Kontinuitet mellem ethvert enkelt Faktum
og andre Fakta udtrykkes i Aarsagsbegrebet, som har
sit Udspring i Subjektets Grundbevidsthed, medens hver
enkelt Aarsag og Virkning maa findes ad Erfaringens
Vej. Den første Spire til Aarsagsloven viser sig i uvilkaarlige
Bevægelser, f. Ex. saadanne, der afværger en
Indgriben; derved forudsættes jo baade, at en Forandring
kan foregaa, og at den kan hindres. Alt, hvad der
saaledes gør sig gældende som Led i en kontinuerlig
Sammenhæng, tillægger vi Existens, det »at være«. Saaledes
staar alt, hvad vi kan erkende, for os som Del af
en Helhed og bliver kun derved Virkelighed for os. Og
al Vorden i det enkelte faar, gennem Sammenhængen i
alt, Indflydelse paa Helheden.

% --- p. 216 ---
\setcounter{footnote}{0}%
\opage{216}Hvad vi kalde muligt eller tilfældigt, er kun det, hvis
samtlige Betingelser vi har ufuldstændig Kundskab om.
Vildfarelse godtgør sig praktisk ved en Modstrid mellem
den Handling, visse Forestillinger fører os til at foretage,
og dens Resultat, idet den Forventning, vi nærede, ikke
bliver opfyldt. I Vildfarelsen, men især i Tvivlen haves
det stærkeste Bevis for Aandslivets Selvstændighed. Anerkendelsen
af noget kan være i Strid med Virkeligheden,
eller vi kan endog helt nægte eller opsætte vor Anerkendelse.

Aarsagsloven gælder ogsaa paa Villiens Omraade. I
enhver Tilstand, i hvilken der handles, er Personligheden
selv en bestemt Faktor, der aldrig bliver Nul overfor de
udefra kommende Paavirkninger. Beslutningens Udslag
bestemmes ved Styrkeforholdet mellem Personligheden og
Paavirkningerne. Men maaske kommer ikke hele min Natur
til at øve sin Indflydelse paa Udslaget; i saa Fald
kan jeg bagefter muligvis reagere mod dette, hvilket udtrykkes
ved, at Samvittigheden melder sig. Den sande
Frihed bestaar i Evne til at kæmpe for den indre Overbevisning
om sit Maal mod Hindringer fra anden Side. ---

Det Aandelige og det Legemlige er forskellige derved,
at hint er iagttagende, dette kun iagttaget. Der er ingen
Overgang fra det Legemlige til det Aandelige, skønt dette
helt igennem er knyttet til hint. Men det Aandelige kommer
ikke udefra ind i den Forbindelse. Aandeligt og Legemligt
er to Sider af det Samme. Ad Analogiens Vej
finder vi os i kontinuerlig Sammenhæng, ikke blot med
andre organiske Væsener, men ogsaa med det, der foregaar
i de kemiske Processer under Stoffernes Modsætning
og Vexelvirkning. I alle kemiske Forbindelser maa der,
sammen med Elementernes Forbindelse, foregaa en Forbindelse
af de i disse Elementer hvilende Aandselementer,
% --- p. 217 ---
\setcounter{footnote}{0}%
\opage{217}og noget saadant foregaar særlig ved Sjælens Tilbliven
under Organismens Dannelse.

Gennem de forskellige Naturformer spores der, især
naar vi lægger Lamarck's Hypotese til Grund, en opstigende
Række af Trin fra den uorganiske Materie gennem
de organiske Væseners Mangfoldighed til det Aandens
og Kærlighedens Liv, der kan udvikle sig i den
menneskelige Verden. Og gennem alt Stof, legemlig Form
og Bevidsthedsliv udfolder der sig en stor Tilværelsesvillie.
Det er ingen Krænkelse af det Aandeliges høje
Værdi som det Inderste i Mennesket, at hele Naturen
betragtes som en Stræben henimod dette Højdepunkt. ---

Til disse erkendelsesteoretiske og naturfilosofiske Ideer
slutter sig en Række Tanker om Grundlaget for det
Etiske. 

Heller ikke paa dette Omraade er den rene Tanke
eneraadig. En Fornuftmoral kan kun have Betydning
som Klaring og Afrunding af de Livsdrifter, som ligger
i Menneskenaturen; i sig selv vilde den intet Indhold
have. Der ytrer sig hos Mennesket en Villen, en Drift
til Eksistens, ikke blot til Ens egen Eksistens, men ogsaa
til medlevende Individers Eksistens. Baade i Medglæde
og i Medlidenhed viser det sig, at Individet føler sig forbundet
med andre Væsener, især de nærmest staaende.
Og naar Forstaaelsen af de andres Væsen og Kaar udvikler
sig, bliver den blotte Medfølelse til en Totalfølelse,
der ytrer sig som Retfærdighed. Den Lov for den enkelte,
som heraf udspringer, er Udtryk for en guddommelig
Villie, for Guds i os indstrømmende Væsen; ti denne
Lov angiver den bestemte Udviklingsretning for den
Helhed, indenfor hvilken Mennesket er blevet til. Kun
i en udviklet social Tilstand kan Pligtfølelsen blive fast
og klar. Den udvikler sig i Vexelvirkning med de forskellige
Livsformer, som Menneskeheden gennem uminde%
% --- p. 218 ---
\setcounter{footnote}{0}%
\opage{218}lige Tider famlende og forsøgsvis har uddannet og stadig
uddanner for et harmonisk Samliv. Enhver moralsk Dom
begrundes ved Individets Forhold til Slægt, Folk og
Menneskehed, altsaa til den Helhed, det tilhører. Det
gælder en Harmoni mellem den Selvopholdelsestrang og
den Sympati, som begge ligger i Menneskets Natur. ---

Disse erkendelsesteoretiske og etiske Reflexioner vidner
i al deres Korthed og Foreløbighed om, at Madvig's
Tanker idethele bevægede sig i Retning af, hvad der i
Slutningen af det 19. Aarhundrede gjorde sig gældende
paa det filosofiske Omraade. Den Interesse og Velvillie,
han ved saa mange Lejligheder viste yngre filosoferende
Forskere, bliver forstaaelig herved. Den Yngre, der i
Samtaler med ham kom ind paa filosofiske Spørgsmaal,
fik det Indtryk at staa overfor en Mand, der tænkte paa
selvstændig Maade. Offentlig udtalte han sig ikke om
Spørgsmaal af denne Art, og særlig ikke om sin Stilling
i religiøse Spørgsmaal, den, som vi nu vil søge at karakterisere.

I Selvbiografien antyder han klart nok, at hans religiøse
Meninger mere og mere afveg, ikke blot fra Kirkens
Lære, men ogsaa fra selve Kristendommen. Som Skolediscipel
i Frederiksborg drøftede han med Kammerater
den Vanskelighed for Aabenbaringstroen, der laa i, at
det aabenbarede Indhold var kommet til os med den
Tilfældighed og Usikkerhed, som er forbunden med al
historisk Overlevering. Særlig fremhæves de kritiske og
filologiske Tvivl angaaende Bibeltexten. Som moden Mand
og Forsker kom han, efterat han (i November 1848)
havde overtaget Kultusministeriet, til at staa overfor at
forberede Opfyldelsen af den nye Grundlovs Løfte om en
% --- p. 219 ---
\setcounter{footnote}{0}%
\opage{219}Forfatning for den danske Kirke. Men han følte, at han
ikke med Tillid, Fasthed og Alvor kunde lede Forhandlinger
i denne Retning, fordi han, skønt paa sin Maade
inderlig religiøst stemt, ikke stod paa hele Aabenbaringstroens
Standpunkt. Det var hans Pligt at befæste og opbygge;
nedrive og svække vilde han ikke; --- og han
frygtede for at vakle eller hykle. Disse Betænkeligheder
er tiltagne med Aarene. Han blev valgt til Medlem af
Kirkekommissionen af 1868, men mener, at han havde
gjort rettest i ikke at tage imod dette Valg, fordi han i
de mellemliggende Aar havde fjernet sig længere og bestemtere
fra den faste kirkelige Grundvold. Følgen var,
at han indenfor Kommissionen optraadte mere konservativt
end tidligere, idet han mente at burde hævde det
Bestaaendes Ret.

Hvad der holdt ham tilbage, har han klart udtalt i
de efterladte Optegnelser. Han opkaster saaledes det
Spørgsmaal, hvorledes et aandeligt Samliv mellem Mand
og Hustru er muligt i Forhold til Religionen, hvor Manden
er Filosof og Videnskabsmand. Og naar han følte
den ulykkelige Stilling mellem den erkendte Sandhed og
Frygten for, at det at udtale den overfor andre kunde
virke nedbrydende uden Mulighed for Erstatning og Nybygning,
saa tænkte han her særlig paa »nære og kære
Væsener, i hvis Følelse der ikke bør gribes ind, fordi
man ikke véd, om man kan føre dem trøstende og tilfredsstillende
ad nye Veje«. Hvad vi her møder hos Madvig,
hænger sammen med det milde og bløde i hans Karakter.
Og han har vel ogsaa følt, at for en Tænker er Grundene
vigtigere end Resultaterne, og at Resultaterne af
dem, der ikke forstaa Grundene, gerne opfattes aldeles
udvortes og dogmatisk og forbindes med de Forestillinger,
der gennem Aarhundreder er blevne knyttede til
Kætterier. Han forholdt sig søgende og kunde kun for%
% --- p. 220 ---
\setcounter{footnote}{0}%
\opage{220}staas af Søgende --- og hvor mange er de? Negtes kan
det dog ikke, at det vilde have været af stor Betydning,
om en Mand med hans videnskabelige og sociale
Autoritet frit og aabent havde udtalt sine Anskuelser.
Overfor de Nærmeste maatte det dog have været ham
muligt at skaffe sin personlige Overbevisning Respekt og
Anerkendelse. Der ytrede sig her hos ham en Ængstelighed,
som ogsaa prægede hans Holdning under den store
politiske Kamp i 80'erne.

I Samtale dølgede han hverken sit religiøse Stade eller,
hvad der holdt ham tilbage fra at udtale sig. Da jeg i
Aaret 1874 havde holdt et Foredrag i Studenterforeningen
om Darwinismen og Filosofien, der blev trykt i »Nær
og Fjern«, udtalte han paa en hjertelig Maade sin Sympati
for min Tankegang, idet han tilføjede, at Tanken
om et Syndefald forekom ham absurd. Derpaa sagde han
med et lunt Smil: »Men jeg tør ikke sige det til mine
Døtre!« Nogle Aar efter henvendte jeg mig til ham for
at foreslaa ham at underskrive en Opfordring til at bidrage
til en Billedstøtte af Spinoza, der med Midler fra
alle Lande skulde rejses i Haag, og jeg mindede ham
om, at han ofte havde omtalt Spinoza med stor Sympati
(en Sympati, der nu, gennem de efterladte Optegnelser,
er fuldt forstaaelig). Han svarede da, at han havde
Betænkelighed paa Grund af den Fordom, der stadig af
mange næredes overfor Spinoza, og at der var Hensyn,
han var nødt til at tage. For ikke at blive misforstaaet
forklarede han, at det var Hensyn til hans Hjem, han
tænkte paa. Senere bestemte han sig dog til at underskrive
og fik betydelige Bidrag samlede.

Hans store Hensynsfuldhed overfor sine Nærmeste har
ogsaa været en af Grundene til, at han stadig gik i Kirke,
hvilket havde til Følge, at han blev regnet for en af Kirkens
Støtter. Dog var Trangen til ikke at skille sig fra
% --- p. 221 ---
\setcounter{footnote}{0}%
\opage{221}sine Nærmeste ikke den eneste Grund til hans Kirkegang.
Han har havt Trang til Fællesskab\footnote[1]{Naar han (ifølge Ligtalen) ogsaa gik til Alters med sin Familie, har det ogsaa været af Trang til aandeligt Fællesskab. Han har kunnet sige, hvad Rousseau sagde, da han gik til Alters i Neufchatel, skønt det for ham var en blot Ceremoni: Je suis avec mes fréres!}, og han har
været stærkt betagen og har følt sig styrket af Orgelspillet
og Kirkesangen, --- just ikke af Prædikenen, som
man ser af følgende Optegnelse.

»I Kirken. Orgelspil og Sang, vel ledet og vel udført,
opløfter mig i en Stemning, hvori der paa én Gang
er Følelse af Menneskehedens Liv, Længsel, Kærlighed og
Mod og en Delagtighed i alle Tiders og Steders menneskelige
Samliv i Aand og Kærlighed. Men saa kommer Prædikenen
med dens saa ofte haltende og modsigende Dogmatik,
dens Uklarhed og skæve Slutninger. Kun den rent
etiske Opfordring tiltaler«.

Ogsaa i en, som det synes, senere nedskreven Udtalelse
fremhæver han skarpt Modsætningen mellem den
stemningsfulde Liturgi og den dogmatiske Prædiken.

Efterat have læst disse Udtalelser faar man et underligt
Indtryk af den Tale, der blev holdt ved hans Begravelse,
og som »Supplementernes« Udgiver har aftrykt
bag i Bogen. Han prises her som Kirkens trofaste Søn.
»Søndag efter Søndag sad han her i Kirken paa sin vante
Plads, medens Evangeliet blev forkyndt«. Prædikanten
har levet i lykkelig Uvidenhed om, hvad hans mest fremragende
Tilhører mente om hans »Forkyndelse«.

Nu er det jo saa, at den store, ofte vældige Poesi,
der er udformet i den kirkelige Kultus, ikke er Kirkens
Ejendom, men tilhører os alle. Det er Menneskehedens
Værk, som al Poesi er det. Enhver har her, som paa
andre aandelige Omraader, Ret til at tage, hvad han
kan bruge i sin aandelige Husholdning. Og det er lige%
% --- p. 222 ---
\setcounter{footnote}{0}%
\opage{222}ledes en naturlig Trang, at man gerne vil dele med sine
Nærmeste, hvad der for dem giver Livet Højhed og
Værdi. Men Mistydning ligger nær, naar man forsøger
at tage Kærnen uden at gøre aabenbart, at man knækker
Skallen; det har til Følge, at man bestyrker andre i, at
Skal og Kærne ikke kan skilles. I hvert Tilfælde maa
man haabe, at Agtelsen for den enkelte Personlighed og
dens Ret til at tage sit eget Stade i Livsspørgsmaalene
vil blive anerkendt for Alvor, ikke blot i større Kredse,
men ogsaa i de inderligste Forhold mellem Mennesker.
Da vil en Lempelse, der strider mod, hvad personlig
Sandhed kræver, blive umulig i alle Parters Øjne. Jo
mere Kvinden udvikler sig til at kunne tage selvstændigt
Stade overfor Livsspørgsmaalene, des mere vil ogsaa
Madvig's bekymrede Spørgsmaal om, hvorledes en Filosof
og Videnskabsmand kan have aandeligt Samliv med
sin Hustru, finde sin Besvarelse --- af sig selv. ---

Lad os nu se, hvilke Tanker der fostredes om det religiøse
Problem i den store Filologs Lønkammer.

Gudsforestillingen skyldes en Trang til Afslutning, til
en sidste Grund eller Aarsag, en Utilfredsstillethed ved
Tilværelsens Mangfoldighed og stadige Forandringer. Men
Tanken indser, at Begreberne Grund og Aarsag kun
finder berettiget Anvendelse indenfor den i Erfaringen
givne Tilværelse. Gud kan da ikke søges udenfor Verden,
men maa opfattes som selve Tilværelsens Enhed og Helhed,
--- som den, der rører sig i alt, i al Tænken og Villen,
som i al Bevægelse. Kun gennem dens Ytringsformer og
dens Virksomheder kender vi Guddommen. At kalde den
personlig, er uberettiget, da Begrebet Personlighed forudsætter
et Vexelvirkningsforhold til noget Ydre og derfor
% --- p. 223 ---
\setcounter{footnote}{0}%
\opage{223}kun passer paa endelige Væsener. Det rette Gudsbegreb
findes kun ved at se den indre Sammenhæng mellem alle
Naturlove og Naturkræfter og ved at agte paa den Udviklingsretning,
der kundgør sig i Tilværelsen. Den, der
er Led i den fremskridende Udvikling, føler sig delagtig
i en stor, universel Stræben, og denne Følelse er den religiøse
Følelse. Gennem den strømmer det hele Livstræs
Saft ind i hver Knop og hver Blomst. Men naar Tanken
forsøger paa anskuelig Maade at udtrykke, hvad der
ligger til Grund for denne Følelse, har den kun Billeder
fra den endelige Verden til Raadighed. Derved forstaas
de modsatte Indtryk, den kirkelige Liturgi og den dogmatiske
Prædiken gør.

Det er muligt, at de fleste Mennesker ikke kan undvære
mystiske og dogmatiske Former for Religion. Men
det er ikke sikkert, at den kristelige Tro i dens dogmatiske
Former i meget høj Grad er trængt ind i Menigheders
og Individers virkelige Bevidsthed; man levede
og lever snarere i en almindelig, mere eller mindre ren
Følelse af Gudhengivenhed og Fortrøstning, og kun til
enkelte Tider skabtes der Fanatisme for uforstaaede
Formler og Navne.

Dette var en stille Mands stille Tanker, udsprungne
af dyb Sandhedstrang, men tilbageholdte af Nænsomhed
overfor Anderledestænkende, der stod ham nær. Det maa
ikke glemmes, at den udmærkede Mands egentlige Kraft
fandt Anvendelse i trofast Arbejde paa andre Omraader,
og det kan ikke undre, at han ikke har villet lade disse
Tankefragmenter komme frem, han, der ellers med kritisk
Alvor overvejede hver Linie, han skrev. Han havde
den Opfattelse, at det er meget faa, der tænker igennem
paa dette Punkt. Ogsaa hos det 19. Aarhundredes Tæn%
% --- p. 224 ---
\setcounter{footnote}{0}%
\opage{224}kere og Digtere --- Kant og Hegel, Schiller og Goethe ---
fandt han Uklarhed i deres Forhold til Kirken. Hvad
har --- spurgte han --- været den virkelige Overbevisning
og Følelse hos de tænkende Mennesker, der overhovedet
i det sidste Aarhundrede har vovet at fæste Blikket paa
dette Punkt? Vi gaar forklædte om mellem hverandre.

Ved denne sidste Sætning (et Citat) antager jeg, han
mener, at naar den store fælles Symbolik, som Kirken
har udviklet, opgives som ufyldestgørende, er der ikke
noget fælles Sprog, paa hvilket alvorlige Mennesker, som
dybt føle Livets Gaader, kan meddele sig til hverandre.
I Negationerne kan man mødes; men de er nu engang
ikke nok til et Sprog. ---

For mit Vedkommende mener jeg, at vi skylder Madvig's
Søn Tak, fordi han har besluttet sig til at udgive
disse »Supplementer«, der er et nyt Vidnesbyrd om, at
der grubles og tænkes mere i Lønkamrene, end man ofte
tror. Bag de ydre Former for det aandelige Liv rører der
sig hos de Bedste et indre Liv, som maaske engang bliver
istand til at skaffe sig Udtryksformer og Livsformer. Dog
er der neppe nogen Mulighed for, at Modsætningen mellem
det ydre og det indre Liv nogensinde vil forsvinde,
og det var heller ikke at ønske.

% --- p. 225 ---
\setcounter{footnote}{0}%

\opage{225}\essayhead{ARISTOTELES OM VENSKAB}

\romanhead{I}

Den mest omtalte og tillige den hyppigst misforstaaede
Tanke i Aristoteles' Etik er hans Bestemmelse af
Dyden, d. v. s. Indbegrebet af værdifulde Karakteregenskaber,
som den rette Midte mellem Modsætninger. Man
har fundet en overfladisk Middelvejsteori deri og ikke
lagt Mærke til, at Aristoteles ved Midte ikke forstaar
et mekanisk Gennemsnit, saaledes at, hvis f. Eks. 6 er
for meget og 2 for lidt, skulde 4 være det rette. Da
menneskelige Individer møder med højst forskellig Trang
og Evne, vil Midten, det vil egentlig sige Harmonien af
Sjælelivets forskellige Elementer, faa en forskellig Karakter
for hvert Individ. Hvad der er Dumdristighed hos
den ene, kan være Fejghed hos den anden. Det kommer
an paa, hvilke indre Betingelser den Enkelte møder
med. Ligesom Forstanderen for Legemsøvelser ikke foreskriver
samme Diæt for alle Deltagere, men tager Hensyn
til deres legemlige Konstitution, saaledes maa man
ved Fastsættelse af værdifulde Karakteregenskaber tage
Hensyn til den sjælelige Konstitution. Tilsidst gælder det
for hver enkelt at finde sin egen rette Midte.

Denne Tanke har sjelden fundet Anerkendelse i den
filosofiske Etik, maaske fordi den gør etiske Spørgsmaal
% --- p. 226 ---
\setcounter{footnote}{0}%
\opage{226}meget vanskeligere, end de fleste holde af at tænke
sig\footnote[1]{Smlg. \emph{Holberg}: \textit{Moralske Tanker} (1744) p. 539: »En, der ingen Galde har, kan ikke berømmes for Koldsindighed. At ride paa et spagt Æsel, er ingen Merite, men at regere en vild og modig Hest, det fortjener Applausum og Hænders Klappen. Dette tager Folk dog gemenligen ikke udi Agt.«}.

Dette Hensyn til individuelle Forskelligheder staar i
ejendommelig Modsætning til den udpræget sociale Karakter
af Aristoteles' Etik. Det er jo ikke sagt, at fordi
det lykkes den enkelte at harmonisere sine Drifter og
Evner, han derfor er rustet til at fyldestgøre de Krav,
Samfund og Stat fra deres Side stiller til ham. Det er
maaske derfor, Aristoteles i den nærmere Gennemførelse
af sin Lære om den rette Midte uvilkaarlig forlader det
individualistiske Synspunkt og anvender et socialt bestemt
Gennemsnit.

Men ogsaa fra en anden Side kan det Spørgsmaal opkastes,
om den enkelte virkelig kan være en afsluttet,
harmonisk lille Verden for sig. Der kunde være Elementer
i selve hans Natur, som pegede ud over denne indre
Verden. Hos Platon, som ogsaa vil hævde en Harmonilære,
skønt ikke med den vidtgaaende Anerkendelse af individuelle
Forskelligheder som hos Aristoteles, rører der
sig afgjort et saadant Element, nemlig Eros, Trangen
til at knytte sig til andre Væsener eller Bestræbelser og
derigennem gaa op i et mere omfattende Hele --- et Elskovsforhold,
et Liv i social Virken eller i en Tankeverden.
Hvor Platon, som i »Symposium«, lægger en saadan
Trang til Grund, fremstiller han en helt anden Type,
end hvor Harmonisynspunktet raader, som i tredie og
fjerde Bog af »Staten«. I mere dæmpet, men ligesaa principiel
Form fremtræder et analogt Element hos Aristoteles.
I ottende og niende Bog af sin Etik taler han om
% --- p. 227 ---
\setcounter{footnote}{0}%
\opage{227}Venskabet (filia), idet hele om et sympatisk Forhold til
andre Mennesker. Det er en noget anden etisk Type, der
fremtræder her, end i de Bøger, der handler om »Midten«.
I Venskabstrangen, den sympatiske Trang er der noget,
som peger ud over Individet og kan hindre Harmonien,
idet den knytter Mennesket til andre Mennesker, deres
Ve og Vel, deres Skæbner. Er det nu sikkert, at den
indre Harmoni kan bevares, naar den enkelte saaledes
drages ud i en større Kreds af Interesser og Skæbner?

I den første Bog af sin Etik, der handler om Lykken,
har Aristoteles mærket Problemet og udtrykt det paa en
naiv Maade, idet han, efterat have fastslaaet, at Lykke
forudsætter en vis Afsluttethed og Selvtilstrækkelighed
(Autarki), tilføjer, at denne Selvtilstrækkelighed ikke betyder
Isolation; den gælder ikke for den enkelte i et
Eneboerliv, men ogsaa for hans Forældre og Børn og
Hustru, fremdeles for Venner og Medborgere, eftersom
Mennesket af Naturen er et Samfundsvæsen. Men --- nu
bliver han betænkelig. Her maa sættes en Grænse, ti ellers
gaar det i det uendelige, idet vi maatte regne Forfædre
og Efterkommere og Venners Venner med! --- Aristoteles
staar her lige ved Problemet. Spørgsmaalet er, efter hvilket
Princip Grænsen skal sættes --- ud fra Individet selv
eller ud fra en større Kreds og dens Livsbetingelser. Der
kommer Uro ind i det afsluttede Lykkebegreb. --- Som
vi skal se, kommer Aristoteles tilbage til Problemet i en
anden Sammenhæng --- og det netop ud fra Sympatiens
egen Natur.

Jeg nævner her i Forbigaaende, at Problemet skærpes
i Stoicismen, som vil gøre den enkelte uafhængig
af andres Slethed og Ulykke, og i Kristendommen, idet
Tro (som Betingelse for den egne Salighed) og Kærlighed
kommer til at staa som Modsætninger. Allerede hos Inderne
rører Problemet sig. Vismanden Bharata var lige
% --- p. 228 ---
\setcounter{footnote}{0}%
\opage{228}ved at naa den højeste Fuldkommenhed --- men han gik
glip af den, fordi han ynkedes over en lille syg Gazelle!
--- Vi behøver ikke at gaa langt for at mærke Problemets
Brod. Hvem kan føle sig lykkelig, naar største Delen af
Menneskene gensidig ødelægger hverandre? ---

Selv om det er rigtigt, at der hos Aristoteles fremtræder
en Modsætning mellem Sympati og individuel Harmoni,
er det dog naturligt at spørge, om der ikke gives
et Synspunkt, fra hvilket Aristoteles' Opfattelse af Harmonien
mellem Sjælelivets Elementer som højeste Karakteregenskab
kan ses i Sammenhæng med den Vægt,
han lægger paa Venskabsforholdet i dets forskellige Grader
og Arter. Og et saadant Synspunkt giver Aristoteles
selv, og det et saadant, der stiller os overfor en
af de dybeste Tanker i hans Filosofi og overfor et af
de smukkeste og fineste Træk i hans Personlighed.

\romanhead{II}
% CHECK roman numeral read

Netop den Lykkelige trænger til Venner, siger Aristoteles.
Og det ikke blot, fordi det at have Venner er
det største af alle ydre Goder, men fordi den dygtige
og alvorlige Mand trænger til nogen, han kan øve Velgerninger
imod. Mennesket er, som Aristoteles atter i
denne Sammenhæng udtaler, et Samfundsvæsen, og Lykke
bestaar i Virksomhed, ikke i Nydelse. Hvad en nulevende
dansk Digter siger: »Jeg besværger Eder, Venner, at I
elsker Eders Lykke!« --- stemmer ganske med Aristoteles'
Tankegang. Lykken var for de Gamle en alvorlig Sag,
en Villiessag, Et med Hævden af Personligheden overfor
de ydre og indre Kræfter, der kunde splitte Livet eller
gøre det ensidigt. Det maa ikke glemmes, at den græske
Harmoni var vunden gennem Kamp med kaotiske Kræfter,
ligesom Olympierne havde maattet kæmpe med Titanerne,
og ligesom Grækerne havde maattet kæmpe med
% --- p. 229 ---
\setcounter{footnote}{0}%
\opage{229}Orientalerne, før den græske Klarhed og Harmoni kunde
blive til. Ogsaa hos de enkelte store Karakterer maa dette
have været saa. I sin Karakteristik af Sokrates har Professor
Heiberg vist, hvorledes meget tyder paa, at den
sokratiske Selvbeherskelse er vunden gennem Kamp.

Og som Lykken bestaar i Virken, i Udfoldelse af Sjælens
Kræfter, saaledes lægger ogsaa Venskabet sig for
Dagen i en Virken, en Aktivitet. Venskab er en Dyd,
eller det forudsætter Dyd. Det viser sig mere i at elske
end i at elskes. Hvad Nytte har man af Rigdom, Magt
og Ære, naar man ikke benytter disse Goder til en for
andre frugtbar Virken? Og dertil giver Venskab den
“ skønneste Lejlighed.

I den Tanke, at Venskab, ligesom Lykke og Dyd,
bestaar i Virksomhed, finder Aristoteles Forklaringen af,
at Velgørere saa ofte elske dem, de have vist en Velgerning,
mere, end de elskes af disse. Grunden hertil
kunde, mener han, ligge deri, at de i den, de have gjort
vel, og i hans Velfærd ser deres eget Værk, i hvilket de
har nedlagt noget af dem selv. Vort Værk er jo en
Del af os selv, en Virkeliggørelse af vore Muligheder. Vi
elsker vort Værk, fordi vi elsker at virke, og dette igen,
fordi det at være er at virke. Derfor, føjer Aristoteles
til, elsker Kunstnerne deres Værker mere, end de vilde
elskes af disse, om de blev levende. Og da især Digterne!
De have en aldeles overvættes Kærlighed til deres Frembringelser.
--- Især naar Værket har kostet Møje, vil det
være Genstand for Kærlighed. Derfor elsker Mødrene
gerne deres Børn mere, end Fædrene gør. Fødselen har
kostet Besvær, og Mødrene har en mere eftertrykkelig
Viden om, at Børnene er deres. Saa højt, siger Aristoteles
i en anden Sammenhæng, kan Mødrene elske deres Børn,
at de kan overlade dem til andre, hvis det er bedst for
Børnene. ---

% --- p. 230 ---
\setcounter{footnote}{0}%
\opage{230}Allerede hele denne Tankegang viser, at der er en
Sammenhæng mellem Kærligheden til andre og Kærligheden
til sig selv. Man øver for sin Ven det Gode,
man kan, og det for Vennens egen Skyld. Men ganske
saaledes staar det dygtige og alvorlige Menneske jo overfor
sig selv, især overfor det Ypperste hos sig, Tankelivet.
Han vil leve i og med sig selv, i sine Erindringer og i
sine Forventninger, og han føler sig hjemme i denne
indre Verden. Og ganske paa samme Vis lever han i sine
Venners Tanker, deres Glæder og Sorger. Vennen er et
andet Jeg. Man kan opfatte Forholdet til sig selv som
en Analogi til Forholdet til sin Ven, eller omvendt Venskabet
som en Forlængelse af Kærligheden til sig selv.

Naar man dadler Kærlighed til sig selv, er det, fordi
man kun tænker paa den Selvhævdelse, der ytrer sig paa
andres Bekostning. Men der gives en Kærlighed til sig
selv, der bestaar i Omsorg for de skønneste Goder, for
Udviklingen af vor Aand og vor Karakter. En saadan
Kærlighed til sig selv kan endog ligge til Grund for Opofrelse
af Liv og Ejendom. Man vil da hellere øve én stor
Gerning end mange smaa, og idet man for sit Land eller
for sine Venner vinder, hvad der for dem er nødvendigt,
vinder man det Skønne for sig selv. (Nødvendigt kalder
Aristoteles, hvad der kun har Værdi som Middel, skønt,
hvad der har Værdi i sig selv). Ja, en brav Mand vil
endog kunne give Afkald paa selv at komme til at udføre
en skøn Bedrift, for at hans Ven kan komme til
at udføre den. Han kan finde det skønnere, at hans Ven
faar Æren, end at han selv faar den.

En nærmere Begrundelse af denne Tankegang finder
Aristoteles deri, at vi bedre kan danne os et anskueligt
Billede af andre end af os selv. Vi ønsker at se Dygtighed
og Energi, fordi Liv og Lykke bestaar i kraftig Udfoldelse,
--- og det kan vi se hos vore Venner under Sam%
% --- p. 231 ---
\setcounter{footnote}{0}%
\opage{231}liv og Samvirken med dem for menneskelige Opgaver.
Det er nemlig ikke nok at græsse paa samme Mark, som
Kvæget gør, ved Siden af hverandre. Samvirken i teoretisk
Syssel er vel nok det, Aristoteles især tænker paa,
men han ser, at menneskelig Samvirken kan have mange
Former, og hvilken fælles Interesse der saa virker, kan
den begrunde et Venskabsforhold. Man kan holde Fest
sammen med sine Venner, eller spille Terning med dem,
eller gymnastisere med dem, eller gaa paa Jagt med
dem, ligesaa vel som man kan filosofere med dem.

Den, som ikke er Ven med sig selv, kan ifølge Aristoteles
heller ikke være Ven med andre i sand Forstand.
Der maa være Harmoni i et Menneske, for at han
skal kunne komme i harmonisk Forhold til andre Mennesker.
Og Harmonien betyder her navnlig Overensstemmelse
mellem, hvad Mennesket er og gør, og hvad han
ønsker og vil. Uden en saadan Overensstemmelse bliver
Livet splittet og uroligt, og et fast Forhold til andre
umuligt. Kun det »dygtige og alvorlige« Menneske kan
indgaa et sandt Venskabsforhold. I Modsætning til en
saadan Karakter stiller Aristoteles den »slette«, som ikke
kan være Ven med sig selv, fordi hans Liv er splittet;
han vil et og gør et andet; af Fejghed og Sløvhed kommer
han aldrig til at udføre, hvad han vil. Derfor finder
han ikke Behag i sig selv og holder ikke af at dvæle i
sit eget Indre, sine Erindringer og sine Forventninger.
Søger han Samkvem med Andre, er det for at glemme sig
selv. Virkeligt Venskab kan han ikke indgaa.

Det er maaske karakteristisk for Aristoteles som Græker,
at han i denne Sammenhæng kun kender to Typer,
den harmoniske og den disharmoniske, og har Tilbøjelighed
til at lade Modsætningen mellem dem falde sammen
med Modsætningen mellem Godhed og Slethed. Han
gaar ud fra, at den Slette altid er disharmonisk, og vender
% --- p. 232 ---
\setcounter{footnote}{0}%
\opage{232}saa uden videre Sætningen om, saa at det disharmoniske
Menneske ogsaa maa være slet. Men der gives disharmoniske,
nedstemte, svingende og ubeslutsomme Naturer,
som vel er ulykkelige, men ikke »slette«, og som meget
vel kan indtræde i et Venskabsforhold, naar den anden
Part har den rette Forstaaelse. Skønt et saadant Menneske
ikke finder noget Elskeligt hos sig selv, ikke har
Sympati med sig selv, ikke kan være Ven med sig selv,
kan han dog, netop paa Grund af Modsætningen, forstaa,
beundre og elske en sluttet og harmonisk Natur,
især naar denne kommer ham imøde med Forstaaelse.
Shakespeare har skildret et saadant Venskabsforhold.
Hamlet elsker Horatio, denne faste, trods al Lidelse
harmoniske Natur, og Horatio forstod sin Ven og havde
opdaget det ædle Hjerte i hans splittede Indre. Derfor
overdrager Hamlet ham ogsaa i sin Dødsstund at forklare
ham og hans Sag for dem, som ej forstaa. Et Eksempel
fra det virkelige Liv er Schleiermacher's Forhold til
Friedrich Schlegel (i dennes unge, før-katolske Tid). I et
Brev udtaler Schleiermacher, at Alt, hvad der synes fejlagtigt,
selvmodsigende og urigtigt hos Schlegel, hænger
sammen med hans Kaar, og at der som det Inderste i
hans Digten og Tragten stadig rører sig noget Stort, Sjældent
og Skønt. --- Forudsætningen for et Venskab mellem
en harmonisk og en disharmonisk Natur vil netop
være, at der midt under Spaltningen og Modsigelserne
hos denne sidste opdages en Stræben, en arbejdende
Villie, selv om denne ikke formaar at naa op til den
Harmoni, der raader i lykkeligere Naturers indre og ydre
Færd.

Havde Aristoteles fastholdt sin Lære om Harmoniens,
den »rette Midte«s Afhængighed af individuelle Forudsætninger,
vilde han kunne have udvidet sin Lære om Venskabet
med et interessant Afsnit. Det er ikke altid de
% --- p. 233 ---
\setcounter{footnote}{0}%
\opage{233}ringeste og fattigste Naturer, der mangler Harmoni. Det
kan netop være Fylden, de mange Muligheder, de rige
Kræfter, der hindrer Enheden og Sluttetheden. Maaske
er den Art Naturer hyppigere i nyere Tid end i Oldtiden.

Aristoteles har iøvrig selv i et andet Afsnit af sin
Etik skildret en disharmonisk Natur, som han i hvert
Fald ikke kalder helt »slet«. Det er Mennesker, som mangler
Selvbeherskelse. Enten følger de enhver Indskydelse,
uden at overveje, eller de fastholder ikke Overvejelsernes
Resultat i Handlingens Øjeblik. Og dog anerkender de
stadig det Rette; Formaalet for deres Stræben er godt,
men kan ikke trænge igennem i de enkelte Øjeblikke.
--- Her var i hvert Tilfælde Udgangspunkt for en Forstaaelse
og et Venskabsforhold tilstede, nemlig for den,
der kunde faa Øje for det atter og atter tilbagetrængte
Formaal.

Desuden kommer Aristoteles selv til at antyde en særegen
Grund til, at ulykkelige Naturer saa ofte ikke egner
sig til Venskab. Mandige Naturer, siger han, vil ikke, at
deres Venner skal dele deres Ulykke og sørge for deres
Skyld; det er dem nok, at de selv er ulykkelige. Saaledes
vil ogsaa, kan vi føje til, disharmoniske Naturer lukke
sig inde i sig selv, uden at kunne eller ville aabne sig
for andre, og det netop af Sympati med dem, de ellers
vilde kunne være komne i det inderligste Forhold til.
I Søren Kierkegaards Papirer kan man studere denne
Type, som forudsætter Selvbeherskelse udadtil, skønt
denne mangler overfor det indre Kaos.

\romanhead{III}
% CHECK roman numeral read

Den ædleste Art af Venskab er efter Aristoteles Karaktervenskabet,
og i det Foregaaende har vi holdt os
til denne Type. Det har den fasteste Grund, i Modsætning
til det Venskab, der skyldes Nytte eller Behag.
% --- p. 234 ---
\setcounter{footnote}{0}%
\opage{234}Nytte og Behag kan skifte, og Venskabet vil falde bort,
naar den ene ikke mere har nogen Fordel af den anden
(ligesom Staters Venskab hører op med den gensidige
Fordel), eller naar Ynden, og med den Forelskelsen er
forsvunden. Dog mener Aristoteles, at det Venskab, der
grunder sig paa Forelskelse, snarere kan gaa over til Karaktervenskab
end det, der grunder sig paa Nytte.

En Betingelse for det fuldkomne Venskab er længere
Samliv og Samvirken, hvorved Karaktererne kan staa
deres Prøve. Man kender ikke hinanden, før man har
spist en vis Mængde Salt\footnote[1]{Aristoteles hentyder til en bekendt Talemaade, Grækerne allerede brugte. Salt var Udtryk for Gæstevenskab. Hos en senere græsk Forfatter forekommer Udtrykket »at have spist en Skæppe Salt sammen.«} med hinanden. Alt, hvad der
kan medføre Samliv og Samvirken kan foranledige et
Venskabsforhold. Saaledes Slægtskab, Rejseselskab, Medborgerforhold
o. s. v. I det enkelte vil Venskabet faa
sin særegne Karakter efter de forskellige sociale og politiske
Forhold. ---

Venskabsforhold mellem Uligestillede kræver en særegen
Undersøgelse. En meget stor Ulighed i Væsen og
Kaar, med overvættes Fortrin paa den ene Side, vil
umuliggøre Venskab. Saaledes mellem Guder og Mennesker.
Skønt man derfor ønsker sine Venner alt Godt, vil
man derfor ikke ønske dem at blive Guder. Hvad Mennesker
angaar, vil den Mægtige let opnaa Venskabsforhold
for Nyttens og Behagets Skyld. Men en dygtig og alvorlig
Mand vil ikke let blive Ven med en i Magt Højerestaaende,
dersom denne ikke besidder ganske særegne
Karakteregenskaber, saa at den Beundring og Kærlighed,
disse vækker, kan veje op mod Uligheden i de ydre
Forhold.

Andre Ulighedsforhold --- mellem Ældre og Yngre,
Forældre og Børn, Mand og Hustru, Herre og Slave ---
% --- p. 235 ---
\setcounter{footnote}{0}%
\opage{235}nuancerer Venskabsforholdet paa mangfoldig Maade, i
Analogi med Forholdet mellem den Mægtige og den Ringe.
Forholdet mellem Mand og Hustru sammenlignes saaledes
med et Aristokrati, hvor Magten øves til Bedste for den
ringere Part. Manden raader i Kraft af sin Naturs Værdighed,
men han yder Hustruen den Indflydelse, der tilkommer
hende; vil han raade blot til egen Fordel, bliver
Ægteskabet at sammenligne med et Oligarki. Hvad Forholdet
til Slaver angaar, lærer Aristoteles, at man kan
være Ven med sin Slave, forsaavidt han er Menneske,
men ikke forsaavidt han er Slave. I sin Statslære forsvarer
han udtrykkelig Slaveriet som en i Naturen grundet
Institution. ---

Ikke mindst i hele dette Afsnit sporer man Aristoteles'
personlige Erfaringer. I den platoniske Skole, som han
har tilhørt i tyve Aar, har han kunnet erfare, hvorledes
Venskab kan opbygges gennem fælles Interesser og Arbejder,
og hvorledes Samliv og Samvirken kan føre til
gensidig Forstaaelse af Karaktererne. Ogsaa Forholdet
mellem Ældre og Yngre har han her kunnet se Betydningen
af, og da han selv blev Leder af en Skole, har han
lært dette Forhold at kende fra den anden Side overfor
sine Disciple. Forholdet til Mægtigere har han gennemlevet
ved sin Stilling til Alexander, hvis Lærer han havde
været, og til Hermeias, Herskeren i Atarneus, som tog
ham under sin Beskyttelse. Af hans Testamente ses, at
han har været en god Ægtemand og en god Herre overfor
sine Slaver.

Som Samlivet kan fremkalde Venskab, saaledes følger
ogsaa Pligter og Rettigheder indenfor Venskabsforholdet
af Samlivsforholdets Natur. De er ikke de samme indenfor
et Rejseselskab, indenfor et Familieliv, indenfor Medborgerforhold,
indenfor Gæstevenskab o. s. v. Retfærdighed
og Venskab bliver begge til indenfor Samlivsforhold,
% --- p. 236 ---
\setcounter{footnote}{0}%
\opage{236}og Retfærdigheden vokser i Finhed med Venskabsforholdet,
baade hvad Grad og Form angaar. Der gælder en
højere Ret mellem Venner end mellem dem, der staar
hinanden fjernere. Og paa sit Højdepunkt faar Retfærdighed
Kærlighedens Præg.

Denne sidste Ytring er sandsynligvis at forstaa saaledes,
at den fuldkomne Forstaaelse af individuel Karakter
og individuelle Forhold, som Venskabsforholdet dels
forudsætter, dels betinger, muliggør en Hensyntagen til
den enkelte Personligheds Karakter og Kaar, som den i
almene Regler formulerede Retfærdighed ikke formaar at
øve.

Som Aristoteles her gaar fra Sympatien over til Retfærdigheden,
saaledes kommer han omvendt i femte Bog
af sin Etik, der særlig handler om Retfærdigheden, til
den Tanke, at de almene Love aldrig kan udtømme alle
en Handlings individuelle Forhold og Betingelser, men
den formelle Retfærdighed suppleres eller korrigeres ved
Forstaaelse af denne individuelle Side af Handlingen, og
den Egenskab, i Kraft af hvilken en saadan Suppleren
eller Korrigeren bliver mulig, kalder han der Billighed
(Epieiki).

\romanhead{IV}
% CHECK roman numeral read

Naar det fuldkomne Venskab viser sig i Samvirken og
i Sympati, maa det være begrænset i Omfang. Vi har allerede
set, hvorledes Aristoteles fra en helt anden Side,
nemlig ud fra den Enkeltes Lykke, kommer til Nødvendigheden
af en Begrænsning af menneskelige Sympatier.
Men Begrundelsen heraf i Bøgerne om Venskabet er af
særlig Interesse. Aristoteles har et klart Blik for, at der
paa Følelseslivets Omraade gælder en Lov om omvendt
Forhold mellem Styrke og Omfang. Derfor maa Vennernes
Antal være begrænset. Den Enkelte kan ikke splitte
% --- p. 237 ---
\setcounter{footnote}{0}%
\opage{237}sig, og paa sit Højdepunkt kan Venskab kun angaa en
enkelt. Dertil kommer, at vore Venner helst skal være
Venner indbyrdes, og dette bliver vanskeligere, jo flere
de er. Denne Betragtning anvender Aristoteles ogsaa overfor
Antallet af Borgere i en Stat, skønt der dér kun er
Tale om politisk Venskab, som har en større Udstrækning
end personligt Venskab. Ogsaa dér er der en Grænse,
og 10 Myriader (100,000) mener Aristoteles er for mange
til en Stat. Han lægger nemlig i sin Statslære stor Vægt
paa, at alle Borgere kender hverandre indbyrdes. Fremdeles
anvender han denne Betragtning til Kritik af Platon's
kommunistiske Statsideal, ifølge hvilket Familieliv
og Ejendom ophæves, for den styrende Klasses Vedkommende,
for at alle, der hører til denne Klasse, kan betragte
hverandre som Brødre. Den Sympati, der kan
raade mellem Menneskene, naar alle Familiebaand falder
bort, vil, mener Aristoteles, blive meget tynd, som naar
man blander en lille Smule Vin i en stor Mængde Vand! ---

Vi ser her den græske Tanke om Begrænsningens Nødvendighed
fra dens smukkeste Side. Og vi forstaar bedre
Platon's og Aristoteles' Iver for at hævde Modsætningen
mellem Hellener og Barbarer. Det er gennem Livet i de
snevreste og inderligste menneskelige Forhold, at Mennesker
kan udvikle sig til at leve med Varme og Kraft i de
mest omfattende Forhold. Hvis det almindelige Humanitetsideal
skal føre til mere end slap Sentimentalitet, maa
det arbejde sig frem Skridt for Skridt. De gamle Tænkere
havde Ret i, at Vejen gaar igennem, ikke udenom
Nationaliteterne.

Dog blev det afgørende Skridt ud over den nationale
Begrænsning allerede gjort i den klassiske Oldtid. I Gæstevenskabet
og i Menneskeligheden overfor Fremmede, som
jo allerede hos Homer staar under Zeus' særlige Beskyttelse,
laa fra gammel Tid en Spire, der senere kunde ud%
% --- p. 238 ---
\setcounter{footnote}{0}%
\opage{238}foldes. Aristoteles anerkender Muligheden af Venskab mellem
Frie og Slaver, og da han især begrunder Slaveriet
paa, at Barbarerne af Naturen er Slaver, fører denne
Anerkendelse ud over national Begrænsning paa et væsentligt
Punkt. Der udviklede sig henimod Oldtidens Slutning,
under Indflydelse af de sociale Forhold, Alexanders
og Romernes Erobringer havde fremkaldt, en almenmenneskelig
Fællesfølelse. Det synes at have været den
stoiske Skole, gennem hvilken denne Udvikling kom til
Bevidsthed, hvorved det blev muligt, at Grækernes latinske
Lærling Cicero kunde tale om caritas generis humani.

\end{document}
